\documentclass[10pt]{article}
\usepackage[margin=1in]{geometry}
\usepackage{amsmath}
\usepackage{amssymb}
\usepackage{mathtools}
\usepackage{stmaryrd}
\usepackage{amsthm}

% the theorem styles and numbering of acmart (sigplan)
\newtheoremstyle{acmplain}
  {.5\baselineskip plus .2\baselineskip minus .2\baselineskip}
  {.5\baselineskip plus .2\baselineskip minus .2\baselineskip}
  {\itshape}{\parindent}{\scshape}{.}{.5em}
  {\thmname{#1}\thmnumber{ #2}\thmnote{ {\upshape(#3)}}}
\theoremstyle{acmplain}
\newtheorem{theorem}{Theorem}[section]
\newtheorem{lemma}[theorem]{Lemma}

\newcommand{\SEM}[1]{\llbracket #1 \rrbracket}
\newcommand{\key}[1]{\mathsf{#1}}
\newcommand{\pass}[1]{\mathit{#1}}
\newcommand{\lam}[1]{\lambda #1.\,}
\newcommand{\clos}[2]{\lambda^{#1}\, #2.\,}
\newcommand{\Let}{\key{let}\;}
\newcommand{\In}{\;\key{in}\;}
\newcommand{\inl}{\key{inl}\,}
\newcommand{\inr}{\key{inr}\,}
\newcommand{\Case}[5]{\key{case}\;#1\;\key{of}\;\key{inl}\,#2 \Rightarrow #3 \mid \key{inr}\,#4 \Rightarrow #5}
\newcommand{\tuple}[1]{\langle #1 \rangle}
\newcommand{\prj}[1]{\pi_{#1}\,}
\newcommand{\fun}[1]{\key{fun}(#1)}
\newcommand{\funref}[1]{\key{fun}_{#1}}
\newcommand{\pair}[2]{\key{pair}(#1, #2)}
\newcommand{\FST}{\key{fst}\,}
\newcommand{\SND}{\key{snd}\,}
\newcommand{\SET}[1]{\mathcal{P}(#1)}
\newcommand{\FSET}[1]{\mathcal{P}_f(#1)}
\newcommand{\V}{\mathbb{V}}
\newcommand{\X}{\mathbb{X}}
\newcommand{\pairv}[2]{\llparenthesis #1, #2 \rrparenthesis}
\newcommand{\tup}[2]{\key{tup}_{#1}(#2)}
\newcommand{\Left}[1]{\key{left}(#1)}
\newcommand{\Right}[1]{\key{right}(#1)}
\newcommand{\emptytup}{\langle\rangle}
\newcommand{\prim}[3]{#1(#2, #3)}
\newcommand{\res}[1]{\lceil #1 \rceil}
\newcommand{\FV}{\mathrm{FV}}
\newcommand{\ext}[3]{#1(#2{:=}#3)}
\newcommand{\app}{\star}
\newcommand{\NE}{\neq\emptyset}
\newcommand{\Rel}[2]{\mathcal{R}(#1, #2)}
\newcommand{\ApproxLR}[2]{\mathsf{ALR}(#1, #2)}
\newcommand{\consis}{\mathit{consis}}
\newcommand{\obseq}{\cong}

% the simulation proof
\newcommand{\ev}[3]{\key{eval}\;#1\;#2\;#3}
\newcommand{\rt}[2]{\key{ret}\;#1\;#2}
\newcommand{\nxt}{\key{next}\;}
\newcommand{\stp}{\key{stop}\;}
\newcommand{\error}{\key{error}}
\newcommand{\retq}{\key{ret?}}
\newcommand{\Mach}[1]{\mathcal{M}^{#1}}
\newcommand{\step}[1]{\mathit{step}^{#1}}
\newcommand{\start}[1]{\mathit{start}^{#1}}
\newcommand{\obs}{\mathit{obs}}
\newcommand{\Final}{\mathit{Final}}
\newcommand{\Term}{\mathit{Terminates}}
\newcommand{\Div}{\mathit{Diverges}}
\newcommand{\Equiv}{\mathit{Equivalent}}
\newcommand{\FSim}{\mathit{FSim}}
\newcommand{\BSim}{\mathit{BSim}}
\newcommand{\cloV}[2]{\key{clo}(#1,\, #2)}
\newcommand{\closV}[2]{\key{clos}(#1;\, #2)}
\newcommand{\tupV}[1]{\key{tup}(#1)}
\newcommand{\pairV}[2]{\key{pair}(#1, #2)}
\newcommand{\funV}[1]{\key{fun}(#1)}
\newcommand{\codeV}[1]{\key{code}_{#1}}
\newcommand{\inlV}{\key{inl}\;}
\newcommand{\inrV}{\key{inr}\;}
\newcommand{\Kf}[2]{\key{#1}(#2)}
\newcommand{\halt}{\key{halt}}
\newcommand{\assign}[3]{#1 \mathrel{:=} #2;\; #3}
\newcommand{\Casei}[4]{\key{case}\;#1\;\key{of}_{#2}\;\key{inl} \Rightarrow #3 \mid \key{inr} \Rightarrow #4}
\newcommand{\evi}[3]{\key{eval}\;#1\;#2\;#3}
\newcommand{\rti}[3]{\key{ret}\;#1\;#2\;#3}
\newcommand{\exc}[4]{\key{exec}\;#1\;#2\;#3\;#4}
\newcommand{\tnext}{\mathit{tuple\text{-}next}}
\newcommand{\cnext}{\mathit{clos\text{-}next}}
\newcommand{\lto}{\longrightarrow}
\newcommand{\approxk}{\approx_k}
\newcommand{\approxq}{\approx^{?}}

% the interpreter proof
\newcommand{\tmo}{\key{timeout}}
\newcommand{\Ret}{\key{ret}\;}
\newcommand{\bind}{\mathbin{>\!\!>\!\!=}}
\newcommand{\Do}{\key{do}\;}
\newcommand{\lift}[1]{\lfloor #1 \rfloor}
\newcommand{\run}{\mathit{run}}
\newcommand{\outc}{\mathit{out}}
\newcommand{\RelRes}[1]{\mathit{RelResult}_{#1}}
\newcommand{\EV}[2]{\key{eval}^{#1}_{#2}}
\newcommand{\EVS}[2]{\key{evals}^{#1}_{#2}}
\newcommand{\APP}[2]{\key{apply}^{#1}_{#2}}
\newcommand{\CALL}[2]{\key{call}^{#1}_{#2}}
\newcommand{\BR}[2]{\key{branch}^{#1}_{#2}}
\newcommand{\Interp}[1]{\mathcal{I}^{#1}}

% X86
\newcommand{\rax}{\%\key{rax}}
\newcommand{\rdi}{\%\key{rdi}}
\newcommand{\rsi}{\%\key{rsi}}
\newcommand{\xrun}{\key{run}}

% the typed pipeline (Part II): types, tags and casts
\newcommand{\Tl}[1]{\textrm{#1}^{\mathsf{T}}}  % a typed language: \Tl{Clos}
\newcommand{\TAny}{\key{Any}}
\newcommand{\Tflat}[1]{\langle\!\langle #1 \rangle\!\rangle}  % the flat type G as a type
\newcommand{\Tcode}[1]{\key{code}\;#1}         % the type of code
\newcommand{\Tfun}{\key{fun}}                  % a flat type, and a tag
\newcommand{\Tsum}{\key{sum}}                  % a flat type, and a tag
\newcommand{\Ttup}[1]{\key{tup}\;#1}           % the flat type of tuples of n components
\newcommand{\Tint}{\key{int}}                  % the tag of integers
\newcommand{\Tttup}[2]{\key{tup}_{#1,#2}}      % the tag of tuples of n components, read at i
\newcommand{\Tflatof}[1]{\mathit{flat}(#1)}    % the flat type a tag tests
\newcommand{\Tout}[1]{\mathit{out}(#1)}        % the flat type of a primitive's result
\newcommand{\Tinject}[1]{\key{inject}_{#1}\,}
\newcommand{\Tproject}[1]{\key{project}_{#1}\,}
% the typed pipeline: terms
\newcommand{\Tcodetm}[3]{\key{code}(#1\,#2.\,#3)}  % closed code: environment, argument, body
\newcommand{\Tpack}[2]{\key{pack}(#1, #2)}     % a closure: code and environment
\newcommand{\Tcallc}[2]{\key{callc}(#1, #2)}   % call a closure
\newcommand{\Tembed}[1]{\ulcorner #1 \urcorner}  % Mon^T into GlobalFun^T
\newcommand{\Tunit}{\Tinject{\Ttup{0}}\emptytup}  % the empty tuple, injected
\newcommand{\Tcheck}[1]{\pass{check}_{#1}}     % the X86 test of a tag
% the typed pipeline: tails of CFG^T, and frames of Imp^T and CFG^T
\newcommand{\Treturn}[1]{\key{return}\;#1}
\newcommand{\Tgoto}[1]{\key{goto}\;#1}
\newcommand{\Tcase}[4]{\key{case}\;#1\;\key{of}_{#2}\;\key{goto}\;#3 \mid \key{goto}\;#4}
\newcommand{\Tcut}[2]{#1|_{#2}}                % the first s slots of a frame
% the typed pipeline: semantics
\newcommand{\Tguard}{\key{guard}^{\mathrm{p}}}            % the guard of a predicate
\newcommand{\Taccepts}[1]{\mathit{accepts}_{#1}}          % the test of a tag, for a kind of function
\newcommand{\Tfuntwo}{\mathit{fun}_2}                     % a function is an entry of a graph
\newcommand{\Tfunfour}{\mathit{fun}_4}                    % a function is a pair
\newcommand{\Tchecked}[1]{\key{checked}^{#1}}             % checked through a map on kinds
\newcommand{\Tkind}{\key{kind}}
\newcommand{\Tkindf}{\key{kind}^{\mathrm{f}}}             % the kind of a flat type
\newcommand{\Tkindfour}{\key{kind}_4}
\newcommand{\Tkeep}[1]{\key{keep}^{#1}}
\newcommand{\Thit}{\key{hit}}
\newcommand{\Tletv}{\key{let}^{\mathsf{T}}}               % bind to the next slot
\newcommand{\Same}{\mathit{Same}}
\newcommand{\TdynV}[2]{\key{dyn}_{#1}(#2)}                % a value of type Any: v tagged with G
\newcommand{\TcodeV}[1]{\key{code}(#1)}                   % code, as a value
\newcommand{\Tres}[1]{\lceil #1 \rceil^{\mathsf{T}}}      % a primitive's result, typed
% the typed pipeline: X86^T
\newcommand{\Tst}[4]{\langle #1;\, #2;\, #3;\, #4 \rangle} % a state type: registers, frame, open closure, flag
\newcommand{\codeA}{\key{code}\,\alpha}                   % the register type of the open closure's code
\newcommand{\tupA}{\key{tup}\,\alpha}                     % the register type of the open closure's environment
\newcommand{\Tup}[1]{#1{\uparrow}}                        % an operand of type <<G>> as one of type Any
\newcommand{\Tchk}[3]{\key{chk}_{#1}\; #2, #3}            % project's check, and the move
\newcommand{\Tentry}[1]{\mathit{entry}(#1)}
\newcommand{\Tcalled}[1]{\mathit{called}(#1)}
\newcommand{\Treturned}[1]{\mathit{returned}(#1)}
\newcommand{\Twipe}[1]{\mathit{wipe}(#1)}
\newcommand{\Tset}{\key{set}}
\newcommand{\Tunset}{\key{unset}}


% ISWIM-mu (Part I): types, terms and X86-mu
\newcommand{\Ml}[1]{\textrm{#1}^{\mu}}                  % a language of the ISWIM-mu pipeline: \Ml{Clos}
\newcommand{\Mmu}[2]{\mu #1.\,#2}                         % a recursive type
\newcommand{\Mtup}[1]{\key{tup}\,#1}                      % a tuple type
\newcommand{\Mcodety}[3]{\key{code}\;#1\;#2\;#3}          % the type of code: environment, argument, result
\newcommand{\Munroll}[2]{#2[\mu #1.\,#2/#1]}              % the unfolding of mu X. A
\newcommand{\Mbool}{\key{bool}}                           % tup() + tup()
\newcommand{\Mfold}{\key{fold}\,}
\newcommand{\Munfold}{\key{unfold}\,}
\newcommand{\codeAB}[2]{\key{code}\alpha\;#1\;#2}         % the register type of the open closure's code
\newcommand{\Tdown}[1]{#1{\downarrow}}                    % (unused since X86mu writes unfold a)
\newcommand{\Mxcase}[4]{\key{case}\;#1\;\key{of}_{#2}\;#3 \mid #4}  % X86-mu's case pseudo-instruction

\newenvironment{syntax}{\[\begin{array}{l@{\quad}rcl}}{\end{array}\]}
\setlength{\emergencystretch}{3em}

\title{Closure Conversion: Definitions}
\author{}
\date{}

\usepackage[colorlinks=true,linkcolor=blue,urlcolor=blue,bookmarksnumbered=true]{hyperref}

\begin{document}
\maketitle
\tableofcontents

\section*{About this document}
\addcontentsline{toc}{section}{About this document}

This is a summary, for review, of the definitions in the Agda development
of the closure-conversion compiler in \texttt{agda}, in paper-ready
notation: the languages, the passes, the semantics of each language in each
of the four styles (graph model, CEK machines, clocked interpreters,
reduction semantics), the relations of each style's proofs, and the theorem
statements. Each definition names its Agda module and definition. It uses
the notation of \texttt{talks/ITRS-2020} and named variables instead of de
Bruijn indices; each part has notes on where it doesn't match the Agda
literally. Relations put the source on the left and the target on the
right.

The closure-conversion languages have the names they have in the paper,
which say what changed; the Agda modules keep their numbered names:
Clos is the Agda \texttt{Clos2} (closures that list their free
variables), EnvClos is \texttt{Clos3} (code that reads its free
variables from an environment tuple), FlatClos is \texttt{Clos4} (a
closure is a pair of code and environment), and GlobalFun is
\texttt{Clos5} (all code is global definitions).

\begin{itemize}
\item Part~\ref{part:mu}, the main experiment: $\Ml{ISWIM}$, a typed ISWIM
  with iso-recursive types, and its compiler, every intermediate language
  intrinsically typed, to the typed assembly language $\Ml{X86}$, proved
  correct in the four styles, with adequacy, the lemmas that nothing errs,
  and one end-to-end theorem for a main program of type $\key{Int}$ common
  to the four styles (\texttt{agda/Mu}, module prefix \texttt{Mu.};
  \texttt{papers/PLDI-2027/notes/iswim-mu-design.md}).
\item Part~\ref{part:dyn}, what dynamic typing and errors add: the
  dynamically typed pipeline from untyped ISWIM through cast insertion to
  $\Tl{X86}$ (\texttt{agda/}), the same compiler specialised to one
  recursive type $\TAny$, with the one operation that fails,
  $\key{project}$, and the graph model with errors that observing errors
  needs.
\item Part~\ref{part:untyped}, history: the untyped pipeline that preceded
  both (\texttt{agda/Untyped}), with the definitions all three share (the
  graph model's domain and operators, machines, simulations, clocked
  programs, reductions as machines; untyped ISWIM and its simplify).
\end{itemize}

\paragraph{Building.} \texttt{latexmk -pdf definitions.tex}
%===============================================================================
\part{ISWIM\texorpdfstring{$^{\mu}$}{-mu}: a Typed ISWIM with Recursive Types}\label{part:mu}

The compiler of the main experiment (\texttt{agda/Mu}, module prefix
\texttt{Mu.}; \texttt{agda/Mu/All.agda}):
\[
\begin{array}{l}
\Ml{ISWIM} \xrightarrow{\pass{simplify}} \Ml{ISWIM}
  \xrightarrow{\pass{enclose}} \Ml{Clos} \xrightarrow{\pass{optimize}} \Ml{Clos}
  \xrightarrow{\pass{concretize}} \Ml{EnvClos} \xrightarrow{\pass{delay}} \Ml{FlatClos} \\
\qquad \xrightarrow{\pass{globalize}} \Ml{GlobalFun} \xrightarrow{\pass{rco}} \Ml{Mon}
  \xrightarrow{\pass{lift}} \Ml{Imp} \xrightarrow{\pass{explicate}} \Ml{CFG}
  \xrightarrow{\pass{select}} \Ml{X86}
\end{array}
\]
The source $\Ml{ISWIM}$ is ISWIM, statically typed, with iso-recursive
types. Every language is intrinsically typed, down to the typed assembly
language $\Ml{X86}$, so a term that misuses a value cannot be written, and
no operation of any language can fail: nothing traps. Divergence and
general recursion come from the recursive types, not from a construct.
The pipeline is a parallel copy of the dynamically typed one of
Part~\ref{part:dyn} (which is, in effect, this pipeline specialised to
the single recursive type $\TAny$): the same passes, re-indexed by
richer types, with $\Mfold$ and $\Munfold$ where that pipeline has
$\key{inject}$ and $\key{project}$, and no cast insertion. The design
and its measurements are in \texttt{papers/PLDI-2027/notes/iswim-mu-design.md}.

The part has the types and languages (Section~\ref{sec:m-langs}), the
passes (Section~\ref{sec:m-passes}), then one section per style, each
with the semantics of the languages and the relations of its proofs: the
graph model (Section~\ref{sec:m-den}), CEK machines
(Section~\ref{sec:m-cek}), clocked interpreters
(Section~\ref{sec:m-interp}) and reduction semantics
(Section~\ref{sec:m-red}); and last adequacy, the no-error lemmas and
the end-to-end theorem that all four styles state
(Section~\ref{sec:m-thm}). The language-independent frameworks (the
graph model's domain and operators, machines as step functions,
simulations, clocked programs, reductions as machines) are those of
\texttt{agda/Framework}, unchanged; they were introduced for the untyped
pipeline and are defined in Part~\ref{part:untyped}
(Sections~\ref{sec:den-ops}, \ref{sec:obs}, \ref{sec:sims},
\ref{sec:u-interp} and~\ref{sec:red}). Of the Agda's names, a module
\texttt{Mu.X} is the file \texttt{agda/Mu/X.agda}, and the files named
in this part are under \texttt{agda/Mu}, except the shared frameworks,
under \texttt{agda/Framework}, \texttt{agda/proof/Framework} and
\texttt{agda/Semantics/Interp/Clock.agda}. Relations put the source on
the left and the target on the right throughout.

%===============================================================================
\section{Types and Languages}\label{sec:m-langs}

\subsection{Types}

The types (\texttt{Language/Type.agda}) are base types, functions, $n$-ary
tuples, binary sums and iso-recursive types $\Mmu{X}{A}$, with type
variables (de~Bruijn indices in the Agda: \texttt{Ty n} has $n$ free type
variables). Terms are indexed by closed types; a context $\Gamma$ or
$\Delta$ is a list of them. The type $\Mcodety{\Delta}{A}{B}$, in
$\Ml{FlatClos}$ and later, is that of closed code that takes an environment,
a tuple of type $\Mtup{\Delta}$, and an argument of type $A$.
\begin{syntax}
  & B &::=& \key{Nat} \mid \key{Int} \mid \key{Bool} \mid \key{Unit} \mid \key{Void} \\
  & A, C &::=& B \mid A \rightarrow C \mid \Mtup{(A_1, \ldots, A_n)} \mid A + C
    \mid \Mmu{X}{A} \mid X \mid \Mcodety{\Delta}{A}{C} \\
  & \Gamma, \Delta &::=& A_1, \ldots, A_n
\end{syntax}
The body $A$ of $\Mmu{X}{A}$ is unrestricted: $\Mfold$ and $\Munfold$
are explicit coercions, so even $\Mmu{X}{X}$ needs no contractiveness
condition.

The unfolding of $\Mmu{X}{A}$ is $\Munroll{X}{A}$, with type substitution
on \texttt{Ty} (\texttt{tren}, \texttt{tsubst}, \texttt{unroll}). The
result type of a primitive is $\Tout{\key{sub}} = \key{Int}$ and
$\Tout{\key{eq}} = \Mbool = \Mtup{()} + \Mtup{()}$: $\key{eq}$ gives
$\inl\emptytup$ for true and $\inr\emptytup$ for false, so $\key{case}$
serves as $\key{if}$. There is no type of tagged values, no flat type, no
tag, and no decidable equality of types: nothing tests a type at run
time.

\subsection{\texorpdfstring{$\Ml{ISWIM}$}{ISWIM-mu}}\label{sec:m-iswim}

ISWIM (Part~\ref{part:untyped}), intrinsically typed
(\texttt{Language/ISWIM.agda}), with $\Mfold$ and $\Munfold$, the two
coercions of the isomorphism between $\Mmu{X}{A}$ and its unfolding,
which do nothing at run time. A term is typed in a context,
$\Gamma \vdash M : A$; a variable is a typed index into $\Gamma$, so every
term is well scoped by construction.
\begin{syntax}
  & L, M, N &::=& x \mid \lam{x} N \mid L\;M \mid c
    \mid \tuple{M_1,\ldots,M_n} \mid \prj{i} M
    \mid \inl M \mid \inr M \\
  &&& \mid \Case{L}{x}{M}{y}{N} \mid \prim{p}{L}{M}
    \mid \Let x = M \In N
    \mid \Mfold M \mid \Munfold M
\end{syntax}
\[
\begin{array}{r@{\;:\;}l@{\qquad}l}
  x & A & x{:}A \in \Gamma \\
  \lam{x} N & A \rightarrow C & \Gamma, x{:}A \vdash N : C \\
  L\;M & C & L : A \rightarrow C,\ M : A \\
  c & B & c \text{ a constant of base type } B \\
  \tuple{M_1,\ldots,M_n} & \Mtup{(A_1, \ldots, A_n)} & M_i : A_i \\
  \prj{i} M & A_i & M : \Mtup{(A_1, \ldots, A_n)},\ 1 \leq i \leq n \\
  \inl M,\ \inr M & A + C & M : A \text{, resp. } M : C \\
  \Case{L}{x}{M}{y}{N} & D & L : A + C,\ \Gamma, x{:}A \vdash M : D,\ \Gamma, y{:}C \vdash N : D \\
  \prim{p}{L}{M} & \Tout{p} & L, M : \key{Int} \\
  \Let x = M \In N & C & M : A,\ \Gamma, x{:}A \vdash N : C \\
  \Mfold M & \Mmu{X}{A} & M : \Munroll{X}{A} \\
  \Munfold M & \Munroll{X}{A} & M : \Mmu{X}{A}
\end{array}
\]
A program is a closed term of any type; the end-to-end theorem is for a
program of type $\key{Int}$ (Section~\ref{sec:m-thm}).

\paragraph{Recursion without a construct.} There is no $\key{rec}$ or
$\key{fix}$. Self-application is typable through a recursive type
(\texttt{proof/Compiler/Examples/RecursionExample.agda}):
\begin{align*}
  \omega &= \lam{x} (\Munfold x)\;x : (\Mmu{X}{X \rightarrow \key{Int}}) \rightarrow \key{Int}
  \qquad \Omega = \omega\;(\Mfold \omega) : \key{Int} , \\
  \mathit{fix} &= \lam{f} g\;(\Mfold g), \quad
    g = \lam{x} f\;(\lam{v} (\Munfold x)\;x\;v), \quad
    x : \Mmu{X}{X \rightarrow (\key{Int} \rightarrow \key{Int})} ,
\end{align*}
where $\Omega$ diverges and $\mathit{fix}$ is the call-by-value fixed-point
combinator at $\key{Int} \rightarrow \key{Int}$; the example compiles a
countdown with $\mathit{fix}$ and runs it, and $\Omega$, on the
$\Ml{ISWIM}$ and $\Ml{FlatClos}$ machines and on the $\Ml{ISWIM}$, $\Ml{FlatClos}$
and $\Ml{X86}$ interpreters, by computation
(\texttt{RecursionInterp.agda}). So no semantics needs a least fixed point.

\subsection{\texorpdfstring{$\Ml{Clos}$}{Clos-mu}}

As $\Ml{ISWIM}$, but each $\lambda$ is a closure
(\texttt{Language/Clos2.agda}): its code $N$ is closed, with the parameter
$y$ and the names $x_1 \ldots x_n$ bound to the values of the captured
terms $M_1 \ldots M_n$, which have types $\Delta = A_1, \ldots, A_n$.
\begin{syntax}
  & L, M, N &::=& x \mid \clos{M_1\ldots M_n}{x_1\ldots x_n\, y} N
    \mid L\;M \mid c \mid \cdots \mid \Mfold M \mid \Munfold M
\end{syntax}
\[
  \frac{\Gamma \vdash M_i : A_i \ (1 \leq i \leq n) \qquad
        x_1{:}A_1, \ldots, x_n{:}A_n, y{:}A \vdash N : C}
       {\Gamma \vdash \clos{M_1\ldots M_n}{x_1\ldots x_n\, y} N : A \rightarrow C}
\]
The closure's type does not show $\Delta$. The other constructs are typed
as in $\Ml{ISWIM}$.

\subsection{\texorpdfstring{$\Ml{EnvClos}$}{EnvClos-mu}}

Each closure stores its captured values in one tuple $e$, from which its
code reads them with projections (\texttt{Language/Clos3.agda}).
\begin{syntax}
  & L, M, N &::=& x \mid \clos{M_1\ldots M_n}{e\, y} N
    \mid L\;M \mid c \mid \cdots \mid \Mfold M \mid \Munfold M
\end{syntax}
\[
  \frac{\Gamma \vdash M_i : A_i \ (1 \leq i \leq n) \qquad
        e{:}\Mtup{\Delta},\, y{:}A \vdash N : C}
       {\Gamma \vdash \clos{M_1\ldots M_n}{e\, y} N : A \rightarrow C}
  \qquad \Delta = A_1, \ldots, A_n
\]

\subsection{\texorpdfstring{$\Ml{FlatClos}$}{FlatClos-mu}}

A closure is a pair of closed code and the tuple of its free variables,
packed so that its type $A \rightarrow C$ hides the types $\Delta$ of the
tuple (\texttt{Language/Clos4.agda}): $A \rightarrow C$ is the existential
$\exists \Delta.\, \Mcodety{\Delta}{A}{C} \times \Mtup{\Delta}$ of typed
closure conversion (Minamide, Morrisett and Harper), with a list of types
as the witness. $\key{pack}$ builds a closure from code and an environment
of the same $\Delta$; $\key{callc}$ opens a closure and calls its code with
its environment and the argument. $\key{callc}$ is the only elimination of
the existential and is parametric in $\Delta$, so it needs no check.
\begin{syntax}
  & L, M, N &::=& x \mid \Tcodetm{e}{y}{N} \mid \Tpack{L}{M} \mid \Tcallc{L}{M}
    \mid c \mid \tuple{M_1,\ldots,M_n} \mid \prj{i} M \mid \inl M \mid \inr M \\
  &&& \mid \Case{L}{x}{M}{y}{N} \mid \prim{p}{L}{M}
    \mid \Let x = M \In N \mid \Mfold M \mid \Munfold M
\end{syntax}
\[
  \frac{e{:}\Mtup{\Delta},\, y{:}A \vdash N : C}
       {\Gamma \vdash \Tcodetm{e}{y}{N} : \Mcodety{\Delta}{A}{C}}
  \qquad
  \frac{\Gamma \vdash L : \Mcodety{\Delta}{A}{C} \quad \Gamma \vdash M : \Mtup{\Delta}}
       {\Gamma \vdash \Tpack{L}{M} : A \rightarrow C}
  \qquad
  \frac{\Gamma \vdash L : A \rightarrow C \quad \Gamma \vdash M : A}
       {\Gamma \vdash \Tcallc{L}{M} : C}
\]

\subsection{\texorpdfstring{$\Ml{GlobalFun}$}{GlobalFun-mu}}

As $\Ml{FlatClos}$, but the code of every function is a global definition
(\texttt{Language/Clos5.agda}). A program has a table
$\Sigma = s_0, \ldots, s_{m-1}$ of the \emph{signatures} of its
definitions $d_0, \ldots, d_{m-1}$, a signature $s_k = (\Delta_k; A_k \to
C_k)$ giving the types of the environment, the argument and the result
(\texttt{DSig}, \texttt{sig}), and a term is typed by the table as well as
by its context, $\Sigma; \Gamma \vdash M : A$: $\funref{k}$ is the code of
$d_k$, of type $\Mcodety{\Delta_k}{A_k}{C_k}$, so a reference to a missing
definition cannot be written. Every definition is typed over the whole
table, and a program has the type of its main term.
\begin{syntax}
  & L, M, N &::=& x \mid \funref{k} \mid \Tpack{L}{M} \mid \Tcallc{L}{M}
    \mid c \mid \cdots \mid \Mfold M \mid \Munfold M \\[1ex]
  & d &::=& \Tcodetm{e}{y}{N} \\
  & P &::=& (d_0, \ldots, d_{m-1};\; M)
\end{syntax}
\[
  \frac{k < m}{\Sigma; \Gamma \vdash \funref{k} : \Mcodety{\Delta_k}{A_k}{C_k}}
  \qquad
  \frac{\Sigma; e{:}\Mtup{\Delta_k}, y{:}A_k \vdash N_k : C_k \ (d_k = \Tcodetm{e}{y}{N_k},\ k < m)
        \quad \Sigma; \emptyset \vdash M : A}
       {\vdash (d_0, \ldots, d_{m-1};\; M) : A}
\]
A term over a smaller table is weakened to a bigger one (\texttt{wk}).

\subsection{\texorpdfstring{$\Ml{Mon}$}{Mon-mu}}

The monadic normal form of $\Ml{GlobalFun}$, with its own syntax
(\texttt{Language/Mon.agda}), so that monadic form needs no proof; its
meaning is that of its embedding $\Tembed{\cdot}$ into $\Ml{GlobalFun}$. An
atom is a variable, a constant or a reference to a definition; the
operands of every operation, $\Mfold$ and $\Munfold$ included, and the
scrutinee of $\key{case}$, are atoms, while the right-hand side and body
of $\key{let}$ and the branches of $\key{case}$ are expressions. Typing is
as in $\Ml{GlobalFun}$.
\begin{syntax}
  & a &::=& x \mid c \mid \funref{k} \\
  & M, N &::=& a \mid \Tpack{a_1}{a_2} \mid \Tcallc{a_1}{a_2}
    \mid \tuple{a_1,\ldots,a_n} \mid \prj{i} a \mid \inl a \mid \inr a
    \mid \prim{p}{a_1}{a_2} \\
  &&& \mid \Mfold a \mid \Munfold a
    \mid \Case{a}{x}{M}{y}{N} \mid \Let x = M \In N \\[1ex]
  & d &::=& \Tcodetm{e}{y}{N} \\
  & P &::=& (d_0, \ldots, d_{m-1};\; M)
\end{syntax}

\subsection{\texorpdfstring{$\Ml{Imp}$}{Imp-mu}}

Imp, typed (\texttt{Language/Imp.agda}). A body runs in a frame of slots,
typed by a context $\Delta = A_0, \ldots, A_{j-1}$, the types of the slots
assigned so far, so that slot $s < j$ holds a value of type $A_s$. A slot
is a typed index into $\Delta$, so a slot not yet assigned cannot be read.
Assignment and $\key{case}$ always write the next slot, $|\Delta|$, which
the rest of the body sees; the $\key{case}$ writes it at the type of the
branch's payload. The simple expressions $r$ are shared with $\Ml{CFG}$;
they and the atoms are those of $\Ml{Mon}$, over slots.
\begin{syntax}
  & a &::=& s \mid c \mid \funref{k} \\
  & r &::=& a \mid \Tpack{a_1}{a_2} \mid \Tcallc{a_1}{a_2}
    \mid \tuple{a_1,\ldots,a_n} \mid \prj{i} a \mid \inl a \mid \inr a
    \mid \prim{p}{a_1}{a_2} \mid \Mfold a \mid \Munfold a \\
  & e &::=& r \mid \assign{s}{e_1}{e_2} \mid \Casei{a}{s}{e_1}{e_2} \\[1ex]
  & P &::=& (e_0, \ldots, e_{m-1};\; e)
\end{syntax}
\[
  \frac{\Sigma; \Delta \vdash e_1 : A \quad \Sigma; \Delta, A \vdash e_2 : C}
       {\Sigma; \Delta \vdash \assign{|\Delta|}{e_1}{e_2} : C}
  \qquad
  \frac{\Sigma; \Delta \vdash a : A + B \quad
        \Sigma; \Delta, A \vdash e_1 : C \quad \Sigma; \Delta, B \vdash e_2 : C}
       {\Sigma; \Delta \vdash \Casei{a}{|\Delta|}{e_1}{e_2} : C}
\]
\[
  \Sigma; \Delta \vdash s : A_s
  \qquad
  \Sigma; \Mtup{\Delta_k}, A_k \vdash e_k : C_k
  \qquad
  \Sigma; \emptyset \vdash e : A
\]
A definition starts with its environment in slot $0$ and its argument in
slot $1$; $\key{main}$ with no slot.

\subsection{\texorpdfstring{$\Ml{CFG}$}{CFG-mu}}

CFG, typed (\texttt{Language/CFG.agda}). A block is a tail typed by the
type $R$ of its body's result, the frame $\Delta$ it starts in, and the
list $\vec{\Delta} = \Delta_0, \ldots, \Delta_{j-1}$ of the frames of the
blocks it may jump to, $\Sigma; R; \vec{\Delta}; \Delta \vdash t$. A label
$l$ is a typed index into $\vec{\Delta}$, so a jump to a missing block, or
to a block that reads a slot not assigned, cannot be written. Block $t_l$
is typed over the blocks older than it, so control flow is acyclic; the
start may jump to any block. An assignment, and the $\key{case}$ that
assigns the payload, write slot $s = |\Delta'|$ for a prefix $\Delta'$ of
$\Delta$: the slots from $s$ on are forgotten (they belong to a finished
scope) and slot $s$ is assigned anew. The two targets of a $\key{case}$
start in different frames: the payload's slot has the left type in one and
the right type in the other.
\begin{syntax}
  & t &::=& \Treturn{r} \mid \assign{s}{r}{t} \mid \Tgoto{l}
    \mid \Tcase{a}{s}{l_1}{l_2} \\[1ex]
  & b &::=& (t_0, \ldots, t_{j-1};\; t) \\
  & P &::=& (b_0, \ldots, b_{m-1};\; b)
\end{syntax}
\[
  \frac{\Sigma; \Delta \vdash r : R}{\Sigma; R; \vec{\Delta}; \Delta \vdash \Treturn{r}}
  \qquad
  \frac{\Delta' \sqsubseteq \Delta \quad \Sigma; \Delta \vdash r : A \quad
        \Sigma; R; \vec{\Delta}; \Delta', A \vdash t}
       {\Sigma; R; \vec{\Delta}; \Delta \vdash \assign{|\Delta'|}{r}{t}}
  \qquad
  \frac{\Delta_l = \Delta}{\Sigma; R; \vec{\Delta}; \Delta \vdash \Tgoto{l}}
\]
\[
  \frac{\Sigma; \Delta \vdash a : A + B \quad \Delta' \sqsubseteq \Delta \quad
        \Delta_{l_1} = \Delta', A \quad \Delta_{l_2} = \Delta', B}
       {\Sigma; R; \vec{\Delta}; \Delta \vdash
        \Tcase{a}{|\Delta'|}{l_1}{l_2}}
\]
\[
  \frac{\Sigma; R; \Delta_0, \ldots, \Delta_{l-1}; \Delta_l \vdash t_l \ (l < j) \quad
        \Sigma; R; \vec{\Delta}; \Delta \vdash t}
       {\Sigma; R; \Delta \vdash (t_0, \ldots, t_{j-1};\; t)}
\]
Here $\Delta' \sqsubseteq \Delta$ says that $\Delta'$ is a prefix of
$\Delta$. A definition's body $b_k$ starts in the frame
$\Mtup{\Delta_k}, A_k$ with result type $C_k$, and main's in the empty
frame. $\key{block}(\vec{t}, l)$ is block $l$ with the blocks older than it
(\texttt{block}, \texttt{Blk}).

\subsection{\texorpdfstring{$\Ml{X86}$}{X86-mu}}\label{sec:m-x86}

The target: a typed assembly language, after Morrisett, Walker, Crary and
Glew's TAL, intrinsically typed (\texttt{Language/X86.agda}). Its
instructions are those of x86 over the values of $\Ml{GlobalFun}$ and the
later languages, with an abstract heap, and one pseudo-instruction; each
is typed by the type of the machine state it starts in and the one it
leaves. A state type $\psi = \Tst{\Theta}{\Delta}{O}{F}$ gives
\begin{itemize}
\item the types $\Theta$ of the registers $\rax$, $\rdi$ and $\rsi$: each
  holds a value of a type $A$, the code ($\codeAB{A}{C}$) or the
  environment ($\tupA$) of the open closure, or is unset; $\emptyset$ has
  every register unset, and $\Theta[r{:}\tau]$ sets $r$;
\item the frame $\Delta$, the types of the slots assigned so far, as in
  $\Ml{CFG}$, in a frame of $n$ slots fixed by the body;
\item the open closure $O$: none ($\cdot$), or a slot $x$ of $\Delta$ of a
  function type $A \rightarrow C$ (below);
\item the flag $F$: $\Tset$ or $\Tunset$.
\end{itemize}
\begin{syntax}
  & \tau &::=& A \mid \codeAB{A}{C} \mid \tupA \\
  & \Theta &::=& \rax{:}\tau_1^?,\ \rdi{:}\tau_2^?,\ \rsi{:}\tau_3^?
    \qquad (\tau^? \text{ a type } \tau \text{ or unset}) \\
  & \psi &::=& \Tst{\Theta}{\Delta}{O}{F}
    \qquad O ::= \cdot \mid x \qquad F ::= \Tset \mid \Tunset \\[1ex]
  & r &::=& \rax \mid \rdi \mid \rsi \\
  & a &::=& s \mid c \mid \funref{k} \mid \%r \mid \Mfold a \mid \Munfold a
    \qquad d \;::=\; \%r \mid s \\
  & o &::=& \key{pair}(a_1, a_2) \mid \key{tuple}(a_1, \ldots, a_m) \mid \key{inl}(a) \mid \key{inr}(a) \\
  & i &::=& \key{movq}\; a, d \mid \key{subq}\; a, \%r \mid \key{cmpq}\; a_2, a_1
    \mid \key{sete}\; d \mid \key{alloc}\; o, d \mid \key{load}\; \pi_i\; a, d \\
  &&& \mid \key{load}\; \key{fst}\; x, \%r \mid \key{load}\; \key{snd}\; x, \%r
    \mid \key{callq}\; a \\
  & \beta &::=& i; \beta \mid \key{retq} \mid \key{jmp}\; l \mid \Mxcase{x}{s}{l_1}{l_2} \\
  & b &::=& (n;\; \beta_0, \ldots, \beta_{j-1};\; \beta) \\
  & P &::=& (b_0, \ldots, b_{m-1};\; b)
\end{syntax}
An operand is an atom of $\Ml{Mon}$ over slots (a slot, an immediate, or
the label of a definition), a register, or an operand of type
$\Munroll{X}{A}$ used at $\Mmu{X}{A}$ ($\Mfold a$) or back
($\Munfold a$): $\Mfold$ and $\Munfold$ are moves. A
destination is a register or a slot. The typing of the atoms is
$\Ml{Imp}$'s, in the frame $\Delta$ of the state; a destination
$n \vdash d : \psi \Rightarrow_\tau \psi'$ is written at type $\tau$ and
leaves the state $\psi'$; a slot write is that of $\Ml{CFG}$, in the frame
of $n$ slots, and closes the open closure:
\[
  \frac{\Sigma; \Delta \vdash a : A}{\Sigma; \Tst{\Theta}{\Delta}{O}{F} \vdash a : A}
  \qquad
  \frac{\Theta(r) = \tau}{\Sigma; \Tst{\Theta}{\Delta}{O}{F} \vdash \%r : \tau}
\]
\[
  \frac{\Sigma; \psi \vdash a : \Munroll{X}{A}}{\Sigma; \psi \vdash \Mfold a : \Mmu{X}{A}}
  \qquad
  \frac{\Sigma; \psi \vdash a : \Mmu{X}{A}}{\Sigma; \psi \vdash \Munfold a : \Munroll{X}{A}}
\]
\[
  n \vdash \%r : \Tst{\Theta}{\Delta}{O}{F} \Rightarrow_\tau \Tst{\Theta[r{:}\tau]}{\Delta}{O}{F}
  \qquad
  \frac{\Delta' \sqsubseteq \Delta \quad |\Delta'| < n}
       {n \vdash |\Delta'| : \Tst{\Theta}{\Delta}{O}{F} \Rightarrow_A \Tst{\Theta}{\Delta', A}{\cdot}{F}}
\]
An object $o$ of $\key{alloc}$ has type $A \rightarrow C$ for
$\key{pair}(a_1, a_2)$ with $a_1 : \Mcodety{\Delta}{A}{C}$ and
$a_2 : \Mtup{\Delta}$ (a closure), $\Mtup{(A_1, \ldots, A_m)}$ for a tuple of
operands of types $A_1, \ldots, A_m$, and $A + B$ for an injection. An
instruction is typed $\Sigma; n \vdash i : \psi \Rightarrow \psi'$, with
$\psi = \Tst{\Theta}{\Delta}{O}{F}$ and operands typed in $\psi$:
\[
\begin{array}{@{}l@{\;:\;}l@{\quad}l@{}}
  \key{movq}\; a, d & \psi \Rightarrow \psi' & a : \tau,\ d : \psi \Rightarrow_\tau \psi' \\
  \key{subq}\; a, \%r & \psi \Rightarrow \psi & a : \key{Int},\ \Theta(r) = \key{Int} \\
  \key{cmpq}\; a_2, a_1 & \psi \Rightarrow \Tst{\Theta}{\Delta}{O}{\Tset} & a_1, a_2 : \key{Int} \\
  \key{sete}\; d & \Tst{\Theta}{\Delta}{O}{\Tset} \Rightarrow \psi' & d : \Tst{\Theta}{\Delta}{O}{\Tset} \Rightarrow_{\Mbool} \psi' \\
  \key{alloc}\; o, d & \psi \Rightarrow \psi' & o : A,\ d : \psi \Rightarrow_A \psi' \\
  \key{load}\; \pi_i\; a, d & \psi \Rightarrow \psi' & a : \Mtup{(A_1, \ldots, A_m)},\ d : \psi \Rightarrow_{A_i} \psi' \\
  \key{load}\; \key{fst}\; x, \%r & \psi \Rightarrow \Tst{\Twipe{\Theta}[r{:}\codeAB{A}{C}]}{\Delta}{x}{F} & x : A \rightarrow C \in \Delta \\
  \key{load}\; \key{snd}\; x, \%r & \Tst{\Theta}{\Delta}{x}{F} \Rightarrow \Tst{\Theta[r{:}\tupA]}{\Delta}{x}{F} & \\
  \key{callq}\; a & \psi \Rightarrow \Treturned{\Delta, C} & a : \codeAB{A}{C}, \\
  \multicolumn{2}{l}{} & \Theta(\rdi) = \tupA,\ \Theta(\rsi) = A
\end{array}
\]
where $\Treturned{\Delta, C} = \Tst{\rax{:}C}{\Delta}{\cdot}{\Tunset}$
(the other registers unset), and $\Twipe{\Theta}$ unsets the registers of
type $\codeAB{A}{C}$ or $\tupA$. A sequence ends in a jump, and is typed by
the type $R$ of its body's result, the state it starts in and the frames
$\vec{\Delta}$ of the blocks it may jump to,
$\Sigma; n; R; \vec{\Delta} \vdash \beta : \psi$:
\[
  \frac{\Sigma; n \vdash i : \psi \Rightarrow \psi' \quad \Sigma; n; R; \vec{\Delta} \vdash \beta : \psi'}
       {\Sigma; n; R; \vec{\Delta} \vdash i; \beta : \psi}
  \qquad
  \frac{\Theta(\rax) = R}{\Sigma; n; R; \vec{\Delta} \vdash \key{retq} : \Tst{\Theta}{\Delta}{O}{F}}
\]
\[
  \frac{\Delta_l = \Delta}{\Sigma; n; R; \vec{\Delta} \vdash \key{jmp}\; l : \Tst{\Theta}{\Delta}{O}{F}}
\]
\[
  \frac{x : A + B \in \Delta \quad \Delta' \sqsubseteq \Delta \quad |\Delta'| < n \quad
        \Delta_{l_1} = \Delta', A \quad \Delta_{l_2} = \Delta', B}
       {\Sigma; n; R; \vec{\Delta} \vdash \Mxcase{x}{|\Delta'|}{l_1}{l_2} : \Tst{\Theta}{\Delta}{O}{F}}
\]
\[
  \frac{\Sigma; n; R; \Delta_0, \ldots, \Delta_{l-1} \vdash \beta_l : \Tentry{\Delta_l}\ (l < j) \quad
        \Sigma; n; R; \vec{\Delta} \vdash \beta : \psi}
       {\Sigma; R \vdash (n;\; \beta_0, \ldots, \beta_{j-1};\; \beta) : \psi}
\]
A block starts with its frame, no register and no flag,
$\Tentry{\Delta} = \Tst{\emptyset}{\Delta}{\cdot}{\Tunset}$; so, as in
$\Ml{CFG}$, the blocks are acyclic and a label is a typed index. A
definition $b_k$ of signature $(\Delta_k; A_k \to C_k)$ starts in
\[ \Tcalled{k} = \Tst{\rdi{:}\Mtup{\Delta_k},\ \rsi{:}A_k}{\emptyset}{\cdot}{\Tunset}\,, \]
with its environment in $\rdi$ and its argument in $\rsi$, and has result
type $C_k$; main starts in $\Tentry{\emptyset}$. So a jump to a missing
block, a read of an unset register, of a slot not assigned or of the flag
when it is not set, and a write of a slot beyond the frame cannot be
written.

\paragraph{The open closure.} The type $A \rightarrow C$ of a closure
hides the types $\Delta$ of its environment ($\exists \alpha.\,
\Mcodety{\alpha}{A}{C} \times \Mtup{\alpha}$), and a call needs the code
and the environment of the \emph{same} closure. $\key{load}\;\key{fst}\;
x, \%r$ opens the closure in slot $x$: $r$ gets the type $\codeAB{A}{C}$ of
its code, $\alpha$ standing for the types of its environment, and the
state records that $x$ is open. $\key{load}\; \key{snd}\; x, \%r$, only of
the open closure $x$, gives $r$ the type $\tupA$ of its environment.
$\key{callq}\; a$ needs $a : \codeAB{A}{C}$, $\rdi : \tupA$ and $\rsi : A$,
and leaves $\rax : C$. The witness is never named in a type: only the
argument and result types, which the type of the closure shows, are in
$\codeAB{A}{C}$. Opening a closure forgets the registers of type
$\codeAB{\cdot}{\cdot}$ or $\tupA$ of an earlier one ($\Twipe{\Theta}$),
and a call or the write of a slot closes the open closure, so $\alpha$
always means one closure's environment.

\paragraph{The case pseudo-instruction.} A $\key{case}$ on an injection
is $\Mxcase{x}{s}{l_1}{l_2}$, a jump that ends a sequence: it writes the
payload of the injection in slot $x$ to slot $s$ and runs block $l_1$ for
an $\key{inl}$ and $l_2$ for an $\key{inr}$. It stands for the three
instructions $\key{testinl}\; x;\ \key{load}\; \key{payload}\; x, s;\
\key{je}\; l_1, l_2$ of the dynamically typed pipeline's X86, but cannot
be split into them here: the payload's slot has the left type in block
$l_1$ and the right type in block $l_2$, and no state type between the
load and the jump can give it a type (TAL refines the type of the tested
value instead). So it is typed as one instruction, as $\key{chk}$ is in
$\Ml{X86}$'s dynamically typed counterpart (Part~\ref{part:dyn}). No
erasure of $\Ml{X86}$ to untyped X86 is defined in the Agda; $\key{case}$
would erase to the three instructions.

\paragraph{Nothing traps.} Every instruction, on operands of the types it
needs, has a result: there is no check, so $\Ml{X86}$ has no counterpart
of $\key{chk}$.

\paragraph{Where it doesn't match the Agda literally.}
\begin{itemize}
\item \textbf{Variables} are typed de Bruijn indices (\texttt{$\Gamma \ni A$},
  \texttt{Framework/Membership.agda}); contexts, frames, tables, tuples'
  lists of types and label lists are listed newest first in the Agda,
  here oldest first, so that slot $s$, definition $k$ and label $l$ are
  the \emph{levels} (\texttt{lvl}) of the indices. A component $\prj{i}$
  is \texttt{get i} for a typed index $i : \Delta \ni A$ into the tuple's
  types, here its position $1 \leq i \leq n$. The context of code is
  \texttt{A $\cdot$ tup $\Delta$ $\cdot$ []} (the argument newest), here
  $e{:}\Mtup{\Delta}, y{:}A$.
\item \textbf{Types}: \texttt{Ty n} with \texttt{base}, \texttt{\_$\Rightarrow$\_},
  \texttt{tup}, the sum (written $+$ here), \texttt{code},
  \texttt{var} and \texttt{$\mu$}; $\Munroll{X}{A}$ is \texttt{unroll A}, $\Tout{p}$ is
  \texttt{out-ty p}, and $\key{Bool}$ is the Agda's blackboard-bold B.
  $\Mfold$ and $\Munfold$ take no equation between types: every definition
  and proof splits terms at a generic type (the design note's ``green
  slime'' discussion, \S5).
\item \textbf{Heterogeneous vectors}: the components of tuples and the
  captured terms of closures are \texttt{All (Tm $\Gamma$) $\Delta$}
  (\texttt{Framework/AllUtil.agda}, with \texttt{lookup}, flags,
  \texttt{select}, \texttt{squeeze}).
\item \textbf{Slots} written in $\Ml{Imp}$ and $\Ml{CFG}$ are implicit in the
  Agda: \texttt{assign} and \texttt{case} of $\Ml{Imp}$ write slot $|\Delta|$;
  \texttt{seq} and \texttt{case} of $\Ml{CFG}$ carry a proof
  $p : \Delta' \sqsubseteq \Delta$ and write slot $|\Delta'|$. The CFG
  assignment $\assign{s}{r}{t}$ is the Agda's \texttt{seq}; a body is
  \texttt{Tail $\Sigma$ R Bs $\Delta$} in \texttt{Body}.
\item \textbf{$\Ml{X86}$}'s register types $\Theta$ are \texttt{Regs}, a
  record of \texttt{Maybe RTy} (\texttt{ty A}, \texttt{code$\alpha$ A B},
  \texttt{tup$\alpha$}); a state type is \texttt{st $\Theta$ $\Delta$ O F}
  (\texttt{St}), with \texttt{O : Maybe (Open $\Delta$)} and the flag a
  \texttt{Bool}; operands are \texttt{Arg} (\texttt{atom}, \texttt{reg},
  \texttt{\_$\uparrow$}, \texttt{\_$\downarrow$}), destinations \texttt{Dst}
  (\texttt{reg}, \texttt{slot p lt} with the irrelevant proof \texttt{lt} of
  $|\Delta'| < n$), objects \texttt{Obj}; $\key{load}\;\pi_i$ is \texttt{load},
  $\key{load}\;\key{fst}$ and $\key{load}\;\key{snd}$ are \texttt{loadfst}
  and \texttt{loadsnd r} (which takes its slot from the state);
  $\Treturned{\Delta, C}$, $\Tentry{\Delta}$ and $\Tcalled{k}$ are
  \texttt{returned}, \texttt{entry} and \texttt{called}; a sequence is
  \texttt{Seq $\Sigma$ n R Bs $\psi$}, a body \texttt{body n bs k}. The
  equations about registers that \texttt{subq}, \texttt{callq} and
  \texttt{retq} carry, and \texttt{reg r e} of an operand, are proofs, not
  operands.
\item \textbf{$\Ml{Mon}$} has its own syntax (\texttt{Atm}, \texttt{Exp});
  its meaning is the embedding $\Tembed{\cdot}$ (\texttt{$\ulcorner\_\urcorner$},
  \texttt{$\ulcorner\_\urcorner_{\mathsf{p}}$}) into $\Ml{GlobalFun}$.
\end{itemize}
%===============================================================================
\section{Compilation}\label{sec:m-passes}

The passes are in \texttt{Compiler/}. They are the passes of the
dynamically typed pipeline (Part~\ref{part:dyn}) on terms of richer
types, and preserve types by construction; $\Mfold$ and $\Munfold$ go
through every pass as one more operation with one operand. No pass looks
at $\mu$: recursive types show up only in the types of values and in the
two coercions. In each pass, the clauses not shown are homomorphic.

\subsection{Simplify (\texorpdfstring{$\Ml{ISWIM} \to \Ml{ISWIM}$}{ISWIM-mu -> ISWIM-mu})}\label{sec:m-simplify}

Simplify (\texttt{Compiler/Simplify.agda}), a simplified version of the
source optimizer cp0 of Chez Scheme, is the first pass. It is the
untyped simplify of Part~\ref{part:untyped}, intrinsically typed, with the
same design, rules and policy, and one more known-value rule, for
$\Munfold$. It carries what it knows of each variable, $E(x) = (T, K)$
(\texttt{Facts}, \texttt{Info}): the term $T$ that replaces $x$ (a variable
or a small value), and what $x$ is known to be, $K$ (\texttt{Known}): a
value, a function $\langle \lam{y} N, E_f \rangle$ (its body before
simplification, with what was known where it was defined), or $\bot$. On a
term it returns the same, $\pass{simp}_E(M) = (M', K)$; we write
$\pass{simp}_E(M)$ for $M'$ and $\mathcal{K}_E(M)$ for $K$. Values are
variables, $\lambda$s, constants, and tuples, injections and $\Mfold$s of
values. A value is \emph{small} when it is a variable, a constant,
$\emptytup$, or $\inl a$ or $\inr a$ of one of those; only small values
are duplicated. Values are known to be themselves, and
$\mathcal{K}_E(\lam{y} N) = \langle \lam{y} N, E \rangle$; under a binder
that stays, its variable is itself, known only as itself ($E'$ below). A
function body $N$ bound to an operand, $\pass{let}_E(x.N, M)$, is shared by
$\key{let}$, $\beta$ and inlining:
\begin{align*}
  \pass{let}_E(x.N, M) &= \pass{simp}_{\ext{E}{x}{(V, \mathcal{K}_E(M))}}(N)
    \quad \text{if } V \text{ is small} \\
  &= \Let x = V \In \pass{simp}_{\ext{E}{x}{(x, \mathcal{K}_E(M))}}(N)
    \quad \text{if } V \text{ is another value} \\
  &= \Let x = V \In \pass{simp}_{E'}(N)
    \quad \text{otherwise, where } V = \pass{simp}_E(M) \\[1ex]
  \pass{simp}_E(x) &= E(x) \\
  \pass{simp}_E(\Let x = M \In N) &= \pass{let}_E(x.N, M) \\
  \pass{simp}_E((\lam{x} N)\;M) &= \pass{let}_E(x.N, M) \\
  \pass{simp}_E(L\;M) &= \pass{let}_{E_f}(x.N, M)
    \quad \text{if } \mathcal{K}_E(L) = \langle \lam{x} N, E_f \rangle \text{ (and small)} \\
  \pass{simp}_E(\prim{p}{L}{M}) &= \res{\delta(p, n, m)}
    \quad \text{if } \mathcal{K}_E(L) = n,\; \mathcal{K}_E(M) = m \\
  \pass{simp}_E(\Case{L}{x}{M}{y}{N}) &= \pass{simp}_{\ext{E}{x}{(K, K)}}(M)
    \quad \text{if } \mathcal{K}_E(L) = \inl K,\; K \text{ small} \\
  \pass{simp}_E(\prj{i} L) &= K_i
    \quad \text{if } \mathcal{K}_E(L) = \tuple{K_1, \ldots, K_n},\; K_i \text{ small} \\
  \mathcal{K}_E(\prj{i} L) &= K_i
    \quad \text{if } \mathcal{K}_E(L) = \tuple{K_1, \ldots, K_n} \\
  \pass{simp}_E(\Munfold L) &= K
    \quad \text{if } \mathcal{K}_E(L) = \Mfold K,\; K \text{ small} \\
  \mathcal{K}_E(\Munfold L) &= K
    \quad \text{if } \mathcal{K}_E(L) = \Mfold K \\[1ex]
  \pass{simplify}(M) &= \pass{simp}_{x \mapsto (x, x)}(M)
\end{align*}
and likewise for $\inr$. Here $\res{o}$ is the term of a primitive's result
(a constant, or $\inl \emptytup$ or $\inr \emptytup$ for $\key{eq}$). The
last two rules are the one that $\mu$ adds (\texttt{simp-unfold}, with
\texttt{unfolded}): an $\Munfold$ of a known $\Mfold$ is the value under
it, as a projection from a known tuple is its component; a $\Mfold$ is not
small, so a $\Mfold$ of a value stays bound by a $\key{let}$ and known. The
inlining clause is as in Part~\ref{part:untyped}: the known function's body
is simplified again, at the call, from the source, under $E_f$ and what is
known of the operand, and replaces the call if it has at most $24$ nodes
(\texttt{inline-size}); $\pass{simp}$ carries fuel ($4$,
\texttt{inline-fuel}), which inlining decreases, so that it is structurally
recursive. It drops only evaluations whose value is known, which have no
effect, and otherwise keeps the order of evaluation: an operand that is
not a value is never substituted, but stays bound by a $\key{let}$.

What typing changes: $E$ maps the variables of the source's context to
facts in the target's, so the pass preserves types by construction; a
known component, branch or unfolding is found by a test that takes the
operand's type generically, with an equation to the expected one
(\texttt{component}, \texttt{branch}, \texttt{unfolded}), and $\lambda$s
are recognized at any type (\texttt{lam?}). Typed terms are closed by
construction, so the untyped pass's lemma that it keeps programs closed is
not needed.

\subsection{CC: Enclose (\texorpdfstring{$\Ml{ISWIM} \to \Ml{Clos}$}{ISWIM-mu -> Clos-mu})}

Enclose (\texttt{Compiler/Enclose.agda}) makes each function a closure
that captures every variable in scope, in order, so its body is unchanged
(pruning is optimize's job):
\begin{align*}
  \pass{enclose}(\lam{y} N) &= \clos{x_1\ldots x_k}{x_1\ldots x_k\, y} \pass{enclose}(N)
    \qquad\text{where } x_1, \ldots, x_k \text{ are the variables in scope}
\end{align*}

\subsection{CC: Optimize (\texorpdfstring{$\Ml{Clos} \to \Ml{Clos}$}{Clos-mu -> Clos-mu})}

Optimize (\texttt{Compiler/Optimize.agda}) removes the captured values that
the code does not use. Closures are strict, so a captured term is dropped
only when the code does not use its name and it is a variable, which is
all enclose captures:
\begin{align*}
  \pass{optimize}(\clos{M_1\ldots M_n}{x_1\ldots x_n\, y} N) &=
    \clos{\pass{optimize}(M_{i_1})\ldots \pass{optimize}(M_{i_k})}{x_{i_1}\ldots x_{i_k}\, y}
    \pass{optimize}(N) \\
  \text{where } \{ i_1 < \cdots < i_k \} &=
    \{ i \mid x_i \in \FV(N) \text{ or } M_i \text{ is not a variable} \}
\end{align*}
Since the captured variables have different types, an unused one cannot
be sent to the argument, as the dynamically typed pass does: the pass's
renaming is defined only on the variables a term uses (\texttt{Ren (uses
M) $\Gamma'$}, given a proof that the variable is used), and a closure's
code is renamed by \texttt{squeeze}, defined on the kept variables, with
the proof (\texttt{covers}) that the code uses only kept ones. The kept
positions are flags (\texttt{keeps}, \texttt{keep?}), and the kept terms
\texttt{select}; a program is \texttt{optimize-program}.

\subsection{CC: Concretize (\texorpdfstring{$\Ml{Clos} \to \Ml{EnvClos}$}{Clos-mu -> EnvClos-mu})}

Each closure stores its captured values in one tuple $e$, of type
$\Mtup{\Delta}$, and inside its code the name $x_i$ becomes $\prj{i} e$
(\texttt{Compiler/Concretize.agda}). The map $\varphi$ sends each
$\Ml{Clos}$ variable to a reference $r ::= x \mid \prj{i} x$ in
$\Ml{EnvClos}$ (\texttt{Ref}: \texttt{var}, \texttt{prj}).
\begin{align*}
  \pass{concretize}_\varphi(x) &= \varphi(x) \\
  \pass{concretize}_\varphi(\clos{M_1\ldots M_n}{x_1\ldots x_n\, y} N) &=
    \clos{\pass{concretize}_\varphi(M_1)\ldots\pass{concretize}_\varphi(M_n)}{e\,y}
       \pass{concretize}_{\varphi_n}(N) \\
  \text{where } \varphi_n &= \{ y \mapsto y,\; x_1 \mapsto \prj{1} e,\; \ldots,\; x_n \mapsto \prj{n} e \}
\end{align*}
($\varphi_n$ is \texttt{fv-ref}); under a $\key{let}$ or $\key{case}$ the
bound variable maps to itself (\texttt{ext-ref}); a program is
concretized with $\varphi = \emptyset$ (\texttt{concretize-program}).

\subsection{CC: Delay (\texorpdfstring{$\Ml{EnvClos} \to \Ml{FlatClos}$}{EnvClos-mu -> FlatClos-mu})}

A closure becomes a packed pair of its code and the tuple of its captured
values, and an application opens the closure and calls its code
(\texttt{Compiler/Delay.agda}); $\key{callc}$ does it all, with no
$\key{let}$, $\key{fst}$ or $\key{snd}$:
\begin{align*}
  \pass{delay}(\clos{M_1\ldots M_n}{e\,y} N) &=
    \Tpack{\Tcodetm{e}{y}{\pass{delay}(N)}}{\tuple{\pass{delay}(M_1),\ldots,\pass{delay}(M_n)}} \\
  \pass{delay}(L\;M) &= \Tcallc{\pass{delay}(L)}{\pass{delay}(M)}
\end{align*}
(It carries a renaming of the variables, the identity for a program,
\texttt{delay-program}.)

\subsection{CC: Globalize (\texorpdfstring{$\Ml{FlatClos} \to \Ml{GlobalFun}$}{FlatClos-mu -> GlobalFun-mu})}

The code of every function becomes a global definition, and the code
itself a reference to it (\texttt{Compiler/Globalize.agda}). The pass
threads the definitions made so far, from left to right; for code, it
first globalizes the body, so the functions nested in it become
definitions first. Each definition is typed over the table of the
definitions older than it (\texttt{Made}), and the term made for an
operand over the table made so far; both are weakened to the bigger table
at the end (\texttt{wk}, \texttt{close}), which has an entry
$(\Delta; A \to C)$ for code of type $\Mcodety{\Delta}{A}{C}$.
\begin{align*}
  \pass{glob}(x, \vec{d}) &= (x, \vec{d}) \\
  \pass{glob}(\Tcodetm{e}{y}{N}, \vec{d}) &= (\funref{|\vec{d}_1|},\; \vec{d}_1, \Tcodetm{e}{y}{N'})
     \quad\text{where } (N', \vec{d}_1) = \pass{glob}(N, \vec{d}) \\
  \pass{glob}(\mathit{op}(\bar{x}_1.A_1, \ldots, \bar{x}_n.A_n), \vec{d}_0)
     &= (\mathit{op}(\bar{x}_1.A'_1, \ldots, \bar{x}_n.A'_n), \vec{d}_n)
     \quad\text{where } (A'_i, \vec{d}_i) = \pass{glob}(A_i, \vec{d}_{i-1}) \\[1ex]
  \pass{globalize}(M) &= (\vec{d};\; M') \quad\text{where } (M', \vec{d}) = \pass{glob}(M, \epsilon)
\end{align*}

\subsection{Remove Complex Operands (\texorpdfstring{$\Ml{GlobalFun} \to \Ml{Mon}$}{GlobalFun-mu -> Mon-mu})}

Every complex operand is bound to a fresh variable by a $\key{let}$, from
left to right, so the operands of every operation become atoms
(\texttt{Compiler/RCO.agda}). The fresh variable has the type of the
operand, which may be any type. A $\key{case}$ or $\key{let}$ is complex:
its branches and body are translated in place, so nothing is duplicated.
An operand that is an atom ($x$, $c$ or $\funref{k}$) stays in place; any
other operand $M$ is translated and bound to a fresh variable:
\begin{align*}
  \pass{atomize}(M) &= (\epsilon;\; M) \qquad M \text{ an atom} \\
  \pass{atomize}(M) &= (x{=}\pass{rco}(M);\; x) \qquad \text{otherwise, } x \text{ fresh} \\
  \pass{atomize}(M_1, \ldots, M_n) &= (\vec{b}_1 \cdots \vec{b}_n;\; a_1, \ldots, a_n)
    \quad\text{where } (\vec{b}_i; a_i) = \pass{atomize}(M_i) \\
  \pass{lets}(x_1{=}N_1 \cdots x_k{=}N_k,\; N) &= \Let x_1 = N_1 \In \cdots \Let x_k = N_k \In N
\end{align*}
The operations with operands are
\begin{multline*}
  \mathit{op}(M_1, \ldots, M_n) \in \{\, \Tpack{M}{N},\ \Tcallc{M}{N},\
     \tuple{M_1, \ldots, M_n},\ \prj{i} M,\ \inl M,\ \inr M, \\
     \prim{p}{M}{N},\ \Mfold M,\ \Munfold M \,\}
\end{multline*}
\begin{align*}
  \pass{rco}(a) &= a \qquad a \text{ an atom} \\
  \pass{rco}(\mathit{op}(M_1, \ldots, M_n)) &= \pass{lets}(\vec{b}, \mathit{op}(a_1, \ldots, a_n)) \\
  &\qquad \text{where } (\vec{b}; a_1, \ldots, a_n) = \pass{atomize}(M_1, \ldots, M_n) \\
  \pass{rco}(\Case{L}{x}{M}{y}{N}) &= \pass{lets}(\vec{b}, \Case{a}{x}{\pass{rco}(M)}{y}{\pass{rco}(N)}) \\
  &\qquad \text{where } (\vec{b}; a) = \pass{atomize}(L) \\
  \pass{rco}(\Let x = M \In N) &= \Let x = \pass{rco}(M) \In \pass{rco}(N) \\
  \pass{rco}((d_0, \ldots, d_{m-1};\; M)) &= (\pass{rco}(d_0), \ldots, \pass{rco}(d_{m-1});\; \pass{rco}(M)) \\
  \pass{rco}(\Tcodetm{e}{y}{N}) &= \Tcodetm{e}{y}{\pass{rco}(N)}
\end{align*}

\subsection{Lift Locals (\texorpdfstring{$\Ml{Mon} \to \Ml{Imp}$}{Mon-mu -> Imp-mu})}

A $\key{let}$ becomes an assignment to the next slot of the frame, and the
variable a $\key{case}$ binds is the slot its payload is assigned to
(\texttt{Compiler/LiftLocals.agda}). The frame of a body is typed by its
context, and a variable is the slot at its level, $\ell(x)$, the number of
variables bound before it (in a definition $\Tcodetm{e}{y}{N}$,
$\ell(e) = 0$ and $\ell(y) = 1$). So the variables and atoms stay as they
are. With $\Gamma$ the variables in scope,
\begin{align*}
  \pass{lift}(x) &= \ell(x) \\
  \pass{lift}(\Let x = M \In N) &= \assign{|\Gamma|}{\pass{lift}(M)}{\pass{lift}(N)} \\
  \pass{lift}(\Case{a}{x}{M}{y}{N}) &=
    \Casei{\pass{lift}(a)}{|\Gamma|}{\pass{lift}(M)}{\pass{lift}(N)} \\
  \pass{lift}(\Tcodetm{e}{y}{N}) &= \pass{lift}(N) \\
  \pass{lift}((d_0, \ldots, d_{m-1};\; M)) &=
    (\pass{lift}(d_0), \ldots, \pass{lift}(d_{m-1});\; \pass{lift}(M))
\end{align*}
and every other construct is a simple expression, unchanged.

\subsection{Explicate Control (\texorpdfstring{$\Ml{Imp} \to \Ml{CFG}$}{Imp-mu -> CFG-mu})}

$\pass{tail}(e)$ translates $e$ in tail position, $\pass{assign}(e, s, k)$
translates ``slot $s := e$, then the tail $k$'' in continuation-passing
style, and a $\key{case}$ on the right of an assignment adds $k$ as a
block, the join point, which both branches jump to
(\texttt{Compiler/Explicate.agda}: \texttt{explicate-tail},
\texttt{explicate-assign}). $\pass{add}(t)$ adds the block $t$ and gives
its label. Both thread the blocks made so far, each typed over the older
ones, so a new block's label is the newest, and a label made earlier stays
valid as more blocks are added. In $\pass{assign}(e, s, k)$, $k$ starts in
the frame $\Delta', A$ with $s = |\Delta'|$, $\Delta'$ a prefix of the
frame $\Delta$ of $e$.
\begin{align*}
  \pass{tail}(r) &= \Treturn{r} \\
  \pass{tail}(\assign{s}{e_1}{e_2}) &= \pass{assign}(e_1, s, \pass{tail}(e_2)) \\
  \pass{tail}(\Casei{a}{s}{e_1}{e_2}) &= \Tcase{a}{s}{l_1}{l_2} \\
  &\qquad \text{where } l_1 = \pass{add}(\pass{tail}(e_1)),\ l_2 = \pass{add}(\pass{tail}(e_2)) \\[1ex]
  \pass{assign}(r, s, k) &= \assign{s}{r}{k} \\
  \pass{assign}(\assign{s'}{e_1}{e_2}, s, k) &= \pass{assign}(e_1, s', \pass{assign}(e_2, s, k)) \\
  \pass{assign}(\Casei{a}{s'}{e_1}{e_2}, s, k) &= \Tcase{a}{s'}{l_1}{l_2} \\
  &\qquad \text{where } l = \pass{add}(k),\
    l_i = \pass{add}(\pass{assign}(e_i, s, \Tgoto{l})) \\[1ex]
  \pass{explicate}(e) &= (\vec{t};\; t)
    \quad\text{where } t = \pass{tail}(e) \text{ and } \vec{t} \text{ are the blocks it adds}
\end{align*}
and $\pass{explicate}$ maps over the definitions and main of a program
(\texttt{explicate-body}). A $\key{case}$ in $\Ml{Imp}$ writes the next slot
of its frame, so the $\key{case}$ tail it becomes writes slot $|\Delta|$.

\subsection{Select Instructions (\texorpdfstring{$\Ml{CFG} \to \Ml{X86}$}{CFG-mu -> X86-mu})}\label{sec:m-select}

A slot is the typed index of $\Ml{CFG}$, which is an $\Ml{X86}$ slot, and an
atom is an $\Ml{X86}$ operand (\texttt{Compiler/Select.agda}); the frame of
a body has as many slots as the deepest level its tails write
($\mathit{depth}$), and each slot write carries the proof that it is in the
frame. The code of a simple expression $r$ is selected in
continuation-passing style, $\pass{sel}(r, o, k)$: its value goes to $o$,
$\rax$ ($\key{ret}$) or the next slot $s$ ($\key{to}\;s$, with
$s = |\Delta'|$ for a prefix $\Delta'$ of the frame and $s < n$), and then
$k$ runs, whatever the registers, the open closure and the flag are. So
$k$ is typed in every state $\key{after}_o(\Theta, O, F)$ that $o$ leaves:
\[
  \key{after}_{\key{ret}}(\Theta, O, F) = \Tst{\Theta[\rax{:}A]}{\Delta}{O}{F}
  \qquad
  \key{after}_{\key{to}\;s}(\Theta, O, F) = \Tst{\Theta}{\Delta', A}{\cdot}{F}
\]
for $r$ of type $A$ in the frame $\Delta$; $\bar{o}$ is the destination,
$\rax$ or $s$. An operand of a function type or of a sum type is a slot
$x$ (no constant has those types), so $\key{callc}$ and $\key{case}$ take a
slot.
\begin{align*}
  \pass{sel}(a, o, k) &= \key{movq}\; a, \bar{o};\ k \\
  \pass{sel}(\Tpack{a_1}{a_2}, o, k) &= \key{alloc}\; \key{pair}(a_1, a_2), \bar{o};\ k \\
  \pass{sel}(\Tcallc{x}{a}, o, k) &= \key{load}\; \key{fst}\; x, \rax;\
    \key{load}\; \key{snd}\; x, \rdi;\ \key{movq}\; a, \rsi; \\
  &\qquad \key{callq}\; \rax;\ \key{movq}\; \rax, \bar{o};\ k \\
  \pass{sel}(\tuple{a_1, \ldots, a_n}, o, k) &= \key{alloc}\; \key{tuple}(a_1, \ldots, a_n), \bar{o};\ k \\
  \pass{sel}(\prj{i} a, o, k) &= \key{load}\; \pi_i\; a, \bar{o};\ k \\
  \pass{sel}(\inl a, o, k) &= \key{alloc}\; \key{inl}(a), \bar{o};\ k
    \qquad\text{and likewise for } \key{inr} \\
  \pass{sel}(\key{sub}(a_1, a_2), o, k) &= \key{movq}\; a_1, \rax;\
    \key{subq}\; a_2, \rax;\ \key{movq}\; \rax, \bar{o};\ k \\
  \pass{sel}(\key{eq}(a_1, a_2), o, k) &= \key{cmpq}\; a_2, a_1;\ \key{sete}\; \bar{o};\ k \\
  \pass{sel}(\Mfold a, o, k) &= \key{movq}\; \Mfold a, \bar{o};\ k \\
  \pass{sel}(\Munfold a, o, k) &= \key{movq}\; \Munfold a, \bar{o};\ k \\[1ex]
  \pass{tail}(\Treturn{r}) &= \pass{sel}(r, \key{ret}, \key{retq}) \\
  \pass{tail}(\assign{s}{r}{t}) &= \pass{sel}(r, \key{to}\;s, \pass{tail}(t)) \\
  \pass{tail}(\Tgoto{l}) &= \key{jmp}\; l \\
  \pass{tail}(\Tcase{x}{s}{l_1}{l_2}) &= \Mxcase{x}{s}{l_1}{l_2}
\end{align*}
In $\key{callc}$, $\key{load}\;\key{fst}$ opens the closure $x$, and
$\key{load}\;\key{snd}$ and the move of the argument give $\rdi$ and $\rsi$
the types $\key{callq}$ needs; the call closes the closure, and the result
is moved from $\rax$. A $\key{case}$ tail is $\Ml{X86}$'s $\key{case}$
pseudo-instruction, one for one. With $\Delta$ the frame a tail starts in,
\begin{align*}
  \mathit{depth}(\Treturn{r}) = \mathit{depth}(\Tgoto{l}) &= |\Delta| \\
  \mathit{depth}(\assign{s}{r}{t}) &= \max(|\Delta|, \mathit{depth}(t)) \\
  \mathit{depth}(\Tcase{a}{s}{l_1}{l_2}) &= \max(|\Delta|, s + 1) \\[1ex]
  \pass{select}((t_0, \ldots, t_{j-1};\; t)) &=
    (n;\ \pass{tail}(t_0), \ldots, \pass{tail}(t_{j-1});\ \pass{tail}(t)) \\
  \pass{select}^{+}((t_0, \ldots, t_{j-1};\; t)) &=
    (n;\ \pass{tail}(t_0), \ldots, \pass{tail}(t_{j-1});\
     \key{movq}\; \rdi, 0;\ \key{movq}\; \rsi, 1;\ \pass{tail}(t))
\end{align*}
for main (\texttt{main-body}) and for a definition (\texttt{def-body}),
with $n = \max(\mathit{depth}(t_0), \ldots, \mathit{depth}(t_{j-1}),
\mathit{depth}(t))$, and $\pass{select}((b_0, \ldots, b_{m-1};\; b)) =
(\pass{select}^{+}(b_0), \ldots, \pass{select}^{+}(b_{m-1});\;
\pass{select}(b))$. The prologue of a definition moves its environment and
argument from $\rdi$ and $\rsi$ to slots $0$ and $1$, the frame its body
starts in. The pass is total on typed programs.

\subsection{The Compiler}

The compiler goes from $\Ml{ISWIM}$ to $\Ml{X86}$
(\texttt{Compiler/Pipeline.agda}), for a closed program of any type $A$,
each compiled program having type $A$:
\[
  \pass{compile}_4 = \pass{delay} \circ \pass{concretize}_{\emptyset}
     \circ \pass{optimize} \circ \pass{enclose} \circ \pass{simplify}
  \qquad
  \pass{compile} = \pass{globalize} \circ \pass{compile}_4
\]
\begin{align*}
  \pass{compile\text{-}mon} &= \pass{rco} \circ \pass{compile}
  &
  \pass{compile\text{-}imp} &= \pass{lift} \circ \pass{compile\text{-}mon} \\
  \pass{compile\text{-}cfg} &= \pass{explicate} \circ \pass{compile\text{-}imp}
  &
  \pass{compile\text{-}x86} &= \pass{select} \circ \pass{compile\text{-}cfg}
\end{align*}
There is no cast insertion and no test that the program is closed: a
closed program is a term of the empty context.

\paragraph{Where it doesn't match the Agda literally.}
\begin{itemize}
\item \textbf{Simplify}'s facts are \texttt{Facts $\Delta$ $\Gamma$} (for each
  variable of the source's context, an \texttt{Info} in the target's:
  \texttt{term} and \texttt{known}), $\ext{E}{x}{I}$ is \texttt{I $\bullet$ E},
  $E'$ is \texttt{under E}, a value kept bound and known is \texttt{known-as};
  $\pass{let}$, the rules and the pass are \texttt{simp-let},
  \texttt{simp-app}, \texttt{simp-call}, \texttt{simp-prim},
  \texttt{simp-case}, \texttt{simp-get}, \texttt{simp-unfold} and
  \texttt{simp}; values and small values are \texttt{value} and
  \texttt{small} (with \texttt{atom}). Renaming is the substitution kit's
  (\texttt{Language/ISWIMSubst.agda}), which the reduction semantics
  shares.
\item \textbf{Optimize} carries a renaming defined on the variables that
  the term uses, \texttt{Ren (uses M) $\Gamma'$}; $x_i \in \FV(N)$ is
  \texttt{uses N (S i)}.
\item \textbf{Concretize}'s references are \texttt{Ref} (\texttt{var},
  \texttt{prj}); \textbf{delay} and \textbf{remove complex operands} carry
  renamings, the identity for a program; rco is in continuation-passing
  style: \texttt{atomize M $\rho$ k} gives the atom for $M$ to the
  continuation $k$, which takes a renaming to the context where the atom is
  used (Kripke style), so there are no lists $\vec{b}$.
\item \textbf{Globalize} threads the definitions made, each over the table
  of the older ones (\texttt{Made}), and returns with each term a proof that
  the table grew (\texttt{Out}); terms are weakened (\texttt{wk}) to the
  bigger table, and at the end every definition is weakened to the whole
  table (\texttt{close}).
\item \textbf{Explicate control}'s results are \texttt{Out} records: a tail
  with the blocks made and a proof that the list of blocks grew; $\pass{add}$
  is $\triangleright$, and a body is \texttt{body bs t}.
\item \textbf{Select instructions}' destinations are \texttt{Out} (\texttt{ret},
  \texttt{to p lt}), $\bar{o}$ is \texttt{dst o}, and the continuation $k$ of
  \texttt{sel s o k} is polymorphic in the registers, open closure and flag;
  $\pass{tail}$ and \texttt{blocks} take the irrelevant proof that the depth
  is at most $n$. The closure of $\key{callc}$ and the scrutinee of
  $\key{case}$ are matched as \texttt{var x}: an atom of a function or sum
  type is a variable.
\end{itemize}
%===============================================================================
\section{Denotational: Graph Model}\label{sec:m-den}

The languages are given meaning in the graph model of
Part~\ref{part:untyped}: the same domain $\V$, the same operators
(Section~\ref{sec:den-ops}) and consistency, and the same clauses
(\texttt{Semantics/Graph/}). $\Mfold$ and $\Munfold$ are the identity: a
value of a recursive type is a value of its unfolding, and the domain
carries no types. Recursion needs no least fixed point: self-application,
through a recursive type, gives it, as in untyped ISWIM. Since the syntax
is intrinsically typed, each semantics is defined by plain structural
recursion on typed terms, with no ABT. The facts from which a language
proves its denotations monotone, consistent and continuous
($\mathit{Ok}$) by induction on the term, one line per construct, are in
\texttt{proof/Framework/GraphModel/TypedOkBase.agda} (shared with the
dynamically typed pipeline), and the lemmas of each language
($\mathit{Ok}$ of $\Ml{ISWIM}$, $\Ml{EnvClos}$, $\Ml{FlatClos}$ and
$\Ml{GlobalFun}$, consistency of a table, monotonicity of a $\Ml{CFG}$ tail in
its frame) are in \texttt{proof/Semantics/Graph/}.

We write variables by name, as in Part~\ref{part:untyped}; in the Agda an
environment $\rho : \mathbb{N} \to \SET{\V}$ gives position $i$ the
meaning of the variable with index $i$, so an environment, like the
domain, carries no types. A variable not in scope means $\{\omega\}$, as
in $\rho_0$; in a typed term there is none, but the body of a closed code
is given the environment $\rho_0$ extended with its parameters.

\subsection{The languages}

$\SEM{-}^{\Ml{ISWIM}}$ (\texttt{Semantics/Graph/ISWIM.agda}) has the clauses
that all the languages of Part~\ref{part:untyped} share (variables,
constants, tuples by $\mathcal{T}$, $\prj{i}$ by $\key{proj}_i$,
$\key{inl}$, $\key{inr}$, $\key{case}$ by $\mathcal{C}$, primitives by
$\key{prim}_p$, and
$\SEM{\Let x = M \In N}\rho = \Lambda(\lambda X.\, \SEM{N}\ext{\rho}{x}{X}) \app \SEM{M}\rho$),
those for $\lambda$ and application, $\Lambda$ and $\app$, and
\begin{align*}
  \SEM{\Mfold M}^{\ell}\rho &= \SEM{M}^{\ell}\rho \qquad
  \SEM{\Munfold M}^{\ell}\rho = \SEM{M}^{\ell}\rho
\end{align*}
in every language $\ell$ (and, in $\Ml{Imp}$ and $\Ml{CFG}$, for the simple
expressions on atoms). $\Ml{Clos}$ and $\Ml{EnvClos}$ are as Clos and EnvClos,
except that a closure captures the values of terms $M_1, \ldots, M_n$, not
only variables:
\begin{align*}
  \SEM{\clos{M_1\ldots M_n}{x_1\ldots x_n\, y} N}^{\Ml{2}}\rho &=
    \key{guard}(D_1, \ldots, D_n;\;
       \Lambda(\lambda Y.\, \SEM{N}^{\Ml{2}}\rho_0(x_1{:=}D_1,\ldots,x_n{:=}D_n, y{:=}Y))) \\
    &\qquad\text{where } D_i = \SEM{M_i}^{\Ml{2}}\rho \\
  \SEM{\clos{M_1\ldots M_n}{e\, y} N}^{\Ml{3}}\rho &=
    \Lambda(\lambda X.\, \Lambda(\lambda Y.\, \SEM{N}^{\Ml{3}}\rho_0(e{:=}X, y{:=}Y)))
    \app \mathcal{T}(\SEM{M_1}^{\Ml{3}}\rho, \ldots, \SEM{M_n}^{\Ml{3}}\rho)
\end{align*}
where $\Ml{2}, \ldots, \Ml{5}$ abbreviate $\Ml{Clos}, \ldots, \Ml{GlobalFun}$.
In $\Ml{FlatClos}$ code is a closed function of two arguments, the
environment and the argument; a closure is the pair of code and
environment, built by $\key{pack}$; and $\key{callc}$ applies the
closure's code ($\key{car}$) to its environment ($\key{cdr}$) and then to
the argument:
\begin{align*}
  \SEM{\Tcodetm{e}{y}{N}}^{\Ml{4}}\rho &= \Lambda(\lambda X.\, \Lambda(\lambda Y.\, \SEM{N}^{\Ml{4}}\rho_0(e{:=}X, y{:=}Y))) \\
  \SEM{\Tpack{L}{M}}^{\Ml{4}}\rho &= \key{pair}(\SEM{L}^{\Ml{4}}\rho, \SEM{M}^{\Ml{4}}\rho) \\
  \SEM{\Tcallc{L}{M}}^{\Ml{4}}\rho &= (\key{car}(D) \app \key{cdr}(D)) \app \SEM{M}^{\Ml{4}}\rho
    = D \odot \SEM{M}^{\Ml{4}}\rho
    \quad\text{where } D = \SEM{L}^{\Ml{4}}\rho
\end{align*}
$\Ml{GlobalFun}$ is $\Ml{FlatClos}$ with code replaced by a reference $\funref{k}$
to the definition $d_k$, given meaning by a table $T$ of the meanings of
the definitions (\texttt{Semantics/Graph/Ops.agda}, \texttt{Table},
\texttt{tables}). Each definition $d_k$ is given meaning with the table of
the older ones, in which a reference to $d_k$ or a newer definition means
$\emptyset$, and main with the table of all of them:
\begin{align*}
  \SEM{\funref{k}}^{\Ml{5}}_T\rho &= T(k) \\
  \key{code}_T(\Tcodetm{e}{y}{N}) &= \Lambda(\lambda X.\, \Lambda(\lambda Y.\, \SEM{N}^{\Ml{5}}_T\rho_0(e{:=}X, y{:=}Y))) \\
  \key{table}(d_0,\ldots,d_{m-1})(k) &=
    \begin{cases}
      \key{code}_{\key{table}(d_0,\ldots,d_{k-1})}(d_k) & k < m \\
      \emptyset & \text{otherwise}
    \end{cases} \\
  \SEM{(d_0, \ldots, d_{m-1};\; M)}^{\Ml{5}} &= \SEM{M}^{\Ml{5}}_{\key{table}(d_0,\ldots,d_{m-1})}\,\rho_0
\end{align*}
with the other clauses those of $\Ml{FlatClos}$ under $T$. The table is
generic in what a definition is, and is also that of $\Ml{Imp}$, $\Ml{CFG}$
and $\Ml{X86}$, with their own $\key{code}$. A $\Ml{Mon}$ program $P$ means
its embedding: $\SEM{P}^{\Ml{Mon}} = \SEM{\Tembed{P}}^{\Ml{5}}$.

\paragraph{$\Ml{Imp}$, continuation-passing} (\texttt{Semantics/Graph/Imp.agda}).
A frame $\sigma$ is an environment of slots, $\sigma(s)$ is the set in slot
$s$, and $(\sigma, X)$ is $\sigma$ extended with its next slot holding
$X$. An expression never changes the slots it can see, so a continuation
$\kappa : \SET{\V} \to \SET{\V}$ takes only the value of the expression.
The atoms and simple expressions, shared with $\Ml{CFG}$, are
$\Ml{GlobalFun}$'s operators on atoms:
\begin{align*}
  \mathcal{A}_\sigma(s) &= \sigma(s) \qquad
  \mathcal{A}_\sigma(c) = \mathcal{B}(c) \qquad
  \mathcal{A}_\sigma(\funref{k}) = T(k) \\[1ex]
  \mathcal{S}_\sigma(a) &= \mathcal{A}_\sigma(a) \qquad
  \mathcal{S}_\sigma(\Tpack{a_1}{a_2}) = \key{pair}(\mathcal{A}_\sigma(a_1), \mathcal{A}_\sigma(a_2)) \\
  \mathcal{S}_\sigma(\Tcallc{a_1}{a_2}) &= \mathcal{A}_\sigma(a_1) \odot \mathcal{A}_\sigma(a_2) \\
  \mathcal{S}_\sigma(\tuple{a_1, \ldots, a_n}) &= \mathcal{T}(\mathcal{A}_\sigma(a_1), \ldots, \mathcal{A}_\sigma(a_n)) \qquad
  \mathcal{S}_\sigma(\prj{i} a) = \key{proj}_i(\mathcal{A}_\sigma(a)) \\
  \mathcal{S}_\sigma(\inl a) &= \key{inl}(\mathcal{A}_\sigma(a)) \qquad
  \mathcal{S}_\sigma(\inr a) = \key{inr}(\mathcal{A}_\sigma(a)) \qquad
  \mathcal{S}_\sigma(\prim{p}{a_1}{a_2}) = \key{prim}_p(\mathcal{A}_\sigma(a_1), \mathcal{A}_\sigma(a_2)) \\
  \mathcal{S}_\sigma(\Mfold a) &= \mathcal{S}_\sigma(\Munfold a) = \mathcal{A}_\sigma(a)
\end{align*}
With $T$ fixed, writing $\SEM{e}\sigma\,\kappa$ for $\SEM{e}^{\Ml{Imp}}_T\sigma\,\kappa$:
\begin{align*}
  \SEM{r}\sigma\,\kappa &= \kappa(\mathcal{S}_\sigma(r)) \\
  \SEM{\assign{s}{e_1}{e_2}}\sigma\,\kappa &= \SEM{e_1}\sigma\,(\lambda D.\,
    \Lambda(\lambda X.\, \SEM{e_2}(\sigma, X)\,\kappa) \app D) \\
  \SEM{\Casei{a}{s}{e_1}{e_2}}\sigma\,\kappa &=
    \mathcal{C}(\mathcal{A}_\sigma(a),\; \lambda X.\, \SEM{e_1}(\sigma, X)\,\kappa,\;
      \lambda X.\, \SEM{e_2}(\sigma, X)\,\kappa) \\[1ex]
  \key{code}^{\mathrm{K}}_T(e) &= \Lambda(\lambda X.\, \Lambda(\lambda Y.\,
    \SEM{e}^{\Ml{Imp}}_T\,(X, Y)\,(\lambda D.\, D))) \\
  \SEM{(e_0, \ldots, e_{m-1};\; e)}^{\Ml{Imp}} &=
    \SEM{e}^{\Ml{Imp}}_{\key{table}^{\mathrm{K}}(e_0,\ldots,e_{m-1})}\,\rho_0\,(\lambda D.\, D)
\end{align*}
where $(X, Y)$ is the frame with the environment $X$ in slot $0$ and the
argument $Y$ in slot $1$, and $\key{table}^{\mathrm{K}}$ is $\Ml{GlobalFun}$'s
table with $\key{code}^{\mathrm{K}}$.

\paragraph{$\Ml{CFG}$} (\texttt{Semantics/Graph/CFG.agda}). A label is an
index into the list of the blocks a tail can see, so a jump runs the
block in the frame as it is. An assignment $\assign{s}{r}{t}$ and a case
write slot $s = |\Delta'|$: they continue in the frame cut to its first
$s$ slots, $\Tcut{\sigma}{s}$ (\texttt{drop}), extended with the new slot.
With $T$ and the meanings $B$ of the visible blocks fixed:
\begin{align*}
  \SEM{\Treturn{r}}\sigma &= \mathcal{S}_\sigma(r) \qquad
  \SEM{\Tgoto{l}}\sigma = B(l)(\sigma) \\
  \SEM{\assign{s}{r}{t}}\sigma &=
    \Lambda(\lambda X.\, \SEM{t}(\Tcut{\sigma}{s}, X)) \app \mathcal{S}_\sigma(r) \\
  \SEM{\Tcase{a}{s}{l_1}{l_2}}\sigma &=
    \mathcal{C}(\mathcal{A}_\sigma(a),\;
      \lambda X.\, B(l_1)(\Tcut{\sigma}{s}, X),\; \lambda X.\, B(l_2)(\Tcut{\sigma}{s}, X)) \\[1ex]
  \key{blocks}_T(t_0, \ldots, t_{j-1})(l) &=
    \SEM{t_l}^{\Ml{CFG}}_{T,\key{blocks}_T(t_0, \ldots, t_{l-1})} \qquad (l < j) \\
  \mathit{run}_T((t_0, \ldots, t_{j-1};\; t)) &= \SEM{t}^{\Ml{CFG}}_{T,\key{blocks}_T(t_0, \ldots, t_{j-1})} \\
  \key{code}^{\mathrm{CFG}}_T(b) &= \Lambda(\lambda X.\, \Lambda(\lambda Y.\, \mathit{run}_T(b)(X, Y))) \\
  \SEM{(b_0, \ldots, b_{m-1};\; b)}^{\Ml{CFG}} &=
    \mathit{run}_{\key{table}^{\mathrm{CFG}}(b_0,\ldots,b_{m-1})}(b)(\rho_0)
\end{align*}

\paragraph{$\Ml{X86}$} (\texttt{Semantics/Graph/X86.agda}). What runs next,
$\kappa$, is given the frame $\sigma$, the registers $R$ and the flag $F$,
$\kappa(\sigma, R, F)$. The frame is the environment, as in $\Ml{CFG}$, and
writing slot $s$ forgets the slots after it; a slot write binds its value
as $\Ml{CFG}$'s assignment does, and a register write, and the flag's, is
a substitution. Only $\key{callq}$ can diverge ($\emptyset$), so only it
writes its result strictly, $\key{guard}(D; E)$ ($E$ provided that
$D \NE$). An unset register and an unset flag hold $\{\omega\}$, as an
unbound variable does ($R_\omega$ has every register $\{\omega\}$); typed
code never reads them. With $\mathcal{A}_\sigma$ of the table $T$:
\begin{align*}
  \mathcal{O}_{\sigma,R}(a) &= \mathcal{A}_\sigma(a) \ \ (a \text{ an atom}) \qquad
  \mathcal{O}_{\sigma,R}(\%r) = R(r) \qquad
  \mathcal{O}_{\sigma,R}(\Mfold a) = \mathcal{O}_{\sigma,R}(\Munfold a) = \mathcal{O}_{\sigma,R}(a) \\
  \mathcal{O}_{\sigma,R}(\key{pair}(a_1, a_2)) &= \key{pair}(\mathcal{O}_{\sigma,R}(a_1), \mathcal{O}_{\sigma,R}(a_2)) \qquad
  \mathcal{O}_{\sigma,R}(\key{tuple}(\vec{a})) = \mathcal{T}(\mathcal{O}_{\sigma,R}(\vec{a})) \\
  \mathcal{O}_{\sigma,R}(\key{inl}(a)) &= \key{inl}(\mathcal{O}_{\sigma,R}(a))
    \qquad\text{and likewise for } \key{inr} \\
  \key{put}_{\%r}(D, \kappa)(\sigma, R, F) &= \kappa(\sigma, R[r := D], F) \qquad
  \key{put}_{s}(D, \kappa)(\sigma, R, F) = \Lambda(\lambda X.\, \kappa((\Tcut{\sigma}{s}, X), R, F)) \app D
\end{align*}
and, writing $\mathcal{O}$ for $\mathcal{O}_{\sigma,R}$,
\begin{align*}
  \SEM{\key{movq}\; a, d}\kappa(\sigma, R, F) &= \key{put}_d(\mathcal{O}(a), \kappa)(\sigma, R, F) \\
  \SEM{\key{subq}\; a, \%r}\kappa(\sigma, R, F) &= \kappa(\sigma, R[r := \key{prim}_{\key{sub}}(R(r), \mathcal{O}(a))], F) \\
  \SEM{\key{cmpq}\; a_2, a_1}\kappa(\sigma, R, F) &= \kappa(\sigma, R, \key{prim}_{\key{eq}}(\mathcal{O}(a_1), \mathcal{O}(a_2))) \\
  \SEM{\key{sete}\; d}\kappa(\sigma, R, F) &= \key{put}_d(F, \kappa)(\sigma, R, F) \\
  \SEM{\key{alloc}\; o, d}\kappa(\sigma, R, F) &= \key{put}_d(\mathcal{O}(o), \kappa)(\sigma, R, F) \\
  \SEM{\key{load}\; \pi_i\; a, d}\kappa(\sigma, R, F) &= \key{put}_d(\key{proj}_i(\mathcal{O}(a)), \kappa)(\sigma, R, F) \\
  \SEM{\key{load}\; \key{fst}\; x, \%r}\kappa(\sigma, R, F) &= \kappa(\sigma, R[r := \key{car}(\sigma(x))], F) \\
  \SEM{\key{load}\; \key{snd}\; x, \%r}\kappa(\sigma, R, F) &= \kappa(\sigma, R[r := \key{cdr}(\sigma(x))], F) \\
  \SEM{\key{callq}\; a}\kappa(\sigma, R, F) &= \key{guard}(C;\; \kappa(\sigma, R_C, \{\omega\}))
    \qquad C = (\mathcal{O}(a) \app R(\rdi)) \app R(\rsi)
\end{align*}
with $R_C$ holding $C$ in $\rax$ and $\{\omega\}$ in $\rdi$ and $\rsi$. A
sequence runs under the meanings $B$ of the blocks it may jump to; the
$\key{case}$ pseudo-instruction is $\Ml{CFG}$'s $\key{case}$:
\begin{align*}
  \SEM{i; \beta}_B &= \SEM{i}\,\SEM{\beta}_B \qquad
  \SEM{\key{retq}}_B(\sigma, R, F) = R(\rax) \qquad
  \SEM{\key{jmp}\; l}_B(\sigma, R, F) = B(l)(\sigma) \\
  \SEM{\Mxcase{x}{s}{l_1}{l_2}}_B(\sigma, R, F) &=
    \mathcal{C}(\sigma(x),\; \lambda X.\, B(l_1)(\Tcut{\sigma}{s}, X),\; \lambda X.\, B(l_2)(\Tcut{\sigma}{s}, X)) \\
  \key{blocks}_T(\beta_0, \ldots, \beta_{j-1})(l)(\sigma) &=
    \SEM{\beta_l}^{\Ml{X86}}_{T,\, \key{blocks}_T(\beta_0, \ldots, \beta_{l-1})}(\sigma, R_\omega, \{\omega\}) \\
  \mathit{run}_T((n;\; \vec{\beta};\; \beta)) &= \SEM{\beta}^{\Ml{X86}}_{T,\, \key{blocks}_T(\vec{\beta})} \\
  \key{code}^{\mathrm{X}}_T(b) &= \Lambda(\lambda X.\, \Lambda(\lambda Y.\, \mathit{run}_T(b)(\rho_0, R_{X,Y}, \{\omega\}))) \\
  \SEM{(b_0, \ldots, b_{m-1};\; b)}^{\Ml{X86}} &= \mathit{run}_{\key{table}^{\mathrm{X}}(b_0, \ldots, b_{m-1})}(b)(\rho_0, R_\omega, \{\omega\})
\end{align*}
with $R_{X,Y}$ holding $\{\omega\}$, $X$ and $Y$ in $\rax$, $\rdi$ and
$\rsi$, and $\key{table}^{\mathrm{X}}$ $\Ml{GlobalFun}$'s table with
$\key{code}^{\mathrm{X}}$.

%-------------------------------------------------------------------------------
\subsection{Relations}\label{sec:m-den-rel}

The proofs are in \texttt{proof/Correctness/Graph/}; the statements of the
theorems for closed programs (of each pass, and of the compiler) are in
\texttt{Theorems/\allowbreak Denotational.agda}, under the names given below, each
defined as its proof. Every pass but delay
preserves the meaning of a program exactly ($\simeq$); delay preserves and
reflects the constants a program produces and whether it produces
anything. A typed term has no unbound variable, so no relation needs to
say what the variables out of scope mean. Each theorem holds for a closed
program of any type.

\subsubsection{Simplify}

Simplify's proof (\texttt{SimplifyCorrect.agda}) is that of the untyped
pass (Part~\ref{part:untyped}, Lemma~\ref{lem:simplify}): the
transformations rest on $\beta$ for a value
(Lemma~\ref{lem:beta-value}, \texttt{let-$\beta$}, \texttt{case-inl},
\texttt{case-inr}), whose operand is nonempty and consistent, as every
value's denotation is (\texttt{sv-ne}). Environments are nonempty and
consistent ($\mathit{Wf}$), and the facts $E$ are true of the source's
environment $\gamma$ in the target's $\gamma'$
($\mathit{Ok}_E(\gamma, \gamma')$): for each $x$ with $E(x) = (T, K)$,
$\gamma(x) \simeq \SEM{T}\gamma'$ and $K$ describes $\gamma(x)$
($\mathit{Desc}$, $\mathit{KDesc}$): if $K$ is a value,
$\gamma(x) \simeq \SEM{K}\gamma'$, and if $K = \langle \lam{y} N, E_f
\rangle$, then $\gamma(x) \simeq \SEM{\lam{y} N}\gamma_f$ for some
$\mathit{Wf}$ $\gamma_f$ with $\mathit{Ok}_{E_f}(\gamma_f, \gamma')$.
\begin{lemma}[Correctness of Simplify]\label{lem:m-den-simplify}
If $\mathit{Wf}(\gamma)$, $\mathit{Wf}(\gamma')$ and $\mathit{Ok}_E(\gamma, \gamma')$, then
\[ \SEM{M}^{\Ml{ISWIM}}\gamma \simeq \SEM{\pass{simp}_E(M)}^{\Ml{ISWIM}}\gamma'\,, \]
and $\mathcal{K}_E(M)$ describes $\SEM{M}\gamma$ (\texttt{simp-ok}, with
\texttt{let-ok} for $\pass{let}$). So
$\SEM{M}\rho_0 \simeq \SEM{\pass{simplify}(M)}\rho_0$ for every closed $M$
(\texttt{simplify-correct}).
\end{lemma}
The new rule is immediate: $\Mfold$ and $\Munfold$ are the identity, so an
$\Munfold$ of a known $\Mfold\,K$ denotes $\SEM{K}$. Unlike untyped ISWIM,
whose renaming and continuity come with the syntax of any signature,
$\Ml{ISWIM}$'s renaming needs its own lemma (\texttt{ren-$\simeq$}), and its
continuity is \texttt{proof/Semantics/Graph/ISWIM.agda}'s.

\subsubsection{CC: Enclose}

\begin{align*}
  \mathit{Rel}(\rho, \rho') &\iff \forall x \text{ in scope}.\; \rho(x) \simeq \rho'(x) \land \rho'(x) \NE
\end{align*}
\begin{lemma}[Correctness of Enclose]\label{lem:m-den-enclose}
If $\mathit{Rel}(\rho, \rho')$, then
\[ \SEM{M}^{\Ml{ISWIM}}\rho \simeq \SEM{\pass{enclose}(M)}^{\Ml{2}}\rho' \]
(\texttt{enc}); in particular $\SEM{M}\rho_0 \simeq \SEM{\pass{enclose}(M)}^{\Ml{2}}\rho_0$
for a closed $M$ (\texttt{enclose-correct}; the proof is
\texttt{EncloseCorrect.agda}).
\end{lemma}
Every case but $\lambda$ is a congruence. A function becomes a closure
that captures all the variables in scope: they are nonempty, so its guard
holds, and the code's environment agrees with $\rho'$ on them.

\subsubsection{CC: Optimize}

The environments agree on the variables that $M$ uses, through the
renaming the pass carries, and $\rho$ is nonempty on all the variables in
scope:
\begin{align*}
  \mathit{NE}(\rho) &\iff \forall x \text{ in scope}.\; \rho(x) \NE \qquad
  \mathit{Agree}_M(\rho, \rho') \iff \forall x \in \FV(M).\; \rho(x) \simeq \rho'(x)
\end{align*}
The captured values of a closure and the kept ones are related by the
selection of the kept positions $i_1 < \cdots < i_k$:
\begin{align*}
  \mathit{Th}(D_1, \ldots, D_n;\; E_1, \ldots, E_k) &\iff
    D_{i_j} \simeq E_j \text{ for each } j, \text{ and } D_i \NE \text{ for each dropped } i
\end{align*}
A dropped captured term is a variable, which is nonempty by
$\mathit{NE}(\rho)$. So the guards of the two closures hold together, and
the environments of the codes agree on the variables the code uses.
\begin{lemma}[Correctness of Optimize]\label{lem:m-den-optimize}
If $\mathit{NE}(\rho)$ and $\mathit{Agree}_M(\rho, \rho')$, then
\[ \SEM{M}^{\Ml{2}}\rho \simeq \SEM{\pass{optimize}(M)}^{\Ml{2}}\rho' \]
(\texttt{OptimizeCorrect.agda}, \texttt{opt-ok}, with \texttt{th-ok}); so
for a closed $M$, $\SEM{M}^{\Ml{2}}\rho_0 \simeq \SEM{\pass{optimize}(M)}^{\Ml{2}}\rho_0$
(\texttt{optimize-correct}).
\end{lemma}

\subsubsection{CC: Concretize}

Each variable means what its reference means:
\begin{align*}
  \SEM{x}_r\rho &= \rho(x) \qquad
  \SEM{\prj{i} x}_r\rho = \key{proj}_i(\rho(x)) \qquad
  \rho_2 \sim_\varphi \rho_3 \iff \forall x \text{ in scope}.\; \rho_2(x) \simeq \SEM{\varphi(x)}_r\rho_3
\end{align*}
A $\Ml{Clos}$ closure is strict in its captured values; a $\Ml{EnvClos}$
closure is its code applied to their tuple, which is nonempty exactly
when they are. Then the code applied to the tuple is the function of the
body on it, and the code finds $x_i$ as $\prj{i} e$ of the tuple:
\begin{lemma}[Code Applied to an Environment]\label{lem:m-den-code-env}
If $\consis(E)$, $E \NE$, $F$ is monotone in its first argument, and an
element of $F(E, Y)$ needs only a finite, nonempty part of $E$, then
\[ \Lambda(\lambda X.\, \Lambda(\lambda Y.\, F(X, Y))) \app E \simeq \Lambda(\lambda Y.\, F(E, Y))\,. \]
If $D_j \NE$ for every $j$, then $\key{proj}_i(\mathcal{T}(D_1, \ldots, D_n)) \simeq D_i$.
\end{lemma}
\begin{lemma}[Correctness of Concretize]\label{lem:m-den-concretize}
If $\consis(\rho_3)$ and $\rho_2 \sim_\varphi \rho_3$, then
\[ \SEM{M}^{\Ml{2}}\rho_2 \simeq \SEM{\pass{concretize}_\varphi(M)}^{\Ml{3}}\rho_3 \]
(\texttt{ConcretizeCorrect.agda}, \texttt{conc}, with \texttt{EnvRel}); so
\[ \SEM{M}^{\Ml{2}}\rho_0 \simeq \SEM{\pass{concretize}_\emptyset(M)}^{\Ml{3}}\rho_0 \]
for a closed $M$ (\texttt{concretize-correct}).
\end{lemma}

\subsubsection{CC: Delay}\label{sec:m-den-delay}

Delay changes how a closure is represented, so its correctness is not an
equation (\texttt{DelayPreserve.agda}, \texttt{DelayReflect.agda},
\texttt{DelayRel.agda}, \texttt{DelayCorrect.agda}): a logical relation
between source values and target sets shows that the pass preserves, and
the approximant logical relation between source sets and finite sets of
target values that it reflects, the constants and nonemptiness. Both are
those of the dynamically typed pipeline's delay (Part~\ref{part:dyn})
without their clauses for $\key{project}$, and $\Mfold$ and $\Munfold$ are
the identity on both sides.

\paragraph{Preservation.}
The relation $v \sqsubseteq D'$ between a $\Ml{EnvClos}$ value $v$ and a set
$D'$ of $\Ml{FlatClos}$ values, by recursion on $v$; a function is exhibited
by a set with a pair, a closure, whose calls ($\odot$, as $\key{callc}$)
exhibit its results:
\begin{align*}
  \nu \sqsubseteq D' &\iff \exists u\, v.\; \pairv{u}{v} \in D' \\
  (U \mapsto w) \sqsubseteq D' &\iff (\exists u\, v.\; \pairv{u}{v} \in D') \land
     \forall E'.\; \consis(E') \land U \sqsubseteq E' \implies w \sqsubseteq D' \odot E' \\
  \tup{i}{d} \sqsubseteq D' &\iff d \sqsubseteq \key{proj}_i(D') \\
  \Left{d} \sqsubseteq D' &\iff d \sqsubseteq \{ d' \mid \Left{d'} \in D' \} \qquad
  \Right{d} \sqsubseteq D' \iff d \sqsubseteq \{ d' \mid \Right{d'} \in D' \} \\
  v \sqsubseteq D' &\iff v \in D' \quad \text{otherwise} \\[1ex]
  U \sqsubseteq E' &\iff \forall u \in U.\; u \sqsubseteq E' \qquad
  \rho \sqsubseteq \rho' \iff \forall x.\; \forall v \in \rho(x).\; v \sqsubseteq \rho'(x)
\end{align*}
A value related to $D'$ makes $D'$ nonempty.
\begin{lemma}[Delay Preserves]\label{lem:m-den-delay-preserve}
If $\rho'(x) \NE$ for every $x$, $\consis(\rho')$, and $\rho \sqsubseteq \rho'$, then for every $v$,
\[ v \in \SEM{M}^{\Ml{3}}\rho \implies v \sqsubseteq \SEM{\pass{delay}(M)}^{\Ml{4}}\rho' \]
(\texttt{preserve}).
\end{lemma}

\paragraph{Reflection.}
The \emph{approximant logical relation} $\ApproxLR{D}{V'}$ relates an
arbitrary set $D$ of $\Ml{EnvClos}$ values to a finite target approximant,
a finite set $V'$ of $\Ml{FlatClos}$ values, by recursion on
$\mathit{depth}(V')$. A target closure is a set of pairs whose calls
combine the codes and environments of several pairs, so a call is a batch
($\mathit{Calls}$); $\tau$ ranges over $\key{tup}_i$, $\key{left}$ and
$\key{right}$:
\begin{align*}
  \mathit{Entry}(V', U_0, w) &\iff \exists F \NE.\; \exists d.\;
     \pairv{F \mapsto (U_0 \mapsto w)}{d} \in V',\; F \subseteq \key{cdr}(V') \\
  \mathit{Calls}(V', U, W) &\iff
  \begin{array}[t]{l}
    \forall w \in W.\; \exists U_0 \NE.\; \mathit{Entry}(V', U_0, w) \land U_0 \subseteq U \\
    \land\; \forall u \in U.\; \exists U_0\, w.\; \mathit{Entry}(V', U_0, w) \land u \in U_0
  \end{array} \\[1ex]
  \ApproxLR{D}{V'} &\iff
  \begin{array}[t]{l}
    \forall c \in V'.\; c \in D \\
    \land\; (V' \NE \implies D \NE) \\
    \land\; \forall U\, W\, E.\;
       \mathit{Calls}(V', U, W) \land \consis(E) \land \consis(U) \\
    \qquad\qquad \land\; \ApproxLR{E}{U} \implies \ApproxLR{D \app E}{W} \\
    \land\; \forall \tau\, W.\; W \subseteq \{ d \mid \tau(d) \in V' \}
       \implies \ApproxLR{\{ d \mid \tau(d) \in D \}}{W}
  \end{array} \\[1ex]
  \rho \rhd \rho' &\iff \forall x\, V'.\; V' \subseteq \rho'(x) \implies \ApproxLR{\rho(x)}{V'}
\end{align*}
The third clause is the logical one: the relation occurs negatively in it,
on the arguments of the calls. It is well founded because the arguments
$U$ and the results $W$ of $\mathit{Calls}(V', U, W)$ are shallower than
$V'$ (\texttt{Bnd-Calls-args}, \texttt{Bnd-Calls-apps}); the Agda iterates
one layer of the relation (\texttt{ApproxLRLayer}) on fuel
(\texttt{ApproxLR}$_\leq$) only to make the recursion structural, and shows
that any fuel bounding the depth gives the same relation
(\texttt{ApproxLR-depth-irrelevant}, with \texttt{unfold} and
\texttt{fold}).

\paragraph{Why a type index cannot replace the depth.} One might try to
define the relation by recursion on the $\Ml{EnvClos}$ type $A$ of the term,
$R_A(D, V')$, as a logical relation usually is. The function case is fine,
since $A$ and $C$ are smaller than $A \rightarrow C$:
$R_{A \rightarrow C}(D, V')$ asks $R_A$ of the arguments and $R_C$ of the
results. But $\Mfold$ and $\Munfold$ are the identity on denotations, so
the relation at $\Mmu{X}{A}$ has to be the one at its unfolding,
$R_{\Mmu{X}{A}} = R_{\Munroll{X}{A}}$, which is not structurally smaller.
At $A_0 = \Mmu{X}{X \rightarrow \key{Int}}$, the type of $\omega$'s
argument, $R_{A_0}$ unfolds to $R_{A_0 \rightarrow \key{Int}}$, whose
clause for calls asks for $R_{A_0}$ on the arguments: the relation at the
same type again. No other order on types repairs this, since unfolding
grows the type, and the relation cannot be an inductive datatype instead,
since it occurs negatively. What decreases is the depth of the finite
approximant: the arguments and results of $\mathit{Calls}(V', U, W)$ are
shallower than $V'$, and $\Mfold$ does not change a value, so the depth
needs no help from the type. This is the denotational counterpart of step
indexing for recursive types (Appel and McAllester; Ahmed): the depth
index is needed because of recursive types, whether $\TAny$ or a
user-written $\mu$, not because of dynamic typing or errors (the comment
of \texttt{DelayRel.agda}; the design note, \S3.2).

\begin{lemma}[Delay Reflects]\label{lem:m-den-delay-reflect}
If $\rho(x) \NE$ for every $x$, $\consis(\rho)$, $\consis(\rho')$, and $\rho \rhd \rho'$,
then for every finite $V'$,
\[ V' \subseteq \SEM{\pass{delay}(M)}^{\Ml{4}}\rho' \implies \ApproxLR{\SEM{M}^{\Ml{3}}\rho}{V'} \]
(\texttt{delay-reflect}).
\end{lemma}
A target call $V' \subseteq L' \odot N'$ is a batch of calls of a finite
$W \subseteq L'$ on a finite $U \subseteq N'$ ($\mathit{Calls}(W, U, V')$,
\texttt{collect}), which the source application reflects; a $\key{case}$
takes, by consistency, all its results from one branch.

\begin{theorem}[Correctness of Delay]\label{thm:m-den-delay}
For every closed $\Ml{EnvClos}$ term $M$ and constant $c$,
\[
  c \in \SEM{M}^{\Ml{3}}\rho_0 \iff c \in \SEM{\pass{delay}(M)}^{\Ml{4}}\rho_0
  \qquad\text{and}\qquad
  \SEM{M}^{\Ml{3}}\rho_0 \NE \iff \SEM{\pass{delay}(M)}^{\Ml{4}}\rho_0 \NE
\]
(\texttt{delay-correct-const}, \texttt{delay-correct-nonempty}).
\end{theorem}

\subsubsection{CC: Globalize}

The definitions made so far, $\vec{d} = d_0, \ldots, d_{k-1}$, each typed
over the older ones, have their own table, $\key{made}(\vec{d})$, which is
$\key{table}(\vec{d})$ before the definitions are weakened to the whole
table. The definitions $\vec{d}'$ extend $\vec{d}$ when $\vec{d}$ is a prefix
of $\vec{d}'$, $\mathit{Extends}(\vec{d}, \vec{d}')$, and then
$\key{made}(\vec{d}')(k) = \key{made}(\vec{d})(k)$ for $k < |\vec{d}|$
(\texttt{GlobalizeCorrect.agda}).
\begin{lemma}[Weakening]\label{lem:m-den-wk}
If $M$ is typed over the table $\Sigma$ and $T'(k) \simeq T(k)$ for every
$k < |\Sigma|$, then $\SEM{M}^{\Ml{5}}_{T'}\rho \simeq \SEM{M}^{\Ml{5}}_T\rho$
(\texttt{wk-sem}).
\end{lemma}
\begin{lemma}[Correctness of Glob]\label{lem:m-den-glob}
If $\pass{glob}(M, \vec{d}) = (M', \vec{d}')$, then $\mathit{Extends}(\vec{d}, \vec{d}')$ and
$\SEM{M}^{\Ml{4}}\rho \simeq \SEM{M'}^{\Ml{5}}_{\key{made}(\vec{d}')}\rho$ for every $\rho$
(\texttt{glob-ok}).
\end{lemma}
\begin{theorem}[Correctness of Globalize]\label{thm:m-den-globalize}
For every closed $\Ml{FlatClos}$ term $M$,
\[ \SEM{M}^{\Ml{4}}\rho_0 \simeq \SEM{\pass{globalize}(M)}^{\Ml{5}} \]
(\texttt{globalize-correct}, with \texttt{table-close}: the table of the
program agrees with $\key{made}$).
\end{theorem}

\subsubsection{Remove Complex Operands}

Fix tables $T$ (of the input) and $T'$ (of the output) with
$T(k) \simeq T'(k)$ and $\consis(T'(k))$ for every $k$. The environment
$\gamma$ of the input agrees with the environment $\sigma$ of the output
through the renaming the pass carries, $\mathit{Ag}(\gamma, \sigma)$
(\texttt{Ag}). Each construct whose operands rco names means a strict
operator $O$ on the meanings of its operands ($\mathit{Strict}(O)$ of
Part~\ref{part:untyped}, \texttt{proof/Framework/GraphModel/Strict.agda}):
$\key{pair}$, $\odot$, $\mathcal{T}$, $\key{proj}_i$, $\key{inl}$,
$\key{inr}$, $\key{prim}_p$, the identity for $\Mfold$ and $\Munfold$, and
$\mathcal{C}(-, \ldots)$ of the branches for $\key{case}$. Binding an
operand by a $\key{let}$ is the $\beta$ law of its body, so it changes
nothing:
\begin{lemma}[Let Chain]\label{lem:m-den-let-chain}
Let $O$ be strict, $(\vec{b};\; a_1, \ldots, a_n) = \pass{atomize}(\pass{rco}(M_1), \ldots, \pass{rco}(M_n))$,
$\sigma$ nonempty and consistent, and $\mathit{Ag}(\gamma, \sigma)$.
If, for every nonempty, consistent $\sigma'$ that extends $\sigma$ with
the variables of $\vec{b}$,
$O(\SEM{a_1}^{\Ml{5}}_{T'}\sigma', \ldots, \SEM{a_n}^{\Ml{5}}_{T'}\sigma') \simeq \SEM{X}^{\Ml{5}}_{T'}\sigma'$, then
$O(\SEM{M_1}^{\Ml{5}}_{T}\gamma, \ldots, \SEM{M_n}^{\Ml{5}}_{T}\gamma) \simeq
  \SEM{\pass{lets}(\vec{b}, X)}^{\Ml{5}}_{T'}\sigma$
(\texttt{atomize-sem}, \texttt{atomizes-sem}, \texttt{atomize$_2$-sem}).
\end{lemma}
\begin{theorem}[Correctness of RCO]\label{thm:m-den-rco}
If $\sigma$ is nonempty and consistent and $\mathit{Ag}(\gamma, \sigma)$, then
\[ \SEM{M}^{\Ml{5}}_{T}\gamma \simeq \SEM{\Tembed{\pass{rco}(M)}}^{\Ml{5}}_{T'}\sigma \]
(\texttt{rco-sem}). For every $\Ml{GlobalFun}$ program $P$,
\[ \SEM{P}^{\Ml{5}} \simeq \SEM{\Tembed{\pass{rco}(P)}}^{\Ml{5}} = \SEM{\pass{rco}(P)}^{\Ml{Mon}} \]
(\texttt{rco-correct}, from \texttt{table-rco} of \texttt{RCOCorrect.agda}).
\end{theorem}

\subsubsection{Lift Locals}

A variable of a $\Ml{Mon}$ body is the slot at its level, so the frame
$\sigma$ of the lifted body is also the environment of the source, and no
relation between them is needed. Fix tables $T$ (of $\Ml{GlobalFun}$) and $T'$
(of $\Ml{Imp}$) with $T(k) \simeq T'(k)$ and $\consis(T(k))$. Lifted code is
run with a continuation that means what the enclosing assignments make of
the value: nothing at the end of a body, and $\Lambda(G) \app D$ for the
right-hand side of an assignment whose rest means $G$,
$K_{\cdot}(D) = D$ and $K_G(D) = \Lambda(G) \app D$. The commuting
conversions are those of Part~\ref{part:untyped}
(Lemmas~\ref{lem:let-law} and~\ref{lem:case-law},
\texttt{proof/Framework/GraphModel/Laws.agda}).
\begin{theorem}[Correctness of Lift Locals]\label{thm:m-den-lift}
If $\consis(\sigma)$, then for every continuation $K$ of the form above,
\[ K(\SEM{\Tembed{M}}^{\Ml{5}}_{T}\sigma) \simeq \SEM{\pass{lift}(M)}^{\Ml{Imp}}_{T'}\sigma\,K \]
(\texttt{lift-sem}); the tables agree definition by definition
(\texttt{table-lift}), and for every $\Ml{Mon}$ program $P$,
\[ \SEM{P}^{\Ml{Mon}} \simeq \SEM{\pass{lift}(P)}^{\Ml{Imp}} \]
(\texttt{lift-locals-correct}; the proof is \texttt{LiftLocalsCorrect.agda}).
\end{theorem}

\subsubsection{Explicate Control}

The pass threads the list of blocks: the blocks $\vec{t}'$ it returns
extend the blocks $\vec{t}$ it was given, $\vec{t} \preceq \vec{t}'$ (a
prefix, \texttt{Ext}), and the blocks keep their meanings, with the
blocks added after them (\texttt{ext-at}). Fix tables $T$ (of $\Ml{Imp}$)
and $T'$ (of $\Ml{CFG}$) with $T(k) \simeq T'(k)$. Assigning slot $s$ and
then running $k$ with the blocks $\vec{t}$ is
$\key{asg}_{s,k,\vec{t}}(\sigma, D) = \Lambda(\lambda X.\,
\SEM{k}^{\Ml{CFG}}_{T',\key{blocks}_{T'}(\vec{t})}(\Tcut{\sigma}{s}, X)) \app D$
(\texttt{asg}). Each case is a congruence, and a $\key{case}$ passes its
continuation to its branches, which reach it through the join point.
\begin{lemma}[Continuation-Passing and Tail]\label{lem:m-den-cps}
If $\pass{assign}(e, s, k, \vec{t}) = (t;\, \vec{t}')$, then $\vec{t} \preceq \vec{t}'$ and
$\SEM{e}^{\Ml{Imp}}_{T}\sigma\,(\lambda D.\, \key{asg}_{s,k,\vec{t}}(\sigma, D)) \simeq
  \SEM{t}^{\Ml{CFG}}_{T',\key{blocks}_{T'}(\vec{t}')}\sigma$ (\texttt{cps}).
If $\pass{tail}(e, \vec{t}) = (t;\, \vec{t}')$, then $\vec{t} \preceq \vec{t}'$ and
$\SEM{e}^{\Ml{Imp}}_{T}\sigma\,(\lambda D.\, D) \simeq
  \SEM{t}^{\Ml{CFG}}_{T',\key{blocks}_{T'}(\vec{t}')}\sigma$ (\texttt{tl}).
\end{lemma}
\begin{theorem}[Correctness of Explicate Control]\label{thm:m-den-explicate}
For every $\Ml{Imp}$ program $P$,
\[ \SEM{P}^{\Ml{Imp}} \simeq \SEM{\pass{explicate}(P)}^{\Ml{CFG}} \]
(\texttt{explicate-correct}; the proof is \texttt{ExplicateCorrect.agda}).
\end{theorem}

\subsubsection{Select Instructions}

The $\Ml{X86}$ frame is the $\Ml{CFG}$ frame, a slot write binds as
$\Ml{CFG}$'s assignment does, a register write is a substitution, and the
$\key{case}$ pseudo-instruction means $\Ml{CFG}$'s $\key{case}$. So the code
selected for a simple expression means it, definitionally, up to the
tables (\texttt{SelectCorrect.agda}): $\key{callq}$ applies the code and
environment that $\key{load}\;\key{fst}$ and $\key{load}\;\key{snd}$ read
from the closure's slot, which is $\mathcal{S}_\sigma$ of the call. What is
left is the start of a definition. Fix tables $T$ (of $\Ml{CFG}$) and $T'$
(of $\Ml{X86}$) with $T(k) \simeq T'(k)$, and the blocks $B'$ a sequence
may jump to. What runs after the value of $r$ is in $o$ means $K$ of it:
\begin{align*}
  \key{then}_{\key{ret}}(K, D) &= D \qquad
  \key{then}_{\key{to}\;s}(K, D) = \Lambda(K) \app D \\
  \mathit{Fits}_{\key{ret}}(K, k, \sigma) &\iff \forall \sigma'\, R\, F.\;
    R(\rax) \simeq \SEM{k}_{B'}(\sigma', R, F) \\
  \mathit{Fits}_{\key{to}\;s}(K, k, \sigma) &\iff \forall V.\; V \text{ consistent} \implies
    \forall R\, F.\; K(V) \simeq \SEM{k}_{B'}((\Tcut{\sigma}{s}, V), R, F)
\end{align*}
for $k$ at any registers, open closure and flag, and $V$ finite. A strict
write before $\key{then}_o$ changes nothing:
\[ \key{guard}(D; \key{then}_o(K, D)) \simeq \key{then}_o(K, D) \]
(\texttt{strict-put}).
\begin{lemma}[Select a Simple Expression]\label{lem:m-den-sel}
If $\mathit{Fits}_o(K, k, \sigma)$, then for every $R$ and $F$,
\[ \key{then}_o(K, \mathcal{S}_\sigma(r)) \simeq \SEM{\pass{sel}(r, o, k)}^{\Ml{X86}}_{T',B'}(\sigma, R, F) \]
(\texttt{sel-$\simeq$}, with \texttt{put-$\simeq$}).
\end{lemma}
The blocks $B$ of $\Ml{CFG}$ and $B'$ of $\Ml{X86}$ agree,
\[ B \mathrel{\mathit{BR}} B' \iff \forall l\, \sigma.\; B(l)(\sigma) \simeq B'(l)(\sigma)\,. \]
\begin{lemma}[Tail and Blocks]\label{lem:m-den-sel-tail}
If $B \mathrel{\mathit{BR}} B'$ and $\mathit{depth}(t) \leq n$, then for every $\sigma$, $R$ and $F$,
\[ \SEM{t}^{\Ml{CFG}}_{T,B}\,\sigma \simeq \SEM{\pass{tail}(t)}^{\Ml{X86}}_{T',B'}(\sigma, R, F) \]
(\texttt{tail-$\simeq$}); and
\[ \key{blocks}^{\Ml{CFG}}_{T}(\vec{t}) \mathrel{\mathit{BR}} \key{blocks}^{\Ml{X86}}_{T'}(\pass{tail}(\vec{t})) \]
(\texttt{blocks-$\simeq$}).
\end{lemma}
The prologue binds finite parts $V$ and $W$ of the parameters (an entry
of $\Lambda$ has a finite input), where $\Ml{CFG}$ runs the body in the
frame $(V, W)$ at once: a bind of a finite consistent set is $\beta$,
$G(V) \simeq \Lambda(G) \app V$ for $G$ monotone (\texttt{fin-$\beta$}),
and a $\Ml{CFG}$ tail is monotone in its frame (\texttt{tail-mono}).
\begin{lemma}[Start]\label{lem:m-den-sel-start}
For the body $b$ of every definition,
\[ \key{code}^{\mathrm{CFG}}_{T}(b) \simeq \key{code}^{\mathrm{X}}_{T'}(\pass{select}^{+}(b)) \]
(\texttt{code-$\simeq$}); so the tables of $P$ and of $\pass{select}(P)$ agree at
every $k$ (\texttt{table-$\simeq$}).
\end{lemma}
\begin{theorem}[Correctness of Select Instructions]\label{thm:m-den-select}
For every $\Ml{CFG}$ program $P$,
\[ \SEM{P}^{\Ml{CFG}} \simeq \SEM{\pass{select}(P)}^{\Ml{X86}} \]
(\texttt{select-correct}).
\end{theorem}
There is no precondition. Unlike the dynamically typed pipeline's proof,
there is no lemma that $\key{testinl}$, the load of the payload and
$\key{je}$ agree with $\Ml{CFG}$'s $\key{case}$ on a slot that never holds
both an $\key{inl}$ and an $\key{inr}$ ($\mathit{Tagged}$, $\mathit{Inv}$,
\texttt{case-core} there): the pseudo-instruction makes the case one for
one.

\subsubsection{End to End}

The relation of the end-to-end theorem is $\Same$ (\texttt{Same},
\texttt{CompilerCorrect.agda}): the same constants, and nonempty
together,
\begin{align*}
  \Same(D, D') &\iff (\forall c.\; c \in D \iff c \in D') \land (D \NE \iff D' \NE) .
\end{align*}
Every pass but delay gives $\simeq$, which composes with $\Same$
(\texttt{then}), and delay gives $\Same$ (Theorem~\ref{thm:m-den-delay}).

\begin{theorem}[Compiler Correctness, Graph Model]\label{thm:m-den-compiler}\sloppy
For every closed $\Ml{ISWIM}$ program $M$ of any type,
\begin{align*}
  &\Same(\SEM{M}\rho_0,\; \SEM{\pass{compile}_4(M)}^{\Ml{4}}\rho_0) &
  &\Same(\SEM{M}\rho_0,\; \SEM{\pass{compile}(M)}^{\Ml{5}}) \\
  &\Same(\SEM{M}\rho_0,\; \SEM{\pass{compile\text{-}mon}(M)}^{\Ml{Mon}}) &
  &\Same(\SEM{M}\rho_0,\; \SEM{\pass{compile\text{-}imp}(M)}^{\Ml{Imp}}) \\
  &\Same(\SEM{M}\rho_0,\; \SEM{\pass{compile\text{-}cfg}(M)}^{\Ml{CFG}}) &
  &\Same(\SEM{M}\rho_0,\; \SEM{\pass{compile\text{-}x86}(M)}^{\Ml{X86}})
\end{align*}
by Lemmas~\ref{lem:m-den-simplify}, \ref{lem:m-den-enclose},
\ref{lem:m-den-optimize} and~\ref{lem:m-den-concretize}, and
Theorems~\ref{thm:m-den-delay}, \ref{thm:m-den-globalize},
\ref{thm:m-den-rco}, \ref{thm:m-den-lift}, \ref{thm:m-den-explicate}
and~\ref{thm:m-den-select} (\texttt{Theorems/\allowbreak Denotational.agda}, each
$\Same$ as its two halves: \texttt{compile-correct-const} and
\texttt{compile-correct-nonempty}, \ldots,
\texttt{compile-x86-correct-const} and
\texttt{compile-x86-correct-nonempty}; for a main of type $\key{Int}$,
\texttt{compile-x86-int-const} and \texttt{compile-x86-int-nonempty}.
The statement for $\pass{compile}_4$ is a step of the proof,
\texttt{compile$_4$-same} of \texttt{CompilerCorrect.agda}, not one of
the public statements).
\end{theorem}
Adequacy turns the last into the operational statement
(Section~\ref{sec:m-thm}).

\paragraph{Where it doesn't match the Agda literally.}
\begin{itemize}
\item \textbf{Environments and frames} are functions from de Bruijn
  positions to sets (\texttt{Env Value}), newest first: variable $x$ is
  position \texttt{idx x}, and the frame of a definition is
  \texttt{Y $\bullet$ X $\bullet$ $\rho_0$}. Here a slot is its level, so the
  frame is $(X, Y)$; $(\sigma, X)$ is \texttt{X $\bullet$ $\sigma$}, and
  $\Tcut{\sigma}{s}$ is \texttt{drop p $\sigma$} for $p : \Delta' \sqsubseteq \Delta$
  with $s = |\Delta'|$.
\item \textbf{Tables} are indexed by typed references $f : \Sigma \ni s$
  (\texttt{Table}, \texttt{tables}), here by their levels $k$; the codes are
  \texttt{code$'$} (of $\Ml{GlobalFun}$, $\Ml{Imp}$ and $\Ml{CFG}$) and
  \texttt{code$_{\mathsf{t}}$} ($\key{code}^{\mathrm{X}}$). The empty frame and
  an unbound variable are \texttt{init} $= \{\omega\}$. Component $i$ of a
  tuple is its position \texttt{pos i}.
\item \textbf{The renamings} that optimize, delay and rco carry appear in
  their relations (\texttt{Agree (uses M) r}, \texttt{Envs r}, \texttt{Ag});
  here the environments are related by name. Optimize's relation
  $\mathit{Th}$ is over flags (\texttt{Th}).
\item \textbf{Delay}'s relations: $\sqsubseteq$ is \texttt{\_}$\sqsubseteq$\texttt{\_} with
  \texttt{PairV}, \texttt{Nth}, \texttt{Lefts}, \texttt{Rights}, lifted to
  environments by \texttt{Envs}; $\mathsf{ALR}$ is \texttt{ApproxLR}, the record
  \texttt{ApproxLRLayer} (\texttt{const-obs}, \texttt{nonempty-obs},
  \texttt{app-obs}, \texttt{tag-obs}) iterated on fuel up to the depth of
  $V'$, a list; $\mathit{Entry}$ is \texttt{CarEntry}, $\mathit{Calls}$ is
  \texttt{Calls}, $\{ d \mid \tau(d) \in D \}$ is \texttt{Untag $\tau$ D},
  and $\rho \rhd \rho'$ is the field \texttt{$\rho\rhd$} of \texttt{Envs}.
  $D \odot E$ is \texttt{D $\odot$ E} and $D \app E$ is \texttt{D $\bullet$ E}.
\item \textbf{Globalize}'s $\key{made}$ is \texttt{made} over \texttt{Made},
  $\mathit{Extends}$ is the record \texttt{Ext} (with \texttt{wk-made}).
\item \textbf{Lift locals}' continuations are \texttt{K$\langle$ k $\rangle$} for
  \texttt{k : Maybe}, \texttt{nothing} for $K_\cdot$ and \texttt{just G} for
  $K_G$.
\item \textbf{Select instructions}: $\key{then}$, $\mathit{Fits}$ and $\mathit{BR}$
  are \texttt{then}, \texttt{Fits} and \texttt{BR}; in $\mathit{Fits}$ a finite
  $V$ is \texttt{mem V} with $V \approx V$.
\end{itemize}
%===============================================================================
\section{CEK Machines}\label{sec:m-cek}

The machines (\texttt{Semantics/CEK/}) are those of Part~\ref{part:untyped}
(Section~\ref{sec:machines}) with typed values, environments and
continuations: machines in the sense of Section~\ref{sec:obs}, step
functions that move to the next state or stop with an observation. A
continuation is indexed by the type of the value it expects, a state
$\ev{M}{\rho}{\kappa}$ or $\rt{v}{\kappa}$ pairs a term or value of type
$A$ with a continuation expecting $A$, and an environment is typed by the
context of its term, so a state is well typed by construction. Every
transition of Part~\ref{part:untyped} that stops with $\error$ on a value
of the wrong kind (an operator that is not a closure, a $\prj{i}$ of a
non-tuple, a $\key{case}$ of a non-injection, a primitive on a
non-integer, an unbound variable, a missing definition or block) has no
counterpart: those states cannot be written. And there is no
$\key{project}$: \emph{no machine of this part has a transition that stops
with $\error$}; a machine stops only when it returns a value to the empty
continuation.

\subsection{Values, continuations and shared transitions}\label{sec:m-cek-shared}

A value of type $\Mmu{X}{A}$ is a $\Mfold$ of a value of type
$\Munroll{X}{A}$; a value of a base type is a constant, one of a tuple type
a tuple, one of a sum type an injection.
\[
\begin{array}{rcl}
  v &::=& \cdots \mid c \mid \tupV{v_1, \ldots, v_n} \mid \inlV v \mid \inrV v \mid \Mfold v \\
  \kappa &::=& \cdots \mid \halt \mid \Kf{tupK}{\bar{v}, \bar{M}, \rho, \kappa} \mid \Kf{getK_i}{\kappa}
    \mid \Kf{inlK}{\kappa} \mid \Kf{inrK}{\kappa} \mid \Kf{caseK}{x.M, y.N, \rho, \kappa} \\
  &\mid& \Kf{primL_p}{M, \rho, \kappa} \mid \Kf{primR_p}{n, \kappa}
    \mid \Kf{letK}{x.N, \rho, \kappa} \mid \Kf{foldK}{\kappa} \mid \Kf{unfoldK}{\kappa}
\end{array}
\]
\begin{align*}
  \obs(\Mfold v) &= \obs(v), \qquad \text{otherwise $\obs$ as in Part~\ref{part:untyped} (closures, code and packs give $\key{fun}$)} \\
  \Tres{n} &= n, \quad
  \Tres{\mathsf{true}} = \inlV \tupV{}, \quad
  \Tres{\mathsf{false}} = \inrV \tupV{} \\
  \Mach{\ell}(M) &= (\mathit{State}^{\ell},\; \step{\ell},\; \ev{M}{\emptyset}{\halt})
    \qquad (\ell \in \{\Ml{ISWIM}, \Ml{Clos}, \Ml{EnvClos}, \Ml{FlatClos}\})
\end{align*}
The transitions shared by the machines of $\Ml{ISWIM}$ to $\Ml{GlobalFun}$ are
those of Part~\ref{part:untyped} for variables, constants, tuples,
$\prj{i}$, $\key{inl}$, $\key{inr}$, $\key{case}$ and primitives, without
the cases that stop with $\error$, together with $\key{let}$ and the
coercions. The values of a tuple are collected newest first, so that the
tuple's type needs no rewriting as they come in.
\[
\begin{array}{@{}l@{\;\mapsto\;}l@{\quad}l@{}}
  \ev{x}{\rho}{\kappa} & \nxt \rt{\rho(x)}{\kappa} \\
  \ev{c}{\rho}{\kappa} & \nxt \rt{c}{\kappa} \\
  \ev{\tuple{M_1, \ldots, M_n}}{\rho}{\kappa} & \nxt \tnext(\epsilon;\, M_1 \ldots M_n;\, \rho;\, \kappa) \\
  \ev{(\prj{i} M)}{\rho}{\kappa} & \nxt \ev{M}{\rho}{\Kf{getK_i}{\kappa}} \\
  \ev{(\inl M)}{\rho}{\kappa} & \nxt \ev{M}{\rho}{\Kf{inlK}{\kappa}}
    \qquad\text{and likewise for } \key{inr} \\
  \ev{(\Case{L}{x}{M}{y}{N})}{\rho}{\kappa} & \nxt \ev{L}{\rho}{\Kf{caseK}{x.M, y.N, \rho, \kappa}} \\
  \ev{\prim{p}{L}{M}}{\rho}{\kappa} & \nxt \ev{L}{\rho}{\Kf{primL_p}{M, \rho, \kappa}} \\
  \ev{(\Let x = M \In N)}{\rho}{\kappa} & \nxt \ev{M}{\rho}{\Kf{letK}{x.N, \rho, \kappa}} \\
  \ev{(\Mfold M)}{\rho}{\kappa} & \nxt \ev{M}{\rho}{\Kf{foldK}{\kappa}} \\
  \ev{(\Munfold M)}{\rho}{\kappa} & \nxt \ev{M}{\rho}{\Kf{unfoldK}{\kappa}} \\[1ex]
  \rt{v}{\halt} & \stp \obs(v) \\
  \rt{v}{\Kf{tupK}{\bar{v}, \bar{M}, \rho, \kappa}} & \nxt \tnext(\bar{v}\,v;\, \bar{M};\, \rho;\, \kappa) \\
  \rt{\tupV{v_1, \ldots, v_n}}{\Kf{getK_i}{\kappa}} & \nxt \rt{v_i}{\kappa} \\
  \rt{v}{\Kf{inlK}{\kappa}} & \nxt \rt{(\inlV v)}{\kappa}
    \qquad\text{and likewise for } \key{inr} \\
  \rt{(\inlV u)}{\Kf{caseK}{x.M, y.N, \rho, \kappa}} & \nxt \ev{M}{\ext{\rho}{x}{u}}{\kappa} \\
  \rt{(\inrV u)}{\Kf{caseK}{x.M, y.N, \rho, \kappa}} & \nxt \ev{N}{\ext{\rho}{y}{u}}{\kappa} \\
  \rt{n}{\Kf{primL_p}{M, \rho, \kappa}} & \nxt \ev{M}{\rho}{\Kf{primR_p}{n, \kappa}} \\
  \rt{m}{\Kf{primR_p}{n, \kappa}} & \nxt \rt{\Tres{\delta(p, n, m)}}{\kappa} \\
  \rt{v}{\Kf{letK}{x.N, \rho, \kappa}} & \nxt \ev{N}{\ext{\rho}{x}{v}}{\kappa} \\
  \rt{v}{\Kf{foldK}{\kappa}} & \nxt \rt{\Mfold v}{\kappa} \\
  \rt{\Mfold v}{\Kf{unfoldK}{\kappa}} & \nxt \rt{v}{\kappa}
\end{array}
\]
The $\key{unfoldK}$ step always succeeds: a value of type $\Mmu{X}{A}$ is a
$\Mfold$. A variable is always bound, a $\prj{i}$ finds a tuple with an
$i$-th component, and a $\key{case}$ or primitive always finds an injection
or an integer, by typing.

\subsection{\texorpdfstring{$\Ml{ISWIM}$ to $\Ml{FlatClos}$}{ISWIM-mu to FlatClos-mu}}

\begin{sloppypar}
The four machines are those of the typed languages of Part~\ref{part:dyn}
(\texttt{Semantics/CEK/\{ISWIM,Clos2,Clos3,Clos4\}.agda}).
$\Ml{ISWIM}$'s closure is $\cloV{\lam{x} N}{\rho}$; an application
evaluates the operator and then the operand ($\key{appL}$, $\key{appR}$)
and runs the body:
\end{sloppypar}
\[
\begin{array}{@{}l@{\;\mapsto\;}l@{\quad}l@{}}
  \ev{(\lam{x} N)}{\rho}{\kappa} & \nxt \rt{\cloV{\lam{x} N}{\rho}}{\kappa} \\
  \ev{(L\;M)}{\rho}{\kappa} & \nxt \ev{L}{\rho}{\Kf{appL}{M, \rho, \kappa}} \\
  \rt{f}{\Kf{appL}{M, \rho, \kappa}} & \nxt \ev{M}{\rho}{\Kf{appR}{f, \kappa}} \\
  \rt{w}{\Kf{appR}{\cloV{\lam{x} N}{\rho}, \kappa}} & \nxt \ev{N}{\ext{\rho}{x}{w}}{\kappa}
\end{array}
\]
A $\Ml{Clos}$ closure evaluates its captured terms left to right, as a
tuple, and keeps their values as the environment of its code; a
$\Ml{EnvClos}$ closure the same, its code running with the tuple:
\[
\begin{array}{@{}l@{\;\mapsto\;}l@{\quad}l@{}}
  \ev{(\clos{M_1 \ldots M_n}{\bar{x}\, y} N)}{\rho}{\kappa} &
    \nxt \tnext(\epsilon;\, M_1 \ldots M_n;\, \rho;\, \Kf{closK}{\bar{x}\, y.\, N, \kappa}) \\
  \rt{\tupV{\bar{v}}}{\Kf{closK}{\bar{x}\, y.\, N, \kappa}} & \nxt \rt{\closV{\bar{x}\, y.\, N}{\bar{v}}}{\kappa} \\
  \rt{w}{\Kf{appR}{\closV{\bar{x}\, y.\, N}{\bar{v}}, \kappa}} & \nxt \ev{N}{\{\bar{x} {\mapsto} \bar{v}, y {\mapsto} w\}}{\kappa}
    & (\Ml{Clos}) \\
  \rt{w}{\Kf{appR}{\closV{e\, y.\, N}{\bar{v}}, \kappa}} & \nxt \ev{N}{\{e {\mapsto} \tupV{\bar{v}}, y {\mapsto} w\}}{\kappa}
    & (\Ml{EnvClos})
\end{array}
\]
In $\Ml{FlatClos}$ code is a value of type $\Mcodety{\Delta}{A}{C}$, and a
closure packs code with an environment of type $\Mtup{\Delta}$;
$\Tcallc{L}{M}$ evaluates the closure and the argument and runs the
closure's code with its environment and the argument, with no check:
\[
\begin{array}{@{}l@{\;\mapsto\;}l@{\quad}l@{}}
  \ev{\Tcodetm{e}{y}{N}}{\rho}{\kappa} & \nxt \rt{\TcodeV{e\, y.\, N}}{\kappa} \\
  \ev{\Tpack{L}{M}}{\rho}{\kappa} & \nxt \ev{L}{\rho}{\Kf{packL}{M, \rho, \kappa}} \\
  \ev{\Tcallc{L}{M}}{\rho}{\kappa} & \nxt \ev{L}{\rho}{\Kf{callL}{M, \rho, \kappa}} \\
  \rt{v}{\Kf{packL}{M, \rho, \kappa}} & \nxt \ev{M}{\rho}{\Kf{packR}{v, \kappa}} \\
  \rt{u}{\Kf{packR}{v, \kappa}} & \nxt \rt{\Tpack{v}{u}}{\kappa} \\
  \rt{f}{\Kf{callL}{M, \rho, \kappa}} & \nxt \ev{M}{\rho}{\Kf{callR}{f, \kappa}} \\
  \rt{w}{\Kf{callR}{\Tpack{\TcodeV{e\, y.\, N}}{u}, \kappa}} & \nxt \ev{N}{\{e {\mapsto} u, y {\mapsto} w\}}{\kappa}
\end{array}
\]

\subsection{\texorpdfstring{$\Ml{GlobalFun}$ and $\Ml{Mon}$}{GlobalFun-mu and Mon-mu}}

A program runs with its definitions $d_0, \ldots, d_{m-1}$
(\texttt{Semantics/CEK/Clos5.agda}); values, environments and
continuations are indexed by the table $\Sigma$. Code is a reference
$\codeV{k}$ to definition $d_k$, so a call always finds its definition.
The machine takes values apart with total functions (the integer of a
value of type $\key{Int}$, the components of a tuple, the payload of an
injection, the value under a $\Mfold$: \texttt{int}, \texttt{elems},
\texttt{choose}, \texttt{unfold$_v$}), since their types say what they are;
these and the values are shared with the machines of $\Ml{Imp}$, $\Ml{CFG}$
and $\Ml{X86}$. The transitions are those of $\Ml{FlatClos}$, with
\[
\begin{array}{@{}l@{\;\mapsto\;}l@{\quad}l@{}}
  \ev{\funref{k}}{\rho}{\kappa} & \nxt \rt{\codeV{k}}{\kappa} \\
  \rt{w}{\Kf{callR}{\Tpack{\codeV{k}}{u}, \kappa}} & \nxt \ev{N_k}{\{e {\mapsto} u, y {\mapsto} w\}}{\kappa}
    & d_k = \Tcodetm{e}{y}{N_k}
\end{array}
\]
in place of those for code. A $\Ml{Mon}$ program runs as its embedding,
$\Mach{\Ml{Mon}}(P) = \Mach{\Ml{GlobalFun}}(\Tembed{P})$.

\subsection{\texorpdfstring{$\Ml{Imp}$}{Imp-mu}}

An expression runs in a frame $\sigma$ of slots typed by its frame context
(\texttt{Semantics/CEK/Imp.agda}); a slot is a typed index, so an atom
always has a value. $\assign{s}{e_1}{e_2}$ and a $\key{case}$ write the
next slot, and only the expression in their scope sees it; $(\sigma, v)$ is
$\sigma$ with its next slot holding $v$. A simple expression takes one step:
its operands are atoms, so it gives a value or calls a closure.
\[
\begin{array}{rcl}
  \mathit{State}^{\Ml{Imp}} &::=& \evi{e}{\sigma}{\kappa} \qquad
  \kappa \;::=\; \halt \mid \Kf{assignK}{e, \sigma, \kappa} \qquad
  \mathit{Res} \;::=\; \key{value}(v) \mid \key{call}(f, w)
\end{array}
\]
\[
\begin{array}{@{}l@{\;=\;}l@{\qquad}l@{\;=\;}l@{}}
  \mathit{run}(a, \sigma) & \key{value}(\sigma(a)) &
  \mathit{run}(\Tpack{a_1}{a_2}, \sigma) & \key{value}(\Tpack{\sigma(a_1)}{\sigma(a_2)}) \\
  \mathit{run}(\Tcallc{a_1}{a_2}, \sigma) & \key{call}(\sigma(a_1), \sigma(a_2)) &
  \mathit{run}(\tuple{a_1, \ldots, a_n}, \sigma) & \key{value}(\tupV{\sigma(a_1), \ldots, \sigma(a_n)}) \\
  \mathit{run}(\prj{i} a, \sigma) & \key{value}(v_i) \quad (\sigma(a) = \tupV{\bar{v}}) &
  \mathit{run}(\inl a, \sigma) & \key{value}(\inlV \sigma(a)) \\
  \mathit{run}(\prim{p}{a_1}{a_2}, \sigma) & \key{value}(\Tres{\delta(p, \sigma(a_1), \sigma(a_2))}) &
  \mathit{run}(\Mfold a, \sigma) & \key{value}(\Mfold \sigma(a)) \\
  \mathit{run}(\Munfold a, \sigma) & \key{value}(v) \quad (\sigma(a) = \Mfold v)
\end{array}
\]
For a program with definitions $e_0, \ldots, e_{m-1}$:
\begin{align*}
  \mathit{cont}(v, \halt) &= \stp \obs(v) \qquad
  \mathit{cont}(v, \Kf{assignK}{e, \sigma, \kappa}) = \nxt \evi{e}{(\sigma, v)}{\kappa} \\
  \mathit{res}(\key{value}(v), \kappa) &= \mathit{cont}(v, \kappa) \qquad
  \mathit{res}(\key{call}(\Tpack{\codeV{k}}{u}, w), \kappa) = \nxt \evi{e_k}{(u, w)}{\kappa}
\end{align*}
\[
\begin{array}{@{}l@{\;\mapsto\;}l@{\quad}l@{}}
  \evi{r}{\sigma}{\kappa} & \mathit{res}(\mathit{run}(r, \sigma), \kappa) \\
  \evi{(\assign{s}{e_1}{e_2})}{\sigma}{\kappa} & \nxt \evi{e_1}{\sigma}{\Kf{assignK}{e_2, \sigma, \kappa}} \\
  \evi{(\Casei{a}{s}{e_1}{e_2})}{\sigma}{\kappa} & \nxt \evi{e_1}{(\sigma, u)}{\kappa} & \sigma(a) = \inlV u \\
  & \nxt \evi{e_2}{(\sigma, u)}{\kappa} & \sigma(a) = \inrV u
\end{array}
\]

\subsection{\texorpdfstring{$\Ml{CFG}$}{CFG-mu}}

The values, frames, $\mathit{Res}$ and $\mathit{run}$ are those of
$\Ml{Imp}$ (\texttt{Semantics/CEK/CFG.agda}). A state runs a tail with the
blocks it may jump to; an assignment or $\key{case}$ that writes slot $s$
continues in $(\Tcut{\sigma}{s}, v)$. A step runs one statement. A call
runs the start of the callee's body under $\key{seqK}$ (the caller's
assignment) or $\key{retK}$ (the caller's $\key{return}$).
\[
  \mathit{State}^{\Ml{CFG}} ::= \exc{t}{\sigma}{\vec{t}}{\kappa} \qquad
  \kappa ::= \halt \mid \Kf{seqK}{s, \sigma, t, \vec{t}, \kappa} \mid \Kf{retK}{\kappa}
\]
\begin{align*}
  \mathit{enter}(\Tpack{\codeV{k}}{u}, w, \kappa) &= \nxt \exc{t}{(u, w)}{\vec{t}}{\kappa}
    \qquad (b_k = (\vec{t};\; t)) \\
  \mathit{cont}(v, \halt) &= \stp \obs(v) \qquad
  \mathit{cont}(v, \Kf{retK}{\kappa}) = \mathit{cont}(v, \kappa) \\
  \mathit{cont}(v, \Kf{seqK}{s, \sigma, t, \vec{t}, \kappa}) &= \nxt \exc{t}{(\Tcut{\sigma}{s}, v)}{\vec{t}}{\kappa}
\end{align*}
\[
\begin{array}{@{}l@{\;\mapsto\;}l@{\quad}l@{}}
  \exc{(\Treturn{r})}{\sigma}{\vec{t}}{\kappa} & \mathit{cont}(v, \kappa) & \mathit{run}(r, \sigma) = \key{value}(v) \\
  & \mathit{enter}(f, w, \Kf{retK}{\kappa}) & \mathit{run}(r, \sigma) = \key{call}(f, w) \\
  \exc{(\assign{s}{r}{t})}{\sigma}{\vec{t}}{\kappa} &
    \nxt \exc{t}{(\Tcut{\sigma}{s}, v)}{\vec{t}}{\kappa} & \mathit{run}(r, \sigma) = \key{value}(v) \\
  & \mathit{enter}(f, w, \Kf{seqK}{s, \sigma, t, \vec{t}, \kappa}) & \mathit{run}(r, \sigma) = \key{call}(f, w) \\
  \exc{(\Tgoto{l})}{\sigma}{\vec{t}}{\kappa} & \nxt \exc{t'}{\sigma}{\vec{t}'}{\kappa} & (t';\; \vec{t}') = \key{block}(\vec{t}, l) \\
  \exc{(\Tcase{a}{s}{l_1}{l_2})}{\sigma}{\vec{t}}{\kappa} &
    \nxt \exc{t'}{(\Tcut{\sigma}{s}, u)}{\vec{t}'}{\kappa} & \sigma(a) = \inlV u,\; (t';\; \vec{t}') = \key{block}(\vec{t}, l_1) \\
  & \nxt \exc{t'}{(\Tcut{\sigma}{s}, u)}{\vec{t}'}{\kappa} & \sigma(a) = \inrV u,\; (t';\; \vec{t}') = \key{block}(\vec{t}, l_2)
\end{array}
\]

\subsection{\texorpdfstring{$\Ml{X86}$}{X86-mu}}\label{sec:m-cek-x86}

A state $\xrun\;\beta\;\vec{\beta}\;\Pi\;\eta\;m$
(\texttt{Semantics/CEK/X86.agda}) runs the sequence $\beta$ with the blocks
$\vec{\beta}$ it may jump to, the call stack $\Pi$, and the machine
$m = \langle R, \sigma, O, f \rangle$ of the type of $\beta$'s start: the
registers $R$, the frame $\sigma$, the environment $O$ of the open closure,
and the flag $f$, a boolean. An unset register, an unset flag and $O$ when
no closure is open hold nothing ($\cdot$). $\eta$ is $\alpha$, the types of
the open closure's environment (the witness that a closure's type hides),
so a register of type $\codeAB{A}{C}$ holds code of type
$\Mcodety{\eta}{A}{C}$ and one of type $\tupA$ a tuple of type
$\Mtup{\eta}$. The stack holds, for each call, the rest of the caller's
sequence, its blocks and its frame:
\[
  \mathit{State}^{\Ml{X86}} ::= \xrun\;\beta\;\vec{\beta}\;\Pi\;\eta\;\langle R, \sigma, O, f \rangle \qquad
  \Pi ::= \halt \mid \key{retK}(\beta, \vec{\beta}, \sigma, \Pi)
\]
An operand has a value of its type; $\Mfold a$ folds and $\Munfold a$ unfolds:
\begin{align*}
  m(a) &= \sigma(a) \ \ (a \text{ an atom}) \qquad
  m(\%r) = R(r) \qquad
  m(\Mfold a) = \Mfold m(a) \qquad
  m(\Munfold a) = v \ \ (m(a) = \Mfold v) \\
  \key{put}(\%r, v, m) &= \langle R[r := v], \sigma, O, f \rangle \qquad
  \key{put}(s, v, m) = \langle R, (\Tcut{\sigma}{s}, v), \cdot, f \rangle
\end{align*}
For a program with bodies $b_0, \ldots, b_{m-1}$:
\begin{align*}
  \mathit{enter}(\codeV{k}, u, w, \Pi) &=
    \nxt \xrun\;\beta\;\vec{\beta}\;\Pi\;\eta\;\langle (\rdi{:}u,\ \rsi{:}w), \emptyset, \cdot, \cdot \rangle
    \qquad (b_k = (n;\; \vec{\beta};\; \beta),\ u : \Mtup{\eta}) \\
  \key{go}((\beta;\; \vec{\beta}), \Pi, \eta, \sigma) &=
    \xrun\;\beta\;\vec{\beta}\;\Pi\;\eta\;\langle \emptyset, \sigma, \cdot, \cdot \rangle \\
  \Mach{\Ml{X86}}(b_0, \ldots, b_{m-1};\; (n;\; \vec{\beta};\; \beta)) &=
    (\mathit{State}^{\Ml{X86}},\; \step{\Ml{X86}}_{b_0 \ldots b_{m-1}},\;
     \xrun\;\beta\;\vec{\beta}\;\halt\;\epsilon\;\langle \emptyset, \emptyset, \cdot, \cdot \rangle)
\end{align*}
Writing $\langle i;\, \beta \rangle$ for $\xrun\;(i;\, \beta)\;\vec{\beta}\;\Pi\;\eta\;m$,
$m = \langle R, \sigma, O, f \rangle$, and $\langle \beta \rangle_{\eta', m'}$ for
$\xrun\;\beta\;\vec{\beta}\;\Pi\;\eta'\;m'$:
\[
\begin{array}{@{}l@{\;\mapsto\;}l@{\quad}l@{}}
  \langle \key{movq}\; a, d;\, \beta \rangle & \nxt \langle \beta \rangle_{\eta, \key{put}(d, m(a), m)} \\
  \langle \key{subq}\; a, \%r;\, \beta \rangle & \nxt \langle \beta \rangle_{\eta, \key{put}(\%r, \delta(\key{sub}, R(r), m(a)), m)} \\
  \langle \key{cmpq}\; a_2, a_1;\, \beta \rangle & \nxt \langle \beta \rangle_{\eta, \langle R, \sigma, O, (m(a_1) = m(a_2)) \rangle} \\
  \langle \key{sete}\; d;\, \beta \rangle & \nxt \langle \beta \rangle_{\eta, \key{put}(d, \Tres{f}, m)} \\
  \langle \key{alloc}\; o, d;\, \beta \rangle & \nxt \langle \beta \rangle_{\eta, \key{put}(d, m(o), m)} \\
  \langle \key{load}\; \pi_i\; a, d;\, \beta \rangle & \nxt \langle \beta \rangle_{\eta, \key{put}(d, v_i, m)}
    & m(a) = \tupV{\bar{v}} \\
  \langle \key{load}\; \key{fst}\; x, \%r;\, \beta \rangle & \nxt \langle \beta \rangle_{\Delta, \langle \Twipe{R}[r := c], \sigma, u, f \rangle}
    & \sigma(x) = \Tpack{c}{u},\ u : \Mtup{\Delta} \\
  \langle \key{load}\; \key{snd}\; x, \%r;\, \beta \rangle & \nxt \langle \beta \rangle_{\eta, \langle R[r := O], \sigma, O, f \rangle} \\
  \langle \key{callq}\; a;\, \beta \rangle & \mathit{enter}(m(a), R(\rdi), R(\rsi), \key{retK}(\beta, \vec{\beta}, \sigma, \Pi)) \\[1ex]
  \xrun\;\key{retq}\;\vec{\beta}\;\halt\;\eta\;m & \stp \obs(R(\rax)) \\
  \xrun\;\key{retq}\;\vec{\beta}\;\key{retK}(\beta', \vec{\beta}', \sigma', \Pi)\;\eta\;m &
    \nxt \xrun\;\beta'\;\vec{\beta}'\;\Pi\;\eta\;\langle (\rax{:}R(\rax)), \sigma', \cdot, \cdot \rangle \\
  \xrun\;(\key{jmp}\; l)\;\vec{\beta}\;\Pi\;\eta\;m & \nxt \key{go}(\key{block}(\vec{\beta}, l), \Pi, \eta, \sigma) \\
  \xrun\;(\Mxcase{x}{s}{l_1}{l_2})\;\vec{\beta}\;\Pi\;\eta\;m & \nxt \key{go}(\key{block}(\vec{\beta}, l_1), \Pi, \eta, (\Tcut{\sigma}{s}, u)) & \sigma(x) = \inlV u \\
  & \nxt \key{go}(\key{block}(\vec{\beta}, l_2), \Pi, \eta, (\Tcut{\sigma}{s}, u)) & \sigma(x) = \inrV u
\end{array}
\]
where $\Twipe{R}$ unsets the registers of type $\codeAB{\cdot}{\cdot}$ or
$\tupA$. There is no step that stops with $\error$: no unset register or
flag, no slot out of the frame, no value of the wrong kind, no missing
block or definition, and no check.

%-------------------------------------------------------------------------------
\subsection{Relations}\label{sec:m-cek-rel}

Each pass is proved correct by a forward simulation
(\texttt{proof/Correctness/CEK/}), $\FSim$ of Section~\ref{sec:sims} with
the answer of the dynamically typed pipeline's CEK column: a source step
that stops with $o$ may also be answered by a target run to a final state
with the same observation. Forward simulation gives equivalence
(Theorem~\ref{thm:equivalent}), and equivalence composes
(Lemma~\ref{lem:equiv-trans}). The statements of the theorems (of each
pass, and of the compiler) are in \texttt{Theorems/\allowbreak Simulation.agda},
under the names given below, each defined as its proof. Besides plus simulations
(Lemma~\ref{lem:plus-sim}), two shapes of simulation are built from a
relation $\sim$ between states
(\texttt{proof/Framework/Simulation.agda}, \texttt{Both.agda}): the
\emph{lagging} one, in which the source is $i$ steps behind a state related
to the target's,
\[ s_1 \sim^{\mathrm{lag}}_i s_2 \iff \exists t.\; s_1 \lto^i t \land t \sim s_2 \]
(\texttt{Lag}, \texttt{lag-fsim}); and $\mathit{Both}_\sim$, in which both
machines move from related states, by one or more steps each, to related
states or to final states with the same observation (\texttt{Both},
\texttt{both-fsim}).

Each pass relates source and target values ($v \approx w$), environments,
continuations ($\kappa_1 \approxk \kappa_2$), and states ($s_1 \sim s_2$).
The clauses of $\approx$ not shown are $c \approx c$,
$\tupV{\bar{v}} \approx \tupV{\bar{w}}$ if $\bar{v} \approx \bar{w}$,
$\inlV v \approx \inlV w$ and $\inrV v \approx \inrV w$ if $v \approx w$,
and $\Mfold v \approx \Mfold w$ if $v \approx w$; environments are related
pointwise unless said otherwise; the frames of $\approxk$ not shown are
related to the same frame whose terms are the pass's translations, with
related values, environments and rest (\ldots,
$\Kf{foldK}{\kappa_1} \approxk \Kf{foldK}{\kappa_2}$,
$\Kf{unfoldK}{\kappa_1} \approxk \Kf{unfoldK}{\kappa_2}$). In every pass
related values are observed alike (\texttt{obs$\approx$}). No step of
either machine stops with an error, so the error cases of the dynamically
typed pipeline's proofs (a $\key{project}$ that fails on both sides) have
no counterpart.

\subsubsection{Simplify}

The simulation goes from the target to the source (\texttt{SimplifySim.agda}),
as in Part~\ref{part:untyped} (Section~\ref{sec:sim-simplify}): the target
does less work, since the reductions that simplify did ahead of time ($\beta$
of a small value, an inlined call, a $\key{case}$ on a known injection, a
projection from a known tuple, an $\Munfold$ of a known $\Mfold$, a
primitive on known integers) are steps of the source only. A target
closure is the simplified body of a source closure, under facts $E$ that
are true of the two environments ($\rho \models_E \rho'$, \texttt{Ok}): for
each $x$ with $E(x) = (T, K)$, $\rho(x)$ is related to a value that the
value term $T$ denotes in $\rho'$ (\texttt{Strong}), and $K$ describes
$\rho(x)$ (\texttt{KDen}): a known value the same way, up to functions
(\texttt{Den}), and a known function $\langle \lam{y} N, E_f \rangle$ the
source closure $\cloV{\lam{y} N}{\rho_f}$ for some $\rho_f \models_{E_f} \rho'$.
\begin{align*}
  \cloV{\lam{y} N}{\rho} \approx \cloV{\lam{y} \pass{simp}_{E'}(N)}{\rho'}
    &\iff \rho \models_E \rho' \qquad (E' = E \text{ under } y) \\
  \ev{M}{\rho}{\kappa} \sim \ev{\pass{simp}_E(M)}{\rho'}{\kappa'}
    &\iff \rho \models_E \rho' \land \kappa \approxk \kappa' \\
  \rt{v}{\kappa} \sim \rt{v'}{\kappa'} &\iff v \approx v' \land \kappa \approxk \kappa'
\end{align*}
Where simplify reduced, the source runs ahead: if $\pass{simp}_E(M)$ is a
value, or $M$ has a known value, then $\ev{M}{\rho}{\kappa} \lto^*
\rt{v}{\kappa}$ for a $v$ it describes (\texttt{run-term},
\texttt{run-known}); at an inlined call, the source runs the operator to
the known function's closure, then continues as for $\beta$. Where the
target does more (it evaluates a duplicated small value, or the value of
a let that simplify kept), it stutters, and a match is indexed by the
number of target steps until a state related to the source's.
\begin{theorem}[Correctness of Simplify, Simulation]\label{thm:m-cek-simplify}
For every closed $\Ml{ISWIM}$ term $M$,
\[ \Equiv(\Mach{\Ml{ISWIM}}(M), \Mach{\Ml{ISWIM}}(\pass{simplify}(M))) \]
(\texttt{simplify-correct}), from the simulation of the target by the
source and the symmetry of $\Equiv$.
\end{theorem}

\subsubsection{CC: Enclose}

A closure captures every variable in scope, so its environment is the
source environment (\texttt{EncloseSim.agda}):
\begin{align*}
  \cloV{\lam{y} N}{\{x_1 {\mapsto} v_1, \ldots, x_k {\mapsto} v_k\}} \approx {}&
    \closV{x_1 \ldots x_k\, y.\, \pass{enclose}(N)}{w_1 \ldots w_k}
    \iff \forall i.\; v_i \approx w_i \\
  \ev{M}{\rho_1}{\kappa_1} \sim {}& \ev{\pass{enclose}(M)}{\rho_2}{\kappa_2}
    \iff \rho_1 \approx \rho_2 \land \kappa_1 \approxk \kappa_2 \\
  \rt{v}{\kappa_1} \sim {}& \rt{w}{\kappa_2} \iff v \approx w \land \kappa_1 \approxk \kappa_2
\end{align*}
The machines run in lock step, except that evaluating a $\lambda$ takes the
target several steps: it evaluates the variables it captures, which
collects exactly $\rho_2$, and builds the closure. So
\[ \mathit{PlusSim}(\Mach{\Ml{ISWIM}}(M), \Mach{\Ml{Clos}}(\pass{enclose}(M)), \sim)\,. \]
\begin{theorem}[Correctness of Enclose, Simulation]\label{thm:m-cek-enclose}
For every closed $M$,
\[ \Equiv(\Mach{\Ml{ISWIM}}(M), \Mach{\Ml{Clos}}(\pass{enclose}(M))) \]
(\texttt{enclose-correct}).
\end{theorem}

\subsubsection{CC: Optimize}

Optimize keeps the free variables $x_i$ of a closure
$\clos{M_1 \ldots M_n}{x_1 \ldots x_n\, y} N$ at the positions
\[ K = \{ i \mid x_i \in \FV(N) \text{ or } M_i \text{ is not a variable} \}\,; \]
$\bar{x}_K$ are the kept names, in order. Environments agree on the free
variables of a term, through the pass's renaming,
\[ \rho_1 \approx_X \rho_2 \iff \forall x \in X.\; \rho_1(x) \approx \rho_2(x) \]
(\texttt{Agree}). A closure is related to one that keeps the values at a
set $K$ of positions that covers the free variables of its code
(\texttt{OptimizeSim.agda}):
\begin{align*}
  \mathit{Thin}_K(v_1 \ldots v_n, \bar{w}) &\iff \bar{w} = (w_i)_{i \in K} \land \forall i \in K.\; v_i \approx w_i \\
  \closV{\bar{x}\, y.\, N}{\bar{v}} \approx \closV{\bar{x}_K\, y.\, \pass{optimize}(N)}{\bar{w}}
    &\iff \{ i \mid x_i \in \FV(N) \} \subseteq K \land \mathit{Thin}_K(\bar{v}, \bar{w})
\end{align*}
While the source evaluates the captured terms of a closure, its frame is
related to the target's when the values still to come give related
closures ($\mathit{Done}$, \texttt{Done}). The states are matched by
$\sim^{\mathrm{lag}}$: the source is $2d$ steps behind a synchronisation
point while it evaluates $d$ dropped free variables in a row (each in an
$\key{eval}$ and a $\key{ret}$ step, matched by no target step).
\begin{theorem}[Correctness of Optimize, Simulation]\label{thm:m-cek-optimize}
For every closed $\Ml{Clos}$ term $M$,
\[ \Equiv(\Mach{\Ml{Clos}}(M), \Mach{\Ml{Clos}}(\pass{optimize}(M))) \]
(\texttt{optimize-correct}).
\end{theorem}
The machines collect a tuple's values newest first, so the closure's
environment is the values so far followed by those still to come, with
no equation about lengths.

\subsubsection{CC: Concretize}

Environments are related through the map $\varphi$ of concretize
(\texttt{EnvRel}, \texttt{ref-val}): $\rho_2 \approx_\varphi \rho_3 \iff
\forall x.\; \rho_2(x) \approx \rho_3[\varphi(x)]$, with
$\rho_3[\prj{i} x] = w_i$ for $\rho_3(x) = \tupV{w_1, \ldots, w_n}$
(\texttt{ConcretizeSim.agda}).
\begin{align*}
  \closV{x_1 \ldots x_n\, y.\, N}{\bar{v}} \approx {}&
    \closV{e\, y.\, \pass{concretize}_{\varphi_n}(N)}{\bar{w}}
    \iff \bar{v} \approx \bar{w} \\
  \ev{M}{\rho_2}{\kappa_2} \sim {}& \ev{\pass{concretize}_\varphi(M)}{\rho_3}{\kappa_3}
    \iff \rho_2 \approx_\varphi \rho_3 \land \kappa_2 \approxk \kappa_3
\end{align*}
with the frames related through $\varphi$ (and $\ext{\varphi}{x}{x}$ under a
binder). The machines run in lock step, except that a variable read in one
step may become a projection read in three, which cannot fail
(\texttt{ref-run}).
\begin{theorem}[Correctness of Concretize, Simulation]\label{thm:m-cek-concretize}
For every closed $\Ml{Clos}$ term $M$,
\[ \Equiv(\Mach{\Ml{Clos}}(M), \Mach{\Ml{EnvClos}}(\pass{concretize}_{\emptyset}(M))) \]
(\texttt{concretize-correct}).
\end{theorem}

\subsubsection{CC: Delay}

\begin{align*}
  \closV{e\, y.\, N}{\bar{v}} \approx {}&
    \Tpack{\TcodeV{e\, y.\, \pass{delay}(N)}}{\tupV{\bar{w}}} \iff \bar{v} \approx \bar{w} \\
  \Kf{appL}{M, \rho_1, \kappa_1} \approxk {}& \Kf{callL}{\pass{delay}(M), \rho_2, \kappa_2} \qquad
  \Kf{appR}{v, \kappa_1} \approxk \Kf{callR}{w, \kappa_2} \\
  \Kf{closK}{e\, y.\, N, \kappa_1} \approxk {}& \Kf{packR}{\TcodeV{e\, y.\, \pass{delay}(N)}, \kappa_2}
\end{align*}
with the environments related through the renaming of the pass
(\texttt{EnvRel}) and the other frames the same frame
(\texttt{DelaySim.agda}). Neither machine checks the operator of a call,
so the machines run in lock step, except that building a closure takes
the target three more steps (the code, then the tuple under
$\key{packR}$): a plus simulation.
\begin{theorem}[Correctness of Delay, Simulation]\label{thm:m-cek-delay}
For every closed $\Ml{EnvClos}$ term $M$,
\[ \Equiv(\Mach{\Ml{EnvClos}}(M), \Mach{\Ml{FlatClos}}(\pass{delay}(M))) \]
(\texttt{delay-correct}).
\end{theorem}

\subsubsection{CC: Globalize}

The relations are those for the definitions
$\vec{d} = d_0, \ldots, d_{m-1}$, $d_k = \Tcodetm{e}{y}{N_k}$, that globalize
makes for the whole program (\texttt{GlobalizeSim.agda}). A $\Ml{FlatClos}$
term is related to a $\Ml{GlobalFun}$ term ($M \leadsto M'$) when they are the
same, except that code is related to a reference to a definition whose
body its body is related to; values likewise:
\[
  \Tcodetm{e}{y}{N} \leadsto \funref{k} \iff N \leadsto N_k \qquad
  \TcodeV{e\, y.\, N} \approx \codeV{k} \iff N \leadsto N_k ,
\]
and continuations and states are related frame by frame. The machines run
in lock step. The pass makes each term over the table made so far and
later weakens it, so the relation is taken relative to an interpretation
of the references of a smaller table ($I \vdash M \leadsto M'$,
\texttt{Interp}, \texttt{\_$\vdash$\_$\leadsto$\_}); the definitions made
before a term are those of $\vec{d}$ at the same positions
(\texttt{Faithful}, \texttt{glob-grows}, \texttt{MExt}).
\begin{theorem}[Correctness of Globalize, Simulation]\label{thm:m-cek-globalize}
For every closed $\Ml{FlatClos}$ term $M$,
\[ \Equiv(\Mach{\Ml{FlatClos}}(M), \Mach{\Ml{GlobalFun}}(\pass{globalize}(M))) \]
(\texttt{globalize-correct}).
\end{theorem}

\subsubsection{Remove Complex Operands}

Both machines are the $\Ml{GlobalFun}$ machine, the source on $P$ and the target
on $\Tembed{\pass{rco}(P)}$ (\texttt{RCOSim.agda}); the two runs compute the
same values. The target environment holds the values of the $\key{let}$s
as well: a source term $M$ in $\rho_1$ is related to $\pass{rco}(M)$ in
$\rho_2$ when $\rho_2$ holds at $r(x)$ the value of $x$ for the pass's
renaming $r$ (\texttt{ER}; $\rho_1 \preceq \rho_2$ by name). For an operand
$N$ of a construct, the rest of the translation is a function $Q$ of the
atom that stands for $N$ (the Kripke continuation of $\pass{atomize}$). A
source continuation $\kappa_1$ continues as $Q$ given the value of the
atom, $\mathit{KRel}(\kappa_1, Q, \rho_2, \kappa_2)$, when for every
$\rho_2' \succeq \rho_2$ and every atom $a$ over $\rho_2'$,
\[ \mathit{Both}_\sim(\rt{\rho_2'(a)}{\kappa_1},\; \ev{\Tembed{Q(a)}}{\rho_2'}{\kappa_2})\,. \]
\begin{align*}
  \halt &\approxk \halt \qquad
  \Kf{letK}{x.N, \rho_1, \kappa_1} \approxk \Kf{letK}{x.\Tembed{\pass{rco}(N)}, \rho_2, \kappa_2}
    \iff \rho_1 \preceq \rho_2 \land \kappa_1 \approxk \kappa_2 \\
  \kappa_1 &\approxk \Kf{letK}{z.\Tembed{Q(z)}, \rho_2, \kappa_2}
    \iff \mathit{KRel}(\kappa_1, Q, \rho_2, \kappa_2) \\
  \ev{M}{\rho_1}{\kappa_1} &\sim \ev{\Tembed{\pass{rco}(M)}}{\rho_2}{\kappa_2}
    \iff \rho_1 \preceq \rho_2 \land \kappa_1 \approxk \kappa_2 \\
  \rt{v}{\kappa_1} &\sim \rt{v}{\kappa_2} \iff \kappa_1 \approxk \kappa_2
\end{align*}
The third clause (\texttt{pend}) relates a source frame waiting for the
value of a complex operand to the target's $\key{let}$ that binds it;
$\mathit{KRel}$ holds of each kind of frame ($\key{getK}$, $\key{inlK}$,
$\key{inrK}$, $\key{foldK}$, $\key{unfoldK}$, $\key{caseK}$, the two
operands of $\key{pack}$, $\key{callc}$ and primitives, and $\key{tupK}$).
An atom evaluates to its value in one step and cannot fail
(\texttt{Clos5Common.agda}). The machines move by $\mathit{Both}$.
\begin{theorem}[Correctness of RCO, Simulation]\label{thm:m-cek-rco}
For every $\Ml{GlobalFun}$ program $P$,
\[ \Equiv(\Mach{\Ml{GlobalFun}}(P), \Mach{\Ml{Mon}}(\pass{rco}(P))) \]
(\texttt{rco-correct}).
\end{theorem}

\subsubsection{Lift Locals}

The source is the $\Ml{GlobalFun}$ machine on $\Tembed{P}$ and the target the
$\Ml{Imp}$ machine on $\pass{lift}(P)$ (\texttt{LiftLocalsSim.agda}); the
simulation goes from $\Ml{Imp}$ to $\Ml{Mon}$, and the theorem follows by
the symmetry of $\Equiv$. The pass is structural, and the frame of the
$\Ml{Imp}$ machine is the environment of the $\Ml{GlobalFun}$ machine, so the two
run on the same values and frames:
\begin{align*}
  \Kf{letK}{x.\Tembed{N}, \sigma, \kappa_1} &\approxk \Kf{assignK}{\pass{lift}(N), \sigma, \kappa_2}
    \iff \kappa_1 \approxk \kappa_2 \\
  \ev{\Tembed{N}}{\sigma}{\kappa_1} &\sim \evi{\pass{lift}(N)}{\sigma}{\kappa_2}
    \iff \kappa_1 \approxk \kappa_2
\end{align*}
One $\Ml{Imp}$ step runs a simple expression on its atoms, which the
$\Ml{GlobalFun}$ machine evaluates one by one, through continuation frames
(\texttt{atm-step}); then both continue with related continuations, or
both call, or both stop with the same observation ($\mathit{Both}$).
\begin{theorem}[Correctness of Lift Locals, Simulation]\label{thm:m-cek-lift}
For every $\Ml{Mon}$ program $P$,
\[ \Equiv(\Mach{\Ml{Mon}}(P), \Mach{\Ml{Imp}}(\pass{lift}(P))) \]
(\texttt{lift-locals-correct}).
\end{theorem}

\subsubsection{Explicate Control}

The machines share the values, the frames and $\mathit{run}$, and an
expression runs in the same frame in both (\texttt{ExplicateSim.agda}).
The pass only adds blocks, $\vec{t} \preceq \vec{t}'$ (\texttt{Ext}), and the
older labels find the same blocks (\texttt{ext-block}). A $\Ml{CFG}$
continuation continues as an $\Ml{Imp}$ one given the value of a body,
$\kappa_1 \approx \kappa_2$; and
$\mathit{Ctl}_{\vec{t}}(\sigma, \kappa_1, k, \kappa_2)$ when $\kappa_1$,
given a value $v$, continues as the tail $k$ in the frame $(\sigma, v)$
with the blocks $\vec{t}$, under $\kappa_2$:
\begin{align*}
  \halt &\approx \halt \qquad
  \kappa_1 \approx \Kf{retK}{\kappa_2} \Leftarrow \kappa_1 \approx \kappa_2 \\
  \kappa_1 &\approx \Kf{seqK}{s, \sigma, k, \vec{t}, \kappa_2} \Leftarrow
    \mathit{Ctl}_{\vec{t}}(\Tcut{\sigma}{s}, \kappa_1, k, \kappa_2) \\[1ex]
  \mathit{Ctl}_{\vec{t}'}(\sigma, \Kf{assignK}{e_2, \sigma, \kappa}, t, \kappa_2) &\Leftarrow
    \kappa \approx \kappa_2 \land \pass{tail}(e_2, \vec{t}) = (t;\, \vec{t}') \\
  \mathit{Ctl}_{\vec{t}'}(\sigma, \Kf{assignK}{e_2, \sigma, \kappa}, t, \kappa_2) &\Leftarrow
    \mathit{Ctl}_{\vec{t}}(\Tcut{\sigma}{s}, \kappa, k, \kappa_2) \land
    \pass{assign}(e_2, s, k, \vec{t}) = (t;\, \vec{t}') \\
  \mathit{Ctl}_{\vec{t}'}(\sigma, \kappa_1, \Tgoto{l}, \kappa_2) &\Leftarrow
    \mathit{Ctl}_{\vec{t}}(\sigma, \kappa_1, k, \kappa_2) \land \key{block}(\vec{t}', l) = (k;\, \vec{t})
\end{align*}
The states are indexed by the measure of the stutter,
$\mathit{depth}(\assign{s}{e_1}{e_2}) = 1 + \mathit{depth}(e_1)$ and
$\mathit{depth}(e) = 0$ otherwise:
\begin{align*}
  \evi{e}{\sigma}{\kappa_1} \sim_{\mathit{depth}(e)} {}& \exc{t}{\sigma}{\vec{t}'}{\kappa_2}
    \Leftarrow \kappa_1 \approx \kappa_2 \land \pass{tail}(e, \vec{t}) = (t;\, \vec{t}') \\
  \evi{e}{\sigma}{\kappa_1} \sim_{\mathit{depth}(e)} {}& \exc{t}{\sigma}{\vec{t}'}{\kappa_2}
    \Leftarrow \mathit{Ctl}_{\vec{t}}(\Tcut{\sigma}{s}, \kappa_1, k, \kappa_2) \land
      \pass{assign}(e, s, k, \vec{t}) = (t;\, \vec{t}')
\end{align*}
An $\Ml{Imp}$ step from $\assign{s}{e_1}{e_2}$ is matched by no $\Ml{CFG}$
step, with a smaller index; any other by one or more $\Ml{CFG}$ steps: a
statement, then the jumps to join points.
\begin{theorem}[Correctness of Explicate Control, Simulation]\label{thm:m-cek-explicate}
For every $\Ml{Imp}$ program $P$,
\[ \Equiv(\Mach{\Ml{Imp}}(P), \Mach{\Ml{CFG}}(\pass{explicate}(P))) \]
(\texttt{explicate-correct}).
\end{theorem}

\subsubsection{Select Instructions}

Both machines hold the same values, and the $\Ml{X86}$ frame is the
$\Ml{CFG}$ frame (\texttt{SelectSim.agda}). So a state of $\Ml{CFG}$ is
matched by the code selected for its tail, in the same frame, with any
registers, open closure and flag ($\mathit{Match}$); and the $\Ml{X86}$ call
stack holds, for each return frame of the $\Ml{CFG}$ continuation, the
caller's frame and the move of $\rax$ to where the caller wants the
result:
\begin{align*}
  \halt &\sim \halt \qquad
  \Kf{seqK}{s, \sigma, t, \vec{t}, \kappa} \sim
    \key{retK}((\key{movq}\; \rax, s;\; \pass{tail}(t)), \pass{tail}(\vec{t}), \sigma, \Pi)
    \Leftarrow \kappa \sim \Pi \\
  \Kf{retK}{\kappa} &\sim \key{retK}((\key{movq}\; \rax, \rax;\; \key{retq}), \vec{\beta}, \sigma, \Pi)
    \Leftarrow \kappa \sim \Pi \\
  \exc{t}{\sigma}{\vec{t}}{\kappa} &\sim
    \xrun\;\pass{tail}(t)\;\pass{tail}(\vec{t})\;\Pi\;\eta\;\langle R, \sigma, O, f \rangle
    \Leftarrow \kappa \sim \Pi
\end{align*}
One $\Ml{CFG}$ step (a statement) is matched by the instructions selected
for it: they write its value to its destination (\texttt{xval}), or call
(\texttt{xcall}): the four steps of $\key{load}\;\key{fst}$,
$\key{load}\;\key{snd}$, the move to $\rsi$ and $\key{callq}$, then the two
of the callee's prologue. A $\key{case}$ is matched by the $\key{case}$
pseudo-instruction, one step to the same block (\texttt{block-sel}).
Neither machine fails, so there is no failing case to match.
\begin{theorem}[Correctness of Select Instructions, Simulation]\label{thm:m-cek-select}
For every $\Ml{CFG}$ program $P$,
\[ \Equiv(\Mach{\Ml{CFG}}(P), \Mach{\Ml{X86}}(\pass{select}(P))) \]
(\texttt{select-correct}).
\end{theorem}

\subsubsection{End to End}

\begin{theorem}[Compiler Correctness, Simulation]\label{thm:m-cek-compile}
For every closed $\Ml{ISWIM}$ program $M$ of any type,
\begin{align*}
  &\Equiv(\Mach{\Ml{ISWIM}}(M), \Mach{\Ml{FlatClos}}(\pass{compile}_4(M))) &
  &\Equiv(\Mach{\Ml{ISWIM}}(M), \Mach{\Ml{GlobalFun}}(\pass{compile}(M))) \\
  &\Equiv(\Mach{\Ml{ISWIM}}(M), \Mach{\Ml{Mon}}(\pass{compile\text{-}mon}(M))) &
  &\Equiv(\Mach{\Ml{ISWIM}}(M), \Mach{\Ml{Imp}}(\pass{compile\text{-}imp}(M))) \\
  &\Equiv(\Mach{\Ml{ISWIM}}(M), \Mach{\Ml{CFG}}(\pass{compile\text{-}cfg}(M))) &
  &\Equiv(\Mach{\Ml{ISWIM}}(M), \Mach{\Ml{X86}}(\pass{compile\text{-}x86}(M)))
\end{align*}
by the theorems for each pass and Lemma~\ref{lem:equiv-trans}
(\texttt{Theorems/\allowbreak Simulation.agda}: \texttt{compile-correct}, \ldots,
\texttt{compile-x86-correct}; for a main of type $\key{Int}$,
\texttt{compile-x86-const} and \texttt{compile-x86-diverges},
Section~\ref{sec:m-thm}; the statement for $\pass{compile}_4$ is a step of
the proof, \texttt{compile$_4$-correct} of \texttt{PipelineSim.agda}).
\end{theorem}

\paragraph{Where it doesn't match the Agda literally.}
\begin{itemize}
\item \textbf{Environments and frames} are typed lists (\texttt{All Val
  $\Gamma$}), newest first, looked up by de Bruijn index; here named. The
  frame $(u, w)$ of a definition is \texttt{w $\cdot$ u $\cdot$ []}, and
  $(\sigma, v)$ is \texttt{v $\cdot$ $\sigma$}. A $\Ml{Clos}$ closure keeps its
  values as \texttt{All Val $\Delta$}, its code over \texttt{A $\cdot$ $\Delta$}.
\item \textbf{Tuples}: $\key{tupK}$ and \texttt{tuple-next} collect the values
  newest first in \texttt{All Val $\Delta'$}, and the tuple's type is
  \texttt{tup ($\Delta'$ <>< $\Delta$)} (reverse-append, \texttt{\_<><\_} and
  \texttt{\_<><$^{\mathsf{a}}$\_} of \texttt{Framework/AllUtil.agda}).
\item \textbf{$\Ml{GlobalFun}$}: the step function is parametrized by the
  definitions (\texttt{step ds}); a code value is \texttt{code f} with $f$ a
  typed index into the table; what an $\key{eval}$ state does is
  \texttt{begin}, and a call is \texttt{call}.
\item \textbf{$\Ml{Imp}$ and $\Ml{CFG}$}: $\Kf{assignK}{e_2, \sigma, \kappa}$
  is \texttt{assignK e$_2$ $\varphi$ $\kappa$}; $\Kf{seqK}{s, \sigma, t,
  \vec{t}, \kappa}$ is \texttt{seqK p $\varphi$ t bs $\kappa$} and
  $\Tcut{\sigma}{s}$ is \texttt{forget p $\varphi$}; $\key{block}$ returns a
  \texttt{Blk}, run by \texttt{go}; $\mathit{run}$'s results are \texttt{Res}
  (\texttt{value}, \texttt{call}).
\item \textbf{$\Ml{X86}$}: a state is \texttt{run k bs $\sigma$ $\eta$ m} with
  the call stack $\sigma$ (here $\Pi$, \texttt{Stack}) and \texttt{m} a
  record \texttt{Mach} $\langle$\texttt{regs}, \texttt{frame},
  \texttt{env$\alpha$}, \texttt{flag}$\rangle$, whose fields' types follow the
  state type (\texttt{RV}, \texttt{MV}, \texttt{Opened}, \texttt{Flag}); $m(a)$
  is \texttt{arg a m}, $\Twipe{R}$ is \texttt{wipe$_r$}, the environment of a
  closure and its types are \texttt{env-of} and \texttt{size}, $\Tres{f}$ is
  \texttt{bit f}.
\item \textbf{The relations} are \texttt{\_$\approx$\_}, \texttt{\_$\approx$k\_}
  and \texttt{\_\char`\~\_} in each \texttt{*Sim.agda}; environments are related
  through the renaming or map that a pass carries (\texttt{Agree},
  \texttt{EnvRel}, \texttt{ER}), here by name. Simplify's are
  \texttt{Ok}, \texttt{Strong}, \texttt{Den}, \texttt{KDen}, with the source's
  run ahead \texttt{run-term}, \texttt{run-known}, \texttt{let-sim}.
  Explicate's two clauses of $\sim$ are \texttt{tail\char`\~} and
  \texttt{asg\char`\~}, those of $\mathit{Ctl}$ \texttt{asgT}, \texttt{asgA}
  and \texttt{jmp}. Select's $\kappa \sim \Pi$ is \texttt{\_\char`\~\_} and
  $s_1 \sim s_2$ is \texttt{Match}.
\end{itemize}
%===============================================================================
\section{Clocked Interpreters}\label{sec:m-interp}

The clocked interpreters (\texttt{Semantics/Interp/}) use the results,
clocked programs and end theorem of Part~\ref{part:untyped}
(\texttt{Framework/Clocked.agda}, \texttt{Semantics/Interp/Clock.agda}): a
clocked program $(\obs, \lambda k.\, \ldots)$, results $\tmo$, $\error$ and
$\Ret v$, $\bind$ (with $\Do x \leftarrow r;\; r'$ for $r \bind \lambda
x.\, r'$), $\Term$, $\Div$ and $\Equiv$. The clock decreases only when a
function is called; every other construct evaluates its subterms with the
same clock, and a call with clock $0$ times out. Values and environments
are typed, so each interpreter is a total function of typed terms that
\emph{never} errs: it only times out or returns a value.

\subsection{Values}\label{sec:m-interp-values}

The values (\texttt{Semantics/Interp/Values.agda}) are shared by all the
interpreters, indexed by their type, and parameterized by what the code of
a closure with an environment of types $\Delta$, an argument of type $A$
and a result of type $C$ is in each language, $C_{\Delta,A,C}$. A closure
is always code packed with the tuple of its environment, even in the
languages before $\Ml{FlatClos}$, where code is not a value of the language.
With $c$ a constant of base type $B$ and $\gamma \in C_{\Delta,A,C}$:
\[
\begin{array}{r@{\;:\;}l@{\qquad}r@{\;:\;}l}
  \TcodeV{\gamma} & \Mcodety{\Delta}{A}{C} &
  \Tpack{v}{u} & A \rightarrow C \quad (v : \Mcodety{\Delta}{A}{C},\ u : \Mtup{\Delta}) \\
  c & B &
  \tupV{v_1, \ldots, v_n} & \Mtup{(A_1, \ldots, A_n)} \quad (v_i : A_i) \\
  \inlV v,\ \inrV v & A + C \quad (v : A, \text{ resp. } C) &
  \Mfold v & \Mmu{X}{A} \quad (v : \Munroll{X}{A})
\end{array}
\]
\[
\begin{array}{l|ccc}
  & \Ml{ISWIM},\ \Ml{Clos} & \Ml{EnvClos},\ \Ml{FlatClos} & \Ml{GlobalFun},\ \Ml{Imp},\ \Ml{CFG},\ \Ml{X86} \\ \hline
  C_{\Delta,A,C} & \Delta, y{:}A \vdash N : C & e{:}\Mtup{\Delta}, y{:}A \vdash N : C &
    k \text{ with } s_k = (\Delta; A \to C)
\end{array}
\]
written $\lam{y} N$ ($\Ml{ISWIM}$), $\bar{x}\, y.\, N$ ($\Ml{Clos}$), $e\,
y.\, N$ and $\codeV{k}$. $\obs$ is the machines' ($\obs(\Mfold v) =
\obs(v)$, code and closures observed as $\key{fun}$). The operations on
values are total on values of their types:
\begin{align*}
  \key{nth}_i(\tupV{v_1, \ldots, v_n}) &= v_i \qquad
  \delta(p, n, m) = \Tres{\delta(p, n, m)} \qquad
  \key{unfold}_v(\Mfold v) = v \\
  \key{either}(\inlV v, f, g) &= f(v) \qquad \key{either}(\inrV v, f, g) = g(v)
\end{align*}

\subsection{The interpreters}

For $\ell$ one of $\Ml{ISWIM}$ to $\Ml{FlatClos}$, and for $\Ml{GlobalFun}$ with
the definitions $\vec{d} = d_0, \ldots, d_{m-1}$ (written $\EV{\vec{d}}{k}$,
$\CALL{\vec{d}}{k}$; \texttt{Semantics/Interp/\{ISWIM,Clos2,Clos3,Clos4,Clos5\}.agda}):
\begin{align*}
  \Interp{\ell}(M) &= (\obs,\; \lambda k.\, \EV{\ell}{k}\,\emptyset\,M) \\
  \Interp{\Ml{GlobalFun}}(\vec{d};\; M) &= (\obs,\; \lambda k.\, \EV{\vec{d}}{k}\,\emptyset\,M) \\[1ex]
  \EV{\ell}{k}\,\rho\,x &= \Ret \rho(x) \qquad
  \EV{\ell}{k}\,\rho\,c = \Ret c \\
  \EV{\ell}{k}\,\rho\,\tuple{M_1, \ldots, M_n} &=
    \Do \bar{v} \leftarrow \EVS{\ell}{k}\,\rho\,(M_1 \ldots M_n);\; \Ret \tupV{\bar{v}} \\
  \EV{\ell}{k}\,\rho\,(\prj{i} M) &= \Do v \leftarrow \EV{\ell}{k}\,\rho\,M;\; \Ret \key{nth}_i(v) \\
  \EV{\ell}{k}\,\rho\,(\inl M) &= \Do v \leftarrow \EV{\ell}{k}\,\rho\,M;\; \Ret (\inlV v)
    \qquad \text{and so for } \key{inr} \\
  \EV{\ell}{k}\,\rho\,(\Case{L}{x}{M}{y}{N}) &=
    \Do u \leftarrow \EV{\ell}{k}\,\rho\,L;\; \key{either}(u,\; \lambda v.\, \EV{\ell}{k}\,\ext{\rho}{x}{v}\,M,\;
    \lambda v.\, \EV{\ell}{k}\,\ext{\rho}{y}{v}\,N) \\
  \EV{\ell}{k}\,\rho\,\prim{p}{L}{M} &=
    \Do v \leftarrow \EV{\ell}{k}\,\rho\,L;\; w \leftarrow \EV{\ell}{k}\,\rho\,M;\; \Ret \delta(p, v, w) \\
  \EV{\ell}{k}\,\rho\,(\Let x = M \In N) &=
    \Do v \leftarrow \EV{\ell}{k}\,\rho\,M;\; \EV{\ell}{k}\,\ext{\rho}{x}{v}\,N \\
  \EV{\ell}{k}\,\rho\,(\Mfold M) &= \Do v \leftarrow \EV{\ell}{k}\,\rho\,M;\; \Ret \Mfold v \\
  \EV{\ell}{k}\,\rho\,(\Munfold M) &= \Do v \leftarrow \EV{\ell}{k}\,\rho\,M;\; \Ret \key{unfold}_v(v) \\[1ex]
  \EVS{\ell}{k}\,\rho\,\epsilon &= \Ret \epsilon \\
  \EVS{\ell}{k}\,\rho\,(M\,\bar{M}) &=
    \Do v \leftarrow \EV{\ell}{k}\,\rho\,M;\; \bar{v} \leftarrow \EVS{\ell}{k}\,\rho\,\bar{M};\; \Ret (v\,\bar{v}) \\
  \CALL{\ell}{0}\,f\,w &= \tmo
\end{align*}
An application evaluates its operator and its operand and then calls,
\[ \Do f \leftarrow \EV{\ell}{k}\,\rho\,L;\; w \leftarrow \EV{\ell}{k}\,\rho\,M;\; \CALL{\ell}{k}\,f\,w\,, \]
for $L\;M$ in $\Ml{ISWIM}$ to $\Ml{EnvClos}$ and for $\Tcallc{L}{M}$ in
$\Ml{FlatClos}$ and $\Ml{GlobalFun}$. Closures and calls:
\begin{align*}
  \EV{\Ml{ISWIM}}{k}\,\rho\,(\lam{y} N) &= \Ret \Tpack{\TcodeV{\lam{y} N}}{\tupV{\rho}} \\
  \CALL{\Ml{ISWIM}}{k+1}\,\Tpack{\TcodeV{\lam{y} N}}{\tupV{\rho}}\,w &= \EV{\Ml{ISWIM}}{k}\,\ext{\rho}{y}{w}\,N \\
  \EV{\Ml{Clos}}{k}\,\rho\,(\clos{M_1 \ldots M_n}{\bar{x}\, y} N) &=
    \Do \bar{v} \leftarrow \EVS{\Ml{Clos}}{k}\,\rho\,(M_1 \ldots M_n); \\
  &\qquad \Ret \Tpack{\TcodeV{\bar{x}\, y.\, N}}{\tupV{\bar{v}}} \\
  \CALL{\Ml{Clos}}{k+1}\,\Tpack{\TcodeV{\bar{x}\, y.\, N}}{\tupV{\bar{v}}}\,w &=
    \EV{\Ml{Clos}}{k}\,\{\bar{x} {\mapsto} \bar{v}, y {\mapsto} w\}\,N \\
  \EV{\Ml{EnvClos}}{k}\,\rho\,(\clos{M_1 \ldots M_n}{e\, y} N) &=
    \Do \bar{v} \leftarrow \EVS{\Ml{EnvClos}}{k}\,\rho\,(M_1 \ldots M_n);\; \Ret \Tpack{\TcodeV{e\, y.\, N}}{\tupV{\bar{v}}} \\
  \CALL{\ell}{k+1}\,\Tpack{\TcodeV{e\, y.\, N}}{u}\,w &=
    \EV{\ell}{k}\,\{e {\mapsto} u, y {\mapsto} w\}\,N \qquad (\ell \in \{\Ml{EnvClos}, \Ml{FlatClos}\}) \\
  \EV{\Ml{FlatClos}}{k}\,\rho\,\Tcodetm{e}{y}{N} &= \Ret \TcodeV{e\, y.\, N} \\
  \EV{\ell}{k}\,\rho\,\Tpack{L}{M} &=
    \Do v \leftarrow \EV{\ell}{k}\,\rho\,L;\; u \leftarrow \EV{\ell}{k}\,\rho\,M;\; \Ret \Tpack{v}{u} \\
  \EV{\vec{d}}{k}\,\rho\,\funref{j} &= \Ret \codeV{j} \\
  \CALL{\vec{d}}{k+1}\,\Tpack{\codeV{j}}{u}\,w &=
    \EV{\vec{d}}{k}\,\{e {\mapsto} u, y {\mapsto} w\}\,N
    \quad (d_j = \Tcodetm{e}{y}{N})
\end{align*}
$\Ml{Mon}$ runs as its embedding, $\Interp{\Ml{Mon}}(P) =
\Interp{\Ml{GlobalFun}}(\Tembed{P})$.

\paragraph{$\Ml{Imp}$ and $\Ml{CFG}$} (\texttt{Semantics/Interp/\{Imp,CFG\}.agda}).
An expression runs in a frame $\sigma$ and returns a value; an atom is read
from the frame and cannot fail. The value of a simple expression, given
how to call a closure ($C$), is shared: $\key{simple}_C(\sigma, \Tcallc{a}{b}) =
C(\sigma(a), \sigma(b))$, $\key{simple}_C(\sigma, \Mfold a) = \Ret \Mfold
\sigma(a)$, $\key{simple}_C(\sigma, \Munfold a) = \Ret
\key{unfold}_v(\sigma(a))$, and the others the operations above on the
atoms' values. For definitions $\vec{e}$ and bodies $\vec{b}$:
\begin{align*}
  \EV{\vec{e}}{k}\,\sigma\,r &= \key{simple}_{\CALL{\vec{e}}{k}}(\sigma, r) \\
  \EV{\vec{e}}{k}\,\sigma\,(\assign{s}{e_1}{e_2}) &=
    \Do v \leftarrow \EV{\vec{e}}{k}\,\sigma\,e_1;\; \EV{\vec{e}}{k}\,(\sigma, v)\,e_2 \\
  \EV{\vec{e}}{k}\,\sigma\,(\Casei{a}{s}{e_1}{e_2}) &=
    \key{either}(\sigma(a),\; \lambda v.\, \EV{\vec{e}}{k}\,(\sigma, v)\,e_1,\;
      \lambda v.\, \EV{\vec{e}}{k}\,(\sigma, v)\,e_2) \\
  \CALL{\vec{e}}{k+1}\,\Tpack{\codeV{j}}{u}\,w &= \EV{\vec{e}}{k}\,(u, w)\,e_j \\[1ex]
  \key{tail}^k_{\vec{b}}\,\vec{t}\,\sigma\,(\Treturn{r}) &= \key{simple}_{\CALL{\vec{b}}{k}}(\sigma, r) \\
  \key{tail}^k_{\vec{b}}\,\vec{t}\,\sigma\,(\assign{s}{r}{t}) &=
    \Do v \leftarrow \key{simple}_{\CALL{\vec{b}}{k}}(\sigma, r);\;
    \key{tail}^k_{\vec{b}}\,\vec{t}\,(\Tcut{\sigma}{s}, v)\,t \\
  \key{tail}^k_{\vec{b}}\,\vec{t}\,\sigma\,(\Tgoto{l}) &= \key{jump}^k_{\vec{b}}\,\vec{t}\,l\,\sigma \\
  \key{tail}^k_{\vec{b}}\,\vec{t}\,\sigma\,(\Tcase{a}{s}{l_1}{l_2}) &=
    \key{either}(\sigma(a),\; \lambda v.\, \key{jump}^k_{\vec{b}}\,\vec{t}\,l_1\,(\Tcut{\sigma}{s}, v),\;
      \lambda v.\, \key{jump}^k_{\vec{b}}\,\vec{t}\,l_2\,(\Tcut{\sigma}{s}, v)) \\
  \key{jump}^k_{\vec{b}}\,(t_0, \ldots, t_{q-1})\,l\,\sigma &= \key{tail}^k_{\vec{b}}\,(t_0, \ldots, t_{l-1})\,\sigma\,t_l \\
  \CALL{\vec{b}}{k+1}\,\Tpack{\codeV{j}}{u}\,w &= \key{tail}^k_{\vec{b}}\,\vec{t}\,(u, w)\,t
    \quad (b_j = (\vec{t};\; t))
\end{align*}
Only calls tick; a jump runs its block with the blocks older than it, so
the interpreter is structurally recursive on the clock, the blocks and the
tail.

\paragraph{$\Ml{X86}$} (\texttt{Semantics/Interp/X86.agda}). The interpreter
runs on the typed states of its machine (Section~\ref{sec:m-cek-x86}),
$m = \langle R, \sigma, O, f \rangle$ and $\eta$; the instructions of a
sequence run in order, each passing its state to the rest, then its jump,
and only calls tick. With $m(a)$, $\key{put}$, $\Twipe{R}$ and $\Tres{f}$
those of the machine, and writing $\key{run}^k\,\eta\,m\,\beta$ for
$\key{run}^k_{\vec{b}}\,\vec{\beta}\,\eta\,m\,\beta$:
\begin{align*}
  \key{run}^k\,\eta\,m\,(i;\, \beta) &= \key{run}^k\,\eta'\,m'\,\beta
    \qquad \text{if } \langle i;\, \beta\rangle \mapsto \langle \beta \rangle_{\eta', m'} \text{ in the machine},\ i \neq \key{callq} \\
  \key{run}^k\,\eta\,m\,(\key{callq}\; a;\, \beta) &=
    \Do v \leftarrow \CALL{\vec{b}}{k}(m(a), R(\rdi), R(\rsi));\;
    \key{run}^k\,\eta\,\langle (\rax{:}v), \sigma, \cdot, \cdot \rangle\,\beta \\
  \key{run}^k\,\eta\,m\,\key{retq} &= \Ret R(\rax) \\
  \key{run}^k\,\eta\,m\,(\key{jmp}\; l) &= \key{jump}^k\,\vec{\beta}\,\eta\,l\,\sigma \\
  \key{run}^k\,\eta\,m\,(\Mxcase{x}{s}{l_1}{l_2}) &=
    \key{either}(\sigma(x),\; \lambda u.\, \key{jump}^k\,\vec{\beta}\,\eta\,l_1\,(\Tcut{\sigma}{s}, u),\;
      \lambda u.\, \key{jump}^k\,\vec{\beta}\,\eta\,l_2\,(\Tcut{\sigma}{s}, u)) \\
  \key{jump}^k\,(\beta_0, \ldots, \beta_{q-1})\,\eta\,l\,\sigma &=
    \key{run}^k_{\vec{b}}\,(\beta_0, \ldots, \beta_{l-1})\,\eta\,\langle \emptyset, \sigma, \cdot, \cdot \rangle\,\beta_l \\
  \CALL{\vec{b}}{k+1}(\codeV{j}, u, w) &= \key{run}^k_{\vec{b}}\,\vec{\beta}\,\eta'\,\langle (\rdi{:}u,\ \rsi{:}w), \emptyset, \cdot, \cdot \rangle\,\beta \\
    &\qquad (b_j = (n;\; \vec{\beta};\; \beta),\ u : \Mtup{\eta'}) \\
  \Interp{\Ml{X86}}(\vec{b};\; (n;\; \vec{\beta};\; \beta)) &=
    (\obs,\; \lambda k.\, \key{run}^k_{\vec{b}}\,\vec{\beta}\,\epsilon\,\langle \emptyset, \emptyset, \cdot, \cdot \rangle\,\beta)
\end{align*}

%-------------------------------------------------------------------------------
\subsection{Relations}\label{sec:m-interp-rel}

Each pass but simplify is proved clock for clock (Theorem~\ref{thm:clocked}):
for every clock, a program and its compilation both time out or return
related values, $\RelRes{\approx}$ (both err never happens: neither
interpreter errs). For the passes from $\Ml{GlobalFun}$ on the values are the
same and the results equal (\texttt{clocked-$\equiv$}). Each proof uses
Lemma~\ref{lem:bind-rel} at every $\bind$. The statements of the theorems
(of each pass, and of the compiler) are in
\texttt{Theorems/\allowbreak Interpreter.agda}, under the names given below, each
defined as its proof.

\paragraph{The value relation between interpreters.}
The passes from $\Ml{ISWIM}$ to $\Ml{GlobalFun}$ relate the values of two
interpreters by one relation (\texttt{proof/Semantics/Interp/Rel.agda}),
given a relation $\mathit{Code}(\gamma, \gamma')$ on codes and a relation
$\mathit{Clo}(\gamma, \bar{v}', \gamma', \bar{w})$ on closures:
\begin{align*}
  \TcodeV{\gamma} \approx \TcodeV{\gamma'} &\iff \mathit{Code}(\gamma, \gamma') \\
  \Tpack{\TcodeV{\gamma}}{\tupV{\bar{v}}} \approx \Tpack{\TcodeV{\gamma'}}{\tupV{\bar{w}}}
    &\iff \exists \bar{v}'.\; \bar{v} \approx \bar{v}' \land \mathit{Clo}(\gamma, \bar{v}', \gamma', \bar{w}) \\
  c \approx c, \qquad
  \tupV{\bar{v}} \approx \tupV{\bar{w}} &\iff \bar{v} \approx \bar{w} \\
  \inlV v \approx \inlV w &\iff v \approx w \ \ (\text{and so for } \key{inr}) \qquad
  \Mfold v \approx \Mfold w \iff v \approx w
\end{align*}
For most passes the closure relation keeps the environment,
$\Same_{\mathit{Code}}(\gamma, \bar{v}', \gamma', \bar{w}) \iff
\mathit{Code}(\gamma, \gamma') \land \bar{w} = \bar{v}'$. Related values are
observed alike, and $\delta$, $\key{nth}_i$, $\key{unfold}_v$ and
$\key{either}$ of related values give related results. Each main lemma is
proved by induction on the term, together with the same statement for
$\EVS{}{}$, and a lemma for calls, by the clock: at clock $0$ both time
out; at $k + 1$ both run the bodies of related closures with clock $k$.

\subsubsection{Simplify}

Simplify removes calls ($\beta$ of a small value, inlining), so its proof
is not clock for clock but clock monotone (Theorem~\ref{thm:clocked-le};
\texttt{SimplifyFuel.agda}), from the monotonicity of the interpreter in
the clock (\texttt{mono}) and two lemmas, with the relations of the CEK
proof (a target closure the simplified body of a source closure, under
facts true of the two environments, \texttt{Ok}, \texttt{KDen}):
\begin{itemize}
\item forward, at the same clock (\texttt{run$\approx$}): if
  $\rho \models_E \rho'$, the source's result, if it does not time out, is
  the target's, related ($\mathit{Le}_{\approx}$);
\item backward (\texttt{back}): if the target does not time out at a clock,
  then the source has a related result at every large enough clock.
\end{itemize}
When $\pass{simp}_E(M)$ returns a value term, or $M$ has a known value,
the source runs $M$ to a value it describes, at every large enough clock
(\texttt{run-term}, \texttt{run-known}; \texttt{run-unf} for the
$\Munfold$ rule).
\begin{theorem}[Correctness of Simplify, Interpreter]\label{thm:m-interp-simplify}
For every closed $M$,
\[ \Equiv(\Interp{\Ml{ISWIM}}(M), \Interp{\Ml{ISWIM}}(\pass{simplify}(M))) \]
(\texttt{simplify-correct}).
\end{theorem}

\subsubsection{CC: Enclose, Optimize, Concretize and Delay}

Each is clock for clock, with the value relation above
(\texttt{EncloseFuel.agda}, \texttt{OptimizeFuel.agda},
\texttt{ConcretizeFuel.agda}, \texttt{DelayFuel.agda}; the lemmas
\texttt{run$\approx$}, \texttt{runs$\approx$} and \texttt{call$\approx$}):
\begin{itemize}
\item \textbf{enclose}: $\Same_{\mathit{Code}}$ with
  \[ \mathit{Code}(\lam{y} N, \gamma') \iff \gamma' = \bar{x}\, y.\, \pass{enclose}(N)\,, \]
  $\bar{x}$ the variables of the environment; environments pointwise. A
  $\lambda$ captures every variable in scope, in order, so its closure
  evaluates the variables to $\rho_2$ itself (\texttt{evals-vars});
\item \textbf{optimize}: $\mathit{Code}$ is empty, and a closure may keep
  only some of its environment (\texttt{Opt}):
  \[
    \mathit{Opt}(\bar{x}\, y.\, N,\; \bar{v}',\; x_{i_1} \ldots x_{i_q}\, y.\,
    \pass{optimize}(N),\; v'_{i_1} \ldots v'_{i_q}) \iff \{ i \mid x_i \in
    \FV(N) \} \subseteq \{ i_1 < \cdots < i_q \} ;
  \]
  environments
  agree on the variables a term uses, through the renaming (\texttt{Agree}),
  and the kept values of a closure are related to the source's by
  $\mathit{Sel}$ (\texttt{Sel}, \texttt{sel$\approx$}): a dropped $M_i$ is a
  variable, which the target does not evaluate;
\item \textbf{concretize}: $\Same_{\mathit{Code}}$ with
  $\mathit{Code}(\bar{x}\, y.\, N,\; e\, y.\, N') \iff N' =
  \pass{concretize}_{\varphi_n}(N)$; environments through $\varphi$,
  $\rho_2 \approx_\varphi \rho_3 \iff \forall x.\; \rho_2(x) \approx
  \rho_3(\varphi(x))$ (\texttt{Env$\approx$}, \texttt{deref});
\item \textbf{delay}: $\Same_{\mathit{Code}}$ with
  \[ \mathit{Code}(e\, y.\, N,\; e\, y.\, N') \iff N' = \pass{delay}(N)\,; \]
  environments through the renaming.
\end{itemize}
\begin{theorem}[Correctness of Enclose, Optimize, Concretize and Delay, Interpreter]\label{thm:m-interp-cc}
For every closed term $M$ of the source language of each pass,
\[ \Equiv(\Interp{\Ml{ISWIM}}(M), \Interp{\Ml{Clos}}(\pass{enclose}(M)))\,, \]
\[ \Equiv(\Interp{\Ml{Clos}}(M), \Interp{\Ml{Clos}}(\pass{optimize}(M)))\,, \]
$\Equiv(\Interp{\Ml{Clos}}(M), \Interp{\Ml{EnvClos}}(\pass{concretize}_\emptyset(M)))$ and
\[ \Equiv(\Interp{\Ml{EnvClos}}(M), \Interp{\Ml{FlatClos}}(\pass{delay}(M))) \]
(\texttt{enclose-correct}, \ldots, \texttt{delay-correct}).
\end{theorem}

\subsubsection{CC: Globalize}

The $\Ml{GlobalFun}$ interpreter runs with the table $D = D_0, \ldots, D_{m-1}$
that the pass makes for the whole program (\texttt{GlobalizeFuel.agda}).
The pass only appends to the definitions it threads, so the definitions
$\vec{d}$ made by any point of the pass are, weakened, a prefix of $D$,
$\vec{d} \sqsubseteq D$ (\texttt{Contains}, \texttt{Ext}). Code is related
to a reference when the referenced definition is the globalized body:
$\mathit{Code}(e\, y.\, N,\; j) \iff \exists \vec{d}.\; \pass{glob}(N,
\vec{d}) = (N', \vec{d}') \land \vec{d}' \sqsubseteq D \land D_j = N'$ (with
$N'$ weakened), and the value relation is that with $\Same_{\mathit{Code}}$.
Globalize maps each $\key{callc}$ to one, so the interpreters call in lock
step.
\begin{theorem}[Correctness of Globalize, Interpreter]\label{thm:m-interp-globalize}
For every closed $\Ml{FlatClos}$ term $M$,
\[ \Equiv(\Interp{\Ml{FlatClos}}(M), \Interp{\Ml{GlobalFun}}(\pass{globalize}(M))) \]
(\texttt{globalize-correct}).
\end{theorem}

\subsubsection{Remove Complex Operands, Lift Locals and Explicate Control}

The values are the same, and the results equal for every clock:
\begin{itemize}
\item \textbf{rco} (\texttt{RCOFuel.agda}): the pass only adds $\key{let}$s,
  which do not tick, and an embedded atom evaluates without calling or
  failing. The target environment holds the source's, and the fresh
  variables (\texttt{Agree}). At a clock where calls agree, a
  continuation $\kappa$ of $\pass{atomize}$ that means $K$ under every
  extension of the environment (\texttt{KOk}) gives
  \[ \EV{\vec{d}}{k}\,\gamma\,M \bind K = \EV{\vec{d}'}{k}\,\varphi\,\Tembed{\pass{atomize}(M, \kappa)} \]
  (\texttt{atomize-ok}), and
  \[ \EV{\vec{d}}{k}\,\gamma\,M = \EV{\vec{d}'}{k}\,\varphi\,\Tembed{\pass{rco}(M)} \]
  (\texttt{rco-ok}); calls agree at every clock (\texttt{calls$\equiv$});
\item \textbf{lift locals} (\texttt{LiftLocalsFuel.agda}): the $\Ml{Imp}$
  frame of an expression is the $\Ml{GlobalFun}$ environment of its embedding,
  and the pass maps each call to one call:
  \[ \EV{\Tembed{\vec{d}}}{k}\,\sigma\,\Tembed{M} = \EV{\pass{lift}(\vec{d})}{k}\,\sigma\,\pass{lift}(M) \]
  (\texttt{lift-ok});
\item \textbf{explicate control} (\texttt{ExplicateFuel.agda}): the pass
  threads the blocks and only adds to them, $\vec{t} \sqsubseteq \vec{t}'$
  (\texttt{Ext}), and a jump to an older block is the same
  (\texttt{jump-at}). The output of the pass from the blocks $\vec{t}$
  means $K$ when $\mathit{Ok}_{\vec{t}}(K, (\vec{t}';\; t)) \iff \vec{t}
  \sqsubseteq \vec{t}' \land \forall \sigma.\; K(\sigma) =
  \key{tail}^k\,\vec{t}'\,\sigma\,t$ (fields \texttt{exts},
  \texttt{means}); $\pass{assign}(e, s, t')$ means $e$ followed by the
  assignment of $s$ and $t'$ (\texttt{assign-ok}), and $\pass{tail}(e)$
  means $e$ (\texttt{tail-ok}).
\end{itemize}
\begin{theorem}[Correctness of RCO, Lift Locals and Explicate Control, Interpreter]\label{thm:m-interp-back}
For every program $P$ of the source language of each pass,
\[ \Equiv(\Interp{\Ml{GlobalFun}}(P), \Interp{\Ml{Mon}}(\pass{rco}(P)))\,, \]
$\Equiv(\Interp{\Ml{Mon}}(P), \Interp{\Ml{Imp}}(\pass{lift}(P)))$ and
\[ \Equiv(\Interp{\Ml{Imp}}(P), \Interp{\Ml{CFG}}(\pass{explicate}(P)))\,, \]
with equal results for every clock (\texttt{rco-correct},
\texttt{lift-locals-correct}, \texttt{explicate-correct}, which state
$\Equiv$).
\end{theorem}

\subsubsection{Select Instructions}

The two interpreters hold the same values in the same frame, so there is
no value relation: for every clock the $\Ml{X86}$ program gives exactly the
result of the $\Ml{CFG}$ one (\texttt{SelectFuel.agda}). At a clock where
calls give the same results, $\mathit{Calls}_k \iff \forall v\, w.\;
\CALL{\vec{b}}{k}\,v\,w = \CALL{\pass{select}(\vec{b})}{k}(\mathit{code}(v),
\mathit{env}(v), w)$ (\texttt{Calls}): the code selected for a simple
expression gives its $\Ml{CFG}$ result bound to what $k$ gives with the
value in $o$, whatever the registers and the flag (\texttt{sel-ok}, every
case definitional but a call, $\key{sete}$, \texttt{equal-bit}, and the
operands of a tuple, \texttt{args-ok}); a tail and a jump give what they
give in $\Ml{CFG}$ (\texttt{tail-ok}, \texttt{jump-ok}); calls agree by
induction on the clock (\texttt{call-ok}).
\begin{theorem}[Correctness of Select Instructions, Interpreter]\label{thm:m-interp-select}
For every $\Ml{CFG}$ program $P$,
$\Equiv(\Interp{\Ml{CFG}}(P), \Interp{\Ml{X86}}(\pass{select}(P)))$, with
equal results for every clock (\texttt{select-correct}).
\end{theorem}

\subsubsection{End to End}

\begin{theorem}[Compiler Correctness, Interpreter]\label{thm:m-interp-compiler}
For every closed $\Ml{ISWIM}$ program $M$ of any type, $\Interp{\Ml{ISWIM}}(M)$
is $\Equiv$ to $\Interp{\Ml{FlatClos}}(\pass{compile}_4(M))$,
$\Interp{\Ml{GlobalFun}}(\pass{compile}(M))$, $\Interp{\Ml{Mon}}(\pass{compile\text{-}mon}(M))$,
$\Interp{\Ml{Imp}}(\pass{compile\text{-}imp}(M))$, $\Interp{\Ml{CFG}}(\pass{compile\text{-}cfg}(M))$
and $\Interp{\Ml{X86}}(\pass{compile\text{-}x86}(M))$, by the theorems for each
pass and Lemma~\ref{lem:equiv-trans-clk} (\texttt{Theorems/\allowbreak Interpreter.agda}:
\texttt{compile-correct}, \ldots, \texttt{compile-x86-correct}; for a
main of type $\key{Int}$, \texttt{compile-x86-const} and
\texttt{compile-x86-diverges}; the statement for $\pass{compile}_4$ is a
step of the proof, \texttt{compile$_4$-correct} of
\texttt{CompilerFuel.agda}).
\end{theorem}

\paragraph{Where it doesn't match the Agda literally.}
\begin{itemize}
\item \textbf{Variables and environments.} Environments are lists typed by
  the context (\texttt{All Val $\Gamma$}), newest first; the frame $(u, w)$
  is \texttt{w :: u :: []}. A code $\gamma$ is a term (\texttt{Tm (A $\cdot$
  $\Delta$) B} in $\Ml{ISWIM}$ and $\Ml{Clos}$, a term over \texttt{A $\cdot$
  tup $\Delta$ $\cdot$ []} in $\Ml{EnvClos}$ and $\Ml{FlatClos}$) or an index
  \texttt{$\Sigma \ni$ sig $\Delta$ A B} into the table (\texttt{Vals}).
\item \textbf{Renamings.} The proofs of optimize and delay relate
  environments through the renaming the pass carries; concretize's
  $\varphi$ maps indices to \texttt{Ref}s, and rco's environments agree
  through a renaming (\texttt{Agree}), here by name.
\item \textbf{$\key{unfold}_v$} is \texttt{unfold-v} (\texttt{unfold$_v$} in
  the machines), $\key{simple}$ is \texttt{eval-simple}, and the $\Ml{X86}$
  interpreter is \texttt{run ck bs $\eta$ m k}, with \texttt{jump} and
  \texttt{call}.
\end{itemize}
%===============================================================================
\section{Reduction Semantics}\label{sec:m-red}

The languages with binders, $\Ml{ISWIM}$ and $\Ml{Clos}$--$\Ml{GlobalFun}$, get a
reduction semantics by substitution (\texttt{Semantics/Reduction/}), as the
untyped ones do in Part~\ref{part:untyped} (Section~\ref{sec:red}): a state
is a closed term, and a step reduces a redex in an evaluation context,
substituting values for bound variables. The framework is the same
(\texttt{Framework/Reduction.agda}): a reduction says of a term whether it
is a value ($\key{val}$), is stuck ($\error$), or reduces ($\key{to}\;M$); it
is one function $\mathit{red}$ written with the congruence
$\key{step\text{-}in}(f, \mathit{red}(X), k)$, and a closed program $M$ runs
as the machine $\mathcal{R}^{\ell}(M)$ of Section~\ref{sec:red}, which stops
with $\obs(V)$ at a value $V$ and with $\error$ when stuck. The operands of
a tuple, and the captured terms of a closure, are reduced left to right by
one combinator for heterogeneous vectors (\texttt{red-::},
\texttt{Framework/ReductionAll.agda}).

What the types change:
\begin{itemize}
\item $\mathit{red}$ is defined on closed terms of every type $A$, and takes
  a term of type $A$ to one of type $A$, so every state is well typed by
  construction.
\item There is no case for a variable (a closed term has none), and no
  reachable case for an operand of the wrong shape: a closed value of a
  function type is a $\lambda$, closure or pack, one of a tuple type a tuple
  of values, one of a sum type an injection, one of type $\key{Int}$ an
  integer, and one of type $\Mmu{X}{A}$ a $\Mfold$. The redexes split their
  operand at a generic type, with an equation to the expected one, since a
  term cannot be split at a concrete type (the type $\Munroll{X}{A}$ of an
  $\Munfold$ does not unify with one), and so they have catch-all cases to
  $\error$, which are never reached (Section~\ref{sec:m-thm}). Nothing is
  stuck: $\Munfold(\Mfold V)$ reduces to $V$.
\item Observations see through folds: $\obs(\Mfold V) = \obs(V)$; otherwise
  $\obs$ is that of Part~\ref{part:untyped}, and $\error$ on a term that is
  not a value (never reached).
\end{itemize}

\paragraph{Renaming and substitution.}
Terms are intrinsically scoped and typed, by a context of types (in
$\Ml{GlobalFun}$ also by the table $\Sigma$). Renaming and substitution are one
traversal $\mathit{trav}_{\mathcal{K}}$, parameterised by a \emph{kit}
$\mathcal{K}$ (McBride's type-preserving renaming and substitution;
\texttt{Framework/Syntax/Kit.agda}; for $\Ml{ISWIM}$,
\texttt{Language/ISWIMSubst.agda}, shared with simplify). On paper we write
$N[x := V]$ and $N[x_1 := V_1, \ldots, x_n := V_n]$, with named variables.
Types are closed, so the kits need no type substitution.

\paragraph{Values} are a predicate $\mathit{Value}$ in the proof halves
(\texttt{proof/Semantics/Reduction/*.agda}), with the lemma that a value
reduces to $\key{val}$, and that renaming and substitution keep values.

\subsection{\texorpdfstring{$\Ml{ISWIM}$}{ISWIM-mu}}

\[
\begin{array}{rcl}
  V, W &::=& \lam{x} N \mid c \mid \tuple{V_1, \ldots, V_n} \mid \inl V \mid \inr V
    \mid \Mfold V \\
  E &::=& \square \mid E\;M \mid V\;E \mid \tuple{V_1, \ldots, V_{i-1}, E, M_{i+1}, \ldots, M_n}
    \mid \prj{i} E \mid \inl E \mid \inr E \\
  &\mid& \Case{E}{x}{M}{y}{N} \mid \prim{p}{E}{M} \mid \prim{p}{V}{E}
    \mid \Let x = E \In N \mid \Mfold E \mid \Munfold E
\end{array}
\]
\[
\begin{array}{@{}r@{\;\to\;}l@{\quad}l@{}}
  (\lam{x} N)\;V & N[x := V] \\
  \prj{i} \tuple{V_1, \ldots, V_n} & V_i \\
  \Case{\inl V}{x}{M}{y}{N} & M[x := V] \\
  \Case{\inr V}{x}{M}{y}{N} & N[y := V] \\
  \Let x = V \In N & N[x := V] \\
  \prim{\key{sub}}{n}{m} & n - m \\
  \prim{\key{eq}}{n}{m} & \inl \emptytup & n = m \\
  \prim{\key{eq}}{n}{m} & \inr \emptytup & n \neq m \\
  \Munfold (\Mfold V) & V
\end{array}
\]
(\texttt{Semantics/Reduction/ISWIM.agda}: $\beta$, \texttt{get?}, \texttt{case?},
\texttt{prim?}, \texttt{unfold?}).

\subsection{\texorpdfstring{$\Ml{Clos}$ to $\Ml{GlobalFun}$}{Clos-mu to GlobalFun-mu}}

As $\Ml{ISWIM}$, with closures in place of $\lambda$. A closure evaluates
its captured terms left to right, and is a value once they are values. Its
code is closed, so no substitution enters it.
\[
  (\clos{V_1\ldots V_n}{x_1\ldots x_n\, y} N)\;W \to N[x_1 := V_1, \ldots, x_n := V_n, y := W]
  \qquad (\Ml{Clos})
\]
\[
  (\clos{V_1\ldots V_n}{e\, y} N)\;W \to N[e := \tuple{V_1, \ldots, V_n}, y := W]
  \qquad (\Ml{EnvClos})
\]
In $\Ml{FlatClos}$ code is a value, and so is a pack of values; a call
evaluates the closure and the argument, left to right, and then calls the
code of the closure with its environment and the argument; in $\Ml{GlobalFun}$
a reference $\funref{k}$ is a value, and the reduction
$\mathit{red}_{\bar{d}}$ of a program with the definitions
$\bar{d} = d_0, \ldots, d_{m-1}$ calls the definition:
\begin{align*}
  \Tcallc{\Tpack{\Tcodetm{e}{y}{N}}{U}}{W} &\to N[e := U, y := W] \\
  \Tcallc{\Tpack{\funref{k}}{U}}{W} &\to N_k[e := U, y := W] \quad (d_k = \Tcodetm{e}{y}{N_k})
\end{align*}
with the contexts
\[
  E ::= \Tpack{E}{M} \mid \Tpack{V}{E} \mid \Tcallc{E}{M} \mid \Tcallc{V}{E} \mid \cdots
\]
(\texttt{Semantics/\allowbreak Reduction/\allowbreak \{Clos2,Clos3,Clos4,Clos5\}.agda}). The program
runs as $\mathcal{R}^{\Ml{GlobalFun}}(d_0, \ldots, d_{m-1};\; M)$.

\subsection{\texorpdfstring{$\Ml{Mon}$, $\Ml{Imp}$, $\Ml{CFG}$ and $\Ml{X86}$}{Mon-mu, Imp-mu, CFG-mu and X86-mu}}

$\Ml{Mon}$ runs as its embedding into $\Ml{GlobalFun}$. $\Ml{Imp}$, $\Ml{CFG}$ and
$\Ml{X86}$ have no binders, so they keep their machines, the CEK column's
(Section~\ref{sec:m-cek}), and the proofs of explicate control and select
instructions are the CEK column's.

%-------------------------------------------------------------------------------
\subsection{Relations}\label{sec:m-red-rel}

The proofs (\texttt{proof/Correctness/Reduction/*Red.agda}) are simulations
between reductions as in Part~\ref{part:untyped}
(Section~\ref{sec:red-sims}, \texttt{proof/Framework/Reduction.agda}): a
relation $\approx$ between closed source and target terms gives a
simulation when $\mathit{Ans}(M, \mathit{red}(M), M')$ whenever $M \approx
M'$, where $\mathit{Ans}$ answers a source value by a target run to a
related value, a stuck source by a target run to a stuck term, and a
source step by one or more target steps to a related term; with
$M_0 \approx M_0'$ for the start and equal observations of related values,
Theorem~\ref{thm:red-equivalent} gives $\Equiv$, and
Lemma~\ref{lem:ans-in} lifts the answers of a subterm through a context.
A relation is between terms of the same type. Each relation is a
congruence ($x \approx x$, $c \approx c$, each construct, $\Mfold$ and
$\Munfold$ included, related to the same construct with related subterms),
closed under substitution for the kits of both renaming and substitution
(Lemma~\ref{lem:red-subst}, \texttt{$\approx$-trav}); we show only the other
clauses. The statements of the theorems (of each pass, and of the
compiler) are in \texttt{Theorems/\allowbreak Reduction.agda}, under the names given
below, each defined as its proof.

\subsubsection{Simplify}

The target, simplified, takes fewer steps than the source, so the
simulation goes from the target to the source (\texttt{SimplifyRed.agda}),
and $\Equiv$ is symmetric. Simplify reduces under binders, so the source
catches up by steps $\to$ on open terms, in which a variable stands for a
value. Writing the source on the left, the relation $\succcurlyeq$ is a
congruence with, besides,
\begin{align*}
  M \succcurlyeq M' &\iff M \to M_1 \land M_1 \succcurlyeq M' \\
  (\lam{x} N)\;W \succcurlyeq (\lam{x} N')\;V' &\iff W \text{ is a value} \land W \succcurlyeq V' \land N[x := W] \succcurlyeq N'[x := V']
\end{align*}
and the like for a $\key{let}$ of a value that simplify kept bound and known,
whose bodies are related only after $\beta$ (the rules \texttt{step},
\texttt{let$\lambda$}, \texttt{let$\lambda$v}, \texttt{letv}).
\begin{lemma}[Fundamental Lemma, Simplify]\label{lem:m-red-simplify-fund}
If $\tau$ and $\upsilon$ map variables to values, and the facts $E$ are true
of $\tau$ after $\upsilon$ (for every $x$ with $E(x) = (T, K)$,
$\tau(x) \succcurlyeq T[\upsilon]$, and $K$ is true of $\tau(x)$: a known value
$K$ with $\tau(x) \succcurlyeq K[\upsilon]$, a known function $\langle \lam{y}
N, E_f \rangle$ reached by $\tau(x) \to^* (\lam{y} N)[\tau_f]$ for a $\tau_f$ of
which $E_f$ is true), then $M[\tau] \succcurlyeq \pass{simp}_E(M)[\upsilon]$,
and $\mathcal{K}_E(M)$ is true of $M[\tau]$ in the same way
(\texttt{simp-ok}, \texttt{Desc}, \texttt{KDesc}).
\end{lemma}
A known value is reached: if $M \succcurlyeq V'$ for a value $V'$, then $M
\to^* W$ for a value $W \succcurlyeq V'$ (\texttt{eval-sv}).
\begin{theorem}[Correctness of Simplify, Reduction]\label{thm:m-red-simplify}
For every closed $M$,
\[ \Equiv(\mathcal{R}^{\Ml{ISWIM}}(M), \mathcal{R}^{\Ml{ISWIM}}(\pass{simplify}(M))) \]
(\texttt{simplify-correct}).
\end{theorem}

\subsubsection{CC: Enclose}

Substitution goes into an $\Ml{ISWIM}$ function, but not into the code of a
$\Ml{Clos}$ closure, only into the terms it captures. So a function, after a
substitution $\sigma$ of the variables $x_1, \ldots, x_n$ in its scope, is
related to the closure of related code that captures terms related to
$\sigma$ (\texttt{EncloseRed.agda}):
\begin{align*}
  &\lam{y} N[x_1 := \sigma(x_1), \ldots, x_n := \sigma(x_n)] \approx
    \clos{F_1\ldots F_n}{x_1\ldots x_n\, y} N' \\
  &\qquad\iff N \approx N' \land \forall i.\; \sigma(x_i) \approx F_i \land F_i \text{ is a variable or a value} .
\end{align*}
The two reductions run in lock step.
\begin{theorem}[Correctness of Enclose, Reduction]\label{thm:m-red-enclose}
For every closed $M$,
\[ \Equiv(\mathcal{R}^{\Ml{ISWIM}}(M), \mathcal{R}^{\Ml{Clos}}(\pass{enclose}(M))) \]
(\texttt{enclose-correct}).
\end{theorem}

\subsubsection{CC: Optimize}

The relation $M \approx_\rho M'$ (\texttt{OptimizeRed.agda}) relates a
variable $x$ to $\rho(x)$, for a partial map $\rho$ from source to target
variables (a free variable that a closure drops has no variable in the
closure's code), and relates a closure to its optimization when its code is
related through the map $\rho_{\mathit{fl}}$ of the kept variables, and its
captured terms as the flags $\mathit{fl}$ keep them:
\begin{align*}
  \clos{M_1\ldots M_n}{x_1\ldots x_n\, y} N &\approx_\rho
    \clos{M'_{i_1}\ldots M'_{i_k}}{x_{i_1}\ldots x_{i_k}\, y} N' \\
  &\iff N \approx_{\rho_{\mathit{fl}}} N' \land \mathit{Caps}_\rho(\mathit{fl}, \bar{M}, \bar{M}') \\
  \mathit{Caps}_\rho(\mathit{fl}, \bar{M}, \bar{M}') &\iff
    \forall i.\; M_i \approx_\rho M_i' \text{ if } \mathit{fl}_i \text{ (kept)}, \\
  &\qquad\quad M_i \text{ is a variable or a value otherwise}
\end{align*}
In a closed term a dropped free variable is a value, which the source
evaluates in no step, so the two reductions run in lock step.
\begin{theorem}[Correctness of Optimize, Reduction]\label{thm:m-red-optimize}
For every closed $\Ml{Clos}$ term $M$,
\[ \Equiv(\mathcal{R}^{\Ml{Clos}}(M), \mathcal{R}^{\Ml{Clos}}(\pass{optimize}(M))) \]
(\texttt{optimize-correct}).
\end{theorem}

\subsubsection{CC: Concretize}

The relation $M \approx_\varphi M'$ (\texttt{ConcretizeRed.agda}) relates a
variable $x$ to where it lives, $\varphi(x)$; in both languages substitution
enters only a closure's captured terms. Applying a $\Ml{EnvClos}$ closure
substitutes the tuple of its captured values for $e$, so where the source
has a value, the target has a projection from that tuple, which takes a
step:
\begin{align*}
  V &\approx_\varphi \prj{i} \tuple{V_1', \ldots, V_n'}
    \iff V \text{ is a value} \land V \approx_\varphi V_i' \land V_1', \ldots, V_n' \text{ are values} \\
  \clos{M_1\ldots M_n}{x_1\ldots x_n\, y} N &\approx_\varphi \clos{M_1'\ldots M_n'}{e\, y} N' \\
    &\iff N \approx_{\varphi_n} N' \land M_i \approx_\varphi M_i' \text{ for every } i
\end{align*}
\begin{theorem}[Correctness of Concretize, Reduction]\label{thm:m-red-concretize}
For every closed $\Ml{Clos}$ term $M$,
\[ \Equiv(\mathcal{R}^{\Ml{Clos}}(M), \mathcal{R}^{\Ml{EnvClos}}(\pass{concretize}_{\emptyset}(M))) \]
(\texttt{concretize-correct}).
\end{theorem}

\subsubsection{CC: Delay}

The relation (\texttt{DelayRed.agda}) relates a closure to the pack of its
code and the tuple of its captured terms, and an application to a call:
\begin{align*}
  \clos{M_1\ldots M_n}{e\, y} N &\approx \Tpack{\Tcodetm{e}{y}{N'}}{\tuple{M_1', \ldots, M_n'}}
    \iff N \approx N' \land M_i \approx M_i' \\
  L\;M &\approx \Tcallc{L'}{M'} \iff L \approx L' \land M \approx M'
\end{align*}
The two reductions run in lock step.
\begin{theorem}[Correctness of Delay, Reduction]\label{thm:m-red-delay}
For every closed $\Ml{EnvClos}$ term $M$,
\[ \Equiv(\mathcal{R}^{\Ml{EnvClos}}(M), \mathcal{R}^{\Ml{FlatClos}}(\pass{delay}(M))) \]
(\texttt{delay-correct}).
\end{theorem}

\subsubsection{CC: Globalize}

The $\Ml{GlobalFun}$ program runs with the definitions $\bar{d}$ that the pass
made for the whole program, and the relation is for $\bar{d}$
(\texttt{GlobalizeRed.agda}):
$\Tcodetm{e}{y}{N} \approx_{\bar{d}} \funref{k} \iff
d_k = \Tcodetm{e}{y}{N_k} \land N \approx_{\bar{d}} N_k$. A subterm is related
to its output, weakened to the whole table, when the definitions made
agree with $\bar{d}$ (\texttt{Agree}). The reductions run in lock step.
\begin{theorem}[Correctness of Globalize, Reduction]\label{thm:m-red-globalize}
For every closed $\Ml{FlatClos}$ term $M$,
\[ \Equiv(\mathcal{R}^{\Ml{FlatClos}}(M), \mathcal{R}^{\Ml{GlobalFun}}(\pass{globalize}(M))) \]
(\texttt{globalize-correct}).
\end{theorem}

\subsubsection{Remove Complex Operands}

The source runs with its definitions $\bar{d}$, the target, the embedding
of the $\Ml{Mon}$ program, with $\Tembed{\pass{rco}(\bar{d})}$
(\texttt{RCORed.agda}). Values contain no code, so source and target
compute the same values, and the relation is built from the code under a
substitution $\sigma$ of values. A complex operand $M$ in the context $h$ of
its construct is reduced by the target under a $\key{let}$; once it is a
value $W$, the target substitutes it, and is then related to $h[W]$:
\begin{align*}
  V &\approx V \\
  h[M] &\approx \Let x = M' \In N
    \iff M \approx M' \land \forall W \text{ value}.\; h[W] \approx N[x := W] \\
  \Case{V}{x}{M[\sigma]}{y}{N[\sigma]} &\approx
    \Case{V}{x}{\Tembed{\pass{rco}(M)}[\sigma]}{y}{\Tembed{\pass{rco}(N)}[\sigma]} \\
  \Let x = M \In N[\sigma] &\approx \Let x = M' \In \Tembed{\pass{rco}(N)}[\sigma]
    \iff M \approx M'
\end{align*}
and a redex whose operands are values to itself (\texttt{letc$\approx$}). The
fundamental lemma: $M[\sigma] \approx \Tembed{\pass{rco}(M)}[\sigma]$
(\texttt{rco$\approx$}, with \texttt{atomize$\approx$} for the continuations of
$\pass{atomize}$, related for every substitution that agrees through the
renaming).
\begin{theorem}[Correctness of RCO, Reduction]\label{thm:m-red-rco}
For every $\Ml{GlobalFun}$ program $P$,
\[ \Equiv(\mathcal{R}^{\Ml{GlobalFun}}(P), \mathcal{R}^{\Ml{GlobalFun}}(\Tembed{\pass{rco}(P)})) \]
(\texttt{rco-correct}).
\end{theorem}

\subsubsection{Lift Locals}

The source, the $\Ml{GlobalFun}$ reduction of $\Tembed{P}$, substitutes values
for the variables; the target, the $\Ml{Imp}$ machine on $\pass{lift}(P)$,
keeps them in its frame (\texttt{LiftLocalsRed.agda}). A machine state
decodes to a term $D(s)$: its expression, with the values of its frame
substituted for the slots, in the stack of $\key{let}$s of its continuation
($\mathit{plug}$); a simple expression that computes a value (a component,
a primitive, an $\Munfold$) does so in the step that hands the value to the
continuation, so a state whose expression is one decodes to its value,
which the source reaches in one step. The source is related to a state $s$
when it is $D(s)$ or one step before it (\texttt{Lag}, the lagging
simulation, \texttt{lag-fsim}). Unlike the dynamically typed pipeline's,
the continuations of the $\Ml{Imp}$ machine do not say the type of the
result, so a state decodes to a pair of a type and a term of that type, and
the source runs on such pairs (\texttt{red$\Sigma$}, equivalent by
\texttt{pairs}, \texttt{Framework/ReductionAll.agda}).
\begin{theorem}[Correctness of Lift Locals, Reduction]\label{thm:m-red-lift}
For every $\Ml{Mon}$ program $P$,
\[ \Equiv(\mathcal{R}^{\Ml{GlobalFun}}(\Tembed{P}), \Mach{\Ml{Imp}}(\pass{lift}(P))) \]
(\texttt{lift-locals-correct}).
\end{theorem}

\subsubsection{Explicate Control, Select Instructions, End to End}

Explicate control and select instructions are between machines, by the
CEK column's proofs (Theorems~\ref{thm:m-cek-explicate}
and~\ref{thm:m-cek-select}; restated in \texttt{Theorems/\allowbreak Reduction.agda}
as \texttt{explicate-correct} and \texttt{select-correct}).
\begin{theorem}[Compiler Correctness, Reduction]\label{thm:m-red-compiler}
For every closed $\Ml{ISWIM}$ program $M$ of any type,
$\mathcal{R}^{\Ml{ISWIM}}(M)$ is $\Equiv$ to
$\mathcal{R}^{\Ml{FlatClos}}(\pass{compile}_4(M))$,
$\mathcal{R}^{\Ml{GlobalFun}}(\pass{compile}(M))$,
\[ \mathcal{R}^{\Ml{GlobalFun}}(\Tembed{\pass{compile\text{-}mon}(M)})\,, \]
$\Mach{\Ml{Imp}}(\pass{compile\text{-}imp}(M))$,
$\Mach{\Ml{CFG}}(\pass{compile\text{-}cfg}(M))$ and
$\Mach{\Ml{X86}}(\pass{compile\text{-}x86}(M))$, by the theorems for each pass
and Lemma~\ref{lem:equiv-trans} (\texttt{Theorems/\allowbreak Reduction.agda}:
\texttt{compile-correct}, \ldots, \texttt{compile-x86-correct}; for a
main of type $\key{Int}$, \texttt{compile-x86-const} and
\texttt{compile-x86-diverges}; the statement for $\pass{compile}_4$ is a
step of the proof, \texttt{compile$_4$-correct} of
\texttt{PipelineRed.agda}).
\end{theorem}

\paragraph{Where it doesn't match the Agda literally.}
\begin{itemize}
\item Variables are de Bruijn indices, newest first; on paper they are
  named. The relations of optimize carry partial maps of variables, delay
  a renaming, and remove complex operands a renaming $\rho$ with the
  target's substitution agreeing with the source's through it; on paper
  names make these the identity.
\item Evaluation contexts are on paper only; the Agda's $\mathit{red}$ is the
  congruence $\key{step\text{-}in}$, and a frame $h$ of RCO's relation is a
  \texttt{Frame}.
\item The redexes take the operand at a generic type $C$ with an equation
  $C \equiv A \rightarrow B$ (and so on), and fall through to $\error$
  otherwise; $\obs$ is \texttt{o-error} on a term that is not a value.
  Neither is reached (\texttt{proof/Semantics/Reduction/NoError.agda}).
\item The relation of simplify is \texttt{\_$\succcurlyeq$\_} (written with
  $\succcurlyeq$ in the Agda too), on the open steps \texttt{\_$\longrightarrow$\_}
  between the values \texttt{SV}; in a closed term they are steps of the
  reduction semantics (\texttt{closed-red}).
\end{itemize}
%===============================================================================
\section{Adequacy, No Errors, and the End-to-End Theorem}\label{sec:m-thm}

The four styles state the same end-to-end theorem, for a main program of
type $\key{Int}$: the source and the compiled $\Ml{X86}$ program stop with the
same constant, or both diverge. The three operational styles get it from
$\Equiv$ for $\pass{compile\text{-}x86}$; the graph model from
$\Same$ (Theorem~\ref{thm:m-den-compiler}) and adequacy for the clocked
interpreters of $\Ml{ISWIM}$ and $\Ml{X86}$, so that it states the
interpreter column's theorem, about the same interpreters. Adequacy is
reported separately in the comparison, since the denotational semantics
can be taken as \emph{the} semantics, and the no-error lemmas are in no
total.

\subsection{Observations}

A machine (or clocked program) $L$ terminates with $o$, $\Term_L(o)$, or
diverges, $\Div_L$ (Section~\ref{sec:obs}); it \emph{converges} when it
terminates with an answer, an observation other than $\error$
(\texttt{Converges}). For a main program of type $\key{Int}$ an answer is a
constant $n$. A clocked program $f$ \emph{returns $v$ from some clock on}
when there is a clock $c$ with $f(k) = \Ret v$ for every $k \geq c$
(\texttt{Stable}, \texttt{proof/Framework/Clocked.agda}); such results
compose (\texttt{st-bind}: a bind returns from the larger of the two
clocks on) and are unique (\texttt{stable-unique}).

\subsection{No errors}

\begin{lemma}[No Errors]\label{lem:m-no-error}\sloppy
For every closed $\Ml{ISWIM}$ term $M$ and $\Ml{X86}$ program $P$ of any type,
\begin{itemize}
\item the CEK machines never stop with an error,
  $\neg\Term_{\Mach{\Ml{ISWIM}}(M)}(\error)$ and
  $\neg\Term_{\Mach{\Ml{X86}}(P)}(\error)$
  (\texttt{proof/\allowbreak Semantics/\allowbreak CEK/\allowbreak NoError.agda}:
  \texttt{iswim-no-error}, \texttt{x86-no-error}), by inspection of the step functions: a machine
  stops only when it returns a value to the empty continuation (stack),
  every other state steps, and no value is observed as $\error$
  (\texttt{iswim-stops}, \texttt{iswim-answer}, \texttt{clos5-answer});
\item the clocked interpreters never err,
  $\neg\Term_{\Interp{\Ml{ISWIM}}(M)}(\error)$ and
  $\neg\Term_{\Interp{\Ml{X86}}(P)}(\error)$
  (\texttt{proof/Semantics/Interp/NoError.agda}: \texttt{iswim-no-error},
  \texttt{x86-no-error}), following their recursion: they have no error
  case, so a result at any clock is a timeout or a value
  (\texttt{bind-safe});
\item the reduction semantics of $\Ml{ISWIM}$ never stops with an error,
  $\neg\Term_{\mathcal{R}^{\Ml{ISWIM}}(M)}(\error)$
  (\texttt{proof/Semantics/Reduction/NoError.agda}: \texttt{iswim-no-error}),
  by progress (\texttt{progress}): a closed term is a value, or reduces; a
  value of a function, tuple, sum, integer or recursive type has the form
  of its type, so each redex reduces and the catch-all cases are never
  reached; and a value is observed as an answer (\texttt{value-answer}).
  The reduction column's $\Ml{X86}$ is the CEK machine.
\end{itemize}
\end{lemma}
The reduction semantics is the only one with error cases left in its
definition; removing them would need redexes on values rather than terms.
A machine that never stops with an error and does not converge diverges
(\texttt{never-final-diverges}, \texttt{proof/Framework/Simulation.agda});
so for these machines divergence is the negation of convergence. No
end-to-end theorem uses the lemma: the soundness halves of adequacy for
the interpreters show again that they never err (below), and the
adequacy for the CEK machines that preceded it, which uses the
machines' lemma, is kept only as an alternative
(\texttt{proof/Adequacy/CEK/}, Section~\ref{sec:m-adequacy}).

\subsection{Adequacy}\label{sec:m-adequacy}

Adequacy relates the graph model to the clocked interpreters of
$\Ml{ISWIM}$ and $\Ml{X86}$ (\texttt{proof/Adequacy/}); it is direct to
$\Ml{X86}$, with no adequacy for $\Ml{GlobalFun}$ or the other intermediate
languages, and does not use the value relations of the interpreter
column. Each half is by induction on the clock and the program, not over
the steps of a machine: soundness has no continuations or call stack,
and the logical relation's evaluation is ``returns $v$ from some clock
on''. The soundness halves only ever produce a timeout or a value whose
denotation is part of the program's, so they also show that the
interpreters never err.

\paragraph{$\Ml{ISWIM}$} (\texttt{ISWIM.agda}). \emph{Soundness}
(\texttt{Soundness.agda}): what the interpreter returns at any clock
denotes a part of the denotation, through denotations $\mathcal{V}(v)$ of
values (a closure denotes its $\lambda$ in the denotation of its
environment; $\mathcal{V}(\Mfold v) = \mathcal{V}(v)$), by induction on
the clock and the term: a result \emph{holds} of a set $D$ when it is a
timeout or a value $v$ with $\mathcal{V}(v) \subseteq D$, never an error
(\texttt{Holds}, \texttt{Sound}), and the binds of the interpreter
preserve it (\texttt{bind-holds}); a call, $\key{let}$ and $\key{case}$
use the continuity of the semantics and that every value denotes a
nonempty, consistent set. \emph{Adequacy} (\texttt{LogRel.agda}): a
logical relation $\mathbb{V}(d, v)$, by recursion on the denotational value
$d$ and on the $\Mfold$s of $v$,
\begin{align*}
  \mathbb{V}(d, \Mfold v) &\iff \mathbb{V}(d, v) \qquad
  \mathbb{V}(c, c') \iff c = c' \qquad
  \mathbb{V}(\nu, \cloV{\lam{x} N}{\gamma}) \qquad
  \mathbb{V}(\emptytup, \tupV{}) \\
  \mathbb{V}(V \mapsto w, \cloV{\lam{x} N}{\gamma}) &\iff \forall v.\; (\forall d \in V.\; \mathbb{V}(d, v)) \\
    &\qquad\quad \implies \exists v'.\; \ext{\gamma}{x}{v} \vdash N \Downarrow v' \land \mathbb{V}(w, v') \\
  \mathbb{V}(\tup{i}{d}, \tupV{\bar{v}}) &\iff \mathbb{V}(d, v_i) \qquad
  \mathbb{V}(\Left{d}, \inlV v) \iff \mathbb{V}(d, v) \qquad
  \mathbb{V}(\Right{d}, \inrV v) \iff \mathbb{V}(d, v)
\end{align*}
and false otherwise, where $\gamma \vdash M \Downarrow v$ is ``the
interpreter, on $M$ in $\gamma$, returns $v$ from some clock on''. The
results of subterms compose by \texttt{st-bind}, so no lemma about the
interpreter's monotonicity in the clock is needed. The fundamental lemma
(\texttt{adequate}) is by induction on the term, with related environments;
an application needs one argument value related to all the (finitely many)
elements of $V$, the value of the argument, since the interpreter returns
one value (\texttt{$\Downarrow$-det}, by \texttt{stable-unique}). The relation is by recursion on
the denotational value $d$, not on the type: as for delay
(Section~\ref{sec:m-den-delay}), a type index would not decrease at $\mu$.
So for every closed $M$ of any type
(\texttt{Theorems/\allowbreak Adequacy.agda}: \texttt{iswim-const-iff},
\texttt{iswim-diverges-iff}),
\begin{align*}
  c \in \SEM{M}\rho_0 &\iff \Term_{\Interp{\Ml{ISWIM}}(M)}(c) \\
  \Div_{\Interp{\Ml{ISWIM}}(M)} &\iff \SEM{M}\rho_0 = \emptyset .
\end{align*}
Divergence needs no lemma about errors: if the denotation is empty, a
sound result at any clock is a timeout; if it is not, the interpreter
returns a value from some clock on. Since the interpreter is typed, there
is no statement about errors, and the shape of an observation is read
off the related value.

\paragraph{$\Ml{X86}$} (\texttt{X86.agda}). \emph{Adequacy}
(\texttt{X86LogRel.agda}), for every program: the logical relation
$\mathbb{V}$ of the $\Ml{GlobalFun}$ interpreter's values, which the $\Ml{X86}$
interpreter shares (\texttt{Clos5Values.agda}; $\mathbb{V}(d, \Mfold v) =
\mathbb{V}(d, v)$), code $\codeV{k}$ being related to $\Lambda X.\, \Lambda
Y.\, w$ when a call of definition $k$ on values related to $X$ and $Y$
returns, from some clock on, a value related to $w$ (\texttt{Run}); the
fundamental lemma is by induction on the sequence, given that the blocks it
may jump to are adequate, with the frame, registers and flag of the
semantics related to the interpreter's at the types the state type gives them
(\texttt{$\mathbb{M}$}), and the open closure that of its slot
(\texttt{OpR}). An instruction other than $\key{callq}$ is one recursive
call of the interpreter on the rest of the sequence, so it needs no step
lemma; a call returns the same value for every element of its
denotation (\texttt{stable-unique}), and the rest of the caller composes
with it by \texttt{st-bind}; $\Mfold$ and $\Munfold$ are the identity on
denotations, and $\key{case}$ takes the injection in its slot apart, as
$\mathcal{C}$ does. \emph{Soundness} (\texttt{X86Soundness.agda}), for a
table that is consistent and contains the code of each definition
interpreted with the whole table, by induction on the clock, the blocks
and the sequence: what the interpreter returns from a sequence run in a
state denotes a part of what the sequence denotes from some
approximation of the state, a location denoting a finite part of the
meaning of the value it holds (\texttt{Apx}), since the semantics binds
a slot to a finite set and calls code on finite parts of $\rdi$ and
$\rsi$. A call combines the callee's approximations of its start, one
for each element of the finite part of the result that the caller uses,
into one (\texttt{merge}), and makes from it the entries $X \mapsto (Y
\mapsto w)$ of the code that the caller applies (\texttt{call-$\subseteq$}).
The programs that the compiler produces have such a
table (\texttt{X86Compiled.agda}, \texttt{compiled}), from globalize's
references to older definitions only
(\texttt{proof/Compiler/GlobalizeOlder.agda}) and the table lemmas of the
graph-model proofs of rco, lift locals, explicate control and select
(\texttt{table-rco}, \texttt{table-lift}, \texttt{table-explicate},
\texttt{table-$\simeq$}). So for every closed $M$ of any type
(\texttt{Theorems/\allowbreak Adequacy.agda}: \texttt{x86-const-iff}, \texttt{x86-diverges-iff}),
\begin{align*}
  c \in \SEM{\pass{compile\text{-}x86}(M)}^{\Ml{X86}} &\iff \Term_{\Interp{\Ml{X86}}(\pass{compile\text{-}x86}(M))}(c) \\
  \Div_{\Interp{\Ml{X86}}(\pass{compile\text{-}x86}(M))} &\iff \SEM{\pass{compile\text{-}x86}(M)}^{\Ml{X86}} = \emptyset .
\end{align*}

\paragraph{The alternative: adequacy for the CEK machines}
(\texttt{proof/Adequacy/CEK/}, not among the statements of
\texttt{Theorems/}). The same two theorems for the CEK machines of
$\Ml{ISWIM}$ and $\Ml{X86}$, with the same logical relations, but
$\gamma \vdash M \Downarrow v$ is ``the machine runs from $\ev{M}{\gamma}{\kappa}$
to $\rt{v}{\kappa}$ for every $\kappa$'', soundness shows that every step can
only shrink the set a state denotes (continuations and call stacks
denoting the contexts they stand for), and the statements are about
constants and convergence; divergence then needs
Lemma~\ref{lem:m-no-error} for the machines. It shares the tuples of
sets, the tables' lemmas and \texttt{X86Compiled.agda}, and is reported
beside the adequacy for the interpreters in the comparison.

\subsection{The end-to-end theorem}

\begin{theorem}[End to End, Graph Model]\label{thm:m-end-to-end-den}\sloppy
For every closed $\Ml{ISWIM}$ program $M : \key{Int}$ and every integer $n$,
\begin{align*}
  \Term_{\Interp{\Ml{ISWIM}}(M)}(n) &\iff \Term_{\Interp{\Ml{X86}}(\pass{compile\text{-}x86}(M))}(n) \\
  \Div_{\Interp{\Ml{ISWIM}}(M)} &\iff \Div_{\Interp{\Ml{X86}}(\pass{compile\text{-}x86}(M))}
\end{align*}
(\texttt{Theorems/\allowbreak Adequacy.agda}: \texttt{end-to-end-x86-const},
\texttt{end-to-end-x86-diverges}; the proof is
\texttt{proof/Adequacy/EndToEnd.agda}).
\end{theorem}
The interpreter of $M$ terminates with $n$ iff $n \in \SEM{M}\rho_0$
(adequacy for $\Ml{ISWIM}$), iff $n \in \SEM{\pass{compile\text{-}x86}(M)}$
(Theorem~\ref{thm:m-den-compiler}), iff the $\Ml{X86}$ interpreter terminates
with $n$ (adequacy for $\Ml{X86}$); likewise for divergence, through the
emptiness of the two denotations, which $\Same$ relates. No graph model
with errors is needed: there are no errors to observe.

\begin{theorem}[End to End, Common to the Four Styles]\label{thm:m-end-to-end}\sloppy
For every closed $\Ml{ISWIM}$ program $M : \key{Int}$, the source and its
compilation $\pass{compile\text{-}x86}(M)$ stop with the same constant,
or both diverge:
\[
  (\forall n.\; \Term_{L(M)}(n) \iff \Term_{L'(\pass{compile\text{-}x86}(M))}(n))
  \land (\Div_{L(M)} \iff \Div_{L'(\pass{compile\text{-}x86}(M))}) ,
\]
where $(L, L')$ is
\begin{itemize}
\item $(\Interp{\Ml{ISWIM}}, \Interp{\Ml{X86}})$ in the graph model
  (Theorem~\ref{thm:m-end-to-end-den}) and in the interpreter column
  (Theorem~\ref{thm:m-interp-compiler}; \texttt{compile-x86-const},
  \texttt{compile-x86-diverges} of \texttt{Theorems/\allowbreak Interpreter.agda});
\item $(\Mach{\Ml{ISWIM}}, \Mach{\Ml{X86}})$ in the CEK column
  (Theorem~\ref{thm:m-cek-compile}; \texttt{compile-x86-const},
  \texttt{compile-x86-diverges} of \texttt{Theorems/\allowbreak Simulation.agda});
\item $(\mathcal{R}^{\Ml{ISWIM}}, \Mach{\Ml{X86}})$ in the reduction column
  (Theorem~\ref{thm:m-red-compiler}; \texttt{compile-x86-const},
  \texttt{compile-x86-diverges} of \texttt{Theorems/\allowbreak Reduction.agda}).
\end{itemize}
\end{theorem}
In the operational columns the theorem is an instance of $\Equiv$, which
holds at any type and also equates errors, of which there are none
(Lemma~\ref{lem:m-no-error}); in the graph model it comes from adequacy,
whose soundness halves show that the interpreters never err. So the four
columns state one theorem, and the comparison has no caveat that the
theorems differ. The intermediate compositions ($\pass{compile}_4$ to
$\pass{compile\text{-}cfg}$) and the passes hold at any type.

\begin{sloppypar}
\paragraph{The statements.} The public statements are in
\texttt{agda/Mu/Theorems/}, one file per style and one for adequacy:
\texttt{Denotational.agda}, \texttt{Simulation.agda},
\texttt{Interpreter.agda}, \texttt{Reduction.agda} and
\texttt{Adequacy.agda} (the latter with adequacy for the two interpreters
and the end-to-end theorem through them). Each theorem
there has its full type written out and is defined as its proof, under
\texttt{agda/Mu/proof}; all are imported by \texttt{agda/Mu/All.agda}.
The statements for $\pass{compile}_4$, the no-error lemmas and the
alternative adequacy for the CEK machines are not among them;
they are the top-level definitions named above, in the proof files.
\end{sloppypar}
%===============================================================================
\part{Dynamic Typing and Errors}\label{part:dyn}

The dynamically typed pipeline (\texttt{agda/}: the typed languages in
\texttt{agda/Language/Typed}, the passes in \texttt{agda/Compiler}, the
semantics in the \texttt{Typed} subdirectories of
\texttt{agda/Semantics/*}, the statements in \texttt{agda/Theorems}) compiles
ISWIM, the untyped source of Part~\ref{part:untyped}, to the typed assembly
language $\Tl{X86}$:
\[
\begin{array}{l}
\mathrm{ISWIM} \xrightarrow{\pass{simplify}} \mathrm{ISWIM} \xrightarrow{\pass{cast}} \Tl{ISWIM}
  \xrightarrow{\pass{enclose}} \Tl{Clos} \xrightarrow{\pass{optimize}} \Tl{Clos}
  \xrightarrow{\pass{concretize}} \Tl{EnvClos} \xrightarrow{\pass{delay}} \Tl{FlatClos} \\
\qquad \xrightarrow{\pass{globalize}} \Tl{GlobalFun} \xrightarrow{\pass{rco}} \Tl{Mon}
  \xrightarrow{\pass{lift}} \Tl{Imp} \xrightarrow{\pass{explicate}} \Tl{CFG}
  \xrightarrow{\pass{select}} \Tl{X86}
\end{array}
\]
It is, in effect, the pipeline of Part~\ref{part:mu} specialised to the
single recursive type $\TAny$, with the case on a type tag fused into a
partial $\key{project}$, which is the only operation that can fail. So the
two pipelines have the same recursion (a recursive type in both) and the
same closure conversion, and differ only in the trapped errors and the tag
checks: the difference in each column is the cost of dynamic typing, with
recursion held constant (the design note, \S4). This part presents only
what dynamic typing adds: the type $\TAny$ and the casts, cast insertion,
$\key{project}$ and $\key{chk}$ in each language and style, and the graph
model with errors, which the plain graph model needs to observe errors.
Everything else (the passes from enclose on, the semantics, the
relations) is that of Part~\ref{part:mu} at the types below.

\section{Types, Casts and Cast Insertion}\label{sec:t-langs}

\paragraph{$\TAny$ and the flat types} (\texttt{Language/Typed/Type.agda}). A
value of type $\TAny$ is a value of some flat type $G$, tagged with it.
The flat types mirror the values of untyped ISWIM: a constant of a base
type, a tuple of $n$ values of type $\TAny$, a sum ($\key{inl}$ or
$\key{inr}$ of a value of type $\TAny$), and a function from $\TAny$ to
$\TAny$. The type $\Tcode{n}$, from $\Tl{FlatClos}$ on, is that of closed code
taking an environment of $n$ values. In the notation of
Part~\ref{part:mu},
\[
  \TAny \cong \mu X.\, \textstyle\sum_{G} G[X], \qquad
  \Tinject{G} = \Mfold \circ \mathit{in}_G, \qquad
  \Tproject{t} = \mathit{case}_t \circ \Munfold ,
\]
where $\mathit{case}_t$ is the case on the tag, with an error in every
branch but $\Tflatof{t}$'s. $\TAny$ is an open-ended sum, which fixes the shape of everything else: the contexts
of $\Tl{ISWIM}$ and $\Tl{Clos}$ are numbers (every variable has type
$\TAny$), $\Tflat{\Tfun}$ is $\TAny \to \TAny$, $\Tflat{\Ttup{n}}$ holds
$\TAny$s, and tables are lists of numbers. The tags $t$ are the flat types
that $\key{project}$ can test; a tuple tag $\Tttup{n}{i}$ also names the
component $i$ that the projection after it reads.
\begin{syntax}
  & B &::=& \key{Nat} \mid \key{Int} \mid \key{Bool} \\
  & G, H &::=& B \mid \Ttup{n} \mid \Tsum \mid \Tfun \\
  & A &::=& \TAny \mid \Tflat{G} \mid \Tcode{n} \\
  & t &::=& \Tfun \mid \Tsum \mid \Tint \mid \Tttup{n}{i}
    \qquad (1 \leq i \leq n)
\end{syntax}
\[
  \Tflatof{\Tfun} = \Tfun \quad
  \Tflatof{\Tsum} = \Tsum \quad
  \Tflatof{\Tint} = \key{Int} \quad
  \Tflatof{\Tttup{n}{i}} = \Ttup{n}
  \qquad
  \Tout{\key{sub}} = \key{Int} \quad
  \Tout{\key{eq}} = \Tsum
\]
A value of type $\TAny$ carries its flat type, and $\key{project}$ compares
it with $\Tflatof{t}$, so flat types have decidable equality (which
$\Ml{ISWIM}$'s types do not need).

\paragraph{The languages.} $\Tl{ISWIM}$ is ISWIM with explicit casts,
\[
  \frac{M : \Tflat{G}}{\Tinject{G} M : \TAny}
  \qquad
  \frac{M : \TAny}{\Tproject{t} M : \Tflat{\Tflatof{t}}} ,
\]
where every variable has type $\TAny$, a $\lambda$ has type
$\Tflat{\Tfun}$ with a body of type $\TAny$, an application takes an
operator of type $\Tflat{\Tfun}$, $\prj{i}$ an operand of type
$\Tflat{\Ttup{n}}$, $\key{case}$ one of type $\Tflat{\Tsum}$, and a
primitive operands of type $\Tflat{\key{Int}}$. $\Tl{Clos}$ to $\Tl{CFG}$
are those of Part~\ref{part:mu} at these types, with $\Tinject{G}$ and
$\Tproject{t}$ in place of $\Mfold$ and $\Munfold$ (and as simple
expressions in $\Tl{Imp}$ and $\Tl{CFG}$); a closure has type
$\Tflat{\Tfun} = \exists n.\, \Tcode{n} \times \Tflat{\Ttup{n}}$ from
$\Tl{FlatClos}$ on, a table is a list $\Sigma = n_0, \ldots, n_{m-1}$, and the two
targets of a $\Tl{CFG}$ $\key{case}$ start in the same frame $\Delta', \TAny$.

\subsection{Cast Insertion (ISWIM \texorpdfstring{$\to$}{->} \texorpdfstring{$\Tl{ISWIM}$}{ISWIM-T})}\label{sec:t-cast}

Cast insertion (\texttt{Compiler/Cast.agda}, after \emph{Essentials of
Compilation}, ch.~9) turns every term of ISWIM into a term of type $\TAny$.
A value is injected with its flat type. An operation projects its operand
to the flat type it needs, at the point where the ISWIM machine checks it:
the operator of an application before the operand is evaluated. A primitive
checks both operands after it has evaluated them both, so they are bound by
$\key{let}$s first.
\begin{align*}
  \pass{cast}(x) &= x \qquad
  \pass{cast}(\lam{x} N) = \Tinject{\Tfun}(\lam{x} \pass{cast}(N)) \\
  \pass{cast}(L\;M) &= (\Tproject{\Tfun} \pass{cast}(L))\; \pass{cast}(M) \\
  \pass{cast}(c) &= \Tinject{B} c \qquad c \text{ of base type } B \\
  \pass{cast}(\tuple{M_1,\ldots,M_n}) &= \Tinject{\Ttup{n}}
    \tuple{\pass{cast}(M_1),\ldots,\pass{cast}(M_n)} \\
  \pass{cast}(\prj{i} M) &= \prj{i} (\Tproject{\Tttup{n}{i}} \pass{cast}(M)) \\
  \pass{cast}(\inl M) &= \Tinject{\Tsum} (\inl \pass{cast}(M))
    \qquad\text{and likewise for } \key{inr} \\
  \pass{cast}(\Case{L}{x}{M}{y}{N}) &=
    \Case{\Tproject{\Tsum} \pass{cast}(L)}{x}{\pass{cast}(M)}{y}{\pass{cast}(N)} \\
  \pass{cast}(\prim{p}{L}{M}) &= \Let x = \pass{cast}(L) \In \Let y = \pass{cast}(M) \In \\
    &\qquad \Tinject{\Tout{p}} \prim{p}{\Tproject{\Tint} x}{\Tproject{\Tint} y} \\
  \pass{cast}(\Let x = M \In N) &= \Let x = \pass{cast}(M) \In \pass{cast}(N)
\end{align*}
The pass is defined on well-scoped terms (\texttt{WS k M}), so it needs a
closed program; the pipeline uses $\pass{cast\text{-}closed}(M) =
\pass{cast}(M)$ if $\FV(M) = \emptyset$, and otherwise the program
$(\Tproject{\Tfun} U)\; U$, $U = \Tunit$, that stops with an error at once.
Simplify, on untyped ISWIM, keeps programs closed
(Lemma~\ref{lem:simplify-closed}), so the theorems, for closed programs,
are about $\pass{cast}$. The compiler is
$\pass{compile}_4 = \pass{delay} \circ \pass{concretize}_{\emptyset} \circ
\pass{optimize} \circ \pass{enclose} \circ \pass{cast\text{-}closed} \circ
\pass{simplify}$, and so on as in Part~\ref{part:mu}
(\texttt{Compiler/Pipeline.agda}).

\subsection{\texorpdfstring{$\Tl{X86}$}{X86-T} and \texorpdfstring{$\key{chk}$}{chk}}\label{sec:t-x86}

$\Tl{X86}$ (\texttt{Language/Typed/X86.agda}) is $\Ml{X86}$ at these types,
with an operand $\Tup{a}$ for $\Tinject{G}$ (a move), registers
$\codeA$, $\tupA$ of the open closure ($\alpha$ standing for the hidden
length), and two differences:
\begin{itemize}
\item a $\key{case}$ is three real instructions, $\key{testinl}\; x;\
  \key{load}\; \key{payload}\; x, s;\ \key{je}\; l_1, l_2$, since the payload
  has type $\TAny$ on both sides;
\item $\key{project}$ is the pseudo-instruction $\Tchk{t}{x}{d}$, the only
  check: it moves slot $x$ of type $\TAny$ to $d$ at the flat type
  $\Tflat{\Tflatof{t}}$, and is the trapped error when $x$ holds a value of
  another kind. It stands for the X86 test of $t$ and the move, which
  clobber $\rax$ or the flag, so it leaves $\rax$ and the flag unset:
  \[
    \Tchk{t}{x}{d} : \psi \Rightarrow \psi' \quad\text{for } x : \TAny \in \Delta \text{ and }
    d : \Tst{\Theta'}{\Delta}{O}{\Tunset} \Rightarrow_{\Tflat{\Tflatof{t}}} \psi' ,
  \]
  with $\Theta'$ the $\Theta$ with $\rax$ unset.
\end{itemize}
Select is $\Ml{X86}$'s, with $\pass{sel}(\Tinject{G} a, o, k) = \key{movq}\;
\Tup{a}, \bar{o};\ k$ and $\pass{sel}(\Tproject{t} x, o, k) = \Tchk{t}{x}{\bar{o}};\ k$.
The erasure $\lfloor - \rfloor$ to untyped X86
(\texttt{Language/Typed/X86Erase.agda}) drops the types and expands
$\lfloor \Tchk{t}{x}{d} \rfloor = \Tcheck{t}(x);\ \key{movq}\; x, d$, the test
being an existing X86 instruction that errs on a value of any other kind
and otherwise only clobbers $\rax$ or the flag:
\begin{align*}
  \Tcheck{\Tint}(o) &= \key{movq}\; o, \rax;\ \key{subq}\; \$0, \rax &
  \Tcheck{\Tsum}(o) &= \key{testinl}\; o \\
  \Tcheck{\Tfun}(o) &= \key{load}\; \key{fst}\; o, \rax &
  \Tcheck{\Tttup{n}{i}}(o) &= \key{load}\; \pi_i\; o, \rax
\end{align*}
\begin{theorem}[Erasure]\label{thm:t-erase}
For every $\Tl{X86}$ program $P$,
\[ \Equiv(\Mach{\Tl{X86}}(P), \Mach{\mathrm{X86}}(\lfloor P \rfloor)) \]
(\texttt{erase-correct}); if the erased program errs, the $\Tl{X86}$ machine
reaches a failing $\key{chk}$ (\texttt{erase-errors}); and
$\lfloor \pass{select}(P) \rfloor = \pass{select}_{\mathrm{X86}}(P)$ for every
$\Tl{CFG}$ program, select into the untyped X86 of Part~\ref{part:untyped}
(\texttt{select-erase}; \texttt{Theorems/X86Erasure.agda}).
\end{theorem}
These are assurance theorems, not passes of the compiler and not among the
comparison's statements.

\section{Project in Each Style}\label{sec:t-styles}

\paragraph{Graph model} (\texttt{Semantics/Graph/Typed/};
\texttt{Framework/GraphModel/TypedGuard.agda}). $\Tinject{G}$ is the
identity, and $\Tproject{t}$ keeps its whole operand when the operand has
a value that the tag accepts, and means $\emptyset$ otherwise (a failed
project means $\emptyset$, as any error does in the plain graph model):
\begin{align*}
  \Tguard(P, D) &= \{ w \in D \mid \exists d \in D.\; P(d) \} \qquad
  \key{project}^{F}_t(D) = \Tguard(\Taccepts{F}(t, -), D) \\
  \Taccepts{F}(\Tfun, d) &\iff F(d) \qquad
  \Taccepts{F}(\Tint, d) \iff d \text{ is an integer} \\
  \Taccepts{F}(\Tsum, d) &\iff d = \Left{u} \text{ or } \Right{u} \\
  \Taccepts{F}(\Tttup{n}{i}, d) &\iff d = \tup{i}{u}, \text{ an entry of a tuple of } n \text{ components}
\end{align*}
where a function is an entry of a graph up to $\Tl{EnvClos}$
($\Tfuntwo(d) \iff d = V \mapsto w$ or $\nu$) and a pair from $\Tl{FlatClos}$
on ($\Tfunfour(d) \iff d = \pairv{u}{v}$). What this costs the proofs:
\begin{itemize}
\item \emph{cast insertion} (\texttt{CastCorrect.agda}): the untyped
  operators use only elements that a project accepts
  ($\key{project}^{\Tfuntwo}_{\Tfun}(D) \app E \simeq D \app E$, and so on),
  and the two $\key{let}$s of a primitive bind sets of which it needs one
  element each, so $\SEM{M}\rho \simeq \SEM{\pass{cast}(M)}^{\Tl{ISWIM}}\sigma$;
\item \emph{delay}: the preservation relation $\sqsubseteq$ and
  $\mathsf{ALR}$ each get a clause for project, $\mathsf{ALR}$ the field
  \texttt{fun-obs} (a pair in $V'$ needs a function in $D$), with the
  lemmas that project preserves and reflects (\texttt{acc-preserve},
  \texttt{acc-reflect}); the depth index is needed as in
  Section~\ref{sec:m-den-delay}, since $\TAny$ is a recursive type;
\item \emph{select}: $\key{testinl}$, the load of the payload and $\key{je}$
  agree with $\Tl{CFG}$'s $\key{case}$ only on a slot that never holds both
  an $\key{inl}$ and an $\key{inr}$ ($\mathit{Tagged}$, $\mathit{Inv}$,
  \texttt{case-core}), which holds of every frame since $\Lambda$ binds
  consistent sets.
\end{itemize}
In this model a program that errs and one that diverges both mean
$\emptyset$, so the end-to-end theorem ($\SEM{M}\rho_0 \obseq
\SEM{\pass{compile\text{-}x86}(M)}^{\Tl{X86}}$, the same constants and
nonempty together, \texttt{Theorems/Denotational.agda}) cannot tell an error
from divergence; the graph model with errors (Section~\ref{sec:t-gerr})
does.

\paragraph{CEK machines} (\texttt{Semantics/CEK/Typed/}). A value of type
$\TAny$ is $\TdynV{G}{v}$; $\key{injK}_G$ tags a value, and
$\rt{\TdynV{H}{v}}{\Kf{projK_t}{\kappa}}$ steps to $\rt{v}{\kappa}$ when $H =
\Tflatof{t}$ and stops with $\error$ otherwise: the only $\error$ of the
machines (in $\Tl{X86}$, a failing $\key{chk}$). In each pass from enclose
to globalize, projects fail together on related values (\texttt{proj$\sim$},
\texttt{proof/Correctness/CEK/Common.agda}). Cast insertion
(\texttt{CastSim.agda}) relates an ISWIM value to a value of type $\TAny$
tagged with the flat type of its kind, and a source frame to the typed
frame wrapped in the $\key{injK}$ and $\key{projK}$ frames of the casts;
where the source stops with $\error$ on an operand of the wrong kind, the
target stops at the corresponding project.

\paragraph{Clocked interpreters} (\texttt{Semantics/Interp/Typed/}).
$\key{project?}_t(\TdynV{H}{v})$ is $\Ret v$ if $H = \Tflatof{t}$ and
$\error$ otherwise, the only error; related values give related results of
$\key{project?}$. Cast insertion is clock for clock (\texttt{CastFuel.agda}).

\paragraph{Reduction semantics} (\texttt{Semantics/Reduction/Typed/}).
$\Tinject{G} V$ is a value, and
$\Tproject{t} (\Tinject{G} V) \to V$ if $G = \Tflatof{t}$, and is stuck
($\error$) otherwise, the only error. Cast insertion's relation
(\texttt{CastRed.agda}) follows the shape of $\pass{cast}$'s output, with the
states after an operator has been projected and after the first operand of
a primitive is bound.

\paragraph{The theorems} (\texttt{agda/Theorems/}). In the three operational
columns each pass, and each prefix of the compiler, gives $\Equiv$
semantics, which equates errors too: the compiled program terminates with
the same constants, errs exactly when the source does, and diverges exactly
when it does (\texttt{compile-const}, \texttt{compile-error},
\texttt{compile-diverges}). Unlike Part~\ref{part:mu}, these theorems and the
plain graph model's are not the same statement, which is what the graph
model with errors repairs.

\section{The Graph Model with Errors}\label{sec:t-gerr}

The typed languages are also given meaning in the graph model with errors
of Part~\ref{part:untyped} (Section~\ref{sec:graph-errors};
\texttt{Semantics/GraphErr/Typed/}, \texttt{Framework/GraphModel/ErrorMonad.agda}):
a set $D$ is a computation that errs when $\omega \in D$, has the values
$\lfloor D \rfloor = D \setminus \{\omega\}$, and diverges when it has
neither; $\mathit{val}(D)$ holds when $\lfloor D \rfloor \NE$, and $D \bind F$
is $\{\omega\}$ when $\omega \in D$, together with $F(\lfloor D \rfloor)$ when
$\mathit{val}(D)$. \emph{Only project checks kinds and errs}: each operator
is the graph model's kernel under the binds of its operands, in the
machines' order, with no check ($D \app^{\mathsf{T}} E = D \bind \lambda F.\,
E \bind \lambda X.\, F \app X$, and so on); an operand may still hold
$\omega$, from a failed project, which the binds propagate. A project to tag
$t$ binds its operand, adds $\omega$ when the operand has a value of another
kind than that of $\Tflatof{t}$, and keeps the values of that kind:
\begin{align*}
  \Tchecked{\varphi}_k(X, A) &= \{ \omega \mid \exists d \in X.\; \varphi(\Tkind(d)) \neq k \} \cup A \\
  \Tkeep{\varphi}_t(X) &= \{ d \in X \mid \Thit(t, X),\; \varphi(\Tkind(d)) = \Tkindf(\Tflatof{t}) \} \\
  \key{project}^{\varphi}_t(D) &= D \bind \lambda X.\, \Tchecked{\varphi}_{\Tkindf(\Tflatof{t})}(X, \Tkeep{\varphi}_t(X))
\end{align*}
where $\Tkind(d)$ is the kind of a value of the domain ($\Tkind(\omega) =
\mathit{err}$), $\Tkindf(G)$ that of the values of a flat type, and
$\Thit(\Tttup{n}{i}, X)$ asks for an entry of component $i$
(\texttt{Framework/GraphModel/TypedCheck.agda}). Up to $\Tl{EnvClos}$ $\varphi$
is the identity; from $\Tl{FlatClos}$ on a closure is a pair, so $\Tkindfour$
sees a pair as a function and code as an error. A $\key{let}$ is the bind
of its body, and a slot is bound as $\Tletv(D, F) = D \bind \lambda X.\,
\Lambda^{\omega}(F) \app X$; $\Tl{X86}$'s $\key{chk}$ is a strict write of
$\key{project}^{\Tkindfour}_t$ of its slot, and $\key{callq}$'s result is
written strictly too.

\paragraph{The proofs} (\texttt{proof/Correctness/GraphErr/},
\texttt{Theorems/GraphErrors.agda}). Every pass but delay preserves
meanings exactly ($\simeq$), by induction on the term with environments of
sets of values ($\mathit{Val}$: nonempty, without $\omega$): the binds are
the monad's left unit, binds associate, and a case and a check commute with
a bind. What errors add:
\begin{itemize}
\item \emph{simplify} holds with environments that also hold no $\omega$,
  so for closed programs only (simplify duplicates and drops only values);
\item \emph{cast insertion} moves the kind errors from the untyped
  operators' kernels to a project of the operand, related by one law
  (\texttt{proj-$\bind$}): if $K$ needs kind $k = \Tkindf(\Tflatof{t})$,
  $\key{project}^{\mathsf{T}}_t(D) \bind K \simeq D \bind \lambda X.\,
  \key{checked}_k(X, K(X))$ for a consistent $D$;
\item \emph{delay} traces every target error to a source error: the
  preservation relation has $\omega \sqsubseteq D' \iff \omega \in D'$, and
  $\mathsf{ALR}$ the fields \texttt{err-obs} ($\omega \in V' \implies \omega
  \in D$) and \texttt{kind-obs} (a target value's kind, through
  $\Tkindfour$, is that of a source value) in place of \texttt{nonempty-obs}
  and \texttt{fun-obs}, with lemmas lifting both relations through the
  monad's binds (\texttt{binds-lift}, \texttt{ApproxLR-bind}); its theorem
  is $\Same$ below;
\item \emph{select} needs $\mathit{Tagged}$ slots of values
  ($\mathit{Atom}$) for the three instructions of a $\key{case}$, and a call
  that binds its operands where $\Tl{X86}$ reads the closure's slot.
\end{itemize}
With $\Same(D, E)$ when $D$ and $E$ err together, have the same constants
and have a value together, for a closed ISWIM program $M$,
$\Same(\SEM{M}^{\mathrm{ISWIM},\omega}\rho_0,
\SEM{\pass{compile\text{-}x86}(M)}^{\Tl{X86},\omega})$, and likewise for each
prefix of the compiler (\texttt{compile-x86-same}).

\paragraph{Adequacy with errors.} The meaning tells the four outcomes of
the $\Tl{X86}$ machine apart, directly
(\texttt{proof/Adequacy/GraphErr/X86.agda}): $\omega$ is in it iff the
machine errs, $c$ iff it terminates with $c$, a value iff it terminates
with an answer, and nothing iff it diverges, for the programs of the
compiler, whose table is consistent and whose definitions' code depends
only on the older ones (\texttt{compile-ok}); likewise for ISWIM. So the
ISWIM machine on $M$ and the $\Tl{X86}$ machine on its compilation stop with
an error together, stop with the same constants, and diverge together
(\texttt{end-to-end-x86-error}, \ldots), the theorem of the operational
columns. This column, about 2500 lines without cast insertion, is what
dynamic typing costs the denotational approach to state the operational
theorem; Part~\ref{part:mu} needs none of it (Section~\ref{sec:m-thm}).

\paragraph{Where it doesn't match the Agda literally.} Tuple tags
$\Tttup{n}{i}$ are the Agda's \texttt{ttup (n-1) (i-1)}; cast insertion
carries a renaming for the variable that the $\key{let}$ of a primitive adds,
and $\FV(M) = \emptyset$ is \texttt{ws? 0 M}; $\key{project}^F_t$ is
\texttt{project$^{\mathtt{g}}$ F t}; the operators of the model with errors
are $\texttt{app}^{\mathsf{t}}$, $\texttt{get}^{\mathsf{t}}$,
$\texttt{case}^{\mathsf{t}}$, $\texttt{prim}^{\mathsf{t}}$ and
$\texttt{project}^{\varphi}$ (\texttt{Semantics/GraphErr/Typed/Ops.agda}),
$\Tletv$ is $\texttt{let}^{\mathsf{t}}$, and $\mathit{val}(D)$ is
\texttt{Conv}.

%===============================================================================
\part{History: The Untyped Pipeline}\label{part:untyped}

This part presents the untyped pipeline, the one that preceded the
dynamically typed pipeline of Part~\ref{part:dyn} and the $\Ml{ISWIM}$
pipeline of Part~\ref{part:mu}. It is preserved, and checked by
\texttt{make check}, in \texttt{agda/Untyped/} (module prefix
\texttt{Untyped.}), whose paths mirror those of \texttt{agda/}; it is no
longer measured. Its intermediate languages are untyped, and their
semantics check the kinds of values themselves:
\[
\begin{array}{l}
\mathrm{ISWIM} \xrightarrow{\pass{simplify}} \mathrm{ISWIM}
  \xrightarrow{\pass{enclose}} \mathrm{Clos} \xrightarrow{\pass{optimize}} \mathrm{Clos}
  \xrightarrow{\pass{concretize}} \mathrm{EnvClos} \xrightarrow{\pass{delay}} \mathrm{FlatClos}
  \xrightarrow{\pass{globalize}} \mathrm{GlobalFun} \\
\qquad \xrightarrow{\pass{rco}} \text{monadic GlobalFun}
  \xrightarrow{\pass{lift}} \mathrm{Imp} \xrightarrow{\pass{explicate}} \mathrm{CFG}
  \xrightarrow{\pass{select}} \mathrm{X86}
\end{array}
\]
It is kept here because the other two parts build on it: the graph
model's domain, operators and consistency (Section~\ref{sec:den-ops}), the
machines and simulations (Sections~\ref{sec:obs} and~\ref{sec:sims}), the
clocked programs (Section~\ref{sec:u-interp}) and reductions as machines
(Section~\ref{sec:red}) are the language-independent frameworks of
\texttt{agda/Framework}, which all three pipelines use; and its definitions
of ISWIM and simplify, and the semantics of ISWIM in each style, are
shared with the dynamically typed pipeline (the $\Ml{ISWIM}$ pipeline has
its own, typed, copies). X86 is this pipeline's target alone: the
dynamically typed pipeline ends at $\Tl{X86}$ (Section~\ref{sec:t-x86}),
whose erasure is X86 (Theorem~\ref{thm:t-erase}). X86's syntax
(\texttt{agda/\allowbreak Language/\allowbreak X86.agda}) and its machine
(\texttt{agda/\allowbreak Semantics/\allowbreak CEK/\allowbreak X86.agda}) stay
outside \texttt{agda/\allowbreak Untyped/}; everything else here, the
semantics of X86 in the graph models and the interpreter and their
adequacy included, is in \texttt{agda/Untyped/}. The part has the languages
(Section~\ref{sec:u-langs}) and the passes (Section~\ref{sec:u-passes}),
and then one section per column, each with its semantics and the
relations of its proofs: the graph model (Section~\ref{sec:u-den}), CEK
machines (Section~\ref{sec:u-cek}), clocked interpreters
(Section~\ref{sec:u-interp}), reduction semantics
(Section~\ref{sec:u-red}), and the graph model with errors
(Section~\ref{sec:graph-errors}).

\section*{Contents and notes of this part}
\addcontentsline{toc}{section}{Contents and notes of this part}

The public theorem statements of this pipeline are in
\texttt{agda/Untyped/Theorems}; they and the definitions they use
(everything outside the \texttt{proof} directories) are the audit surface,
and the proofs are in \texttt{agda/Untyped/proof}, which mirrors the paths
of the surface. The meaning of the primitives is in
\texttt{Semantics/Graph/Prim.agda}.

\paragraph{Contents.}
\begin{enumerate}
\item \textbf{The compiler:}
  \begin{itemize}
  \item the syntax of ISWIM and Clos--GlobalFun, including GlobalFun programs,
    and the primitives $\key{sub}$ and $\key{eq}$;
  \item enclose, optimize, concretize, delay (with the let), and globalize;
  \item remove complex operands (rco) into monadic GlobalFun, $\pass{compile\text{-}mon}$,
    and well-scoped GlobalFun programs ($\mathit{Scoped}_5$);
  \item the imperative language Imp, lift locals into it, and
    $\pass{compile\text{-}imp}$;
  \item the control-flow-graph language CFG, explicate control into it, and
    $\pass{compile\text{-}cfg}$;
  \item x86 with variables over abstract values (X86), select
    instructions into it, and $\pass{compile\text{-}x86}$.
  \end{itemize}
\item \textbf{Denotational: graph model:}
  \begin{itemize}
  \item the value domain and environments, with $\rho_0(x) = \{\omega\}$;
  \item the operators $\star$, $\Lambda$, $\key{pair}$/$\key{car}$/$\key{cdr}$,
    strict $\mathcal{T}$, $\key{proj}$, $\key{inl}$/$\key{inr}$/$\mathcal{C}$,
    $\key{guard}$, $\key{prim}_p$, and $\odot$;
  \item consistency ($\sim$, $\approx$, $\consis$), the binders $\Lambda$ and
    $\mathcal{C}$ binding only consistent sets, the preservation of a relation
    by a semantics, and continuity;
  \item the semantics: the clauses all languages share, then each language's own
    clauses, including GlobalFun's table-indexed semantics, $\key{code}$, and $\key{table}$;
  \item what the semantics of Imp and CFG share: stores, $\mathit{read}_n$,
    $\mathit{checked}_n$, atoms $\mathcal{A}$, simple expressions $\mathcal{S}$,
    $\key{frame}$, and $\key{table}$;
  \item the continuation-passing semantics of Imp (the main one), and the direct
    one, with ${\downarrow}$, $\key{let}_s$, and $\key{case}_s$;
  \item the semantics of CFG, with $\key{blocks}$ and $\mathit{run}$;
  \item the semantics of X86: the frame, registers, and flag, strict
    writes ($\key{put}$, $\key{flag}$), the instructions in
    continuation-passing style, $\key{callq}$, $\key{code}$, and $\key{table}$;
  \item the enclose and optimize equalities;
  \item for concretize, the $\mathit{Ref}$ semantics and $\sim_\varphi$;
  \item for delay, the let equation, the element-wise relation $\sqsubseteq$
    for preservation, and $\mathit{depth}$, $\mathsf{ALR}$, and $\rhd$
    for reflection;
  \item for globalize, $\mathit{Extends}$, $\mathit{RB}_n$, and $\mathit{Agree}_n$;
  \item for rco, the let laws, strict operators, the let chain, and the
    rco equation;
  \item for lift locals, $\mathit{Agree}$ and $\mathit{Scope}$, and, against the
    continuation-passing semantics, the enclosing lets $K$ and the let and
    case laws; the lemma against the direct semantics;
  \item for explicate control, $\preceq$ on block lists, $\mathit{Stable}$,
    $\key{assignK}$, the continuation-passing lemma, and the theorem; the
    theorem against the direct semantics;
  \item for select instructions, $\mathit{Tagged}$, $\mathit{Inv}$,
    $\mathit{Fits}$, and $\key{after}$, the lemmas for simple expressions,
    case, tails, monotone CFG tails, and the start of a definition, and the
    theorem for definitions of at least 2 slots;
  \item the end-to-end theorems, for $\pass{compile}$, $\pass{compile\text{-}mon}$,
    $\pass{compile\text{-}imp}$, $\pass{compile\text{-}cfg}$, and $\pass{compile\text{-}x86}$.
  \end{itemize}
\item \textbf{CEK machine + simulation:}
  \begin{itemize}
  \item observations, machines as step functions, termination, divergence,
    and determinism;
  \item the machines of ISWIM, Clos--GlobalFun, Imp, and CFG: values, continuation frames,
    and transitions; the evaluation of simple expressions ($\mathit{first}$,
    $\mathit{resume}$) that Imp and CFG share;
  \item the X86 machine: states with a call stack, flag, frame, and
    registers, and its transitions, with $\key{callq}$, $\key{retq}$, and jumps;
  \item forward simulations, $\mathit{Equivalent}$, and the theorem that a
    forward simulation gives $\mathit{Equivalent}$;
  \item for each pass, the relations $\approx$, $\approx_k$, and $\sim$ on
    values, continuations, and states, and its correctness theorem;
  \item for rco, $\preceq$ on environments, the pending let chain
    $\mathit{Pend}$, and its frames;
  \item for lift locals, frames related to environments ($\approx^n_\varphi$), and
    continuations given such a frame (head, weakened, and calls);
  \item for explicate control, the modes, $\preceq$, the relations $\approx$,
    $\mathit{Ctl}$, and $\mathit{Br}$ on continuations, the states indexed by
    $\mathit{depth}$, and the join lemma;
  \item for select instructions, $\approx$ between continuations and call
    stacks, $\mathit{Dest}$, the synchronisation points $\sim$, and the
    scripts $\mathit{Fate}$;
  \item the end-to-end theorems, for $\pass{compile}$, $\pass{compile\text{-}mon}$,
    $\pass{compile\text{-}imp}$, $\pass{compile\text{-}cfg}$, and $\pass{compile\text{-}x86}$;
  \item backward simulations, and the direct backward simulation for delay.
  \end{itemize}
\item \textbf{Clocked interpreter + value relations:}
  \begin{itemize}
  \item results, clocked programs, termination, divergence, and $\mathit{Equivalent}$;
  \item $\mathit{RelResult}$, and the theorem that related results for every
    clock give $\mathit{Equivalent}$;
  \item the clocked interpreters of ISWIM, Clos--GlobalFun, Imp, CFG, and X86, with
    $\key{eval\text{-}simple}$, $\key{inj?}$, and $\key{set}$ shared by Imp and CFG;
  \item for each pass, the lemma relating the interpreters clock for clock,
    and its correctness theorem; for rco, with equal results, and the let chain;
    for lift locals, with $\mathit{Out}$; for explicate control, with $\sqsubseteq$,
    $\mathit{Means}$, and $\mathit{Ok}$; for select instructions, with equal results;
  \item the end-to-end theorems, for $\pass{compile}$, $\pass{compile\text{-}mon}$,
    $\pass{compile\text{-}imp}$, $\pass{compile\text{-}cfg}$, and $\pass{compile\text{-}x86}$.
  \end{itemize}
\item \textbf{Reduction semantics + simulation:}
  \begin{itemize}
  \item reductions as machines ($\mathit{Red}$, $\mathit{toNext}$), and the
    reduction semantics of ISWIM and Clos--GlobalFun by substitution, with
    evaluation contexts; Imp, CFG, and X86 keep their machines;
  \item simulations between reductions: answering a stop by a run to a final
    state, $\mathit{Ans}$, congruence, the theorem that a simulation gives
    $\mathit{Equivalent}$, and closure under substitution;
  \item for each pass, the relation on terms, its fundamental lemma, and its
    correctness theorem; for rco, the let chain related for every value; for
    lift locals, terms in a context of lets related to Imp states;
  \item the end-to-end theorems, for $\pass{compile}$, $\pass{compile\text{-}mon}$,
    $\pass{compile\text{-}imp}$, $\pass{compile\text{-}cfg}$, and $\pass{compile\text{-}x86}$.
  \end{itemize}
\end{enumerate}

\paragraph{Where it doesn't match the Agda literally.}
Using names instead of de Bruijn indices changes a few definitions:
\begin{itemize}
\item \textbf{Enclose} lists $\FV(\lam{y} N)$ as the closure's free variables.
  The Agda lists the first $m$ surrounding variables, enough to cover every
  variable the body uses.
\item \textbf{Closure free variables} are variables in Clos and
  references in EnvClos, which is all the passes produce. The Agda allows
  arbitrary terms there.
\item \textbf{Optimize} has no renaming parameter. It keeps $x_i$ when
  $x_i \in \FV(N)$.
\item \textbf{Concretize}'s map has no fallback for variables past the free
  variables. Closed code can't refer to them.
\item \textbf{Tuple indices} run from $1$ to $n$.
\item \textbf{Global definitions} are listed oldest first, so $\funref{k}$ names $d_k$.
\item \textbf{GlobalFun definitions} are closed by the grammar. The Agda's
  syntax allows any definition; $\mathit{Scoped}_5$ is the Agda's
  \texttt{Scoped} (\texttt{Compiler/PipelineScope.agda}).
\item \textbf{Fresh variables} of rco are de Bruijn indices in the Agda:
  each let shifts the operands after it, and the atoms under the chain.
\item \textbf{Lift locals} has no map $\varphi$ in the Agda: variable $x$
  at depth $d$ is slot $d - (x + 1)$, its de Bruijn level, and slot $n$ when
  $x \geq d$. The Agda defines the pass on derivations that its input is
  in monadic normal form (\texttt{Mon}).
\item \textbf{Lift locals}' relations ($\mathit{Agree}$, $\approx^n_\varphi$,
  $\mathit{Out}$) relate the frame to the environment through $\varphi$.
  In the Agda, at depth $d$ (the length of the environment $\rho$ in the
  operational proofs) variable $x$ is slot $d - (x + 1)$, so the frame
  is $\rho$ reversed followed by the slots above it ($\rho \mathbin{{}^{r}{+}\!\!{+}} \tau$),
  and $\mathit{Scope}$ asks $\rho(x) \simeq \{\omega\}$ for $x \geq d$.
\item \textbf{Explicate control} threads the list of blocks explicitly
  in the Agda, newest first: block $t_l$ is the one with $l$ blocks after
  it in the list, and is looked up with the blocks older than it
  (\texttt{block}). So a list that extends $\vec{t}$ with newer blocks
  ($\vec{t} \preceq \vec{t}'$, $\vec{t} \sqsubseteq \vec{t}'$) has $\vec{t}$
  as a suffix in the Agda, not a prefix.
\item \textbf{Consistency} of a set, $\consis(D)$, is the Agda's \texttt{Con};
  of an environment or a store, $\consis(\rho)$, is \texttt{Con-env}; and
  $D \approx E$ on sets is \texttt{Every \_\char`~\_ D E}.
\item \begin{sloppypar}\textbf{The direct and continuation-passing semantics of Imp} are
  \texttt{ImpDirect.agda} (an experiment, under
  \texttt{proof/Semantics/Graph}) and \texttt{ImpK.agda} (under
  \texttt{Semantics/Graph}); what they share with CFG is in
  \texttt{Imp.agda}.\end{sloppypar}
\item \begin{sloppypar}\textbf{Compiler correctness} is stated with $\obseq$.
  The Agda proves its two halves as separate theorems,
  \texttt{compile-correct-const} and \texttt{compile-correct-nonempty},
  and likewise for $\pass{compile\text{-}imp}$ and $\pass{compile\text{-}cfg}$:
  \texttt{compile-imp-correct-const}, \texttt{compile-imp-correct-nonempty},
  \texttt{compile-cfg-correct-const},
  \texttt{compile-cfg-correct-nonempty}, \texttt{compile-x86-correct-const},
  and \texttt{compile-x86-correct-nonempty}.\end{sloppypar}
\item \textbf{The X86 semantics}' strict write of a register or the flag,
  $\key{guard}(D; E)$, is the Agda's \texttt{strict D E}; the flag is
  \texttt{left}~$\emptytup$ or \texttt{right}~$\emptytup$. Its blocks and
  table are listed newest first, as in CFG.
\item \begin{sloppypar}\textbf{Select instructions}' precondition is \texttt{Defs}, and
  that $\pass{compile\text{-}cfg}$ meets it is \texttt{defs-2}; the
  lemmas are \texttt{sel}, \texttt{case-core}, \texttt{tail},
  \texttt{blocks-}$\simeq$, \texttt{Mono.tail-mono}, \texttt{frame-}$\simeq$,
  and \texttt{start-}$\simeq$ (\texttt{proof/Correctness/Graph/SelectCorrect.agda}).
  $B \mathrel{\mathit{BR}} B'$ is \texttt{BR B B'}.\end{sloppypar}
\end{itemize}
In the simulation proof:
\begin{itemize}
\item \textbf{Machine environments} are finite maps from variables.
  The Agda uses lists indexed by de Bruijn indices, and a Clos closure
  holds the values of its free variables last first.
\item \textbf{Values collected} in $\key{tupK}$ and $\key{closK}$ frames are
  listed in order. The Agda lists them in reverse.
\item \textbf{Primitives} test their operands through their observations
  ($\delta^?$ on $\obs(v)$ and $\obs(w)$), which are integers
  exactly when the values are.
\item \textbf{Projections} $\prj{i} M$ carry the arity $n$ of the tuple,
  as \texttt{get} does in the Agda ($i : \mathtt{Fin}\;n$).
\item \textbf{Enclose}'s $\approx_X$ needs no bound on the Clos environment.
  The Agda's \texttt{Agree} also asks that the Clos environment covers
  the variables the closure captures.
\item \textbf{Optimize}'s relations need no renaming. The Agda relates
  variable $x$ to $r(x)$ for the renaming $r$ that optimize carries, and
  the values collected for a closure through the renumbering of positions
  that \texttt{prune} computes. $\mathit{Scoped}$ is \texttt{WS}, and
  $\mathit{WS}(M)$ is \texttt{WS 0 M}.
\item \textbf{Closed} ISWIM terms, $\FV(M) = \emptyset$, are the Agda's
  \texttt{Closed}.
\item \textbf{Concretize}'s $\approx_\varphi$ ranges over the domain of
  $\varphi$. The Agda's $\varphi$ is total, mapping a variable past a
  closure's free variables to an unbound one, so its theorem has no
  precondition.
\item \textbf{Delay}'s $\approx^{?}$ on environments skips the fresh variables $f$.
  The Agda relates environments through a renaming that skips the
  let-bound variables.
\item \textbf{Global definitions} in the Agda's GlobalFun machine and globalize
  relation are listed newest first.
\item \textbf{RCO}'s $\rho_1 \preceq \rho_2$ is the Agda's \texttt{Ren}: the
  source term $M$ is related to \texttt{rename r (rco M)} for a renaming
  \texttt{r} with $\rho_2(r(x)) = \rho_1(x)$. The Agda's frames are indexed
  by the kind of construct (\texttt{Frame}, \texttt{Kind}), and its
  $\mathit{Pend}$ takes the values and operands as lists.
\item \textbf{Lift locals}' relations assume that the variables bound in a
  body are distinct from each other and from its parameters, so extending
  $\rho$ never shadows. The Agda's frame $\rho \mathbin{{}^{r}{+}\!\!{+}} \tau$ keeps a
  shadowed value in its slot, so it needs no such assumption.
  $\vdash_n$ and $\vdash^{\mathsf{h}}_n$ are \texttt{Kont} and \texttt{KontH},
  and $\approx^n_\varphi$ is \texttt{Fr}.
\item \textbf{The frames of simple expressions} $F$ in the Imp and CFG
  machines are wrapped in a continuation by \texttt{frm}: $F \cdot \kappa$ is
  \texttt{frm F $\kappa$}.
\item \begin{sloppypar}\textbf{Explicate}'s Imp states $\evi{a}{\sigma}{\kappa}$, for an
  atom $a$, are the Agda's \texttt{eval (simple (atm a)) $\varphi$ $\kappa$}.
  The machines and the simulation are in \texttt{Semantics/CEK/CFG.agda} and
  \texttt{proof/Correctness/CEK/ExplicateSim.agda}.\end{sloppypar}
\item \textbf{The X86 machine}'s state $\xrun\;\vec{\imath}\;j\;\vec{\beta}\;\Sigma\;f\;\sigma\;\rho$
  is the Agda's \texttt{run is j bs} $\sigma$ \texttt{f} $\varphi$ $\rho$, where
  $\sigma$ is the call stack and $\varphi$ the frame.
\item \textbf{Select instructions}' $\kappa \approx \Sigma$ is the Agda's
  \texttt{K \char`\~{}} $\sigma$, $s_1 \sim s_2$ is \texttt{Match} (with
  \texttt{exec\char`\~}, \texttt{ret\char`\~}, and \texttt{now}), and the
  scripts' cases are \texttt{stop}, \texttt{lag}, and \texttt{go}
  (\texttt{proof/Correctness/CEK/SelectSim.agda}). Lemmas~\ref{lem:to-dest},
  \ref{lem:sel-simple-sim}, and~\ref{lem:fate} are \texttt{to-dest},
  \texttt{simple}, and \texttt{fate}.
\end{itemize}
In the interpreter proof:
\begin{itemize}
\item \textbf{Environments, values, and relations} are those of the
  simulation proof, so the items above about them apply here too.
\item \textbf{Values collected} by $\key{evals}$ are listed in order. The
  Agda's \texttt{eval-args} conses them onto an accumulator, last first,
  and a tuple reverses them.
\item \textbf{Optimize}'s lemma for the free variables of a closure is
  stated with $\mathit{Built}$. The Agda's \texttt{prune}$\approx$ relates the
  values collected so far through the renumbering of positions that
  \texttt{prune} computes (\texttt{Collected}).
\item \textbf{Delay}'s lemma relates $M$ to $\pass{delay}(M)$. The Agda
  relates $M$ to \texttt{rename r (delay M)} for a renaming \texttt{r}
  of the environments, since the operand of an application is evaluated
  under the let-bound variable.
\item \textbf{The GlobalFun interpreter}'s table is listed newest first in the
  Agda, which looks up definition $j$ in its reverse (\texttt{global}).
\item \textbf{Clocked programs} are the Agda's \texttt{Semantics}, and
  $\outc$ is \texttt{outcome}, which returns a \texttt{Maybe Obs}.
\item \textbf{The Imp interpreter} collects the values of the atoms of a
  tuple onto an accumulator, last first, as GlobalFun's does.
\item \textbf{The CFG interpreter}'s $\key{tail}$ and $\key{jump}$ are
  structurally recursive on the clock, then the list of blocks, then the
  tail: a jump walks the list, newest first, to its block.
\item \textbf{The X86 interpreter}'s registers are a record
  (\texttt{Regs}) of \texttt{Maybe}s for $\rax$, $\rdi$, $\rsi$, and the
  flag; its $\key{cmpq}$ applies \texttt{bool?} to
  $\delta^?(\key{eq}, \obs(v), \obs(w))$, which is $\key{cmp}(v, w)$.
  It is structurally recursive on the clock, the list of blocks, and the
  instructions of a block.
\item \textbf{Select instructions}' lemmas are \texttt{simple-ok},
  \texttt{tail-ok}, and \texttt{jump-ok}, and calls agree by
  \texttt{calls}$\equiv$ and \texttt{call-pw}
  (\texttt{proof/Correctness/Interp/SelectFuel.agda}).
\end{itemize}

In the reduction proof:
\begin{itemize}
\item \textbf{Substitution} is on de Bruijn terms in the Agda
  (\texttt{Framework/Syntax/Subst.agda}), and leaves the cleared code of a
  closure untouched. The reduction semantics are in
  \texttt{Semantics/Reduction}, written as functions with congruence rules
  (\texttt{step-in}) rather than with evaluation contexts.
\item \textbf{Values} are a predicate on terms (\texttt{Value}), and a
  simulation answers with a target value (\texttt{Ans}, \texttt{val}).
\item \textbf{Enclose}'s closure clause relates the variables that the
  body uses (\texttt{Uses}) and keeps the target closed (\texttt{WS 0}).
\item \textbf{The preconditions} $\FV(M) = \emptyset$ of optimize and
  concretize are the Agda's \texttt{WS 0 M}.
\item \textbf{RCO}'s contexts $h$ are frames of the GlobalFun reduction
  (\texttt{Frame}), and the operands of a construct are reduced left to right
  (\texttt{LR}); the let chains are those of \texttt{lift}.
\item \textbf{Lift locals}' frame is $\rho$ reversed followed by the slots
  above it, as in the CEK proof (\texttt{Fr}), with $\sigma$ related to
  $\rho$ (\texttt{Ag}) and the context to the continuation (\texttt{Kont}).
\end{itemize}

\section{Languages}\label{sec:u-langs}

\subsection{ISWIM}

The source language: a call-by-value $\lambda$-calculus with constants, tuples, sums, primitive operations $p$ on integers $n, m$, and $\key{let}$, which evaluates $M$ and then $N$ with $x$ bound to its value.
There are no booleans: $\key{eq}$ returns $\inl\tuple{}$ for true and $\inr\tuple{}$ for false, so $\key{case}$ serves as $\key{if}$.
\begin{syntax}
  & p &::=& \key{sub} \mid \key{eq} \\
  & L, M, N &::=& x \mid \lam{x} N \mid L\;M \mid c
    \mid \tuple{M_1,\ldots,M_n} \mid \prj{i} M \\
  &&& \mid \inl M \mid \inr M \mid \Case{L}{x}{M}{y}{N} \mid \prim{p}{L}{M}
    \mid \Let x = M \In N
\end{syntax}
\[
  \FV(\lam{y} N) = \FV(N) \setminus \{y\}
  \qquad
  \delta(\key{sub}, n, m) = n - m
  \qquad
  \delta(\key{eq}, n, m) = (n = m)
\]

\subsection{Clos}

Like ISWIM, but each $\lambda$ becomes a closure that binds the values of its free variables $z_1\ldots z_n$ to names $x_1\ldots x_n$, so its body refers only to those names and its parameter.
\begin{syntax}
  & L, M, N &::=& x \mid \clos{z_1\ldots z_n}{x_1\ldots x_n\, y} N \mid L\;M \mid c
    \mid \tuple{M_1,\ldots,M_n} \mid \prj{i} M \\
  &&& \mid \inl M \mid \inr M \mid \Case{L}{x}{M}{y}{N} \mid \prim{p}{L}{M}
    \mid \Let x = M \In N
\end{syntax}
\[
  \FV(\clos{z_1\ldots z_n}{x_1\ldots x_n\, y} N) = \{z_1, \ldots, z_n\}
\]

\subsection{EnvClos}

Each closure now packs references $r_1\ldots r_n$ to its free variables into one tuple bound to $e$, so its body reaches them by projections from $e$.
\begin{syntax}
  & r &::=& x \mid \prj{i} x \\
  & L, M, N &::=& x \mid \clos{r_1\ldots r_n}{e\, y} N \mid L\;M \mid c
    \mid \tuple{M_1,\ldots,M_n} \mid \prj{i} M \\
  &&& \mid \inl M \mid \inr M \mid \Case{L}{x}{M}{y}{N} \mid \prim{p}{L}{M}
    \mid \Let x = M \In N
\end{syntax}
\[
  \FV(\clos{r_1\ldots r_n}{e\, y} N) = \FV(r_1) \cup \cdots \cup \FV(r_n)
\]

\subsection{FlatClos}

Closures are explicit pairs of closed code $\fun{e\,y.\,N}$ and an environment, built with $\key{pair}$, taken apart with $\key{fst}$ and $\key{snd}$, and applied to an environment and an argument with $L(M, N)$; $\key{let}$, from the source, is also used by the delayed application.
\begin{syntax}
  & L, M, N &::=& x \mid \fun{e\,y.\,N} \mid L(M, N) \mid c
    \mid \pair{M}{N} \mid \FST M \mid \SND M \mid \prim{p}{L}{M} \\
  &&& \mid \tuple{M_1,\ldots,M_n} \mid \prj{i} M
    \mid \inl M \mid \inr M \mid \Case{L}{x}{M}{y}{N}
    \mid \Let x = M \In N
\end{syntax}
\[
  \FV(\fun{e\,y.\,N}) = \emptyset
\]

\subsection{GlobalFun}

Like FlatClos, but the code is hoisted into a program's list of global definitions $d_k$, which terms refer to as $\funref{k}$.
\begin{syntax}
  & L, M, N &::=& x \mid \funref{k} \mid L(M, N) \mid c
    \mid \pair{M}{N} \mid \FST M \mid \SND M \mid \prim{p}{L}{M} \\
  &&& \mid \tuple{M_1,\ldots,M_n} \mid \prj{i} M
    \mid \inl M \mid \inr M \mid \Case{L}{x}{M}{y}{N}
    \mid \Let x = M \In N \\[1ex]
  & d &::=& \key{fun}\; e\,y.\,N \qquad\qquad (\FV(\key{fun}\; e\,y.\,N) = \emptyset) \\
  & P &::=& (d_0, \ldots, d_{m-1};\; M)
\end{syntax}

\subsection{Monadic GlobalFun}

The sublanguage of GlobalFun in monadic normal form: the operands of every construct but $\key{let}$ are atoms $a$, while the right-hand side and body of $\key{let}$ and the branches of $\key{case}$ may be complex.
\begin{syntax}
  & a &::=& x \mid \funref{k} \mid c \\
  & M, N &::=& a \mid a_1(a_2, a_3) \mid \pair{a_1}{a_2} \mid \FST a \mid \SND a
    \mid \prim{p}{a_1}{a_2} \\
  &&& \mid \tuple{a_1,\ldots,a_n} \mid \prj{i} a
    \mid \inl a \mid \inr a \mid \Case{a}{x}{M}{y}{N}
    \mid \Let x = M \In N \\[1ex]
  & d &::=& \key{fun}\; e\,y.\,N \\
  & P &::=& (d_0, \ldots, d_{m-1};\; M)
\end{syntax}

\subsection{Imp}

An imperative language with structured control flow and no binders. A
body $(n; e)$ runs $e$ in a frame of $n$ local slots $s \in \{0, \ldots,
n-1\}$, and an atom names a slot of the current frame. The assignment
$\assign{s}{e_1}{e_2}$ runs $e_1$, stores its value in slot $s$, and runs
$e_2$; $\Casei{a}{s}{e_1}{e_2}$ stores the payload of the injection in slot
$s$ before running a branch. Assignments are seen by everything that
runs after them in the same frame. A call $a_1(a_2, a_3)$ runs the body
of the definition $a_1$ refers to in a fresh frame with $a_2$ (the free
variables) in slot $0$ and $a_3$ (the argument) in slot $1$; reading or
assigning a slot beyond the frame errors.
\begin{syntax}
  & a &::=& s \mid \funref{k} \mid c \\
  & r &::=& a \mid a_1(a_2, a_3) \mid \pair{a_1}{a_2} \mid \FST a \mid \SND a
    \mid \prim{p}{a_1}{a_2}
    \mid \tuple{a_1,\ldots,a_n} \mid \prj{i} a
    \mid \inl a \mid \inr a \\
  & e &::=& r \mid \Casei{a}{s}{e_1}{e_2}
    \mid \assign{s}{e_1}{e_2} \\[1ex]
  & b &::=& (n;\; e) \\
  & P &::=& (b_0, \ldots, b_{m-1};\; b)
\end{syntax}
The simple expressions $r$ are shared with CFG.

\subsection{CFG}

A C-like language: each body is a control-flow graph of blocks, tails
$t$ that end in $\key{return}$, a $\key{goto}$, or a $\key{case}$ that
assigns the payload to slot $s$ and jumps to block $l_1$ or $l_2$. Frames,
slots and calls are as in Imp, and a call runs the start $t$ of the
callee's body. The blocks $t_0, \ldots, t_{j-1}$ of a body are numbered
in order of creation; block $t_l$ may jump only to blocks $l' < l$, and
the start to any block, so control flow within a body is acyclic. A
jump to any other label errors.
\begin{syntax}
  & t &::=& \key{return}\; r \mid \assign{s}{r}{t} \mid \key{goto}\; l
    \mid \key{case}\;a\;\key{of}_{s}\;\key{goto}\;l_1 \mid \key{goto}\;l_2 \\[1ex]
  & b &::=& (n;\; t_0, \ldots, t_{j-1};\; t) \\
  & P &::=& (b_0, \ldots, b_{m-1};\; b)
\end{syntax}

\subsection{X86}\label{sec:u-x86}

x86 with variables over the same values as before. The locations $d$ are
the slots of the frame and the registers; an operand $o$ is a location,
an immediate, or the label of a global definition. A flag is set by
$\key{cmpq}$ and $\key{testinl}$ and read by $\key{sete}$ and $\key{je}$.
The heap is abstract: $\key{alloc}$ builds a pair, tuple or injection
from its operands, and $\key{load}$ takes a field of one. $\key{callq}\;o$
reads $\%\key{rdi}$ and $\%\key{rsi}$ and runs the body of the definition
$o$ refers to, with them as they are, $\%\key{rax}$ unset and no flag, in
a fresh frame, and continues with the result in $\%\key{rax}$; it
clobbers $\%\key{rdi}$, $\%\key{rsi}$ and the flag. $\key{retq}$ reads
$\%\key{rax}$. Reading an unset register or flag errors. Frames,
blocks and jumps are as in CFG.
\begin{syntax}
  & d &::=& s \mid \%\key{rax} \mid \%\key{rdi} \mid \%\key{rsi} \\
  & o &::=& d \mid c \mid \funref{k} \\
  & i &::=& \key{movq}\; o, d \mid \key{subq}\; o, d \mid \key{cmpq}\; o_2, o_1
    \mid \key{sete}\; d \mid \key{testinl}\; o \\
  &&& \mid \key{alloc}\; \mathit{sh}\; \vec{o}, d \mid \key{load}\; f\; o, d
    \mid \key{callq}\; o \\
  & \mathit{sh} &::=& \key{pair} \mid \key{tuple} \mid \key{inl} \mid \key{inr}
  \qquad f \;::=\; \key{fst} \mid \key{snd} \mid \pi_i \mid \key{payload} \\
  & j &::=& \key{retq} \mid \key{jmp}\; l \mid \key{je}\; l_1, l_2 \\
  & \beta &::=& i_1; \ldots; i_p;\; j \\
  & b &::=& (n;\; \beta_0, \ldots, \beta_{q-1};\; \beta) \\
  & P &::=& (b_0, \ldots, b_{m-1};\; b)
\end{syntax}

%===============================================================================
\section{Compilation}\label{sec:u-passes}

\[
  \pass{compile} = \pass{globalize} \circ \pass{delay} \circ \pass{concretize}_{\mathit{id}}
     \circ \pass{optimize} \circ \pass{enclose} \circ \pass{simplify}
\]

\begin{align*}
  \pass{compile\text{-}mon} &= \pass{rco} \circ \pass{compile}
  &
  \pass{compile\text{-}imp} &= \pass{lift} \circ \pass{compile\text{-}mon} \\
  \pass{compile\text{-}cfg} &= \pass{explicate} \circ \pass{compile\text{-}imp}
  &
  \pass{compile\text{-}x86} &= \pass{select} \circ \pass{compile\text{-}cfg}
\end{align*}

In each pass, the clauses not shown are homomorphic.

\subsection{Simplify (ISWIM \texorpdfstring{$\to$}{->} ISWIM)}

Simplify, a simplified version of the source optimizer cp0 of Chez Scheme,
is the first pass of $\pass{compile}$. It carries what
it knows of each variable, $E(x) = (T, K)$: the term $T$ that replaces $x$ (a
variable or a small value), and what $x$ is known to be, $K$: a value, a
function $\langle \lam{y} N, E_f \rangle$ (its body before simplification, with
what was known where it was defined), or $\bot$. On a term it returns the
same, $\pass{simp}_E(M) = (M', K)$; we write $\pass{simp}_E(M)$ for $M'$ and
$\mathcal{K}_E(M)$ for $K$. A value is
\emph{small} when it is a variable, a constant, $\emptytup$, or $\inl a$ or
$\inr a$ of one of those; only small values are duplicated. Values are known
to be themselves, and $\mathcal{K}_E(\lam{y} N) = \langle \lam{y} N, E \rangle$;
under a binder that stays, its variable is itself, known only as itself
($E'$ below). A function body $N$ bound to an operand, $\pass{let}_E(x.N, M)$,
is shared by $\key{let}$, $\beta$ and inlining:
\begin{align*}
  \pass{let}_E(x.N, M) &= \pass{simp}_{\ext{E}{x}{(V, \mathcal{K}_E(M))}}(N)
    \quad \text{if } V \text{ is small} \\
  &= \Let x = V \In \pass{simp}_{\ext{E}{x}{(x, \mathcal{K}_E(M))}}(N)
    \quad \text{if } V \text{ is another value} \\
  &= \Let x = V \In \pass{simp}_{E'}(N)
    \quad \text{otherwise, where } V = \pass{simp}_E(M) \\[1ex]
  \pass{simp}_E(x) &= E(x) \\
  \pass{simp}_E(\Let x = M \In N) &= \pass{let}_E(x.N, M) \\
  \pass{simp}_E((\lam{x} N)\;M) &= \pass{let}_E(x.N, M) \\
  \pass{simp}_E(L\;M) &= \pass{let}_{E_f}(x.N, M)
    \quad \text{if } \mathcal{K}_E(L) = \langle \lam{x} N, E_f \rangle \text{ and the result is small enough} \\
  \pass{simp}_E(\prim{p}{L}{M}) &= \res{\delta(p, n, m)}
    \quad \text{if } \mathcal{K}_E(L) = n,\; \mathcal{K}_E(M) = m \\
  \pass{simp}_E(\Case{L}{x}{M}{y}{N}) &= \pass{simp}_{\ext{E}{x}{(K, K)}}(M)
    \quad \text{if } \mathcal{K}_E(L) = \inl K,\; K \text{ small} \\
  \pass{simp}_E(\prj{i} L) &= K_i
    \quad \text{if } \mathcal{K}_E(L) = \tuple{K_1, \ldots, K_n},\; K_i \text{ small} \\
  \mathcal{K}_E(\prj{i} L) &= K_i
    \quad \text{if } \mathcal{K}_E(L) = \tuple{K_1, \ldots, K_n} \\[1ex]
  \pass{simplify}(M) &= \pass{simp}_{x \mapsto (x, x)}(M)
\end{align*}
and likewise for $\inr$. Here $\res{o}$ is the term of a primitive's result
(a constant, or $\inl \emptytup$ or $\inr \emptytup$ for $\key{eq}$). The third
clause of $\pass{simp}$ is inlining: the known function's body is simplified
again, at the call, from the source, under $E_f$ and what is known of the
operand, and replaces the call if it has at most $24$ nodes; $\pass{simp}$
carries fuel ($4$ in $\pass{simplify}$), which inlining decreases, so that it
is structurally recursive. It drops only evaluations whose value is known,
which have no effect (the operator of an inlined call), and otherwise keeps
the order of evaluation: as in cp0, an operand that is not a value is never
substituted, but stays bound by a $\key{let}$, and is evaluated before the
body even if the body does not use it. Enclose keeps that $\key{let}$, so a
residual binding allocates no closure.

\begin{lemma}[Simplify Keeps Programs Closed]\label{lem:simplify-closed}
For every ISWIM term $M$,
\[
  \FV(M) = \emptyset \implies \FV(\pass{simplify}(M)) = \emptyset .
\]
\end{lemma}

\subsection{CC: Enclose (ISWIM \texorpdfstring{$\to$}{->} Clos)}

Enclose decides which variables each closure captures (its free variables) and closes off its body, so that the body sees the captured values only through its own bound names.
\begin{align*}
  \pass{enclose}(\lam{y} N) &= \clos{x_1\ldots x_n}{x_1\ldots x_n\, y} \pass{enclose}(N)
    \qquad\text{where } x_1, \ldots, x_n = \FV(\lam{y} N)
\end{align*}

\subsection{CC: Optimize (Clos \texorpdfstring{$\to$}{->} Clos)}

Optimize removes the captured variables that the body does not use.
\begin{align*}
  \pass{optimize}(\clos{z_1\ldots z_n}{x_1\ldots x_n\, y} N) &=
    \clos{z_{i_1}\ldots z_{i_k}}{x_{i_1}\ldots x_{i_k}\, y} \pass{optimize}(N) \\
  \text{where } \{ i_1 < \cdots < i_k \} &= \{ i \mid x_i \in \FV(N) \}
\end{align*}

\subsection{CC: Concretize (Clos \texorpdfstring{$\to$}{->} EnvClos)}

Concretize stores each closure's captured values in a single tuple and turns their uses in the body into projections from that tuple.
\[
  \varphi : \X \to \mathit{Ref}
  \qquad (\mathit{Ref} \text{ is the set of EnvClos references } r)
\]
\begin{align*}
  \pass{concretize}_\varphi(x) &= \varphi(x) \\
  \pass{concretize}_\varphi(\clos{z_1\ldots z_n}{x_1\ldots x_n\, y} N) &=
    \clos{\varphi(z_1)\ldots\varphi(z_n)}{e\,y}
       \pass{concretize}_{\varphi_n}(N) \\
  \text{where } \varphi_n &= \{ y \mapsto y,\; x_1 \mapsto \prj{1} e,\; \ldots,\; x_n \mapsto \prj{n} e \},
    \quad e \text{ fresh} \\
  \pass{concretize}_\varphi(\Case{L}{x}{M}{y}{N}) &=
    \key{case}\;\pass{concretize}_\varphi(L)\;\key{of} \\
    &\qquad \key{inl}\,x \Rightarrow \pass{concretize}_{\ext{\varphi}{x}{x}}(M)
     \mid \key{inr}\,y \Rightarrow \pass{concretize}_{\ext{\varphi}{y}{y}}(N) \\
  \pass{concretize}_\varphi(\Let x = M \In N) &=
    \Let x = \pass{concretize}_\varphi(M) \In \pass{concretize}_{\ext{\varphi}{x}{x}}(N)
\end{align*}

\subsection{CC: Delay (EnvClos \texorpdfstring{$\to$}{->} FlatClos)}

Delay represents each closure as an explicit pair of closed code and its environment tuple, and makes application pass the environment to the code.
\begin{align*}
  \pass{delay}(\clos{r_1\ldots r_n}{e\,y} N) &=
    \pair{\fun{e\,y.\,\pass{delay}(N)}}{\tuple{r_1,\ldots,r_n}} \\
  \pass{delay}(L\;M) &=
    \Let f = \pass{delay}(L) \In (\FST f)(\SND f, \pass{delay}(M))
    \qquad f \text{ fresh}
\end{align*}

\subsection{CC: Globalize (FlatClos \texorpdfstring{$\to$}{->} GlobalFun)}

Globalize hoists all code into a table of global definitions and replaces each piece of code by a reference to its entry.
\[
  \vec{d} = d_0, \ldots, d_{m-1}
  \qquad
  \pass{glob}(M, \vec{d}) = (M', \vec{d}')
\]
\begin{align*}
  \pass{glob}(x, \vec{d}) &= (x, \vec{d}) \\
  \pass{glob}(\fun{e\,y.\,N}, \vec{d}) &= (\funref{|\vec{d}_1|},\; \vec{d}_1, (\key{fun}\; e\,y.\,N'))
     \quad\text{where } (N', \vec{d}_1) = \pass{glob}(N, \vec{d}) \\
  \pass{glob}(\mathit{op}(\bar{x}_1.A_1, \ldots, \bar{x}_n.A_n), \vec{d}_0)
     &= (\mathit{op}(\bar{x}_1.A'_1, \ldots, \bar{x}_n.A'_n), \vec{d}_n)
     \quad\text{where } (A'_i, \vec{d}_i) = \pass{glob}(A_i, \vec{d}_{i-1}) \\[1ex]
  \pass{globalize}(M) &= (\vec{d};\; M') \quad\text{where } (M', \vec{d}) = \pass{glob}(M, \epsilon)
\end{align*}

\subsection{Remove Complex Operands (GlobalFun \texorpdfstring{$\to$}{->} Monadic GlobalFun)}

Remove complex operands binds each complex operand to a fresh variable with a $\key{let}$, from left to right, so that every operand becomes an atom.
\[
  \vec{b} = x_1{=}M_1, \ldots, x_k{=}M_k
  \qquad
  \pass{atomize}(M_1, \ldots, M_n) = (\vec{b};\; a_1, \ldots, a_n)
\]
\begin{align*}
  \pass{atomize}(\epsilon) &= (\epsilon;\; \epsilon) \\
  \pass{atomize}(a, \vec{M}) &= (\vec{b};\; a, \vec{a})
     \quad\text{where } (\vec{b}; \vec{a}) = \pass{atomize}(\vec{M}) \\
  \pass{atomize}(M, \vec{M}) &= (x{=}M, \vec{b};\; x, \vec{a})
     \quad\text{where } (\vec{b}; \vec{a}) = \pass{atomize}(\vec{M}),\
     M \text{ not an atom},\ x \text{ fresh} \\[1ex]
  \pass{lets}(\epsilon, N) &= N \\
  \pass{lets}((x{=}M, \vec{b}), N) &= \Let x = M \In \pass{lets}(\vec{b}, N)
\end{align*}
\[
  \mathit{op}(M_1, \ldots, M_n) \in \{\, L(M, N),\ \pair{M}{N},\ \FST M,\ \SND M,\ \prim{p}{M}{N},\
     \tuple{M_1, \ldots, M_n},\ \prj{i} M,\ \inl M,\ \inr M \,\}
\]
\begin{align*}
  \pass{rco}(\mathit{op}(M_1, \ldots, M_n)) &= \pass{lets}(\vec{b}, \mathit{op}(a_1, \ldots, a_n)) \\
  &\qquad \text{where } (\vec{b}; a_1, \ldots, a_n) = \pass{atomize}(\pass{rco}(M_1), \ldots, \pass{rco}(M_n)) \\
  \pass{rco}(\Case{L}{x}{M}{y}{N}) &= \pass{lets}(\vec{b}, \Case{a}{x}{M'}{y}{N'}) \\
  &\qquad \text{where } (\vec{b}; a) = \pass{atomize}(\pass{rco}(L)),\ M' = \pass{rco}(M),\ N' = \pass{rco}(N) \\
  \pass{rco}((d_0, \ldots, d_{m-1};\; M)) &= (\pass{rco}(d_0), \ldots, \pass{rco}(d_{m-1});\; \pass{rco}(M)) \\
  \pass{rco}(\key{fun}\; e\,y.\,N) &= \key{fun}\; e\,y.\,\pass{rco}(N)
\end{align*}
An atom left in place is evaluated after the complex operands that follow
it. This preserves errors in the machines and interpreters, which
check that the operator of $L(M, N)$ is code only at the call, but an
unbound variable would error later: $\pair{x}{L}$, with $L$ a loop,
errors in the empty environment, while
$\Let y = L \In \pair{x}{y}$ diverges. So the operational proofs of rco
assume well-scoped programs.

\paragraph{Well-scoped programs.}
\begin{align*}
  \mathit{Scoped}_5((d_0, \ldots, d_{m-1};\; M)) &\iff \FV(M) = \emptyset \land
    \forall k < m.\; \FV(d_k) = \emptyset \\
  \FV(\key{fun}\; e\,y.\,N) &= \FV(N) \setminus \{e, y\}
\end{align*}

\begin{lemma}[The Compiler Produces Well-Scoped Programs]\label{lem:compile-scoped}
For every ISWIM term $M$,
\[
  \FV(M) = \emptyset \implies \mathit{Scoped}_5(\pass{compile}(M)) .
\]
\end{lemma}

\subsection{Lift Locals (Monadic GlobalFun \texorpdfstring{$\to$}{->} Imp)}

Lift locals gives every variable of a body a slot of the body's frame,
declared up front, and turns $\key{let}$ into assignment. A variable gets
the slot numbered by its binding depth: a $\key{let}$ or $\key{case}$ at
depth $d$ assigns slot $d$, and the code under it is at depth $d+1$.
A slot at depth $d$ is assigned only by code at depth $d$ or more, where
no variable in scope names it, so no assignment clobbers a variable in
scope. The map $\varphi : \X \rightharpoonup \mathbb{N}$ gives the slots
of the variables in scope, and $n$ is the size of the frame.
\begin{align*}
  \pass{lift}_{n,\varphi,d}(x) &= \varphi(x)
    \qquad (n \text{ if } x \notin \mathrm{dom}(\varphi)) \\
  \pass{lift}_{n,\varphi,d}(\Let x = M \In N) &=
    \assign{d}{\pass{lift}_{n,\varphi,d}(M)}{\pass{lift}_{n,\ext{\varphi}{x}{d},d+1}(N)} \\
  \pass{lift}_{n,\varphi,d}(\Case{a}{x}{M}{y}{N}) &=
    \key{case}\;\pass{lift}_{n,\varphi,d}(a)\;\key{of}_{d} \\
  &\qquad \key{inl} \Rightarrow \pass{lift}_{n,\ext{\varphi}{x}{d},d+1}(M)
    \mid \key{inr} \Rightarrow \pass{lift}_{n,\ext{\varphi}{y}{d},d+1}(N) \\[1ex]
  \mathit{height}(\Let x = M \In N) &= \max(\mathit{height}(M),\, 1 + \mathit{height}(N)) \\
  \mathit{height}(\Case{a}{x}{M}{y}{N}) &= 1 + \max(\mathit{height}(M), \mathit{height}(N)) \\
  \mathit{height}(M) &= 0 \qquad \text{otherwise} \\[1ex]
  \pass{lift}(\key{fun}\; e\,y.\,N) &= (n;\; \pass{lift}_{n,\{e \mapsto 0,\, y \mapsto 1\},2}(N))
    \quad\text{where } n = 2 + \mathit{height}(N) \\
  \pass{lift}((d_0, \ldots, d_{m-1};\; M)) &=
    (\pass{lift}(d_0), \ldots, \pass{lift}(d_{m-1});\; b) \\
  &\qquad\text{where } b = (n;\; \pass{lift}_{n,\emptyset,0}(M)),\ n = \mathit{height}(M)
\end{align*}
A variable out of scope gets the slot $n$, beyond the frame, so that it
errors, as an unbound variable does in GlobalFun.

\subsection{Explicate Control (Imp \texorpdfstring{$\to$}{->} CFG)}

Explicate control flattens the nested assignments and cases of a body
into blocks, in continuation-passing style. $\pass{tail}(e)$ translates
$e$ in tail position, and $\pass{assign}(e, s, k)$ translates the
assignment of $e$ to slot $s$ followed by the tail $k$. Both add blocks
to the body's list; $\pass{add}(t)$ adds the block $t$ and returns its
label, the number of blocks before it. A case on the right of an
assignment adds $k$ as a block, the join point, which its branches jump
to, so $k$ is not duplicated.
\begin{align*}
  \pass{tail}(r) &= \key{return}\; r \\
  \pass{tail}(\assign{s}{e_1}{e_2}) &= \pass{assign}(e_1, s, \pass{tail}(e_2)) \\
  \pass{tail}(\Casei{a}{s}{e_1}{e_2}) &=
    \key{case}\;a\;\key{of}_{s}\;\key{goto}\;l_1 \mid \key{goto}\;l_2 \\
  &\qquad \text{where } l_1 = \pass{add}(\pass{tail}(e_1)),\ l_2 = \pass{add}(\pass{tail}(e_2)) \\[1ex]
  \pass{assign}(r, s, k) &= \assign{s}{r}{k} \\
  \pass{assign}(\assign{s'}{e_1}{e_2}, s, k) &= \pass{assign}(e_1, s', \pass{assign}(e_2, s, k)) \\
  \pass{assign}(\Casei{a}{s'}{e_1}{e_2}, s, k) &=
    \key{case}\;a\;\key{of}_{s'}\;\key{goto}\;l_1 \mid \key{goto}\;l_2 \\
  &\qquad \text{where } l = \pass{add}(k),\
    l_i = \pass{add}(\pass{assign}(e_i, s, \key{goto}\;l)) \\[1ex]
  \pass{explicate}((n;\; e)) &= (n;\; \vec{t};\; t)
    \quad\text{where } t = \pass{tail}(e) \text{ and } \vec{t} \text{ are the blocks it adds}
\end{align*}
and $\pass{explicate}$ maps over the definitions and main of a program.
Every block is added after the blocks it jumps to.

\subsection{Select Instructions (CFG \texorpdfstring{$\to$}{->} X86)}

Select instructions turns each statement into a few instructions; the
blocks keep their labels. $\pass{sel}(r, d)$ computes $r$ into the
location $d$, with $\bar{a}$ the operand for the atom $a$.
\begin{align*}
  \pass{sel}(a, d) &= \key{movq}\; \bar{a}, d \\
  \pass{sel}(a_1(a_2, a_3), d) &= \key{movq}\; \bar{a}_2, \%\key{rdi};\
    \key{movq}\; \bar{a}_3, \%\key{rsi};\ \key{callq}\; \bar{a}_1;\
    \key{movq}\; \%\key{rax}, d \\
  \pass{sel}(\key{sub}(a_1, a_2), d) &= \key{movq}\; \bar{a}_1, \%\key{rax};\
    \key{subq}\; \bar{a}_2, \%\key{rax};\ \key{movq}\; \%\key{rax}, d \\
  \pass{sel}(\key{eq}(a_1, a_2), d) &= \key{cmpq}\; \bar{a}_2, \bar{a}_1;\ \key{sete}\; d \\
  \pass{sel}(\pair{a_1}{a_2}, d) &= \key{alloc}\; \key{pair}\; \bar{a}_1, \bar{a}_2, d
    \qquad\text{and likewise for tuples and injections} \\
  \pass{sel}(\FST a, d) &= \key{load}\; \key{fst}\; \bar{a}, d
    \qquad\text{and likewise for } \SND,\ \prj{i} \\[1ex]
  \pass{sel}(\key{return}\; r) &= \pass{sel}(r, \%\key{rax});\ \key{retq} \\
  \pass{sel}(\assign{s}{r}{t}) &= \pass{sel}(r, s);\ \pass{sel}(t) \\
  \pass{sel}(\key{goto}\; l) &= \key{jmp}\; l \\
  \pass{sel}(\key{case}\;a\;\key{of}_{s}\;\key{goto}\;l_1 \mid \key{goto}\;l_2) &=
    \key{testinl}\; \bar{a};\ \key{load}\; \key{payload}\; \bar{a}, s;\ \key{je}\; l_1, l_2
\end{align*}
Subtraction goes through $\%\key{rax}$, so it is right also when $s$ is
$a_2$. The start of each global definition begins with the prologue
$\key{movq}\; \%\key{rdi}, 0;\ \key{movq}\; \%\key{rsi}, 1$, which moves its
free variables and argument to the slots where its body expects them.

%===============================================================================
\section{Denotational: Graph Model}\label{sec:u-den}

\subsection{Values and Operators}

\subsubsection{Values}

\[
\begin{array}{rcl}
  c &\in& \text{constants} \qquad x, y, z, e, f \in \X \\[1ex]
  \V \ni d, u, v, w &::=& c \mid V \mapsto w \mid \nu \mid \omega
     \mid \pairv{u}{v}
     \mid \tup{i}{v} \mid \Left{v} \mid \Right{v} \mid \emptytup \\
  U, V, W &\in& \FSET{\V} \\
  D, E, X, Y &\in& \SET{\V} \\
  F, G &:& \SET{\V} \to \SET{\V}
\end{array}
\]

\[
\begin{array}{rcl}
  \rho &:& \X \to \SET{\V} \\
  \rho_0(x) &=& \{\omega\} \\
  D \simeq E &\iff& D \subseteq E \text{ and } E \subseteq D \\
  \rho \simeq \rho' &\iff& \forall x.\; \rho(x) \simeq \rho'(x)
\end{array}
\]

\subsubsection{Operators}\label{sec:den-ops}

\begin{align*}
  D \app E &= \{ w \mid \exists V \NE.\; (V \mapsto w) \in D,\; V \subseteq E \} \\
  \Lambda(F) &= \{ V \mapsto w \mid V \NE,\; V \approx V,\; w \in F(V) \} \cup \{ \nu \} \\
  \mathcal{B}(c) &= \{ c \} \\
  \key{pair}(D, E) &= \{ \pairv{u}{v} \mid u \in D,\; v \in E \} \\
  \key{car}(D) &= \{ u \mid \exists v.\; \pairv{u}{v} \in D \} \\
  \key{cdr}(D) &= \{ v \mid \exists u.\; \pairv{u}{v} \in D \} \\
  \mathcal{T}() &= \{ \emptytup \} \\
  \mathcal{T}(D_1,\ldots,D_n) &= \{ \tup{i}{d} \mid d \in D_i,\; \forall j.\; D_j \NE \}
     \qquad (n > 0) \\
  \key{proj}_i(D) &= \{ d \mid \tup{i}{d} \in D \} \\
  \key{inl}(D) &= \{ \Left{d} \mid d \in D \} \\
  \key{inr}(D) &= \{ \Right{d} \mid d \in D \} \\
  \mathcal{C}(D, F, G) &=
    \{ w \mid \exists V \NE.\; V \approx V,\; (\forall v \in V.\; \Left{v} \in D),\; w \in F(V) \} \\
    &\cup\; \{ w \mid \exists V \NE.\; V \approx V,\; (\forall v \in V.\; \Right{v} \in D),\; w \in G(V) \} \\
  \key{guard}(D_1,\ldots,D_n; E) &= \{ w \mid (\forall i.\; D_i \NE),\; w \in E \} \\
  \key{prim}_p(D, E) &= \{ \res{\delta(p, n, m)} \mid n \in D,\; m \in E \}
\end{align*}
where $\res{n} = n$, $\res{\mathsf{true}} = \Left{\emptytup}$, and $\res{\mathsf{false}} = \Right{\emptytup}$.
The binders $\Lambda$ and $\mathcal{C}$ bind only finite sets that are
consistent ($V \approx V$, Section~\ref{sec:consistency}), so a denotation
under a consistent environment has no inconsistent elements.

\begin{align*}
  D \odot E &= (\key{car}(D) \app \key{cdr}(D)) \app E
\end{align*}

\subsubsection{Consistency}\label{sec:consistency}

\begin{align*}
  c \sim c' &\iff c = c' \\
  (V \mapsto w) \sim (V' \mapsto w') &\iff (V \approx V' \implies w \sim w') \\
  (V \mapsto w) \sim \nu, \quad \nu \sim (V \mapsto w), \quad \nu \sim \nu &\iff \mathsf{true} \\
  \omega \sim \omega, \quad \emptytup \sim \emptytup &\iff \mathsf{true} \\
  \pairv{u}{v} \sim \pairv{u'}{v'} &\iff u \sim u' \land v \sim v' \\
  \tup{i}{d} \sim \tup{j}{d'} &\iff n = n' \land (i \neq j \lor d \sim d')
     \qquad (i \in \{1..n\},\; j \in \{1..n'\}) \\
  \Left{d} \sim \Left{d'} &\iff d \sim d' \\
  \Right{d} \sim \Right{d'} &\iff d \sim d' \\
  u \sim v &\iff \mathsf{false} \quad \text{otherwise} \\[1ex]
  D \approx E &\iff \forall u \in D, v \in E.\; u \sim v \\
  \consis(D) &\iff D \approx D \\
  \consis(\rho) &\iff \forall x.\; \consis(\rho(x))
\end{align*}
Every value is consistent with itself ($v \sim v$), so $\consis(\{v\})$.

\paragraph{Preservation.}
A relation $R$ between sets is preserved by the operators of a semantics
when each operator maps related arguments to related results, an
argument that binds a variable being a function that maps related sets
to related results.

\begin{lemma}[A Semantics Preserves Relations]\label{lem:pres}
Let $R$ be preserved by the operators of $\SEM{-}^{\ell}$, with
$\{\omega\} \mathrel{R} \{\omega\}$. If $\rho(x) \mathrel{R} \rho'(x)$ for
every $x$, then $\SEM{M}^{\ell}\rho \mathrel{R} \SEM{M}^{\ell}\rho'$.
\end{lemma}
With $R = {\subseteq}$, each semantics is monotone. With $R = {\approx}$,
which the operators preserve, each is consistent:
\begin{lemma}[Consistency]\label{lem:con}
If $\consis(\rho)$, then $\consis(\SEM{M}^{\ell}\rho)$ for
$\ell \in \{\mathrm{ISWIM}, 3, 4\}$, and $\consis(\SEM{M}^{\mathrm{5}}_T\rho)$
when $\consis(T(k))$ for every $k$. For all GlobalFun definitions
$d_0, \ldots, d_{m-1}$ and every $k$,
\[
  \consis(\key{table}(d_0, \ldots, d_{m-1})(k)) .
\]
\end{lemma}

\subsubsection{Continuity}

\begin{align*}
  \rho' \subseteq \rho &\iff \forall x.\; \rho'(x) \subseteq \rho(x) \\
  \mathit{finiteNE}(\rho') &\iff \forall x.\; \exists V \NE.\; \rho'(x) = V \\
  \mathit{continuous}(\mathcal{E}) &\iff
    \forall \rho\, v.\; (\forall x.\; \rho(x) \NE) \land v \in \mathcal{E}(\rho) \\
    &\qquad \implies \exists \rho'.\; \mathit{finiteNE}(\rho'),\; \rho' \subseteq \rho,\; v \in \mathcal{E}(\rho')
\end{align*}

%===============================================================================
\subsection{Semantics}

\begin{align*}
  \SEM{-}^{\mathrm{ISWIM}}, \SEM{-}^{\mathrm{2}}, \ldots, \SEM{-}^{\mathrm{4}}
    &: \mathit{Term} \to (\X \to \SET{\V}) \to \SET{\V} \\
  \SEM{-}^{\mathrm{5}}_T
    &: \mathit{Term} \to (\X \to \SET{\V}) \to \SET{\V}
    \qquad (T : \mathbb{N} \to \SET{\V}) \\
  \SEM{-}^{\mathrm{5}} &: \mathit{Program} \to \SET{\V}
\end{align*}
For each language $\ell \in \{\mathrm{ISWIM}, 2, 3, 4\}$, $\SEM{-}^{\ell}$ is
defined by the common clauses below together with the clauses for $\ell$.

\subsubsection{Clauses common to all languages}

\begin{align*}
  \SEM{x}^{\ell}\rho &= \rho(x) \\
  \SEM{c}^{\ell}\rho &= \mathcal{B}(c) \\
  \SEM{\tuple{M_1,\ldots,M_n}}^{\ell}\rho &= \mathcal{T}(\SEM{M_1}^{\ell}\rho, \ldots, \SEM{M_n}^{\ell}\rho) \\
  \SEM{\prj{i} M}^{\ell}\rho &= \key{proj}_i(\SEM{M}^{\ell}\rho) \\
  \SEM{\inl M}^{\ell}\rho &= \key{inl}(\SEM{M}^{\ell}\rho) \\
  \SEM{\inr M}^{\ell}\rho &= \key{inr}(\SEM{M}^{\ell}\rho) \\
  \SEM{\Case{L}{x}{M}{y}{N}}^{\ell}\rho &=
     \mathcal{C}(\SEM{L}^{\ell}\rho,\; \lambda X.\, \SEM{M}^{\ell}\ext{\rho}{x}{X},\; \lambda Y.\, \SEM{N}^{\ell}\ext{\rho}{y}{Y}) \\
  \SEM{\prim{p}{L}{M}}^{\ell}\rho &= \key{prim}_p(\SEM{L}^{\ell}\rho, \SEM{M}^{\ell}\rho) \\
  \SEM{\Let x = M \In N}^{\ell}\rho &= \Lambda(\lambda X.\, \SEM{N}^{\ell}\ext{\rho}{x}{X}) \app \SEM{M}^{\ell}\rho
\end{align*}

\subsubsection{ISWIM}

\begin{align*}
  \SEM{\lam{x} N}^{\mathrm{ISWIM}}\rho &= \Lambda(\lambda X.\, \SEM{N}^{\mathrm{ISWIM}}\ext{\rho}{x}{X}) \\
  \SEM{L\;M}^{\mathrm{ISWIM}}\rho &= \SEM{L}^{\mathrm{ISWIM}}\rho \app \SEM{M}^{\mathrm{ISWIM}}\rho
\end{align*}

\subsubsection{Clos}

\begin{align*}
  \SEM{\clos{z_1\ldots z_n}{x_1\ldots x_n\, y} N}^{\mathrm{2}}\rho &=
    \key{guard}(D_1, \ldots, D_n;\;
       \Lambda(\lambda Y.\, \SEM{N}^{\mathrm{2}}\rho_0(x_1{:=}D_1,\ldots,x_n{:=}D_n, y{:=}Y)))
    \\ &\qquad\text{where } D_i = \rho(z_i) \\
  \SEM{L\;M}^{\mathrm{2}}\rho &= \SEM{L}^{\mathrm{2}}\rho \app \SEM{M}^{\mathrm{2}}\rho
\end{align*}

\subsubsection{EnvClos}

\begin{align*}
  \SEM{\clos{r_1\ldots r_n}{e\, y} N}^{\mathrm{3}}\rho &=
    \Lambda(\lambda X.\, \Lambda(\lambda Y.\, \SEM{N}^{\mathrm{3}}\rho_0(e{:=}X, y{:=}Y)))
    \app \mathcal{T}(\SEM{r_1}^{\mathrm{3}}\rho, \ldots, \SEM{r_n}^{\mathrm{3}}\rho) \\
  \SEM{L\;M}^{\mathrm{3}}\rho &= \SEM{L}^{\mathrm{3}}\rho \app \SEM{M}^{\mathrm{3}}\rho
\end{align*}

\subsubsection{FlatClos}

\begin{align*}
  \SEM{\fun{e\,y.\,N}}^{\mathrm{4}}\rho &= \Lambda(\lambda X.\, \Lambda(\lambda Y.\, \SEM{N}^{\mathrm{4}}\rho_0(e{:=}X, y{:=}Y))) \\
  \SEM{L(M, N)}^{\mathrm{4}}\rho &= (\SEM{L}^{\mathrm{4}}\rho \app \SEM{M}^{\mathrm{4}}\rho) \app \SEM{N}^{\mathrm{4}}\rho \\
  \SEM{\pair{M}{N}}^{\mathrm{4}}\rho &= \key{pair}(\SEM{M}^{\mathrm{4}}\rho, \SEM{N}^{\mathrm{4}}\rho) \\
  \SEM{\FST M}^{\mathrm{4}}\rho &= \key{car}(\SEM{M}^{\mathrm{4}}\rho) \\
  \SEM{\SND M}^{\mathrm{4}}\rho &= \key{cdr}(\SEM{M}^{\mathrm{4}}\rho)
\end{align*}

\subsubsection{GlobalFun}

\begin{align*}
  \SEM{\funref{k}}^{\mathrm{5}}_T\rho &= T(k) \\
  \SEM{L(M, N)}^{\mathrm{5}}_T\rho &= (\SEM{L}^{\mathrm{5}}_T\rho \app \SEM{M}^{\mathrm{5}}_T\rho) \app \SEM{N}^{\mathrm{5}}_T\rho \\
  \SEM{\pair{M}{N}}^{\mathrm{5}}_T\rho &= \key{pair}(\SEM{M}^{\mathrm{5}}_T\rho, \SEM{N}^{\mathrm{5}}_T\rho) \\
  \SEM{\FST M}^{\mathrm{5}}_T\rho &= \key{car}(\SEM{M}^{\mathrm{5}}_T\rho) \\
  \SEM{\SND M}^{\mathrm{5}}_T\rho &= \key{cdr}(\SEM{M}^{\mathrm{5}}_T\rho)
\end{align*}
together with the common clauses, with $\SEM{-}^{\mathrm{5}}_T$ in place of $\SEM{-}^{\ell}$.

\begin{align*}
  \key{code}_T(\key{fun}\; e\,y.\,N) &= \Lambda(\lambda X.\, \Lambda(\lambda Y.\, \SEM{N}^{\mathrm{5}}_T\rho_0(e{:=}X, y{:=}Y))) \\[1ex]
  \key{table}(d_0,\ldots,d_{m-1})(k) &=
    \begin{cases}
      \key{code}_{\key{table}(d_0,\ldots,d_{k-1})}(d_k) & k < m \\
      \emptyset & \text{otherwise}
    \end{cases} \\[1ex]
  \SEM{(d_0, \ldots, d_{m-1};\; M)}^{\mathrm{5}} &= \SEM{M}^{\mathrm{5}}_{\key{table}(d_0,\ldots,d_{m-1})}\,\rho_0
\end{align*}

\subsubsection{Imp and CFG: what they share}

A store $\sigma$ gives each slot a meaning. An expression or tail runs in
a frame of $n$ slots. Reading a slot beyond the frame gives $\{\omega\}$,
as an unbound variable does in GlobalFun; assigning one errors, so the
construct means $\emptyset$ ($\mathit{checked}_n$). A simple expression
leaves the store as it was. With the table $T$ and $n$ fixed:
\begin{align*}
  \sigma &: \mathbb{N} \to \SET{\V} \\
  \mathit{read}_n(\sigma, s) &=
    \begin{cases} \sigma(s) & s < n \\ \{\omega\} & \text{otherwise} \end{cases} \qquad
  \mathit{checked}_n(s, D) =
    \begin{cases} D & s < n \\ \emptyset & \text{otherwise} \end{cases} \\[1ex]
  \mathcal{A}_\sigma(s) = \mathit{read}_n(\sigma, s), \quad
  \mathcal{A}_\sigma(c) &= \mathcal{B}(c), \quad
  \mathcal{A}_\sigma(\funref{k}) = T(k) \\[1ex]
  \mathcal{S}_\sigma(a) &= \mathcal{A}_\sigma(a) \\
  \mathcal{S}_\sigma(a_1(a_2, a_3)) &= (\mathcal{A}_\sigma(a_1) \app \mathcal{A}_\sigma(a_2)) \app \mathcal{A}_\sigma(a_3) \\
  \mathcal{S}_\sigma(\pair{a_1}{a_2}) &= \key{pair}(\mathcal{A}_\sigma(a_1), \mathcal{A}_\sigma(a_2)) \\
  \mathcal{S}_\sigma(\FST a) &= \key{car}(\mathcal{A}_\sigma(a)) \qquad
  \mathcal{S}_\sigma(\SND a) = \key{cdr}(\mathcal{A}_\sigma(a)) \\
  \mathcal{S}_\sigma(\tuple{a_1, \ldots, a_l}) &= \mathcal{T}(\mathcal{A}_\sigma(a_1), \ldots, \mathcal{A}_\sigma(a_l)) \qquad
  \mathcal{S}_\sigma(\prj{i} a) = \key{proj}_i(\mathcal{A}_\sigma(a)) \\
  \mathcal{S}_\sigma(\inl a) &= \key{inl}(\mathcal{A}_\sigma(a)) \qquad
  \mathcal{S}_\sigma(\inr a) = \key{inr}(\mathcal{A}_\sigma(a)) \\
  \mathcal{S}_\sigma(\prim{p}{a_1}{a_2}) &= \key{prim}_p(\mathcal{A}_\sigma(a_1), \mathcal{A}_\sigma(a_2))
\end{align*}
A function body runs in a fresh frame, with the free variables in slot
$0$ and the argument in slot $1$; main runs in a frame of zeros. Given
the code of a body, the table of a list of bodies gives each the table
of the older ones:
\begin{align*}
  \key{zeros}(t) &= \{0\} \\
  \key{frame}(X, Y)(0) = X, \quad \key{frame}(X, Y)(1) &= Y, \quad
  \key{frame}(X, Y)(t) = \{0\} \quad (t \geq 2) \\[1ex]
  \key{table}(b_0,\ldots,b_{m-1})(k) &=
    \begin{cases}
      \key{code}_{\key{table}(b_0,\ldots,b_{k-1})}(b_k) & k < m \\
      \emptyset & \text{otherwise}
    \end{cases}
\end{align*}

\subsubsection{Imp, continuation-passing}

The semantics the main proof uses. An expression runs in a frame of $n$
slots, given the store $\sigma$ as it enters and a continuation $\kappa$,
the meaning of what runs after it: $\kappa(D, \sigma')$ is the meaning of
the rest of the body, given the value $D$ of the expression and the
store $\sigma'$ it leaves. So each run of an expression passes its own
store on.
\begin{align*}
  \kappa \in \mathit{Kont} &= \SET{\V} \to (\mathbb{N} \to \SET{\V}) \to \SET{\V} \\
  \SEM{-}^{\mathrm{ImpK}}_{T,n} &: \mathit{Exp} \to (\mathbb{N} \to \SET{\V}) \to \mathit{Kont} \to \SET{\V}
    \qquad (T : \mathbb{N} \to \SET{\V}) \\
  \SEM{-}^{\mathrm{ImpK}} &: \mathit{Program} \to \SET{\V}
\end{align*}
With $T$ and $n$ fixed, writing $\SEM{e}\sigma\,\kappa$ for $\SEM{e}^{\mathrm{ImpK}}_{T,n}\sigma\,\kappa$:
\begin{align*}
  \SEM{r}\sigma\,\kappa &= \kappa(\mathcal{S}_\sigma(r), \sigma) \\
  \SEM{\Casei{a}{s}{e_1}{e_2}}\sigma\,\kappa &= \mathit{checked}_n(s,
    \mathcal{C}(\mathcal{A}_\sigma(a),\; \lambda X.\, \SEM{e_1}\ext{\sigma}{s}{X}\,\kappa,\;
      \lambda X.\, \SEM{e_2}\ext{\sigma}{s}{X}\,\kappa)) \\
  \SEM{\assign{s}{e_1}{e_2}}\sigma\,\kappa &= \SEM{e_1}\sigma\,(\lambda D\, \sigma'.\,
    \mathit{checked}_n(s, \Lambda(\lambda X.\, \SEM{e_2}\ext{\sigma'}{s}{X}\,\kappa) \app D)) \\[1ex]
  \key{halt}(D, \sigma) &= D \\
  \key{code}^{\mathrm{K}}_T((n;\; e)) &= \Lambda(\lambda X.\, \Lambda(\lambda Y.\,
    \SEM{e}^{\mathrm{ImpK}}_{T,n}\,\key{frame}(X, Y)\,\key{halt})) \\
  \SEM{(b_0, \ldots, b_{m-1};\; (n;\; e))}^{\mathrm{ImpK}} &=
    \SEM{e}^{\mathrm{ImpK}}_{\key{table}^{\mathrm{K}}(b_0,\ldots,b_{m-1}),n}\,\key{zeros}\,\key{halt}
\end{align*}
where $\key{table}^{\mathrm{K}}$ is $\key{table}$ with $\key{code}^{\mathrm{K}}$.

\subsubsection{Imp, direct}

An alternative, kept alongside: given the store as it enters, an
expression means a set of values and the store as it leaves.
\begin{align*}
  \SEM{-}^{\mathrm{ImpD}}_{T,n} &: \mathit{Exp} \to (\mathbb{N} \to \SET{\V}) \to \SET{\V} \times (\mathbb{N} \to \SET{\V}) \\
  \SEM{-}^{\mathrm{ImpD}} &: \mathit{Program} \to \SET{\V}
\end{align*}
\begin{align*}
  \mathit{checked}^{\mathrm{D}}_n(s, \sigma, r) &=
    \begin{cases} r & s < n \\ (\emptyset, \sigma) & \text{otherwise} \end{cases} \\[1ex]
  r{\downarrow}t &=
    \begin{cases} \pi_2(r)(t) & \pi_1(r) \NE \\ \emptyset & \text{otherwise} \end{cases} \\
  \key{let}_s(D, E) &= (\Lambda(\lambda X.\, \pi_1(E(X))) \app D,\;
    \lambda t.\, \Lambda(\lambda X.\, E(X){\downarrow}t) \app D) \\
  \key{case}_s(A, E_1, E_2) &= (\mathcal{C}(A,\; \lambda X.\, \pi_1(E_1(X)),\; \lambda X.\, \pi_1(E_2(X))), \\
    &\qquad \lambda t.\, \mathcal{C}(A,\; \lambda X.\, E_1(X){\downarrow}t,\; \lambda X.\, E_2(X){\downarrow}t))
\end{align*}
where $\pi_1$ and $\pi_2$ take the value and the store of a pair. The
slot $t$ of the store a let or case leaves collects slot $t$ of the runs
of its body that have a value ($r{\downarrow}t$): a run that errors
or diverges leaves no store.
With $T$ and $n$ fixed, writing $\SEM{e}\sigma$ for $\SEM{e}^{\mathrm{ImpD}}_{T,n}\sigma$:
\begin{align*}
  \SEM{r}\sigma &= (\mathcal{S}_\sigma(r), \sigma) \\
  \SEM{\Casei{a}{s}{e_1}{e_2}}\sigma &= \mathit{checked}^{\mathrm{D}}_n(s, \sigma,
    \key{case}_s(\mathcal{A}_\sigma(a),\; \lambda X.\, \SEM{e_1}\ext{\sigma}{s}{X},\; \lambda X.\, \SEM{e_2}\ext{\sigma}{s}{X})) \\
  \SEM{\assign{s}{e_1}{e_2}}\sigma &= \mathit{checked}^{\mathrm{D}}_n(s, \sigma,
    \key{let}_s(D_1,\; \lambda X.\, \SEM{e_2}\ext{\sigma_1}{s}{X}))
    \quad\text{where } (D_1, \sigma_1) = \SEM{e_1}\sigma \\[1ex]
  \key{code}^{\mathrm{D}}_T((n;\; e)) &= \Lambda(\lambda X.\, \Lambda(\lambda Y.\, \pi_1(\SEM{e}^{\mathrm{ImpD}}_{T,n}\key{frame}(X, Y)))) \\
  \SEM{(b_0, \ldots, b_{m-1};\; (n;\; e))}^{\mathrm{ImpD}} &=
    \pi_1(\SEM{e}^{\mathrm{ImpD}}_{\key{table}^{\mathrm{D}}(b_0,\ldots,b_{m-1}),n}\,\key{zeros})
\end{align*}
where $\key{table}^{\mathrm{D}}$ is $\key{table}$ with $\key{code}^{\mathrm{D}}$.

\subsubsection{CFG}

A tail runs in a frame of $n$ slots, given the meanings $B$ of the blocks
it can see, and means a set of values. The blocks of a body are given
meaning as the global definitions are: each with the meanings of the
blocks older than it, which are all it can jump to, so no fixed point is
needed. A jump to a label that is not among them means $\emptyset$.
\begin{align*}
  B &: \mathbb{N} \to (\mathbb{N} \to \SET{\V}) \to \SET{\V} \\
  \SEM{-}^{\mathrm{CFG}}_{T,n,B} &: \mathit{Tail} \to (\mathbb{N} \to \SET{\V}) \to \SET{\V} \\
  \SEM{-}^{\mathrm{CFG}} &: \mathit{Program} \to \SET{\V}
\end{align*}
With $T$, $n$, and $B$ fixed, writing $\SEM{t}\sigma$ for $\SEM{t}^{\mathrm{CFG}}_{T,n,B}\sigma$:
\begin{align*}
  \SEM{\key{return}\; r}\sigma &= \mathcal{S}_\sigma(r) \\
  \SEM{\assign{s}{r}{t}}\sigma &= \mathit{checked}_n(s,
    \Lambda(\lambda X.\, \SEM{t}\ext{\sigma}{s}{X}) \app \mathcal{S}_\sigma(r)) \\
  \SEM{\key{goto}\; l}\sigma &= B(l)(\sigma) \\
  \SEM{\key{case}\;a\;\key{of}_{s}\;\key{goto}\;l_1 \mid \key{goto}\;l_2}\sigma &=
    \mathit{checked}_n(s, \mathcal{C}(\mathcal{A}_\sigma(a),\;
      \lambda X.\, B(l_1)(\ext{\sigma}{s}{X}),\; \lambda X.\, B(l_2)(\ext{\sigma}{s}{X}))) \\[1ex]
  \key{blocks}_{T,n}(t_0, \ldots, t_{j-1})(l) &=
    \begin{cases}
      \SEM{t_l}^{\mathrm{CFG}}_{T,n,\key{blocks}_{T,n}(t_0, \ldots, t_{l-1})} & l < j \\
      \lambda \sigma.\, \emptyset & \text{otherwise}
    \end{cases} \\
  \mathit{run}_T((n;\; \vec{t};\; t)) &= \SEM{t}^{\mathrm{CFG}}_{T,n,\key{blocks}_{T,n}(\vec{t})} \\[1ex]
  \key{code}^{\mathrm{CFG}}_T(b) &= \Lambda(\lambda X.\, \Lambda(\lambda Y.\, \mathit{run}_T(b)(\key{frame}(X, Y)))) \\
  \SEM{(b_0, \ldots, b_{m-1};\; b)}^{\mathrm{CFG}} &=
    \mathit{run}_{\key{table}^{\mathrm{CFG}}(b_0,\ldots,b_{m-1})}(b)(\key{zeros})
\end{align*}
where $\key{table}^{\mathrm{CFG}}$ is $\key{table}$ with $\key{code}^{\mathrm{CFG}}$.

\subsubsection{X86}\label{sec:den-x86}

Continuation-passing, as Imp. The state of a body is its frame $\sigma$,
a store as in CFG; its registers $\rho$; and the flag $\varphi$, a set of
values: $\Left{\emptytup}$ is set, $\Right{\emptytup}$ is not set, and
$\emptyset$ is no flag. What runs from a point of a body on, $\kappa$,
means a set of values given the state. Every write is strict: a slot is
bound as in CFG, and a register or the flag holds a set $D$ only when
$D \NE$, so an instruction whose value is $\emptyset$ stops the body.
Reading an unset register or a clobbered flag gives $\emptyset$.
\begin{align*}
  \rho &: \{\rax, \rdi, \rsi\} \to \SET{\V} \qquad \varphi \in \SET{\V} \\
  \kappa \in \mathit{Kont}^{\mathrm{X86}} &= (\mathbb{N} \to \SET{\V}) \to (\{\rax, \rdi, \rsi\} \to \SET{\V}) \to \SET{\V} \to \SET{\V} \\
  B &: \mathbb{N} \to \mathit{Kont}^{\mathrm{X86}} \\
  \SEM{-}^{\mathrm{X86}}_{T,n} &: \mathit{Instr}^* \to \mathit{Kont}^{\mathrm{X86}} \to \mathit{Kont}^{\mathrm{X86}} \qquad
  \SEM{-}^{\mathrm{X86}}_{T,n,B} : \mathit{Block} \to \mathit{Kont}^{\mathrm{X86}} \\
  \SEM{-}^{\mathrm{X86}} &: \mathit{Program} \to \SET{\V}
\end{align*}
With $T$ and $n$ fixed, the operands, the strict writes, and the values
the instructions compute:
\begin{align*}
  \mathcal{O}_{\sigma\rho}(s) = \mathit{read}_n(\sigma, s), \quad
  \mathcal{O}_{\sigma\rho}(\%g) &= \rho(\%g), \quad
  \mathcal{O}_{\sigma\rho}(c) = \mathcal{B}(c), \quad
  \mathcal{O}_{\sigma\rho}(\funref{k}) = T(k) \\[1ex]
  \key{put}_s(D, \kappa)(\sigma, \rho, \varphi) &=
    \mathit{checked}_n(s, \Lambda(\lambda X.\, \kappa(\ext{\sigma}{s}{X}, \rho, \varphi)) \app D) \\
  \key{put}_{\%g}(D, \kappa)(\sigma, \rho, \varphi) &=
    \key{guard}(D;\; \kappa(\sigma, \ext{\rho}{\%g}{D}, \varphi)) \\
  \key{flag}(F, \kappa)(\sigma, \rho) &= \key{guard}(F;\; \kappa(\sigma, \rho, F)) \\[1ex]
  \key{isinl}(D) &= \{ \Left{\emptytup} \mid \exists u.\; \Left{u} \in D \}
    \cup \{ \Right{\emptytup} \mid \exists u.\; \Right{u} \in D \} \\
  \key{payload}(D) &= \{ v \mid \Left{v} \in D \lor \Right{v} \in D \} \\
  \key{object}_{\key{pair}}(D_1, D_2) &= \key{pair}(D_1, D_2), \quad
  \key{object}_{\key{tuple}}(\vec{D}) = \mathcal{T}(\vec{D}), \quad
  \key{object}_{\key{inl}}(D) = \key{inl}(D), \\
  \key{object}_{\key{inr}}(D) &= \key{inr}(D), \quad
  \key{object}_{\mathit{sh}}(\vec{D}) = \emptyset \quad \text{otherwise (arity)} \\
  \key{field}_{\key{fst}} &= \key{car}, \quad
  \key{field}_{\key{snd}} = \key{cdr}, \quad
  \key{field}_{\pi_i} = \key{proj}_i, \quad
  \key{field}_{\key{payload}} = \key{payload} \\[1ex]
  \key{called}(R)(\rax) &= R, \quad \key{called}(R)(\%g) = \emptyset \quad (\%g \neq \rax) \\
  \key{args}(X, Y)(\rax) &= \emptyset, \quad \key{args}(X, Y)(\rdi) = X, \quad \key{args}(X, Y)(\rsi) = Y
\end{align*}
Writing $\mathcal{O}$ for $\mathcal{O}_{\sigma\rho}$ and $\SEM{-}$ for $\SEM{-}^{\mathrm{X86}}_{T,n}$:
\begin{align*}
  \SEM{\key{movq}\; o, d}\kappa\,\sigma\,\rho\,\varphi &= \key{put}_d(\mathcal{O}(o), \kappa)(\sigma, \rho, \varphi) \\
  \SEM{\key{subq}\; o, d}\kappa\,\sigma\,\rho\,\varphi &=
    \key{put}_d(\key{prim}_{\key{sub}}(\mathcal{O}(d), \mathcal{O}(o)), \kappa)(\sigma, \rho, \varphi) \\
  \SEM{\key{cmpq}\; o_2, o_1}\kappa\,\sigma\,\rho\,\varphi &=
    \key{flag}(\key{prim}_{\key{eq}}(\mathcal{O}(o_1), \mathcal{O}(o_2)), \kappa)(\sigma, \rho) \\
  \SEM{\key{sete}\; d}\kappa\,\sigma\,\rho\,\varphi &= \key{put}_d(\varphi, \kappa)(\sigma, \rho, \varphi) \\
  \SEM{\key{testinl}\; o}\kappa\,\sigma\,\rho\,\varphi &=
    \key{flag}(\key{isinl}(\mathcal{O}(o)), \kappa)(\sigma, \rho) \\
  \SEM{\key{alloc}\; \mathit{sh}\; o_1, \ldots, o_p, d}\kappa\,\sigma\,\rho\,\varphi &=
    \key{put}_d(\key{object}_{\mathit{sh}}(\mathcal{O}(o_1), \ldots, \mathcal{O}(o_p)), \kappa)(\sigma, \rho, \varphi) \\
  \SEM{\key{load}\; f\; o, d}\kappa\,\sigma\,\rho\,\varphi &=
    \key{put}_d(\key{field}_f(\mathcal{O}(o)), \kappa)(\sigma, \rho, \varphi) \\
  \SEM{\key{callq}\; o}\kappa\,\sigma\,\rho\,\varphi &= \key{guard}(R;\; \kappa(\sigma, \key{called}(R), \emptyset))
    \quad\text{where } R = (\mathcal{O}(o) \app \rho(\rdi)) \app \rho(\rsi) \\[1ex]
  \SEM{i_1; \ldots; i_p}\kappa &= \SEM{i_1}(\cdots(\SEM{i_p}\kappa)) \\
  \SEM{\key{retq}}_B\,\sigma\,\rho\,\varphi &= \rho(\rax) \\
  \SEM{\key{jmp}\; l}_B &= B(l) \\
  \SEM{\key{je}\; l_1, l_2}_B\,\sigma\,\rho\,\varphi &=
    \mathcal{C}(\varphi,\; \lambda X.\, B(l_1)(\sigma, \rho, \varphi),\; \lambda X.\, B(l_2)(\sigma, \rho, \varphi)) \\
  \SEM{i_1; \ldots; i_p;\; j}^{\mathrm{X86}}_{T,n,B} &= \SEM{i_1; \ldots; i_p}(\SEM{j}_B)
\end{align*}
A call applies the code pointer to $\rdi$ and then to $\rsi$, as Imp's
call does, puts the result in $\rax$, and clobbers $\rdi$, $\rsi$ and
the flag; the caller's frame is as it was. The blocks and the table are
as in CFG. A global definition runs in a fresh frame of $\max(2, n)$
slots holding $0$, with its parameters in $\rdi$ and $\rsi$; main runs
in a frame of $n$ zeros with $\{0\}$ in $\rax$.
\begin{align*}
  \key{blocks}_{T,n}(\beta_0, \ldots, \beta_{q-1})(l) &=
    \begin{cases}
      \SEM{\beta_l}^{\mathrm{X86}}_{T,n,\key{blocks}_{T,n}(\beta_0, \ldots, \beta_{l-1})} & l < q \\
      \lambda \sigma\, \rho\, \varphi.\, \emptyset & \text{otherwise}
    \end{cases} \\
  \mathit{run}_{T,n'}((n;\; \vec{\beta};\; \beta)) &= \SEM{\beta}^{\mathrm{X86}}_{T,n',\key{blocks}_{T,n'}(\vec{\beta})} \\[1ex]
  \key{code}^{\mathrm{X86}}_T((n;\; \vec{\beta};\; \beta)) &=
    \Lambda(\lambda X.\, \Lambda(\lambda Y.\, \mathit{run}_{T,\max(2,n)}((n;\; \vec{\beta};\; \beta))(\key{zeros}, \key{args}(X, Y), \emptyset))) \\
  \SEM{(b_0, \ldots, b_{m-1};\; (n;\; \vec{\beta};\; \beta))}^{\mathrm{X86}} &=
    \mathit{run}_{\key{table}^{\mathrm{X86}}(b_0,\ldots,b_{m-1}),n}((n;\; \vec{\beta};\; \beta))(\key{zeros}, \key{called}(\{0\}), \emptyset)
\end{align*}
where $\key{table}^{\mathrm{X86}}$ is $\key{table}$ with $\key{code}^{\mathrm{X86}}$.

%===============================================================================
\subsection{Relations}

\subsubsection{Simplify}

\begin{lemma}[$\beta$ for a Value]\label{lem:beta-value}
If $\rho(x) \NE$ for every $x$, and $D$ is nonempty and consistent, then
\[
  \SEM{\lam{x} N}^{\mathrm{ISWIM}}\rho \star D \simeq \SEM{N}^{\mathrm{ISWIM}}\ext{\rho}{x}{D} ,
\]
and likewise for $\key{case}$ on an $\inl$ or $\inr$ of a nonempty, consistent set. By
continuity an element of the right needs only a finite part of $\SEM{M}\rho$,
which is consistent and nonempty, as $\Lambda$ requires.
\end{lemma}

\begin{lemma}[Correctness of Simplify]\label{lem:simplify}
Let $\rho$ and $\rho'$ be nonempty and consistent, and for every $x$ with $E(x) = (T, K)$, let
$\rho(x) \simeq \SEM{T}\rho'$ and $K$ describe $\rho(x)$: if $K$ is a value,
$\rho(x) \simeq \SEM{K}\rho'$, and if $K = \langle \lam{y} N, E_f \rangle$, then
$\rho(x) \simeq \SEM{\lam{y} N}\rho_f$ for some nonempty, consistent $\rho_f$ of
which the same holds for $E_f$. Then
\[
  \SEM{M}^{\mathrm{ISWIM}}\rho \simeq \SEM{\pass{simp}_E(M)}^{\mathrm{ISWIM}}\rho' ,
\]
and $\mathcal{K}_E(M)$ describes $\SEM{M}\rho$. In particular
$\SEM{M}\rho_0 \simeq \SEM{\pass{simplify}(M)}\rho_0$.
\end{lemma}

\subsubsection{CC: Enclose}

\begin{lemma}[Correctness of Enclose]\label{lem:enclose}
If $\rho_1(x) \simeq \rho_2(x)$ for every $x \in \FV(M)$
and $\rho_2(x) \NE$ for every $x$, then
\[
  \SEM{M}^{\mathrm{ISWIM}}\rho_1 \simeq \SEM{\pass{enclose}(M)}^{\mathrm{2}}\rho_2 .
\]
\end{lemma}

\subsubsection{CC: Optimize}

\begin{lemma}[Correctness of Optimize]\label{lem:optimize}
If $\rho(x) \simeq \rho'(x)$ for every $x \in \FV(M)$
and $\rho(x) \NE$ for every $x$, then
\[
  \SEM{M}^{\mathrm{2}}\rho \simeq \SEM{\pass{optimize}(M)}^{\mathrm{2}}\rho' .
\]
\end{lemma}

\subsubsection{CC: Concretize}

\begin{align*}
  \SEM{x}_r\rho &= \rho(x) \\
  \SEM{\prj{i} x}_r\rho &= \key{proj}_i(\rho(x)) \\[1ex]
  \rho_2 \sim_\varphi \rho_3 &\iff \forall x.\; \rho_2(x) \simeq \SEM{\varphi(x)}_r\rho_3
\end{align*}
\begin{lemma}[Correctness of Concretize]\label{lem:concretize}
If $\consis(\rho_3)$ and $\rho_2 \sim_\varphi \rho_3$, then
\[
  \SEM{M}^{\mathrm{2}}\rho_2 \simeq \SEM{\pass{concretize}_\varphi(M)}^{\mathrm{3}}\rho_3 .
\]
\end{lemma}

\subsubsection{CC: Delay}

\begin{lemma}[Continuity]\label{lem:continuity}
$\SEM{M}^{\mathrm{3}}$ and $\SEM{M}^{\mathrm{4}}$ are continuous.
\end{lemma}

\begin{lemma}[Delayed Application]\label{lem:let-app}
If $f \notin \FV(N)$ and $\consis(\SEM{L}^{\mathrm{4}}\rho)$, then
\[
  \SEM{\Let f = L \In (\FST f)(\SND f, N)}^{\mathrm{4}}\rho
  \simeq \SEM{L}^{\mathrm{4}}\rho \odot \SEM{N}^{\mathrm{4}}\rho .
\]
\end{lemma}

\paragraph{Preservation.}
The relation $v \sqsubseteq D'$ between a EnvClos value $v$ and a set $D'$ of
FlatClos values is defined by recursion on $v$.
\begin{align*}
  \nu \sqsubseteq D' &\iff D' \NE \\
  (U \mapsto w) \sqsubseteq D' &\iff D' \NE \land
     \forall E'.\; U \sqsubseteq E' \implies w \sqsubseteq D' \odot E' \\
  \tup{i}{d} \sqsubseteq D' &\iff d \sqsubseteq \key{proj}_i(D') \\
  \Left{d} \sqsubseteq D' &\iff d \sqsubseteq \{ d' \mid \Left{d'} \in D' \} \\
  \Right{d} \sqsubseteq D' &\iff d \sqsubseteq \{ d' \mid \Right{d'} \in D' \} \\
  v \sqsubseteq D' &\iff v \in D' \quad \text{otherwise} \\[1ex]
  U \sqsubseteq E' &\iff \forall u \in U.\; u \sqsubseteq E' \\
  \rho \sqsubseteq \rho' &\iff \forall x.\; \forall v \in \rho(x).\; v \sqsubseteq \rho'(x)
\end{align*}

\begin{lemma}[Delay Preserves]\label{lem:delay-preserve}
If $\rho'(x) \NE$ for every $x$, $\consis(\rho')$, and $\rho \sqsubseteq \rho'$, then for every $v$,
\[
  v \in \SEM{M}^{\mathrm{3}}\rho \implies v \sqsubseteq \SEM{\pass{delay}(M)}^{\mathrm{4}}\rho' .
\]
\end{lemma}

\paragraph{Depth.}
\begin{align*}
  \mathit{depth}(c) = \mathit{depth}(\nu) = \mathit{depth}(\omega) = \mathit{depth}(\emptytup) &= 0 \\
  \mathit{depth}(V \mapsto w) &= 1 + \max(\mathit{depth}(V), \mathit{depth}(w)) \\
  \mathit{depth}(\pairv{u}{v}) &= 1 + \max(\mathit{depth}(u), \mathit{depth}(v)) \\
  \mathit{depth}(\tup{i}{d}) = \mathit{depth}(\Left{d}) = \mathit{depth}(\Right{d})
    &= 1 + \mathit{depth}(d) \\
  \mathit{depth}(\emptyset) &= 0 \\
  \mathit{depth}(V) &= \max_{v \in V} \mathit{depth}(v)
\end{align*}

\paragraph{Reflection.}
The relation $\Rel{V'}{D}$ between a finite set $V'$ of FlatClos values and
a set $D$ of EnvClos values, defined by recursion on $\mathit{depth}(V')$,
where $t$ ranges over the tags $\key{tup}_i$, $\key{left}$, and $\key{right}$,
and $U$ and $W$ range over finite sets of FlatClos values.
\begin{align*}
  \mathit{Entry}(V', U_0, w) &\iff \exists F \NE.\; \exists d.\;
     \pairv{F \mapsto (U_0 \mapsto w)}{d} \in V',\; F \subseteq \key{cdr}(V') \\
  \mathit{Calls}(V', U, W) &\iff
  \begin{array}[t]{l}
    \forall w \in W.\; \exists U_0 \NE.\; \mathit{Entry}(V', U_0, w) \land U_0 \subseteq U \\
    \land\; \forall u \in U.\; \exists U_0\, w.\; \mathit{Entry}(V', U_0, w) \land u \in U_0
  \end{array} \\[1ex]
  \Rel{V'}{D} &\iff
  \begin{array}[t]{l}
    \forall c \in V'.\; c \in D \\
    \land\; (V' \NE \implies D \NE) \\
    \land\; \forall U\, W\, E.\;
       \mathit{Calls}(V', U, W) \land \consis(U) \land \Rel{U}{E} \implies \Rel{W}{D \app E} \\
    \land\; \forall t\, W.\; W \subseteq \{ d \mid t(d) \in V' \}
       \implies \Rel{W}{\{ d \mid t(d) \in D \}}
  \end{array} \\[1ex]
  \rho' \rhd \rho &\iff \forall x\, V'.\; V' \subseteq \rho'(x) \implies \Rel{V'}{\rho(x)}
\end{align*}

\begin{lemma}[Delay Reflects]\label{lem:delay-reflect}
If $\rho(x) \NE$ for every $x$, $\consis(\rho)$, $\consis(\rho')$, and $\rho' \rhd \rho$,
then for every finite $V'$,
\[
  V' \subseteq \SEM{\pass{delay}(M)}^{\mathrm{4}}\rho' \implies \Rel{V'}{\SEM{M}^{\mathrm{3}}\rho} .
\]
\end{lemma}
For $\key{let}$, both directions use $\beta$ for a $\key{let}$ by continuity
(an element of $\SEM{N}\ext{\rho}{x}{D}$ needs only a finite, nonempty,
consistent part of $D$); reflection first gathers the finite parts of the
operand that the elements of $V'$ use into one.

\subsubsection{CC: Globalize}

\begin{align*}
  \mathit{Extends}(\vec{d}, \vec{d}') &\iff \vec{d} \text{ is a prefix of } \vec{d}' \\
  \mathit{RB}_n(M) &\iff \text{every } \funref{k} \text{ in } M \text{ has } k < n \\
  \mathit{Agree}_n(T, T') &\iff \forall k < n.\; T(k) \simeq T'(k)
\end{align*}
\begin{lemma}[Well-Formed Tables]\label{lem:glob-wf}
If $\pass{glob}(M, \vec{d}) = (M', \vec{d}')$, then
$\mathit{Extends}(\vec{d}, \vec{d}')$ and $\mathit{RB}_{|\vec{d}'|}(M')$.
\end{lemma}

\begin{lemma}[Table Agreement]\label{lem:agree}
If $\mathit{Agree}_n(T, T')$, $\mathit{RB}_n(M)$, and $\rho \simeq \rho'$, then
$\SEM{M}^{\mathrm{5}}_T\rho \simeq \SEM{M}^{\mathrm{5}}_{T'}\rho'$.
\end{lemma}

\begin{lemma}[Correctness of Glob]\label{lem:glob}
If $\pass{glob}(M, \vec{d}) = (M', \vec{d}')$ and $\rho \simeq \rho'$, then
$\SEM{M}^{\mathrm{4}}\rho \simeq \SEM{M'}^{\mathrm{5}}_{\key{table}(\vec{d}')}\rho'$.
\end{lemma}

\begin{lemma}[Correctness of Globalize]\label{lem:globalize}
$\SEM{M}^{\mathrm{4}}\rho_0 \simeq \SEM{\pass{globalize}(M)}^{\mathrm{5}}$.
\end{lemma}

\subsubsection{Remove Complex Operands}

\begin{lemma}[Continuity of GlobalFun]\label{lem:continuity5}
For every table $T$, $\SEM{M}^{\mathrm{5}}_T$ is monotone and continuous.
\end{lemma}

\begin{lemma}[Let Is Strict]\label{lem:let-strict}
If $\SEM{\Let x = M \In N}^{\mathrm{5}}_T\rho \NE$, then $\SEM{M}^{\mathrm{5}}_T\rho \NE$.
\end{lemma}

\begin{lemma}[Let]\label{lem:let-beta}
If $\rho(x) \NE$ for every $x$, $\SEM{M}^{\mathrm{5}}_T\rho \NE$, and
$\consis(\SEM{M}^{\mathrm{5}}_T\rho)$, then
\[
  \SEM{\Let x = M \In N}^{\mathrm{5}}_T\rho \simeq
  \SEM{N}^{\mathrm{5}}_T\ext{\rho}{x}{\SEM{M}^{\mathrm{5}}_T\rho} .
\]
\end{lemma}

\paragraph{Strict operators.}
An operator $O$ on $n$ sets is strict when
\begin{align*}
  \mathit{Strict}(O) &\iff
    (\forall \bar{D}\, \bar{E}.\; \bar{D} \simeq \bar{E} \implies O(\bar{D}) \simeq O(\bar{E}))
    \land (\forall \bar{D}.\; O(D_1, \ldots, D_n) \NE \implies \forall i.\; D_i \NE) .
\end{align*}
Each construct whose operands rco names means a strict operator on the
meanings of its operands:
\begin{align*}
  O_{L(M, N)}(D, E, F) &= (D \app E) \app F &
  O_{\pair{M}{N}}(D, E) &= \key{pair}(D, E) &
  O_{\FST M}(D) &= \key{car}(D) \\
  O_{\SND M}(D) &= \key{cdr}(D) &
  O_{\tuple{M_1, \ldots, M_n}}(\bar{D}) &= \mathcal{T}(\bar{D}) &
  O_{\prj{i} M}(D) &= \key{proj}_i(D) \\
  O_{\inl M}(D) &= \key{inl}(D) &
  O_{\inr M}(D) &= \key{inr}(D) &
  O_{\prim{p}{L}{M}}(D, E) &= \key{prim}_p(D, E)
\end{align*}
and $O(D) = \mathcal{C}(D,\; \lambda X.\, \SEM{M}^{\mathrm{5}}_T\ext{\rho}{x}{X},\;
\lambda Y.\, \SEM{N}^{\mathrm{5}}_T\ext{\rho}{y}{Y})$ for
$\Case{L}{x}{M}{y}{N}$.

\begin{lemma}[Let Chain]\label{lem:let-chain}
Let $O$ be strict, $(\vec{b};\; a_1, \ldots, a_n) = \pass{atomize}(M_1, \ldots, M_n)$,
$\consis(T(k))$ for every $k$, $\rho(x) \NE$ for every $x$, and $\consis(\rho)$.
If, for every $\rho'$ that extends
$\rho$ with the variables of $\vec{b}$,
$\SEM{X}^{\mathrm{5}}_T\rho' \simeq O(\SEM{a_1}^{\mathrm{5}}_T\rho', \ldots, \SEM{a_n}^{\mathrm{5}}_T\rho')$, then
\[
  \SEM{\pass{lets}(\vec{b}, X)}^{\mathrm{5}}_T\rho \simeq
  O(\SEM{M_1}^{\mathrm{5}}_T\rho, \ldots, \SEM{M_n}^{\mathrm{5}}_T\rho) .
\]
\end{lemma}

\begin{theorem}[Correctness of RCO]\label{thm:rco}
If $\consis(T(k))$ for every $k$, $\rho(x) \NE$ for every $x$, and $\consis(\rho)$, then
$\SEM{\pass{rco}(M)}^{\mathrm{5}}_T\rho \simeq \SEM{M}^{\mathrm{5}}_T\rho$.
For every GlobalFun program $P$,
\[
  \SEM{\pass{rco}(P)}^{\mathrm{5}} \simeq \SEM{P}^{\mathrm{5}} .
\]
\end{theorem}
The theorem needs no scoping: an error and divergence both mean
$\emptyset$, an unbound variable means $\rho_0(x) = \{\omega\}$, and the
order in which operands are evaluated is invisible.

\subsubsection{Lift Locals}

The store agrees with the environment on the slots that $\varphi$ gives
the variables in scope, and the variables out of scope mean
$\{\omega\}$, as slot $n$, beyond the frame, does:
\begin{align*}
  \mathit{Agree}_{\varphi}(\sigma, \rho) &\iff
    \forall x \in \mathrm{dom}(\varphi).\; \sigma(\varphi(x)) \simeq \rho(x) \\
  \mathit{Scope}_{\varphi}(\rho) &\iff
    \forall x \notin \mathrm{dom}(\varphi).\; \rho(x) \simeq \{\omega\}
\end{align*}

\paragraph{Against the continuation-passing semantics.}
Lifted code is run with a continuation that means what the enclosing
lets make of the value, given any store that agrees with the
environment: nothing at the end of a body, and $\Lambda(G) \app D$ for
the right-hand side of a let whose body and continuation mean $G$.
\begin{align*}
  K_{\cdot}(D) &= D \qquad K_G(D) = \Lambda(G) \app D
\end{align*}
The commuting conversions are then on the GlobalFun side. Each needs the
bound set to be consistent, so that the finitely many runs that
$\Lambda(G) \app -$ uses can be joined into one.

\begin{lemma}[Let Law]\label{lem:let-law}
If $\consis(D)$ and $F$ is monotone, then
\[
  K(\Lambda(F) \app D) \simeq \Lambda(\lambda X.\, K(F(X))) \app D .
\]
\end{lemma}

\begin{lemma}[Case Law]\label{lem:case-law}
If $\consis(A)$ and $F_1$ and $F_2$ are monotone, then
\[
  K(\mathcal{C}(A, F_1, F_2)) \simeq \mathcal{C}(A, \lambda X.\, K(F_1(X)), \lambda X.\, K(F_2(X))) .
\]
\end{lemma}

\begin{lemma}[Correctness of Lift]\label{lem:lift-sem}
Let $T(k) \simeq T'(k)$ and $\consis(T'(k))$ for every $k$, and $M$ be a
monadic GlobalFun term. If $d + \mathit{height}(M) \leq n$, $\varphi(x) < d$
for every $x \in \mathrm{dom}(\varphi)$, $\mathit{Scope}_{\varphi}(\rho)$,
$\mathit{Agree}_{\varphi}(\sigma, \rho)$, $\consis(\rho)$, and
$\kappa(D, \sigma') \simeq K(D)$ for every $D$ and every $\sigma'$ with
$\mathit{Agree}_{\varphi}(\sigma', \rho)$, then
\[
  \SEM{\pass{lift}_{n,\varphi,d}(M)}^{\mathrm{ImpK}}_{T,n}\sigma\,\kappa
  \simeq K(\SEM{M}^{\mathrm{5}}_{T'}\rho) .
\]
\end{lemma}

\begin{lemma}[Lifted Tables]\label{lem:lift-table}
For every $k$,
$\key{table}^{\mathrm{K}}(\pass{lift}(d_0), \ldots, \pass{lift}(d_{m-1}))(k) \simeq \key{table}(d_0, \ldots, d_{m-1})(k)$.
\end{lemma}

\begin{theorem}[Correctness of Lift Locals]\label{thm:lift}
For every monadic GlobalFun program $P$,
\[
  \SEM{\pass{lift}(P)}^{\mathrm{ImpK}} \simeq \SEM{P}^{\mathrm{5}} .
\]
\end{theorem}

\paragraph{Against the direct semantics.}
The store that lift leaves agrees with the same environment, when there
is a value. As in the operational proofs, this is relative to $\rho$,
not to the store $\sigma$, so after a binder, whose slot is in
$\varphi[x \mapsto d]$ but not in $\varphi$, only the slot is forgotten.

\begin{lemma}[Correctness of Lift, Direct]\label{lem:lift-sem-d}
Let $T(k) \simeq T'(k)$ for every $k$, and $M$ be a monadic GlobalFun term.
If $d + \mathit{height}(M) \leq n$, $\varphi(x) < d$ for every
$x \in \mathrm{dom}(\varphi)$, $\mathit{Scope}_{\varphi}(\rho)$, and
$\mathit{Agree}_{\varphi}(\sigma, \rho)$, then
\[
  D \simeq \SEM{M}^{\mathrm{5}}_{T'}\rho
  \quad\text{and}\quad
  (D \NE \implies \mathit{Agree}_{\varphi}(\sigma', \rho)),
  \qquad\text{where } (D, \sigma') = \SEM{\pass{lift}_{n,\varphi,d}(M)}^{\mathrm{ImpD}}_{T,n}\sigma .
\]
\end{lemma}

\begin{theorem}[Correctness of Lift Locals, Direct]\label{thm:lift-d}
For every monadic GlobalFun program $P$,
\[
  \SEM{\pass{lift}(P)}^{\mathrm{ImpD}} \simeq \SEM{P}^{\mathrm{5}} .
\]
\end{theorem}

\subsubsection{Explicate Control}\label{sec:den-explicate}

The pass threads the list of blocks: from the blocks $\vec{t}$,
$\pass{tail}(e, \vec{t}) = (t;\, \vec{t}')$ and
$\pass{assign}(e, s, k, \vec{t}) = (t;\, \vec{t}')$ give the tail $t$
and the blocks $\vec{t}'$, $\vec{t}$ followed by the blocks the pass
adds. A list extends another with newer blocks,
\begin{align*}
  \vec{t} \preceq \vec{t}' &\iff \vec{t} \text{ is a prefix of } \vec{t}' ,
\end{align*}
and the blocks keep their labels: if $\vec{t}, t \preceq \vec{t}'$, then
$\key{blocks}_{T,n}(\vec{t}')(|\vec{t}|) = \SEM{t}^{\mathrm{CFG}}_{T,n,\key{blocks}_{T,n}(\vec{t})}$.
The pass only adds blocks: if $\pass{assign}(e, s, k, \vec{t}) = (t;\, \vec{t}')$
or $\pass{tail}(e, \vec{t}) = (t;\, \vec{t}')$, then $\vec{t} \preceq \vec{t}'$.
So a tail made at the blocks $\vec{t}$ means the same with newer blocks
when it is stable:
\begin{align*}
  \mathit{Stable}_{\vec{t}}(k) &\iff \forall \vec{t}'' \succeq \vec{t}.\; \forall \sigma.\;
    \SEM{k}^{\mathrm{CFG}}_{T,n,\key{blocks}_{T,n}(\vec{t}'')}\sigma \simeq
    \SEM{k}^{\mathrm{CFG}}_{T,n,\key{blocks}_{T,n}(\vec{t})}\sigma \\
  \key{assignK}_s(K)(D, \sigma) &= \mathit{checked}_n(s, \Lambda(\lambda X.\, K(\ext{\sigma}{s}{X})) \app D)
\end{align*}
$\key{assignK}_s(K)$ assigns $s$ the value and then runs $K$, as
$\assign{s}{r}{k}$ does in CFG. The pass is in continuation-passing
style, and so is the semantics of Imp: each case below is a congruence.

\begin{lemma}[Continuation-Passing]\label{lem:cps}
Let $T(k) \simeq T'(k)$ for every $k$, and $\pass{assign}(e, s, k, \vec{t}) = (t;\, \vec{t}')$.
If $\mathit{Stable}_{\vec{t}}(k)$, $\vec{t}' \preceq \vec{t}''$, and
$\kappa(D, \sigma') \simeq \key{assignK}_s(\SEM{k}^{\mathrm{CFG}}_{T,n,\key{blocks}_{T,n}(\vec{t})})(D, \sigma')$
for every $D$ and $\sigma'$, then
\[
  \SEM{t}^{\mathrm{CFG}}_{T,n,\key{blocks}_{T,n}(\vec{t}'')}\sigma \simeq
  \SEM{e}^{\mathrm{ImpK}}_{T',n}\sigma\,\kappa .
\]
\end{lemma}

\begin{lemma}[Tail]\label{lem:explicate-tail}
Let $T(k) \simeq T'(k)$ for every $k$, and $\pass{tail}(e, \vec{t}) = (t;\, \vec{t}')$.
If $\vec{t}' \preceq \vec{t}''$, then
\[
  \SEM{t}^{\mathrm{CFG}}_{T,n,\key{blocks}_{T,n}(\vec{t}'')}\sigma \simeq
  \SEM{e}^{\mathrm{ImpK}}_{T',n}\sigma\,\key{halt} .
\]
\end{lemma}

\begin{theorem}[Correctness of Explicate Control]\label{thm:explicate}
For every Imp program $P$,
\[
  \SEM{\pass{explicate}(P)}^{\mathrm{CFG}} \simeq \SEM{P}^{\mathrm{ImpK}} .
\]
\end{theorem}

\paragraph{Against the direct semantics.}
Here the pass is a commuting conversion: in Imp, the store after
$\assign{s}{e}{k}$ collects the runs of $e$, while in CFG $k$ runs after
each run of $e$ on its own. The tail that $\pass{assign}(e, s, k, \vec{t})$
makes means
$\mathit{checked}_n(s, \Lambda(\lambda X.\, K(\ext{\pi_2(r)}{s}{X})) \app \pi_1(r))$
for $r = \SEM{e}^{\mathrm{ImpD}}\sigma$ and $K$ the meaning of $k$,
by laws for nested assignment and for case whose proofs join the
finitely many runs $K$ uses into one. That needs consistent stores, and
CFG tails that preserve directed unions of stores.

\begin{theorem}[Correctness of Explicate Control, Direct]\label{thm:explicate-d}
For every Imp program $P$,
\[
  \SEM{\pass{explicate}(P)}^{\mathrm{CFG}} \simeq \SEM{P}^{\mathrm{ImpD}} .
\]
\end{theorem}

\subsubsection{Select Instructions}\label{sec:den-select}

The X86 frame is the CFG store, with any registers and flag. Write
$\pass{select}^{+}(b)$ for the body $b$ selected with the prologue, as
for a global definition. A case branches on the flag that $\key{testinl}$
sets; that is $\mathcal{C}$ on the tested set when no slot holds both an
$\key{inl}$ and an $\key{inr}$:
\begin{align*}
  \mathit{Tagged}(D) &\iff \neg\exists u\, v.\; \Left{u} \in D \land \Right{v} \in D \\
  \mathit{Inv}(\sigma) &\iff \forall s.\; \mathit{Tagged}(\sigma(s))
\end{align*}
Fix tables $T$ (of X86) and $T'$ (of CFG) with $T(k) \simeq T'(k)$ and
$\mathit{Tagged}(T(k))$ for every $k$, and $n$. The X86 $\kappa$ that
runs after $d$ is written fits the CFG $K$, and $\key{after}_d(K)$
writes $d$ in CFG:
\begin{align*}
  \mathit{Fits}_s(\kappa, K) &\iff \forall \sigma.\; \mathit{Inv}(\sigma) \implies
    \forall \rho\, \varphi.\; \kappa(\sigma, \rho, \varphi) \simeq K(\sigma) \\
  \mathit{Fits}_{\%g}(\kappa, K) &\iff \forall \sigma\, \rho\, \varphi.\; \kappa(\sigma, \rho, \varphi) \simeq \rho(\%g) \\
  \key{after}_s(K)(\sigma, D) &= \mathit{checked}_n(s, \Lambda(\lambda X.\, K(\ext{\sigma}{s}{X})) \app D) \qquad
  \key{after}_{\%g}(K)(\sigma, D) = D
\end{align*}
Each case below is a congruence. The instructions compute $r$ through
registers and only the final write to $d$ binds, as CFG does; the writes
before it are strict, but the value of $r$ is nonempty only when theirs
are, so they change nothing.

\begin{lemma}[Select a Simple Expression]\label{lem:sel}
If $\mathit{Fits}_d(\SEM{\vec{\imath}}^{\mathrm{X86}}_{T,n}\kappa, K)$ and $\mathit{Inv}(\sigma)$, then
\[
  \SEM{\pass{sel}(r, d); \vec{\imath}}^{\mathrm{X86}}_{T,n}\kappa\,\sigma\,\rho\,\varphi
  \simeq \key{after}_d(K)(\sigma, \mathcal{S}_\sigma(r)) ,
\]
with $\mathcal{S}$ that of $T'$.
\end{lemma}

\begin{lemma}[Case]\label{lem:case-core}
If $A \simeq A'$, $\mathit{Tagged}(A)$, and $F_i(V) \simeq G_i(V)$ for
every $V$ with $V \approx V$ ($i = 1, 2$), then
\[
  \Lambda(\lambda X.\, \mathcal{C}(\key{isinl}(A),\; \lambda Z.\, F_1(X),\; \lambda Z.\, F_2(X))) \app \key{payload}(A)
  \simeq \mathcal{C}(A', G_1, G_2) .
\]
\end{lemma}
The blocks $B$ of X86 and $B'$ of CFG agree,
\[
  B \mathrel{\mathit{BR}} B' \iff \forall l\, \sigma.\; \mathit{Inv}(\sigma) \implies
    \forall \rho\, \varphi.\; B(l)(\sigma, \rho, \varphi) \simeq B'(l)(\sigma) .
\]

\begin{lemma}[Tail and Blocks]\label{lem:sel-tail}
If $B \mathrel{\mathit{BR}} B'$ and $\mathit{Inv}(\sigma)$, then
\[
  \SEM{\pass{sel}(t)}^{\mathrm{X86}}_{T,n,B}\,\sigma\,\rho\,\varphi \simeq \SEM{t}^{\mathrm{CFG}}_{T',n,B'}\,\sigma .
\]
For all tails $\vec{t}$,
$\key{blocks}_{T,n}(\pass{sel}(\vec{t})) \mathrel{\mathit{BR}} \key{blocks}_{T',n}(\vec{t})$.
\end{lemma}

The prologue binds finite parts of the parameters, where CFG runs the
body with them whole, so it needs that the meaning of a CFG tail is
monotone in the store:

\begin{lemma}[Monotone]\label{lem:cfg-mono}
If $\sigma(s) \subseteq \sigma'(s)$ for every $s$, then
\[
  \SEM{t}^{\mathrm{CFG}}_{T',n,\key{blocks}_{T',n}(\vec{t})}\sigma \subseteq
  \SEM{t}^{\mathrm{CFG}}_{T',n,\key{blocks}_{T',n}(\vec{t})}\sigma' .
\]
\end{lemma}

\begin{lemma}[Start]\label{lem:sel-start}
Let $b = (n;\; \vec{t};\; t)$ with $2 \leq n$. If $V \NE$, $U \NE$,
$V \approx V$, and $U \approx U$, then
\[
  \mathit{run}_{T,n}(\pass{select}^{+}(b))(\key{zeros}, \key{args}(V, U), \emptyset)
  \simeq \mathit{run}_{T'}(b)(\key{frame}(V, U)) .
\]
So $\key{code}^{\mathrm{X86}}_T(\pass{select}^{+}(b)) \simeq \key{code}^{\mathrm{CFG}}_{T'}(b)$.
\end{lemma}

\begin{theorem}[Correctness of Select Instructions]\label{thm:select}
If every global definition $(n_k;\; \vec{t}_k;\; t_k)$ of the CFG
program $P$ has $2 \leq n_k$, then
\[
  \SEM{\pass{select}(P)}^{\mathrm{X86}} \simeq \SEM{P}^{\mathrm{CFG}} .
\]
\end{theorem}
The global definitions of $\pass{compile\text{-}cfg}(M)$ have at least
2 slots: lift locals gives each $2 + h$ slots.

\subsubsection{End to End}

\begin{align*}
  D \obseq D' &\iff (\forall c.\; c \in D \iff c \in D') \land (D \NE \iff D' \NE)
\end{align*}

\begin{theorem}[Compiler Correctness]\label{thm:compiler}
For every ISWIM program $M$,
\[
  \SEM{M}^{\mathrm{ISWIM}}\rho_0 \obseq \SEM{\pass{compile}(M)}^{\mathrm{5}} .
\]
\end{theorem}

\begin{theorem}[Correctness of Compile-Mon]\label{thm:compile-mon}
For every ISWIM program $M$,
\[
  \SEM{M}^{\mathrm{ISWIM}}\rho_0 \obseq \SEM{\pass{compile\text{-}mon}(M)}^{\mathrm{5}} .
\]
\end{theorem}

\begin{theorem}[Correctness of Compile-Imp]\label{thm:compile-imp}
For every ISWIM program $M$,
\[
  \SEM{M}^{\mathrm{ISWIM}}\rho_0 \obseq \SEM{\pass{compile\text{-}imp}(M)}^{\mathrm{ImpK}} ,
\]
by Theorems~\ref{thm:compile-mon} and~\ref{thm:lift}, and likewise for
$\SEM{-}^{\mathrm{ImpD}}$, by Theorem~\ref{thm:lift-d}.
\end{theorem}

\begin{theorem}[Correctness of Compile-CFG]\label{thm:compile-cfg}
For every ISWIM program $M$,
\[
  \SEM{M}^{\mathrm{ISWIM}}\rho_0 \obseq \SEM{\pass{compile\text{-}cfg}(M)}^{\mathrm{CFG}} ,
\]
by Theorems~\ref{thm:compile-imp} and~\ref{thm:explicate}, and likewise
by Theorems~\ref{thm:lift-d} and~\ref{thm:explicate-d}.
\end{theorem}

\begin{theorem}[Correctness of Compile-X86]\label{thm:compile-x86}
For every ISWIM program $M$,
\[
  \SEM{M}^{\mathrm{ISWIM}}\rho_0 \obseq \SEM{\pass{compile\text{-}x86}(M)}^{\mathrm{X86}} ,
\]
by Theorems~\ref{thm:compile-cfg} and~\ref{thm:select}.
\end{theorem}

%===============================================================================
\section{CEK Machine + Simulation}\label{sec:u-cek}

\subsection{Observations and Machines}

\subsubsection{Observations}\label{sec:obs}

\[
  o ::= c \mid \key{fun} \mid \key{tuple}(o_1, \ldots, o_n) \mid \inl o \mid \inr o \mid \error
\]

\subsubsection{Machines}

A machine is a step function, which either moves to the next state or stops with an observation.
\begin{align*}
  \mathit{Next}(S) &::= \nxt s \mid \stp o \qquad (s \in S) \\
  L &= (S_L,\; \step{L} : S_L \to \mathit{Next}(S_L),\; \start{L} \in S_L)
\end{align*}
\begin{align*}
  s \lto_L s' &\iff \step{L}(s) = \nxt s' \\
  \Final_L(s, o) &\iff \step{L}(s) = \stp o \\
  \Term_L(o) &\iff \exists s.\; \start{L} \lto_L^* s \land \Final_L(s, o) \\
  \Div_L &\iff \forall n.\; \exists s.\; \start{L} \lto_L^n s
\end{align*}
where $\lto_L^*$, $\lto_L^+$, and $\lto_L^n$ are zero or more, one or more, and exactly $n$ steps.

\begin{lemma}[Determinism]\label{lem:det}
If $s \lto_L s_1$ and $s \lto_L s_2$, then $s_1 = s_2$.
If $\Final_L(s, o)$ and $\Final_L(s, o')$, then $o = o'$.
If $\Final_L(s, o)$, then there is no $s'$ with $s \lto_L s'$.
\end{lemma}

\begin{lemma}[Divergence]\label{lem:div}
If $\Div_L$, then not $\Term_L(o)$.
If $\Final_L(s, o)$ for no $o$ and no $s$ with $\start{L} \lto_L^* s$, then $\Div_L$.
\end{lemma}

%-------------------------------------------------------------------------------
\subsection{The Machines}\label{sec:machines}

\subsubsection{Definitions shared by all machines}\label{sec:machine-shared}

\[
\begin{array}{rcl}
  s &::=& \ev{M}{\rho}{\kappa} \mid \rt{v}{\kappa} \\
  \rho &:& \X \rightharpoonup \mathit{Val} \quad \text{(finite)} \qquad
  \emptyset, \quad \ext{\rho}{x}{v}, \quad \{x_1 {\mapsto} v_1, \ldots, x_n {\mapsto} v_n\}
\end{array}
\]
\begin{align*}
  \retq(v, \kappa) &= \nxt \rt{v}{\kappa} \\
  \retq(\bot, \kappa) &= \stp \error \\[1ex]
  \key{proj}^n_i(v_1, \ldots, v_m) &=
    \begin{cases} v_i & m = n \\ \bot & \text{otherwise} \end{cases}
    \qquad (\prj{i} \text{ of arity } n) \\[1ex]
  \tnext(\bar{v};\, \epsilon;\, \rho;\, \kappa) &= \rt{\tupV{\bar{v}}}{\kappa} \\
  \tnext(\bar{v};\, M\,\bar{M};\, \rho;\, \kappa) &= \ev{M}{\rho}{\Kf{tupK}{\bar{v}, \bar{M}, \rho, \kappa}} \\[1ex]
  \obs(c) &= c \\
  \obs(\tupV{v_1, \ldots, v_n}) &= \key{tuple}(\obs(v_1), \ldots, \obs(v_n)) \\
  \obs(\inlV v) &= \inl \obs(v) \\
  \obs(\inrV v) &= \inr \obs(v) \\
  \obs(v) &= \key{fun} \quad \text{otherwise (closures, code, and pairs)} \\[1ex]
  \res{n} &= n, \quad \res{\mathsf{true}} = \inlV \tupV{}, \quad \res{\mathsf{false}} = \inrV \tupV{} \\[1ex]
  \Mach{\ell}(M) &= (\mathit{State}^{\ell},\; \step{\ell},\; \ev{M}{\emptyset}{\halt})
\end{align*}
We write $s \mapsto n$ for $\step{\ell}(s) = n$.

\subsubsection{Rules common to all machines}

\[
\begin{array}{rcl}
  v &::=& \cdots \mid c \mid \tupV{v_1, \ldots, v_n} \mid \inlV v \mid \inrV v \\
  \kappa &::=& \cdots \mid \halt \mid \Kf{tupK}{\bar{v}, \bar{M}, \rho, \kappa} \mid \Kf{getK_i}{\kappa}
    \mid \Kf{inlK}{\kappa} \mid \Kf{inrK}{\kappa} \mid \Kf{caseK}{x.M, y.N, \rho, \kappa} \\
  &\mid& \Kf{primL_p}{M, \rho, \kappa} \mid \Kf{primR_p}{v, \kappa}
\end{array}
\]
\[
\begin{array}{@{}l@{\;\mapsto\;}l@{\quad}l@{}}
  \ev{x}{\rho}{\kappa} & \retq(\rho(x), \kappa) \\
  \ev{c}{\rho}{\kappa} & \nxt \rt{c}{\kappa} \\
  \ev{\tuple{M_1, \ldots, M_n}}{\rho}{\kappa} & \nxt \tnext(\epsilon;\, M_1 \ldots M_n;\, \rho;\, \kappa) \\
  \ev{(\prj{i} M)}{\rho}{\kappa} & \nxt \ev{M}{\rho}{\Kf{getK_i}{\kappa}} \\
  \ev{(\inl M)}{\rho}{\kappa} & \nxt \ev{M}{\rho}{\Kf{inlK}{\kappa}} \\
  \ev{(\inr M)}{\rho}{\kappa} & \nxt \ev{M}{\rho}{\Kf{inrK}{\kappa}} \\
  \ev{(\Case{L}{x}{M}{y}{N})}{\rho}{\kappa} & \nxt \ev{L}{\rho}{\Kf{caseK}{x.M, y.N, \rho, \kappa}} \\
  \ev{\prim{p}{L}{M}}{\rho}{\kappa} & \nxt \ev{L}{\rho}{\Kf{primL_p}{M, \rho, \kappa}} \\[1ex]
  \rt{v}{\halt} & \stp \obs(v) \\
  \rt{v}{\Kf{tupK}{\bar{v}, \bar{M}, \rho, \kappa}} & \nxt \tnext(\bar{v}\,v;\, \bar{M};\, \rho;\, \kappa) \\
  \rt{v}{\Kf{getK_i}{\kappa}} & \retq(\key{proj}^n_i(\bar{v}), \kappa) & v = \tupV{\bar{v}} \\
  & \stp \error & \text{otherwise} \\
  \rt{v}{\Kf{inlK}{\kappa}} & \nxt \rt{(\inlV v)}{\kappa} \\
  \rt{v}{\Kf{inrK}{\kappa}} & \nxt \rt{(\inrV v)}{\kappa} \\
  \rt{v}{\Kf{caseK}{x.M, y.N, \rho, \kappa}} & \nxt \ev{M}{\ext{\rho}{x}{u}}{\kappa} & v = \inlV u \\
  & \nxt \ev{N}{\ext{\rho}{y}{u}}{\kappa} & v = \inrV u \\
  & \stp \error & \text{otherwise} \\
  \rt{v}{\Kf{primL_p}{M, \rho, \kappa}} & \nxt \ev{M}{\rho}{\Kf{primR_p}{v, \kappa}} \\
  \rt{w}{\Kf{primR_p}{v, \kappa}} & \nxt \rt{\res{\delta(p, n, m)}}{\kappa} & v = n,\; w = m \\
  & \stp \error & \text{otherwise}
\end{array}
\]

\subsubsection{ISWIM}

\[
\begin{array}{rcl}
  v &::=& \cloV{\lam{x} N}{\rho} \mid \cdots \\
  \kappa &::=& \Kf{appL}{M, \rho, \kappa} \mid \Kf{appR}{v, \kappa} \mid \cdots
\end{array}
\]
\[
\begin{array}{@{}l@{\;\mapsto\;}l@{\quad}l@{}}
  \ev{(\lam{x} N)}{\rho}{\kappa} & \nxt \rt{\cloV{\lam{x} N}{\rho}}{\kappa} \\
  \ev{(L\;M)}{\rho}{\kappa} & \nxt \ev{L}{\rho}{\Kf{appL}{M, \rho, \kappa}} \\[1ex]
  \rt{v}{\Kf{appL}{M, \rho, \kappa}} & \nxt \ev{M}{\rho}{\Kf{appR}{v, \kappa}} & v = \cloV{\lam{x} N}{\rho'} \\
  & \stp \error & \text{otherwise} \\
  \rt{w}{\Kf{appR}{v, \kappa}} & \nxt \ev{N}{\ext{\rho}{x}{w}}{\kappa} & v = \cloV{\lam{x} N}{\rho} \\
  & \stp \error & \text{otherwise}
\end{array}
\]

\subsubsection{Clos}

\[
\begin{array}{rcl}
  v &::=& \closV{x_1 \ldots x_n\, y.\, N}{v_1 \ldots v_n} \mid \cdots \\
  \kappa &::=& \Kf{appL}{M, \rho, \kappa} \mid \Kf{appR}{v, \kappa}
    \mid \Kf{closK}{x_1 \ldots x_n\, y.\, N, \bar{v}, \bar{M}, \rho, \kappa} \mid \cdots
\end{array}
\]
\begin{align*}
  \cnext(\bar{x}\, y.\, N;\, \bar{v};\, \epsilon;\, \rho;\, \kappa) &= \rt{\closV{\bar{x}\, y.\, N}{\bar{v}}}{\kappa} \\
  \cnext(\bar{x}\, y.\, N;\, \bar{v};\, M\,\bar{M};\, \rho;\, \kappa) &=
    \ev{M}{\rho}{\Kf{closK}{\bar{x}\, y.\, N, \bar{v}, \bar{M}, \rho, \kappa}}
\end{align*}
\[
\begin{array}{@{}l@{\;\mapsto\;}l@{\quad}l@{}}
  \ev{(\clos{z_1 \ldots z_n}{x_1 \ldots x_n\, y} N)}{\rho}{\kappa} &
    \nxt \cnext(x_1 \ldots x_n\, y.\, N;\, \epsilon;\, z_1 \ldots z_n;\, \rho;\, \kappa) \\
  \ev{(L\;M)}{\rho}{\kappa} & \nxt \ev{L}{\rho}{\Kf{appL}{M, \rho, \kappa}} \\[1ex]
  \rt{v}{\Kf{closK}{\bar{x}\, y.\, N, \bar{v}, \bar{M}, \rho, \kappa}} &
    \nxt \cnext(\bar{x}\, y.\, N;\, \bar{v}\,v;\, \bar{M};\, \rho;\, \kappa) \\
  \rt{v}{\Kf{appL}{M, \rho, \kappa}} & \nxt \ev{M}{\rho}{\Kf{appR}{v, \kappa}} & v = \closV{\ldots}{\ldots} \\
  & \stp \error & \text{otherwise} \\
  \rt{w}{\Kf{appR}{v, \kappa}} & \nxt \ev{N}{\{\bar{x} {\mapsto} \bar{v}, y {\mapsto} w\}}{\kappa}
    & v = \closV{\bar{x}\, y.\, N}{\bar{v}} \\
  & \stp \error & \text{otherwise}
\end{array}
\]

\subsubsection{EnvClos}

\[
\begin{array}{rcl}
  v &::=& \closV{e\, y.\, N}{v} \mid \cdots \\
  \kappa &::=& \Kf{appL}{M, \rho, \kappa} \mid \Kf{appR}{v, \kappa} \mid \Kf{closK}{e\, y.\, N, \kappa} \mid \cdots
\end{array}
\]
\[
\begin{array}{@{}l@{\;\mapsto\;}l@{\quad}l@{}}
  \ev{(\clos{r_1 \ldots r_n}{e\, y} N)}{\rho}{\kappa} &
    \nxt \tnext(\epsilon;\, r_1 \ldots r_n;\, \rho;\, \Kf{closK}{e\, y.\, N, \kappa}) \\
  \ev{(L\;M)}{\rho}{\kappa} & \nxt \ev{L}{\rho}{\Kf{appL}{M, \rho, \kappa}} \\[1ex]
  \rt{v}{\Kf{closK}{e\, y.\, N, \kappa}} & \nxt \rt{\closV{e\, y.\, N}{v}}{\kappa} \\
  \rt{v}{\Kf{appL}{M, \rho, \kappa}} & \nxt \ev{M}{\rho}{\Kf{appR}{v, \kappa}} & v = \closV{\ldots}{\ldots} \\
  & \stp \error & \text{otherwise} \\
  \rt{w}{\Kf{appR}{v, \kappa}} & \nxt \ev{N}{\{e {\mapsto} u, y {\mapsto} w\}}{\kappa}
    & v = \closV{e\, y.\, N}{u} \\
  & \stp \error & \text{otherwise}
\end{array}
\]

\subsubsection{FlatClos}

An application checks that its operator is code when it calls it, after
evaluating all three operands.
\[
\begin{array}{rcl}
  v, f, u &::=& \funV{e\, y.\, N} \mid \pairV{v}{w} \mid \cdots \\
  \kappa &::=& \Kf{appF}{M, N, \rho, \kappa} \mid \Kf{appE}{f, N, \rho, \kappa} \mid \Kf{appA}{f, u, \kappa}
    \mid \Kf{pairL}{N, \rho, \kappa} \mid \Kf{pairR}{v, \kappa} \\
  &\mid& \Kf{fstK}{\kappa} \mid \Kf{sndK}{\kappa} \mid \Kf{letK}{x.N, \rho, \kappa} \mid \cdots
\end{array}
\]
\[
\begin{array}{@{}l@{\;\mapsto\;}l@{\quad}l@{}}
  \ev{\fun{e\,y.\,N}}{\rho}{\kappa} & \nxt \rt{\funV{e\, y.\, N}}{\kappa} \\
  \ev{L(M, N)}{\rho}{\kappa} & \nxt \ev{L}{\rho}{\Kf{appF}{M, N, \rho, \kappa}} \\
  \ev{\pair{M}{N}}{\rho}{\kappa} & \nxt \ev{M}{\rho}{\Kf{pairL}{N, \rho, \kappa}} \\
  \ev{(\FST M)}{\rho}{\kappa} & \nxt \ev{M}{\rho}{\Kf{fstK}{\kappa}} \\
  \ev{(\SND M)}{\rho}{\kappa} & \nxt \ev{M}{\rho}{\Kf{sndK}{\kappa}} \\
  \ev{(\Let x = M \In N)}{\rho}{\kappa} & \nxt \ev{M}{\rho}{\Kf{letK}{x.N, \rho, \kappa}} \\[1ex]
  \rt{f}{\Kf{appF}{M, N, \rho, \kappa}} & \nxt \ev{M}{\rho}{\Kf{appE}{f, N, \rho, \kappa}} \\
  \rt{u}{\Kf{appE}{f, N, \rho, \kappa}} & \nxt \ev{N}{\rho}{\Kf{appA}{f, u, \kappa}} \\
  \rt{w}{\Kf{appA}{f, u, \kappa}} & \nxt \ev{N}{\{e {\mapsto} u, y {\mapsto} w\}}{\kappa} & f = \funV{e\, y.\, N} \\
  & \stp \error & \text{otherwise} \\
  \rt{v}{\Kf{pairL}{N, \rho, \kappa}} & \nxt \ev{N}{\rho}{\Kf{pairR}{v, \kappa}} \\
  \rt{w}{\Kf{pairR}{v, \kappa}} & \nxt \rt{\pairV{v}{w}}{\kappa} \\
  \rt{v}{\Kf{fstK}{\kappa}} & \nxt \rt{v_1}{\kappa} & v = \pairV{v_1}{v_2} \\
  & \stp \error & \text{otherwise} \\
  \rt{v}{\Kf{sndK}{\kappa}} & \nxt \rt{v_2}{\kappa} & v = \pairV{v_1}{v_2} \\
  & \stp \error & \text{otherwise} \\
  \rt{v}{\Kf{letK}{x.N, \rho, \kappa}} & \nxt \ev{N}{\ext{\rho}{x}{v}}{\kappa}
\end{array}
\]

\subsubsection{GlobalFun}

The machine of a program runs with its global definitions.
\[
\begin{array}{rcl}
  v, f, u &::=& \codeV{k} \mid \pairV{v}{w} \mid \cdots \\
  \Mach{5}(d_0, \ldots, d_{m-1};\; M) &=& (\mathit{State}^{5},\; \step{5}_{d_0 \ldots d_{m-1}},\; \ev{M}{\emptyset}{\halt})
\end{array}
\]
The continuations and the rules of $\step{5}_{d_0 \ldots d_{m-1}}$ are those of FlatClos, with $\codeV{k}$ in place of $\funV{e\,y.\,N}$:
\[
\begin{array}{@{}l@{\;\mapsto\;}l@{\quad}l@{}}
  \ev{\funref{k}}{\rho}{\kappa} & \nxt \rt{\codeV{k}}{\kappa} \\[1ex]
  \rt{w}{\Kf{appA}{f, u, \kappa}} & \nxt \ev{N}{\{e {\mapsto} u, y {\mapsto} w\}}{\kappa}
    & f = \codeV{k},\; k < m,\; d_k = \key{fun}\; e\,y.\,N \\
  & \stp \error & \text{otherwise}
\end{array}
\]

\subsubsection{Imp}

Both kinds of state carry the frame $\sigma$ of the running body, a list
of values, so an assignment is seen by everything after it. The values
are those of GlobalFun, and the continuation frames hold no environment.
The operands of a simple expression are atoms, evaluated left to right
in a step each, collecting their values in frames $F$ as in GlobalFun. What
a simple expression does at its start ($\mathit{first}$) and when the
value of an operand comes back ($\mathit{resume}$) is shared with the CFG
machine: evaluate another operand, return a value (or stop with an error), or call
definition $k$.
\[
\begin{array}{rcl}
  \mathit{State}^{\mathrm{Imp}} &::=& \evi{e}{\sigma}{\kappa} \mid \rti{v}{\sigma}{\kappa} \qquad
  \sigma = v_0 \ldots v_{n-1} \\
  F &::=& \key{appF}(a_2, a_3) \mid \key{appE}(f, a_3) \mid \key{appA}(f, u)
    \mid \key{pairL}(a) \mid \key{pairR}(v) \mid \key{fstK} \mid \key{sndK} \\
  &\mid& \key{tupK}(\bar{v}, \bar{a}) \mid \key{getK}_i \mid \key{inlK} \mid \key{inrK}
    \mid \key{primL}_p(a) \mid \key{primR}_p(v) \\
  \kappa &::=& \halt \mid F \cdot \kappa \mid \Kf{caseK}{s, e_1, e_2, \kappa}
    \mid \Kf{assignK}{s, e, \kappa} \mid \Kf{callK}{\sigma, \kappa} \\
  \mathit{Res} &::=& \key{operand}(a, F) \mid \key{value}(v^?) \mid \key{apply}(k, u, w)
\end{array}
\]
\begin{align*}
  \sigma(s) &= \begin{cases} v_s & s < n \\ \bot & \text{otherwise} \end{cases}
  \qquad \sigma(c) = c \qquad \sigma(\funref{k}) = \codeV{k} \\
  \sigma[s := v] &= \begin{cases} v_0 \ldots v_{s-1}\, v\, v_{s+1} \ldots v_{n-1} & s < n \\ \bot & \text{otherwise} \end{cases} \\[1ex]
  \tnext(\bar{v};\, \epsilon) &= \key{value}(\tupV{\bar{v}}) \qquad
  \tnext(\bar{v};\, a\,\bar{a}) = \key{operand}(a, \key{tupK}(\bar{v}, \bar{a}))
\end{align*}
\[
\begin{array}{@{}l@{\;=\;}l@{\qquad}l@{\;=\;}l@{}}
  \mathit{first}(\sigma, a) & \key{value}(\sigma(a)) &
  \mathit{first}(\sigma, a_1(a_2, a_3)) & \key{operand}(a_1, \key{appF}(a_2, a_3)) \\
  \mathit{first}(\sigma, \pair{a_1}{a_2}) & \key{operand}(a_1, \key{pairL}(a_2)) &
  \mathit{first}(\sigma, \FST a) & \key{operand}(a, \key{fstK}) \\
  \mathit{first}(\sigma, \SND a) & \key{operand}(a, \key{sndK}) &
  \mathit{first}(\sigma, \tuple{a_1, \ldots, a_l}) & \tnext(\epsilon;\, a_1 \ldots a_l) \\
  \mathit{first}(\sigma, \prj{i} a) & \key{operand}(a, \key{getK}_i) &
  \mathit{first}(\sigma, \inl a) & \key{operand}(a, \key{inlK}) \\
  \mathit{first}(\sigma, \inr a) & \key{operand}(a, \key{inrK}) &
  \mathit{first}(\sigma, \prim{p}{a_1}{a_2}) & \key{operand}(a_1, \key{primL}_p(a_2))
\end{array}
\]
\[
\begin{array}{@{}l@{\;=\;}l@{\quad}l@{}}
  \mathit{resume}(f, \key{appF}(a_2, a_3)) & \key{operand}(a_2, \key{appE}(f, a_3)) \\
  \mathit{resume}(u, \key{appE}(f, a_3)) & \key{operand}(a_3, \key{appA}(f, u)) \\
  \mathit{resume}(w, \key{appA}(f, u)) & \key{apply}(k, u, w) & f = \codeV{k} \\
  & \key{value}(\bot) & \text{otherwise} \\
  \mathit{resume}(v, \key{pairL}(a)) & \key{operand}(a, \key{pairR}(v)) \\
  \mathit{resume}(w, \key{pairR}(v)) & \key{value}(\pairV{v}{w}) \\
  \mathit{resume}(v, \key{fstK}) & \key{value}(v_1) & v = \pairV{v_1}{v_2} \\
  & \key{value}(\bot) & \text{otherwise} \\
  \mathit{resume}(v, \key{sndK}) & \key{value}(v_2) & v = \pairV{v_1}{v_2} \\
  & \key{value}(\bot) & \text{otherwise} \\
  \mathit{resume}(v, \key{tupK}(\bar{v}, \bar{a})) & \tnext(\bar{v}\,v;\, \bar{a}) \\
  \mathit{resume}(v, \key{getK}_i) & \key{value}(\key{proj}^n_i(\bar{v})) & v = \tupV{\bar{v}} \\
  & \key{value}(\bot) & \text{otherwise} \\
  \mathit{resume}(v, \key{inlK}) & \key{value}(\inlV v) \qquad
  \mathit{resume}(v, \key{inrK}) = \key{value}(\inrV v) \\
  \mathit{resume}(v, \key{primL}_p(a)) & \key{operand}(a, \key{primR}_p(v)) \\
  \mathit{resume}(w, \key{primR}_p(v)) & \key{value}(\res{\delta(p, n, m)}) & v = n,\; w = m \\
  & \key{value}(\bot) & \text{otherwise}
\end{array}
\]
As in GlobalFun, a call checks that its operator is code after evaluating its
three operands. It runs the body $(n;\; e)$ in a fresh frame of
$\max(2, n)$ slots, and $\key{callK}$ returns to the caller's frame.
For a program with bodies $b_0, \ldots, b_{m-1}$:
\begin{align*}
  \retq(v, \sigma, \kappa) &= \nxt \rti{v}{\sigma}{\kappa}, \quad \retq(\bot, \sigma, \kappa) = \stp \error \\
  \key{goto}(\sigma, e, \kappa) &= \nxt \evi{e}{\sigma}{\kappa}, \quad \key{goto}(\bot, e, \kappa) = \stp \error \\
  \mathit{call}(k, u, w, \sigma, \kappa) &=
    \begin{cases}
      \nxt \evi{e}{(u\, w\, 0^{n-2})}{\Kf{callK}{\sigma, \kappa}} & k < m,\; b_k = (n;\; e) \\
      \stp \error & \text{otherwise}
    \end{cases} \\[1ex]
  \mathit{res}(\key{operand}(a, F), \sigma, \kappa) &= \nxt \evi{a}{\sigma}{F \cdot \kappa} \\
  \mathit{res}(\key{value}(v^?), \sigma, \kappa) &= \retq(v^?, \sigma, \kappa) \\
  \mathit{res}(\key{apply}(k, u, w), \sigma, \kappa) &= \mathit{call}(k, u, w, \sigma, \kappa) \\[1ex]
  \Mach{\mathrm{Imp}}(b_0, \ldots, b_{m-1};\; (n;\; e)) &= (\mathit{State}^{\mathrm{Imp}},\; \step{\mathrm{Imp}}_{b_0 \ldots b_{m-1}},\; \evi{e}{0^n}{\halt})
\end{align*}
where $0^{n-2}$ is $n - 2$ zeros, none when $n < 2$.
\[
\begin{array}{@{}l@{\;\mapsto\;}l@{\quad}l@{}}
  \evi{r}{\sigma}{\kappa} & \mathit{res}(\mathit{first}(\sigma, r), \sigma, \kappa) \\
  \evi{(\Casei{a}{s}{e_1}{e_2})}{\sigma}{\kappa} & \nxt \evi{a}{\sigma}{\Kf{caseK}{s, e_1, e_2, \kappa}} \\
  \evi{(\assign{s}{e_1}{e_2})}{\sigma}{\kappa} & \nxt \evi{e_1}{\sigma}{\Kf{assignK}{s, e_2, \kappa}} \\[1ex]
  \rti{v}{\sigma}{\halt} & \stp \obs(v) \\
  \rti{v}{\sigma}{F \cdot \kappa} & \mathit{res}(\mathit{resume}(v, F), \sigma, \kappa) \\
  \rti{v}{\sigma}{\Kf{caseK}{s, e_1, e_2, \kappa}} & \key{goto}(\sigma[s := u], e_1, \kappa) & v = \inlV u \\
  & \key{goto}(\sigma[s := u], e_2, \kappa) & v = \inrV u \\
  & \stp \error & \text{otherwise} \\
  \rti{v}{\sigma}{\Kf{assignK}{s, e, \kappa}} & \key{goto}(\sigma[s := v], e, \kappa) \\
  \rti{v}{\sigma}{\Kf{callK}{\sigma', \kappa}} & \nxt \rti{v}{\sigma'}{\kappa}
\end{array}
\]

\subsubsection{CFG}

The values, frames, frames $F$ of simple expressions, $\mathit{first}$,
and $\mathit{resume}$ are those of Imp. A tail runs in a frame with the
blocks it may jump to: those of its body for the start of a body, and
those older than block $l$ for block $l$. An operand is evaluated in a
state of its own.
\[
\begin{array}{rcl}
  \mathit{State}^{\mathrm{CFG}} &::=& \exc{t}{\sigma}{\vec{t}}{\kappa}
    \mid \evi{a}{\sigma}{\kappa} \mid \rti{v}{\sigma}{\kappa} \\
  \kappa &::=& \halt \mid F \cdot \kappa \mid \Kf{seqK}{s, t, \vec{t}, \kappa}
    \mid \Kf{caseK}{s, l_1, l_2, \vec{t}, \kappa} \mid \Kf{callK}{\sigma, \kappa}
\end{array}
\]
For a program with bodies $b_0, \ldots, b_{m-1}$, $\mathit{res}$ is that
of Imp with $\mathit{call}$ running the start of the callee's body with its blocks:
\begin{align*}
  \key{block}(t_0 \ldots t_{j-1}, l) &=
    \begin{cases} (t_l;\; t_0 \ldots t_{l-1}) & l < j \\ \bot & \text{otherwise} \end{cases} \\
  \key{jump}(\sigma, (t;\, \vec{t}), \kappa) &= \nxt \exc{t}{\sigma}{\vec{t}}{\kappa},
    \quad \key{jump}(\sigma^?, p^?, \kappa) = \stp \error \text{ if } \sigma^? = \bot \text{ or } p^? = \bot \\
  \mathit{call}(k, u, w, \sigma, \kappa) &=
    \begin{cases}
      \nxt \exc{t}{(u\, w\, 0^{n-2})}{\vec{t}}{\Kf{callK}{\sigma, \kappa}} & k < m,\; b_k = (n;\; \vec{t};\; t) \\
      \stp \error & \text{otherwise}
    \end{cases} \\[1ex]
  \Mach{\mathrm{CFG}}(b_0, \ldots, b_{m-1};\; (n;\; \vec{t};\; t)) &=
    (\mathit{State}^{\mathrm{CFG}},\; \step{\mathrm{CFG}}_{b_0 \ldots b_{m-1}},\; \exc{t}{0^n}{\vec{t}}{\halt})
\end{align*}
\[
\begin{array}{@{}l@{\;\mapsto\;}l@{\quad}l@{}}
  \exc{(\key{return}\; r)}{\sigma}{\vec{t}}{\kappa} & \mathit{res}(\mathit{first}(\sigma, r), \sigma, \kappa) \\
  \exc{(\assign{s}{r}{t})}{\sigma}{\vec{t}}{\kappa} &
    \mathit{res}(\mathit{first}(\sigma, r), \sigma, \Kf{seqK}{s, t, \vec{t}, \kappa}) \\
  \exc{(\key{goto}\; l)}{\sigma}{\vec{t}}{\kappa} & \key{jump}(\sigma, \key{block}(\vec{t}, l), \kappa) \\
  \exc{(\key{case}\;a\;\key{of}_{s}\;\key{goto}\;l_1 \mid \key{goto}\;l_2)}{\sigma}{\vec{t}}{\kappa} &
    \nxt \evi{a}{\sigma}{\Kf{caseK}{s, l_1, l_2, \vec{t}, \kappa}} \\
  \evi{a}{\sigma}{\kappa} & \mathit{res}(\key{value}(\sigma(a)), \sigma, \kappa) \\[1ex]
  \rti{v}{\sigma}{\halt} & \stp \obs(v) \\
  \rti{v}{\sigma}{F \cdot \kappa} & \mathit{res}(\mathit{resume}(v, F), \sigma, \kappa) \\
  \rti{v}{\sigma}{\Kf{seqK}{s, t, \vec{t}, \kappa}} & \key{jump}(\sigma[s := v], (t;\, \vec{t}), \kappa) \\
  \rti{v}{\sigma}{\Kf{caseK}{s, l_1, l_2, \vec{t}, \kappa}} & \key{jump}(\sigma[s := u], \key{block}(\vec{t}, l_1), \kappa) & v = \inlV u \\
  & \key{jump}(\sigma[s := u], \key{block}(\vec{t}, l_2), \kappa) & v = \inrV u \\
  & \stp \error & \text{otherwise} \\
  \rti{v}{\sigma}{\Kf{callK}{\sigma', \kappa}} & \nxt \rti{v}{\sigma'}{\kappa}
\end{array}
\]

\subsubsection{X86}\label{sec:machine-x86}

One instruction or jump per step. The values are those of GlobalFun and
the frames those of Imp. A state runs the instructions $\vec{\imath}$
and then the jump $j$, with the blocks $\vec{\beta}$ that $j$ may jump
to (as in CFG), the call stack $\Sigma$, the flag $f$, the frame
$\sigma$, and the registers $\rho$, each holding a value or unset ($\bot$).
\[
\begin{array}{rcl}
  \mathit{State}^{\mathrm{X86}} &::=& \xrun\;\vec{\imath}\;j\;\vec{\beta}\;\Sigma\;f\;\sigma\;\rho \\
  \Sigma &::=& \halt \mid \key{retK}(\vec{\imath}, j, \vec{\beta}, \sigma, \Sigma) \\
  f &::=& \bot \mid \mathsf{true} \mid \mathsf{false} \qquad
  \rho : \{\rax, \rdi, \rsi\} \to \mathit{Val}^{?}
\end{array}
\]
An instruction errors on an operand that is unset or beyond the
frame, on a value of the wrong kind, or on a destination beyond the
frame. With $\sigma$ and $\rho$ those of the state:
\begin{align*}
  \mathit{arg}(s) = \sigma(s), \quad \mathit{arg}(\%g) &= \rho(\%g), \quad
  \mathit{arg}(c) = c, \quad \mathit{arg}(\funref{k}) = \codeV{k} \\[1ex]
  \key{put}(s, v, K) &= \nxt K(\sigma[s := v], \rho) \quad (\sigma[s := v] \neq \bot) \\
  \key{put}(\%g, v, K) &= \nxt K(\sigma, \ext{\rho}{\%g}{v}) \\
  \key{put}(d, v^?, K) &= \stp \error \quad \text{otherwise} \\
  \key{test}(b, K) &= \nxt K(b), \quad \key{test}(\bot, K) = \stp \error \\[1ex]
  \key{prim}^?_p(v, w) &= \begin{cases} \res{\delta(p, n, m)} & v = n,\; w = m \\ \bot & \text{otherwise} \end{cases} \qquad
  \key{cmp}(v, w) = \begin{cases} n = m & v = n,\; w = m \\ \bot & \text{otherwise} \end{cases} \\
  \key{isinl}(v) &= \begin{cases} \mathsf{true} & v = \inlV u \\ \mathsf{false} & v = \inrV u \\ \bot & \text{otherwise} \end{cases} \\
  \key{build}_{\key{pair}}(v, w) &= \pairV{v}{w}, \quad
  \key{build}_{\key{tuple}}(\bar{v}) = \tupV{\bar{v}}, \quad
  \key{build}_{\key{inl}}(v) = \inlV v, \quad
  \key{build}_{\key{inr}}(v) = \inrV v \\
  \key{field}_{\key{fst}}(\pairV{v}{w}) &= v, \quad
  \key{field}_{\key{snd}}(\pairV{v}{w}) = w, \quad
  \key{field}_{\pi_i}(\tupV{\bar{v}}) = \key{proj}^n_i(\bar{v}), \\
  \key{field}_{\key{payload}}(\inlV v) &= v, \quad
  \key{field}_{\key{payload}}(\inrV v) = v
\end{align*}
and $\key{build}$, $\key{field}$, $\key{prim}^?$ and $\key{cmp}$ are
$\bot$ on any other arguments, or on an operand that is $\bot$.
$\key{callq}\; o$, with $o$ the code of definition $k$ and $\rdi$ and
$\rsi$ set, pushes the caller's remaining instructions, jump, blocks and
frame, and runs the start of $b_k$ in a fresh frame of $\max(2, n)$
zeros, with $\rdi$ and $\rsi$ as they are, $\rax$ unset and no flag.
$\key{retq}$ reads $\rax$: to $\halt$ it stops with its observation;
otherwise it pops the caller, which continues with $\rax$ and with
$\rdi$, $\rsi$ and the flag unset. $\key{block}$ is that of CFG. For a program with bodies
$b_0, \ldots, b_{m-1}$:
\begin{align*}
  \key{goto}((\vec{\imath};\, j;\, \vec{\beta}), \Sigma, f, \sigma, \rho) &= \nxt \xrun\;\vec{\imath}\;j\;\vec{\beta}\;\Sigma\;f\;\sigma\;\rho, \quad
  \key{goto}(\bot, \Sigma, f, \sigma, \rho) = \stp \error \\
  \mathit{callq}(\codeV{k}, \Sigma, \rho) &=
    \nxt \xrun\;\vec{\imath}\;j\;\vec{\beta}\;\Sigma\;\bot\;0^{\max(2,n)}\;(\bot, \rho(\rdi), \rho(\rsi)) \\
    &\qquad (k < m,\; b_k = (n;\; \vec{\beta};\; \vec{\imath};\, j),\; \rho(\rdi) \neq \bot \neq \rho(\rsi)) \\
  \mathit{callq}(v^?, \Sigma, \rho) &= \stp \error \quad \text{otherwise} \\[1ex]
  \rho_0 &= (0, \bot, \bot) \quad \text{(the values of $\rax$, $\rdi$, $\rsi$)} \\
  \Mach{\mathrm{X86}}(b_0, \ldots, b_{m-1};\; (n;\; \vec{\beta};\; \vec{\imath};\, j)) &=
    (\mathit{State}^{\mathrm{X86}},\; \step{\mathrm{X86}}_{b_0 \ldots b_{m-1}},\; \xrun\;\vec{\imath}\;j\;\vec{\beta}\;\halt\;\bot\;0^n\;\rho_0)
\end{align*}
Writing $\langle \vec{\imath} \rangle$ for $\xrun\;\vec{\imath}\;j\;\vec{\beta}\;\Sigma\;f\;\sigma\;\rho$
and $\langle \vec{\imath} \rangle_{f'}$ for it with the flag $f'$:
\[
\begin{array}{@{}l@{\;\mapsto\;}l@{\quad}l@{}}
  \langle \key{movq}\; o, d;\; \vec{\imath} \rangle & \key{put}(d, \mathit{arg}(o), \xrun\;\vec{\imath}\;j\;\vec{\beta}\;\Sigma\;f) \\
  \langle \key{subq}\; o, d;\; \vec{\imath} \rangle &
    \key{put}(d, \key{prim}^?_{\key{sub}}(\mathit{arg}(d), \mathit{arg}(o)), \xrun\;\vec{\imath}\;j\;\vec{\beta}\;\Sigma\;f) \\
  \langle \key{cmpq}\; o_2, o_1;\; \vec{\imath} \rangle &
    \key{test}(\key{cmp}(\mathit{arg}(o_1), \mathit{arg}(o_2)), \lambda f'.\, \langle \vec{\imath} \rangle_{f'}) \\
  \langle \key{sete}\; d;\; \vec{\imath} \rangle & \key{put}(d, \res{f}, \xrun\;\vec{\imath}\;j\;\vec{\beta}\;\Sigma\;f) & f \neq \bot \\
  & \stp \error & f = \bot \\
  \langle \key{testinl}\; o;\; \vec{\imath} \rangle &
    \key{test}(\key{isinl}(\mathit{arg}(o)), \lambda f'.\, \langle \vec{\imath} \rangle_{f'}) \\
  \langle \key{alloc}\; \mathit{sh}\; o_1, \ldots, o_p, d;\; \vec{\imath} \rangle &
    \key{put}(d, \key{build}_{\mathit{sh}}(\mathit{arg}(o_1), \ldots, \mathit{arg}(o_p)), \xrun\;\vec{\imath}\;j\;\vec{\beta}\;\Sigma\;f) \\
  \langle \key{load}\; \mathit{fl}\; o, d;\; \vec{\imath} \rangle &
    \key{put}(d, \key{field}_{\mathit{fl}}(\mathit{arg}(o)), \xrun\;\vec{\imath}\;j\;\vec{\beta}\;\Sigma\;f) \\
  \langle \key{callq}\; o;\; \vec{\imath} \rangle &
    \mathit{callq}(\mathit{arg}(o), \key{retK}(\vec{\imath}, j, \vec{\beta}, \sigma, \Sigma), \rho) \\[1ex]
  \xrun\;\epsilon\;\key{retq}\;\vec{\beta}\;\halt\;f\;\sigma\;\rho & \stp \obs(\rho(\rax)) & \rho(\rax) \neq \bot \\
  & \stp \error & \rho(\rax) = \bot \\
  \xrun\;\epsilon\;\key{retq}\;\vec{\beta}\;\key{retK}(\vec{\imath}, j, \vec{\beta}', \sigma', \Sigma)\;f\;\sigma\;\rho &
    \nxt \xrun\;\vec{\imath}\;j\;\vec{\beta}'\;\Sigma\;\bot\;\sigma'\;(\rho(\rax), \bot, \bot) & \rho(\rax) \neq \bot \\
  & \stp \error & \rho(\rax) = \bot \\
  \xrun\;\epsilon\;(\key{jmp}\; l)\;\vec{\beta}\;\Sigma\;f\;\sigma\;\rho & \key{goto}(\key{block}(\vec{\beta}, l), \Sigma, f, \sigma, \rho) \\
  \xrun\;\epsilon\;(\key{je}\; l_1, l_2)\;\vec{\beta}\;\Sigma\;f\;\sigma\;\rho & \key{goto}(\key{block}(\vec{\beta}, l_1), \Sigma, f, \sigma, \rho) & f = \mathsf{true} \\
  & \key{goto}(\key{block}(\vec{\beta}, l_2), \Sigma, f, \sigma, \rho) & f = \mathsf{false} \\
  & \stp \error & f = \bot
\end{array}
\]
where $\res{f}$ is $\inlV \tupV{}$ or $\inrV \tupV{}$, as $\key{eq}$ produces.

%-------------------------------------------------------------------------------
\subsection{Simulations}\label{sec:sims}

A forward simulation from $L_1$ to $L_2$ is a family of relations
$s_1 \sim_i s_2$ between their states, indexed by $i \in \mathbb{N}$.
\begin{align*}
  \FSim(L_1, L_2, \sim) &\iff
  \begin{array}[t]{l}
    \exists i.\; \start{L_1} \sim_i \start{L_2} \\
    \land\; \forall i\, s_1\, s_2.\; s_1 \sim_i s_2 \implies \mathit{Answers}_i(\step{L_1}(s_1), s_2)
  \end{array} \\[1ex]
  \mathit{Answers}_i(\stp o, s_2) &\iff \Final_{L_2}(s_2, o) \\
  \mathit{Answers}_i(\nxt s_1', s_2) &\iff
    (\exists j\, s_2'.\; s_2 \lto_{L_2}^+ s_2' \land s_1' \sim_j s_2')
    \lor (\exists j < i.\; s_1' \sim_j s_2)
\end{align*}
A plus simulation from $L_1$ to $L_2$ is a relation $s_1 \sim s_2$ between their states.
\begin{align*}
  \mathit{PlusSim}(L_1, L_2, \sim) &\iff
    \start{L_1} \sim \start{L_2} \land
    \forall s_1\, s_2.\; s_1 \sim s_2 \implies \mathit{Answers}^+(\step{L_1}(s_1), s_2) \\[1ex]
  \mathit{Answers}^+(\stp o, s_2) &\iff \Final_{L_2}(s_2, o) \\
  \mathit{Answers}^+(\nxt s_1', s_2) &\iff \exists s_2'.\; s_2 \lto_{L_2}^+ s_2' \land s_1' \sim s_2'
\end{align*}
\begin{align*}
  \Equiv(L_1, L_2) &\iff
    (\forall o.\; \Term_{L_1}(o) \iff \Term_{L_2}(o)) \land (\Div_{L_1} \iff \Div_{L_2})
\end{align*}

\begin{lemma}[Plus Simulation]\label{lem:plus-sim}
If $\mathit{PlusSim}(L_1, L_2, \sim)$, then $\FSim(L_1, L_2, \sim')$
where $s_1 \sim'_i s_2 \iff s_1 \sim s_2$.
\end{lemma}

\begin{lemma}[Forward Termination]\label{lem:fwd-term}
If $\FSim(L_1, L_2, \sim)$ and $\Term_{L_1}(o)$, then $\Term_{L_2}(o)$.
\end{lemma}

\begin{lemma}[Backward Termination]\label{lem:bwd-term}
If $\FSim(L_1, L_2, \sim)$ and $\Term_{L_2}(o)$, then $\Term_{L_1}(o)$.
\end{lemma}

\begin{theorem}[Forward Simulation Gives Equivalence]\label{thm:equivalent}
If $\FSim(L_1, L_2, \sim)$, then $\Equiv(L_1, L_2)$.
\end{theorem}

\begin{lemma}[Equivalence Composes]\label{lem:equiv-trans}
If $\Equiv(L_1, L_2)$ and $\Equiv(L_2, L_3)$, then
\[
  \Equiv(L_1, L_3) .
\]
\end{lemma}

%-------------------------------------------------------------------------------
\subsection{Relations}\label{sec:sim-relations}

Each pass relates source and target values ($v \approx w$), continuations
($\kappa_1 \approxk \kappa_2$), and states ($s_1 \sim s_2$).
The clauses of $\approx$ not shown are
$c \approx c$, $\tupV{\bar{v}} \approx \tupV{\bar{w}}$ if $\bar{v} \approx \bar{w}$,
and $\inlV v \approx \inlV w$ and $\inrV v \approx \inrV w$ if $v \approx w$.
The frames of $\approxk$ not shown are related when they are the same frame
with related values and related rest ($\halt \approxk \halt$,
$\Kf{appR}{v, \kappa_1} \approxk \Kf{appR}{w, \kappa_2}$ if $v \approx w$ and
$\kappa_1 \approxk \kappa_2$, \ldots).
\begin{align*}
  \bar{v} \approx \bar{w} &\iff |\bar{v}| = |\bar{w}| \land \forall i.\; v_i \approx w_i \\
  v^? \approxq w^? &\iff (v^? = \bot \land w^? = \bot) \lor (v^? = v \land w^? = w \land v \approx w) \\
  \rho_1 \approx_X \rho_2 &\iff \forall x \in X.\; \exists v\, w.\;
    \rho_1(x) = v \land \rho_2(x) = w \land v \approx w \\
  \FV(M_1 \ldots M_n) &= \FV(M_1) \cup \cdots \cup \FV(M_n)
\end{align*}

\begin{lemma}[Related Values Are Observed Alike]\label{lem:obs}
In each pass, if $v \approx w$, then $\obs(v) = \obs(w)$.
\end{lemma}

\subsubsection{Simplify}\label{sec:sim-simplify}

A target closure is the simplified body of a source closure, under facts
$E$ that are true of the two environments ($\rho \models_E \rho'$): for each
$x$ with $E(x) = (T, K)$, $\rho(x)$ is related to a value that the value term
$T$ denotes in $\rho'$, and, if $K \neq \bot$, $K$ describes $\rho(x)$: a value the same
up to functions, and a known function $\langle \lam{y} N, E_f \rangle$ the source
closure $\cloV{\lam{y} N}{\rho_f}$ for some $\rho_f$ with $\rho_f \models_{E_f} \rho'$.
States are related frame by frame, the code by $\pass{simp}$:
\begin{align*}
  \cloV{\lam{y} N}{\rho} \approx \cloV{\lam{y} \pass{simp}_{E'}(N)}{\rho'}
    &\iff \rho \models_E \rho' \qquad (E' = E \text{ under } y,\; y \text{ unknown}) \\
  \ev{M}{\rho}{\kappa} \sim \ev{\pass{simp}_E(M)}{\rho'}{\kappa'}
    &\iff \rho \models_E \rho' \land \kappa \approxk \kappa' \\
  \rt{v}{\kappa} \sim \rt{v'}{\kappa'} &\iff v \approx v' \land \kappa \approxk \kappa'
\end{align*}
with the source on the left. Where simplify reduced, the source runs ahead:
if $\pass{simp}_E(M)$ is a value, or $M$ has a known value, then
$\ev{M}{\rho}{\kappa} \lto^* \rt{v}{\kappa}$ for a $v$ it describes; at an
inlined call, the source runs the operator to the known function's closure,
then continues as for $\beta$. Where the
target does more (it evaluates a duplicated small value, or the value of a
let that simplify kept), it stutters. The simulation goes from the target
to the source, so its matches are of the target's state $s'$ and the
source's $s$: $s' \mathrel{\sim^{-1}_i} s$ says that the target is $i$ steps from a
state $t'$ with $s \sim t'$.

\begin{lemma}[Simplify Simulation]\label{lem:simplify-fsim}
For every ISWIM term $M$,
\[
  \FV(M) = \emptyset \implies
  \FSim(\Mach{\mathrm{ISWIM}}(\pass{simplify}(M)), \Mach{\mathrm{ISWIM}}(M), \sim^{-1}) .
\]
\end{lemma}

\begin{theorem}[Correctness of Simplify]\label{thm:simplify-correct}
For every ISWIM term $M$,
\[
  \FV(M) = \emptyset \implies
  \Equiv(\Mach{\mathrm{ISWIM}}(M), \Mach{\mathrm{ISWIM}}(\pass{simplify}(M))) ,
\]
by Lemma~\ref{lem:simplify-fsim} and the symmetry of $\Equiv$.
\end{theorem}

\subsubsection{CC: Enclose}\label{sec:sim-enclose}

\begin{align*}
  \cloV{\lam{y} N}{\rho_1} \approx \closV{x_1 \ldots x_n\, y.\, \pass{enclose}(N)}{v_1 \ldots v_n}
    &\iff \rho_1 \approx_{\FV(\lam{y} N)} \{x_1 {\mapsto} v_1, \ldots, x_n {\mapsto} v_n\} \\[1ex]
  \Kf{appL}{M, \rho_1, \kappa_1} \approxk \Kf{appL}{\pass{enclose}(M), \rho_2, \kappa_2}
    &\iff \rho_1 \approx_{\FV(M)} \rho_2 \land \kappa_1 \approxk \kappa_2 \\
  \Kf{tupK}{\bar{v}, \bar{M}, \rho_1, \kappa_1} \approxk \Kf{tupK}{\bar{w}, \pass{enclose}(\bar{M}), \rho_2, \kappa_2}
    &\iff \bar{v} \approx \bar{w} \land \rho_1 \approx_{\FV(\bar{M})} \rho_2 \land \kappa_1 \approxk \kappa_2 \\
  \Kf{caseK}{x.M, y.N, \rho_1, \kappa_1} \approxk {}&
    \Kf{caseK}{x.\pass{enclose}(M), y.\pass{enclose}(N), \rho_2, \kappa_2} \\
    &\iff \rho_1 \approx_{\FV(\lam{x} M) \cup \FV(\lam{y} N)} \rho_2 \land \kappa_1 \approxk \kappa_2 \\
  \Kf{primL_p}{M, \rho_1, \kappa_1} \approxk \Kf{primL_p}{\pass{enclose}(M), \rho_2, \kappa_2}
    &\iff \rho_1 \approx_{\FV(M)} \rho_2 \land \kappa_1 \approxk \kappa_2 \\[1ex]
  \ev{M}{\rho_1}{\kappa_1} \sim \ev{\pass{enclose}(M)}{\rho_2}{\kappa_2}
    &\iff \rho_1 \approx_{\FV(M)} \rho_2 \land \kappa_1 \approxk \kappa_2 \\
  \rt{v}{\kappa_1} \sim \rt{w}{\kappa_2} &\iff v \approx w \land \kappa_1 \approxk \kappa_2
\end{align*}

\begin{lemma}[Enclose Simulation]\label{lem:enclose-fsim}
For every ISWIM term $M$,
\[
  \FV(M) = \emptyset \implies
  \mathit{PlusSim}(\Mach{\mathrm{ISWIM}}(M), \Mach{2}(\pass{enclose}(M)), \sim) .
\]
\end{lemma}

\begin{theorem}[Correctness of Enclose]\label{thm:enclose-correct}
For every ISWIM term $M$,
\[
  \FV(M) = \emptyset \implies
  \Equiv(\Mach{\mathrm{ISWIM}}(M), \Mach{2}(\pass{enclose}(M))) .
\]
\end{theorem}

\subsubsection{CC: Optimize}\label{sec:sim-optimize}

For a closure $\clos{z_1 \ldots z_n}{x_1 \ldots x_n\, y} N$, let
$\{ i_1 < \cdots < i_k \} = \{ i \mid x_i \in \FV(N) \}$, as in optimize, and
\[
  \bar{x}'_{<j} = x_{i_l} \ldots \quad (i_l < j) \qquad
  \bar{z}'_{>j} = z_{i_l} \ldots \quad (i_l > j) \qquad
  \bar{x}' = \bar{x}'_{<n+1} .
\]
\begin{align*}
  \mathit{Scoped}(M) &\iff \text{every closure } \clos{\bar{z}}{\bar{x}\, y} N \text{ in } M
    \text{ has } \FV(N) \subseteq \{\bar{x}, y\} \\
  \mathit{WS}(M) &\iff \mathit{Scoped}(M) \land \FV(M) = \emptyset \\
  \mathit{Built}^j_N(v_1 \ldots v_{j-1}, \bar{w}) &\iff \mathit{Scoped}(N) \land
    \{x_1 {\mapsto} v_1, \ldots, x_{j-1} {\mapsto} v_{j-1}\}
    \approx_{\FV(N) \cap \{x_1, \ldots, x_{j-1}\}} \{\bar{x}'_{<j} {\mapsto} \bar{w}\}
\end{align*}
\begin{align*}
  \closV{\bar{x}\, y.\, N}{\bar{v}} \approx {}& \closV{\bar{x}'\, y.\, \pass{optimize}(N)}{\bar{w}} \\
    &\iff \mathit{Scoped}(N) \land
      \{\bar{x} {\mapsto} \bar{v}\} \approx_{\FV(N) \setminus \{y\}} \{\bar{x}' {\mapsto} \bar{w}\} \\[1ex]
  \Kf{appL}{M, \rho_1, \kappa_1} \approxk {}& \Kf{appL}{\pass{optimize}(M), \rho_2, \kappa_2} \\
    &\iff \mathit{Scoped}(M) \land \rho_1 \approx_{\FV(M)} \rho_2 \land \kappa_1 \approxk \kappa_2 \\
  \Kf{primL_p}{M, \rho_1, \kappa_1} \approxk {}& \Kf{primL_p}{\pass{optimize}(M), \rho_2, \kappa_2} \\
    &\iff \mathit{Scoped}(M) \land \rho_1 \approx_{\FV(M)} \rho_2 \land \kappa_1 \approxk \kappa_2 \\
  \Kf{tupK}{\bar{v}, \bar{M}, \rho_1, \kappa_1} \approxk {}&
    \Kf{tupK}{\bar{w}, \pass{optimize}(\bar{M}), \rho_2, \kappa_2} \\
    &\iff \bar{v} \approx \bar{w} \land (\forall i.\; \mathit{Scoped}(M_i)) \land
      \rho_1 \approx_{\FV(\bar{M})} \rho_2 \land \kappa_1 \approxk \kappa_2 \\
  \Kf{caseK}{x.M, y.N, \rho_1, \kappa_1} \approxk {}&
    \Kf{caseK}{x.\pass{optimize}(M), y.\pass{optimize}(N), \rho_2, \kappa_2} \\
    &\iff \mathit{Scoped}(M) \land \mathit{Scoped}(N) \land
      \rho_1 \approx_{\FV(\lam{x} M) \cup \FV(\lam{y} N)} \rho_2 \\
    &\qquad \land\; \kappa_1 \approxk \kappa_2 \\
  \Kf{closK}{\bar{x}\, y.\, N, \bar{v}, z_{j+1} \ldots z_n, \rho_1, \kappa_1} \approxk {}&
    \Kf{closK}{\bar{x}'\, y.\, \pass{optimize}(N), \bar{w}, \bar{z}'_{>j}, \rho_2, \kappa_2} \\
    &\iff x_j \in \FV(N) \land \mathit{Built}^j_N(\bar{v}, \bar{w}) \land
      \rho_1 \approx_{\{z_{j+1}, \ldots, z_n\}} \rho_2 \land \kappa_1 \approxk \kappa_2
\end{align*}
The states are related with an index, which is $0$ except while the source
evaluates a free variable $z_j$ that optimize drops.
\begin{align*}
  \ev{M}{\rho_1}{\kappa_1} \sim_0 \ev{\pass{optimize}(M)}{\rho_2}{\kappa_2}
    &\iff \mathit{Scoped}(M) \land \rho_1 \approx_{\FV(M)} \rho_2 \land \kappa_1 \approxk \kappa_2 \\
  \rt{v}{\kappa_1} \sim_0 \rt{w}{\kappa_2} &\iff v \approx w \land \kappa_1 \approxk \kappa_2 \\
  \ev{z_j}{\rho_1}{\kappa'_1} \sim_{2(n-j)+2} s_2
    &\iff \mathit{Drop}_j \land \rho_1(z_j) \text{ is defined} \\
  \rt{v}{\kappa'_1} \sim_{2(n-j)+1} s_2 &\iff \mathit{Drop}_j
\end{align*}
where
\begin{align*}
  \kappa'_1 &= \Kf{closK}{\bar{x}\, y.\, N, \bar{v}, z_{j+1} \ldots z_n, \rho_1, \kappa_1} \\
  s_2 &= \cnext(\bar{x}'\, y.\, \pass{optimize}(N);\, \bar{w};\, \bar{z}'_{>j};\, \rho_2;\, \kappa_2) \\
  \mathit{Drop}_j &\iff x_j \notin \FV(N) \land \mathit{Built}^j_N(\bar{v}, \bar{w}) \land
      \rho_1 \approx_{\{z_{j+1}, \ldots, z_n\}} \rho_2 \land \kappa_1 \approxk \kappa_2
\end{align*}

\begin{lemma}[Enclose Is Well Scoped]\label{lem:enclose-ws}
For every ISWIM term $M$,
\[
  \FV(M) = \emptyset \implies \mathit{WS}(\pass{enclose}(M)) .
\]
\end{lemma}

\begin{lemma}[Optimize Simulation]\label{lem:optimize-fsim}
For every Clos term $M$,
\[
  \mathit{WS}(M) \implies \FSim(\Mach{2}(M), \Mach{2}(\pass{optimize}(M)), \sim) .
\]
\end{lemma}

\begin{theorem}[Correctness of Optimize]\label{thm:optimize-correct}
For every Clos term $M$,
\[
  \mathit{WS}(M) \implies \Equiv(\Mach{2}(M), \Mach{2}(\pass{optimize}(M))) .
\]
\end{theorem}

\subsubsection{CC: Concretize}\label{sec:sim-concretize}

\begin{align*}
  \rho_2 \approx_\varphi \rho_3 &\iff \forall x \in \mathrm{dom}(\varphi).\;
    \begin{cases}
      \rho_2(x) \approxq \rho_3(z) & \varphi(x) = z \\
      \exists \bar{t}.\; \rho_3(z) = \tupV{\bar{t}} \land \rho_2(x) \approxq \key{proj}^n_i(\bar{t})
        & \varphi(x) = \prj{i} z
    \end{cases}
\end{align*}
\begin{align*}
  \closV{x_1 \ldots x_n\, y.\, N}{\bar{v}} \approx {}&
    \closV{e\, y.\, \pass{concretize}_{\varphi_n}(N)}{\tupV{\bar{t}}}
    \iff \bar{v} \approx \bar{t} \\[1ex]
  \Kf{appL}{M, \rho_2, \kappa_2} \approxk {}&
    \Kf{appL}{\pass{concretize}_\varphi(M), \rho_3, \kappa_3}
    \iff \rho_2 \approx_\varphi \rho_3 \land \kappa_2 \approxk \kappa_3 \\
  \Kf{primL_p}{M, \rho_2, \kappa_2} \approxk {}&
    \Kf{primL_p}{\pass{concretize}_\varphi(M), \rho_3, \kappa_3}
    \iff \rho_2 \approx_\varphi \rho_3 \land \kappa_2 \approxk \kappa_3 \\
  \Kf{tupK}{\bar{v}, \bar{M}, \rho_2, \kappa_2} \approxk {}&
    \Kf{tupK}{\bar{w}, \pass{concretize}_\varphi(\bar{M}), \rho_3, \kappa_3}
    \iff \bar{v} \approx \bar{w} \land \rho_2 \approx_\varphi \rho_3 \land \kappa_2 \approxk \kappa_3 \\
  \Kf{caseK}{x.M, y.N, \rho_2, \kappa_2} \approxk {}&
    \Kf{caseK}{x.\pass{concretize}_{\ext{\varphi}{x}{x}}(M),\,
               y.\pass{concretize}_{\ext{\varphi}{y}{y}}(N), \rho_3, \kappa_3} \\
    &\iff \rho_2 \approx_\varphi \rho_3 \land \kappa_2 \approxk \kappa_3 \\
  \Kf{closK}{x_1 \ldots x_n\, y.\, N, \bar{v}, \bar{M}, \rho_2, \kappa_2} \approxk {}&
    \Kf{tupK}{\bar{w}, \pass{concretize}_\varphi(\bar{M}), \rho_3,
      \Kf{closK}{e\, y.\, \pass{concretize}_{\varphi_n}(N), \kappa_3}} \\
    &\iff \bar{v} \approx \bar{w} \land \rho_2 \approx_\varphi \rho_3
      \land |\bar{v}| + 1 + |\bar{M}| = n \land \kappa_2 \approxk \kappa_3 \\[1ex]
  \ev{M}{\rho_2}{\kappa_2} \sim {}& \ev{\pass{concretize}_\varphi(M)}{\rho_3}{\kappa_3}
    \iff M \notin \X \land \rho_2 \approx_\varphi \rho_3 \land \kappa_2 \approxk \kappa_3 \\
  \ev{x}{\rho_2}{\kappa_2} \sim {}& s_3
    \iff \exists w^?.\; \rho_2(x) \approxq w^? \land \kappa_2 \approxk \kappa_3
      \land \step{3}(s_3) = \retq(w^?, \kappa_3) \\
  \rt{v}{\kappa_2} \sim {}& \rt{w}{\kappa_3} \iff v \approx w \land \kappa_2 \approxk \kappa_3
\end{align*}

\begin{lemma}[Concretize Simulation]\label{lem:concretize-fsim}
For every Clos term $M$,
\[
  \mathit{PlusSim}(\Mach{2}(M), \Mach{3}(\pass{concretize}_{\mathit{id}}(M)), \sim) .
\]
\end{lemma}

\begin{theorem}[Correctness of Concretize]\label{thm:concretize-correct}
For every Clos term $M$,
\[
  \Equiv(\Mach{2}(M), \Mach{3}(\pass{concretize}_{\mathit{id}}(M))) .
\]
\end{theorem}

\subsubsection{CC: Delay}\label{sec:sim-delay}

\begin{align*}
  \rho_1 \approxq \rho_2 &\iff \forall x.\; x \text{ is not a fresh variable } f \text{ of delay}
    \implies \rho_1(x) \approxq \rho_2(x)
\end{align*}
\begin{align*}
  \closV{e\, y.\, N}{v} \approx \pairV{\funV{e\, y.\, \pass{delay}(N)}}{w} &\iff v \approx w
\end{align*}
\begin{align*}
  \Kf{appL}{M, \rho_1, \kappa_1} \approxk {}& \Kf{letK}{f.\,L_M, \rho_2, \kappa_2} \\
    &\iff \rho_1 \approxq \rho_2 \land \kappa_1 \approxk \kappa_2 \\
  \Kf{appR}{\closV{e\, y.\, N}{v}, \kappa_1} \approxk {}& \Kf{appA}{\funV{e\, y.\, \pass{delay}(N)}, w, \kappa_2} \\
    &\iff v \approx w \land \kappa_1 \approxk \kappa_2 \\
  \Kf{tupK}{\bar{v}, \bar{M}, \rho_1, \kappa_1} \approxk {}& \Kf{tupK}{\bar{w}, \pass{delay}(\bar{M}), \rho_2, \kappa_2} \\
    &\iff \bar{v} \approx \bar{w} \land \rho_1 \approxq \rho_2 \land \kappa_1 \approxk \kappa_2 \\
  \Kf{closK}{e\, y.\, N, \kappa_1} \approxk {}& \Kf{pairR}{\funV{e\, y.\, \pass{delay}(N)}, \kappa_2} \\
    &\iff \kappa_1 \approxk \kappa_2 \\
  \Kf{caseK}{x.M, y.N, \rho_1, \kappa_1} \approxk {}& \Kf{caseK}{x.\pass{delay}(M), y.\pass{delay}(N), \rho_2, \kappa_2} \\
    &\iff \rho_1 \approxq \rho_2 \land \kappa_1 \approxk \kappa_2 \\
  \Kf{primL_p}{M, \rho_1, \kappa_1} \approxk {}& \Kf{primL_p}{\pass{delay}(M), \rho_2, \kappa_2}
    \iff \rho_1 \approxq \rho_2 \land \kappa_1 \approxk \kappa_2
\end{align*}
\begin{align*}
  \ev{M}{\rho_1}{\kappa_1} \sim {}& \ev{\pass{delay}(M)}{\rho_2}{\kappa_2}
    \iff \rho_1 \approxq \rho_2 \land \kappa_1 \approxk \kappa_2 \\
  \rt{v}{\kappa_1} \sim {}& \rt{w}{\kappa_2}
    \iff v \approx w \land \kappa_1 \approxk \kappa_2 \land \kappa_1 \neq \Kf{appL}{\ldots} \\
  \rt{v}{\Kf{appL}{M, \rho_1, \kappa_1}} \sim {}& \rt{w}{\Kf{fstK}{\kappa_F}}
    \iff v \approx w \land \rho_1 \approxq \rho_2 \land \kappa_1 \approxk \kappa_2
\end{align*}
where $L_M = (\FST f)(\SND f, \pass{delay}(M))$, as in delay, and
$\kappa_F = \Kf{appF}{\SND f, \pass{delay}(M), \ext{\rho_2}{f}{w}, \kappa_2}$.

\begin{lemma}[Delay Simulation]\label{lem:delay-fsim}
For every EnvClos term $M$,
\[
  \mathit{PlusSim}(\Mach{3}(M), \Mach{4}(\pass{delay}(M)), \sim) .
\]
\end{lemma}

\begin{theorem}[Correctness of Delay]\label{thm:delay-correct}
For every EnvClos term $M$,
\[
  \Equiv(\Mach{3}(M), \Mach{4}(\pass{delay}(M))) .
\]
\end{theorem}

\subsubsection{CC: Globalize}\label{sec:sim-globalize}

The relations are those for the table $D$ that globalize makes for the whole program.
\begin{align*}
  M \leadsto M' &\iff \exists \vec{d}\, \vec{d}'.\;
    \pass{glob}(M, \vec{d}) = (M', \vec{d}') \land \mathit{Extends}(\vec{d}', D) \\[1ex]
  \funV{e\, y.\, N} \approx \codeV{k} &\iff \exists \vec{d}\, \vec{d}'\, N'.\;
    \pass{glob}(N, \vec{d}) = (N', \vec{d}') \land
    \mathit{Extends}((\vec{d}', \key{fun}\; e\,y.\,N'), D) \land k = |\vec{d}'| \\
  \pairV{v}{w} \approx \pairV{v'}{w'} &\iff v \approx v' \land w \approx w' \\
  \rho_1 \approx \rho_2 &\iff \mathrm{dom}(\rho_1) = \mathrm{dom}(\rho_2) \land
    \rho_1 \approx_{\mathrm{dom}(\rho_1)} \rho_2
\end{align*}
Each frame is related to the same frame when its terms are related by
$\leadsto$, its values and environments by $\approx$, and its rest by $\approxk$;
for example,
\begin{align*}
  \Kf{appF}{M, N, \rho_1, \kappa_1} \approxk \Kf{appF}{M', N', \rho_2, \kappa_2}
    &\iff M \leadsto M' \land N \leadsto N' \land \rho_1 \approx \rho_2 \land \kappa_1 \approxk \kappa_2 \\[1ex]
  \ev{M}{\rho_1}{\kappa_1} \sim \ev{M'}{\rho_2}{\kappa_2}
    &\iff M \leadsto M' \land \rho_1 \approx \rho_2 \land \kappa_1 \approxk \kappa_2 \\
  \rt{v}{\kappa_1} \sim \rt{w}{\kappa_2} &\iff v \approx w \land \kappa_1 \approxk \kappa_2
\end{align*}

\begin{lemma}[Definitions Stay Put]\label{lem:global-extends}
If $\mathit{Extends}((\vec{d}, d), D)$, then $D_{|\vec{d}|} = d$.
\end{lemma}

\begin{lemma}[Globalize Simulation]\label{lem:globalize-fsim}
For every FlatClos term $M$, if $\pass{globalize}(M) = (D;\; M')$, then
\[
  \mathit{PlusSim}(\Mach{4}(M), \Mach{5}(D;\; M'), \sim) .
\]
\end{lemma}

\begin{theorem}[Correctness of Globalize]\label{thm:globalize-correct}
For every FlatClos term $M$,
\[
  \Equiv(\Mach{4}(M), \Mach{5}(\pass{globalize}(M))) .
\]
\end{theorem}

\subsubsection{Remove Complex Operands}\label{sec:sim-rco}

The source runs with the definitions $D$ and the target with
$\pass{rco}(D)$, so values are related by equality ($\codeV{k}$ names
$D_k$ in the source and $\pass{rco}(D_k)$ in the target). The target
environment has the values of the lets as well:
\begin{align*}
  \rho_1 \preceq \rho_2 &\iff \forall x \in \mathrm{dom}(\rho_1).\; \rho_2(x) = \rho_1(x)
\end{align*}
Let $c$ be a construct with operands $M_1, \ldots, M_n$ (an application,
pair, projection, tuple, injection, primitive, or the scrutinee of a
case). We write $F^c(v_1 \ldots v_{j-1};\, M_j \ldots M_n;\, \rho;\, \kappa)$
for the frame in which the machine waits for the value of $M_j$ after the values
$v_1, \ldots, v_{j-1}$ (for $L(M, N)$: $\Kf{appF}{M, N, \rho, \kappa}$,
$\Kf{appE}{v_1, N, \rho, \kappa}$, $\Kf{appA}{v_1, v_2, \kappa}$), and
$F^c(v_1 \ldots v_{n-1};\, \epsilon;\, \rho;\, \kappa)$ for the frame
that waits for the last one. The target's code $Q$ that finishes $c$ is
pending, $\mathit{Pend}^c_{\rho_1, \kappa_1, \kappa_2}(\bar{v};\, \bar{M};\, Q, \rho_2)$,
by one rule per operand:
\begin{align*}
  \mathit{Pend}(\bar{v};\, M\,\bar{M};\, \Let x = \pass{rco}(M) \In Q, \rho_2)
    &\Leftarrow M \text{ complex} \land \rho_1 \preceq \rho_2 \land \FV(M) \subseteq \mathrm{dom}(\rho_1) \\
    &\qquad \land \forall w.\; \mathit{Pend}(\bar{v}\, w;\, \bar{M};\, Q, \ext{\rho_2}{x}{w}) \\
  \mathit{Pend}(\bar{v};\, a\,\bar{M};\, Q, \rho_2)
    &\Leftarrow a \text{ an atom} \land \rho_1(a) = u \land \mathit{Pend}(\bar{v}\, u;\, \bar{M};\, Q, \rho_2) \\
  \mathit{Pend}(\bar{v}\, v;\, \epsilon;\, Q, \rho_2)
    &\Leftarrow \ev{Q}{\rho_2}{\kappa_2} \lto^+ \rt{v}{F^c(\bar{v};\, \epsilon;\, \rho_2';\, \kappa_2)}
    \land \kappa_1 \approxk \kappa_2
\end{align*}
where $\rho_1(a)$ is the value of an atom ($\rho_1(x)$, $c$, or $\codeV{k}$).
The continuations:
\begin{align*}
  \Kf{letK}{x.N, \rho_1, \kappa_1} \approxk {}& \Kf{letK}{x.\pass{rco}(N), \rho_2, \kappa_2} \\
    &\iff (\forall v.\; \ext{\rho_1}{x}{v} \preceq \ext{\rho_2}{x}{v}) \land
      \FV(N) \subseteq \mathrm{dom}(\rho_1) \cup \{x\} \land \kappa_1 \approxk \kappa_2 \\
  F^c(\bar{v};\, \bar{M};\, \rho_1;\, \kappa_1) \approxk {}& \Kf{letK}{x.Q, \rho_2, \kappa_2} \\
    &\iff \forall w.\; \mathit{Pend}^c_{\rho_1, \kappa_1, \kappa_2}(\bar{v}\, w;\, \bar{M};\, Q, \ext{\rho_2}{x}{w}) \\
  F^c(\bar{v};\, \epsilon;\, \rho_1;\, \kappa_1) \approxk {}& F^c(\bar{v};\, \epsilon;\, \rho_2;\, \kappa_2)
    \iff \kappa_1 \approxk \kappa_2
\end{align*}
(for a case, the branches $x.M$ and $x.\pass{rco}(M)$ under the lets, and
$\rho_1 \preceq \rho_2$). In the states, the source evaluates an atom
operand while the target waits, so the relation is indexed by a bound
$|M|$ on the source steps before the target moves:
\begin{align*}
  \ev{M}{\rho_1}{\kappa_1} \sim_{|M|} {}& \ev{\pass{rco}(M)}{\rho_2}{\kappa_2}
    \iff \rho_1 \preceq \rho_2 \land \FV(M) \subseteq \mathrm{dom}(\rho_1) \land \kappa_1 \approxk \kappa_2 \\
  \rt{v}{\kappa_1} \sim_0 {}& \key{retT}(v, \kappa_2) \iff \kappa_1 \approxk \kappa_2 \\
  \ev{a}{\rho_1}{F^c(\bar{v};\, \bar{M};\, \rho_1;\, \kappa_1)} \sim_{2 + |\bar{M}|} {}& \ev{Q}{\rho_2}{\kappa_2}
    \iff \rho_1(a) = u \land \mathit{Pend}(\bar{v}\, u;\, \bar{M};\, Q, \rho_2) \\
  \rt{u}{F^c(\bar{v};\, M\,\bar{M};\, \rho_1;\, \kappa_1)} \sim_{1 + |M\,\bar{M}|} {}& \ev{Q}{\rho_2}{\kappa_2}
    \iff \mathit{Pend}(\bar{v}\, u;\, M\,\bar{M};\, Q, \rho_2)
\end{align*}
where $\key{retT}(v, \kappa_2)$ is the target state related to returning
$v$: when $\kappa_2 = \Kf{letK}{x.Q, \rho_2, \kappa}$ waits for the last
complex operand of a construct, the state that $\mathit{Pend}$ runs ahead
to, $\rt{v}{F^c(\bar{v};\, \epsilon;\, \rho_2';\, \kappa)}$, and otherwise
$\rt{v}{\kappa_2}$. So both machines take the construct's own step (a call,
a projection, \ldots) together, and error together.

\begin{lemma}[RCO Simulation]\label{lem:rco-fsim}
For every GlobalFun program $P$,
\[
  \mathit{Scoped}_5(P) \implies \FSim(\Mach{5}(P), \Mach{5}(\pass{rco}(P)), \sim) .
\]
\end{lemma}

\begin{theorem}[Correctness of RCO, Simulation]\label{thm:rco-sim}
For every GlobalFun program $P$,
\[
  \mathit{Scoped}_5(P) \implies \Equiv(\Mach{5}(P), \Mach{5}(\pass{rco}(P))) .
\]
\end{theorem}

\subsubsection{Lift Locals}\label{sec:sim-lift}

The machines run in lock step and compute the same values, so values are
related by equality ($\codeV{k}$ names $D_k$ in the source and
$\pass{lift}(D_k)$ in the target). A frame $\sigma$ of $n$ slots is
related to an environment $\rho$ when $\varphi$ gives the variables of
$\rho$ the slots below $|\rho|$, which hold their values:
\begin{align*}
  \rho \approx^n_\varphi \sigma &\iff |\sigma| = n \land
    \varphi \text{ maps } \mathrm{dom}(\rho) \text{ one to one onto } \{0, \ldots, |\rho| - 1\}
    \land \forall x \in \mathrm{dom}(\rho).\; \sigma(\varphi(x)) = \rho(x)
\end{align*}
A target continuation continues as a source one given any related frame,
$(\rho, \varphi) \vdash_n \kappa_1 \approx \kappa_2$, when its head frame
is related, it was made before a binder (weak), or it returns from a call:
\begin{align*}
  (\rho, \varphi) \vdash_n \kappa_1 \approx \kappa_2
    &\Leftarrow (\rho, \varphi) \vdash^{\mathsf{h}}_n \kappa_1 \approx \kappa_2 \\
  (\ext{\rho}{x}{v}, \ext{\varphi}{x}{|\rho|}) \vdash_n \kappa_1 \approx \kappa_2
    &\Leftarrow (\rho, \varphi) \vdash_n \kappa_1 \approx \kappa_2 \\
  (\rho, \varphi) \vdash_n \kappa_1 \approx \Kf{callK}{\sigma', \kappa_2}
    &\Leftarrow (\rho', \varphi') \vdash_{n'} \kappa_1 \approx \kappa_2
      \land \rho' \approx^{n'}_{\varphi'} \sigma'
\end{align*}
With $d = |\rho|$, $a' = \pass{lift}_{n,\varphi,d}(a)$, and
$(\rho, \varphi) \vdash_n \kappa_1 \approx \kappa_2$ for the rest of
each frame, the head frames:
\begin{align*}
  (\rho, \varphi) \vdash^{\mathsf{h}}_n {}& \halt \approx \halt \\
  (\rho, \varphi) \vdash^{\mathsf{h}}_n {}& \Kf{letK}{x.N, \rho, \kappa_1} \approx
    \Kf{assignK}{d, \pass{lift}_{n,\ext{\varphi}{x}{d},d+1}(N), \kappa_2}
    \quad\text{if } d + 1 + \mathit{height}(N) \leq n \\
  (\rho, \varphi) \vdash^{\mathsf{h}}_n {}& \Kf{caseK}{x.M, y.N, \rho, \kappa_1} \approx
    \Kf{caseK}{d, \pass{lift}_{n,\ext{\varphi}{x}{d},d+1}(M), \pass{lift}_{n,\ext{\varphi}{y}{d},d+1}(N), \kappa_2} \\
    &\qquad \text{if } d + 1 + \max(\mathit{height}(M), \mathit{height}(N)) \leq n \\
  (\rho, \varphi) \vdash^{\mathsf{h}}_n {}& \Kf{appF}{a_2, a_3, \rho, \kappa_1} \approx \key{appF}(a_2', a_3') \cdot \kappa_2 \qquad
  (\rho, \varphi) \vdash^{\mathsf{h}}_n \Kf{appE}{f, a_3, \rho, \kappa_1} \approx \key{appE}(f, a_3') \cdot \kappa_2 \\
  (\rho, \varphi) \vdash^{\mathsf{h}}_n {}& \Kf{pairL}{a, \rho, \kappa_1} \approx \key{pairL}(a') \cdot \kappa_2 \qquad
  (\rho, \varphi) \vdash^{\mathsf{h}}_n \Kf{primL_p}{a, \rho, \kappa_1} \approx \key{primL}_p(a') \cdot \kappa_2 \\
  (\rho, \varphi) \vdash^{\mathsf{h}}_n {}& \Kf{tupK}{\bar{v}, \bar{a}, \rho, \kappa_1} \approx \key{tupK}(\bar{v}, \bar{a}') \cdot \kappa_2
\end{align*}
and the frames that hold no term ($\key{appA}$, $\key{pairR}$, $\key{fstK}$,
$\key{sndK}$, $\key{getK}_i$, $\key{inlK}$, $\key{inrK}$, $\key{primR}_p$)
are related to the same frame $F$ with the same values, $F \cdot \kappa_2$. The states, for some
$n$ and $\varphi$:
\begin{align*}
  \ev{M}{\rho}{\kappa_1} \sim {}& \evi{\pass{lift}_{n,\varphi,|\rho|}(M)}{\sigma}{\kappa_2}
    \iff |\rho| + \mathit{height}(M) \leq n \land \rho \approx^n_\varphi \sigma
      \land (\rho, \varphi) \vdash_n \kappa_1 \approx \kappa_2 \\
  \rt{v}{\kappa_1} \sim {}& \rti{v}{\sigma}{\kappa_2}
    \iff \rho \approx^n_\varphi \sigma \land (\rho, \varphi) \vdash^{\mathsf{h}}_n \kappa_1 \approx \kappa_2
    \quad \text{for some } \rho
\end{align*}
The target's only extra steps are the returns from calls: when the source
returns a value, the target pops its $\key{callK}$ frames.

\begin{lemma}[Lift Simulation]\label{lem:lift-fsim}
For every monadic GlobalFun program $P$,
\[
  \mathit{PlusSim}(\Mach{5}(P), \Mach{\mathrm{Imp}}(\pass{lift}(P)), \sim) .
\]
\end{lemma}

\begin{theorem}[Correctness of Lift Locals, Simulation]\label{thm:lift-sim}
For every monadic GlobalFun program $P$,
\[
  \Equiv(\Mach{5}(P), \Mach{\mathrm{Imp}}(\pass{lift}(P))) .
\]
\end{theorem}
There is no precondition: a variable out of scope errors in both
machines, as its slot $n$ is beyond the frame.

\subsubsection{Explicate Control}\label{sec:sim-explicate}

The machines share the frame and the evaluation of simple expressions,
and run in lock step except for the control flow that the pass makes
explicit; values are related by equality. As in
Section~\ref{sec:den-explicate}, the pass threads the list of blocks,
and $\vec{t} \preceq \vec{t}'$ when $\vec{t}$ is a prefix of $\vec{t}'$.
The blocks of $\vec{t}$ keep their labels among newer blocks: if
$\vec{t}, t \preceq \vec{t}'$, then
$\key{block}(\vec{t}', |\vec{t}|) = (t;\, \vec{t})$. An Imp expression
runs as the tail the pass makes of it in a mode $m$: in tail position,
or assigned to $s$ before the tail $k$ runs.
\begin{align*}
  m &::= \key{tail} \mid \key{asg}(s, k) \\
  \pass{ex}_{\key{tail}}(e, \vec{t}) &= \pass{tail}(e, \vec{t}) \qquad
  \pass{ex}_{\key{asg}(s, k)}(e, \vec{t}) = \pass{assign}(e, s, k, \vec{t}) \\[1ex]
  \mathit{depth}(\assign{s}{e_1}{e_2}) &= 1 + \mathit{depth}(e_1) \qquad
  \mathit{depth}(e) = 0 \quad \text{otherwise}
\end{align*}
If $\pass{ex}_m(e, \vec{t}) = (t;\, \vec{t}')$, then $\vec{t} \preceq \vec{t}'$.
The stack of Imp continuations $\Kf{assignK}{s, e_2, \kappa}$ is a CFG
tail. A CFG continuation continues as an Imp one given a value,
$\kappa_1 \approx \kappa_2$; $\mathit{Ctl}_{\vec{t}}(\kappa_1, m, \kappa_2)$
when $\kappa_1$ is the continuation of an expression that runs in mode
$m$, made at the blocks $\vec{t}$, under $\kappa_2$: either the tail
made of $e_2$, or a jump to the join point that runs it; and
$\mathit{Br}_{\vec{t}', l}(\kappa_1, e, \kappa_2)$ when block $l$ of
$\vec{t}'$ runs $e$ under $\kappa_1$:
\begin{align*}
  \halt &\approx \halt \\
  F \cdot \kappa_1 &\approx F \cdot \kappa_2 \Leftarrow \kappa_1 \approx \kappa_2 \\
  \Kf{callK}{\sigma, \kappa_1} &\approx \Kf{callK}{\sigma, \kappa_2} \Leftarrow \kappa_1 \approx \kappa_2 \\
  \kappa_1 &\approx \Kf{seqK}{s, k, \vec{t}', \kappa_2} \Leftarrow
    \mathit{Ctl}_{\vec{t}}(\kappa_1, \key{asg}(s, k), \kappa_2) \land \vec{t} \preceq \vec{t}' \\
  \Kf{caseK}{s, e_1, e_2, \kappa_1} &\approx \Kf{caseK}{s, l_1, l_2, \vec{t}, \kappa_2} \Leftarrow
    \mathit{Br}_{\vec{t}, l_1}(\kappa_1, e_1, \kappa_2) \land \mathit{Br}_{\vec{t}, l_2}(\kappa_1, e_2, \kappa_2) \\[1ex]
  \mathit{Ctl}_{\vec{t}}(\kappa_1, \key{tail}, \kappa_2) &\Leftarrow \kappa_1 \approx \kappa_2 \\
  \mathit{Ctl}_{\vec{t}'}(\Kf{assignK}{s, e, \kappa_1}, \key{asg}(s, t), \kappa_2) &\Leftarrow
    \mathit{Ctl}_{\vec{t}}(\kappa_1, m, \kappa_2) \land \pass{ex}_m(e, \vec{t}) = (t;\, \vec{t}'') \land \vec{t}'' \preceq \vec{t}' \\
  \mathit{Ctl}_{\vec{t}'}(\kappa_1, \key{asg}(s, \key{goto}\; |\vec{t}|), \kappa_2) &\Leftarrow
    \mathit{Ctl}_{\vec{t}}(\kappa_1, \key{asg}(s, k), \kappa_2) \land \vec{t}, k \preceq \vec{t}' \\[1ex]
  \mathit{Br}_{\vec{t}', l}(\kappa_1, e, \kappa_2) &\iff \exists \vec{t}\, m.\;
    \mathit{Ctl}_{\vec{t}}(\kappa_1, m, \kappa_2) \land \key{block}(\vec{t}', l) = \pass{ex}_m(e, \vec{t})
\end{align*}
The states, indexed by the measure of the stutter:
\begin{align*}
  \evi{e}{\sigma}{\kappa_1} \sim_{\mathit{depth}(e)} {}& \exc{t}{\sigma}{\vec{t}'}{\kappa_2}
    \Leftarrow \mathit{Ctl}_{\vec{t}}(\kappa_1, m, \kappa_2) \land \pass{ex}_m(e, \vec{t}) = (t;\, \vec{t}'')
      \land \vec{t}'' \preceq \vec{t}' \\
  \evi{a}{\sigma}{\kappa_1} \sim_0 {}& \evi{a}{\sigma}{\kappa_2} \Leftarrow \kappa_1 \approx \kappa_2 \\
  \rti{v}{\sigma}{\kappa_1} \sim_0 {}& \rti{v}{\sigma}{\kappa_2} \Leftarrow \kappa_1 \approx \kappa_2
\end{align*}
An Imp step from $\assign{s}{e_1}{e_2}$ is matched by no CFG step: $e_1$
runs as the tail that $\assign{s}{e_1}{e_2}$ runs as, with a smaller
index. A value returned to $\Kf{assignK}{s, e_2, \kappa}$ is matched by
the assignment of $\key{seqK}$ and the jumps to join points:

\begin{lemma}[Join]\label{lem:join}
If $\mathit{Ctl}_{\vec{t}}(\kappa_1, \key{asg}(s, k), \kappa_2)$ and
$\vec{t} \preceq \vec{t}'$, then $\kappa_1 = \Kf{assignK}{s, e, \kappa}$
for some $e$ and $\kappa$, and for every $\sigma$ there is an $s_2$ with
\[
  \exc{k}{\sigma}{\vec{t}'}{\kappa_2} \lto^*_{\mathrm{CFG}} s_2
  \quad\text{and}\quad \evi{e}{\sigma}{\kappa} \sim_{\mathit{depth}(e)} s_2 .
\]
\end{lemma}

\begin{lemma}[Explicate Simulation]\label{lem:explicate-fsim}
For every Imp program $P$,
\[
  \FSim(\Mach{\mathrm{Imp}}(P), \Mach{\mathrm{CFG}}(\pass{explicate}(P)), \sim) .
\]
\end{lemma}

\begin{theorem}[Correctness of Explicate Control, Simulation]\label{thm:explicate-sim}
For every Imp program $P$,
\[
  \Equiv(\Mach{\mathrm{Imp}}(P), \Mach{\mathrm{CFG}}(\pass{explicate}(P))) .
\]
\end{theorem}
There is no precondition.

\subsubsection{Select Instructions}\label{sec:sim-select}

The frames are the same and values are related by equality; the
registers matter only where a value is in flight ($\rax$ after a simple
expression or a return, $\rdi$ and $\rsi$ at a call). Write
$\pass{sel}(t) = \mathit{ins}(t);\, \mathit{end}(t)$ for the instructions
and jump of the block of $t$. The call stack holds, for each
$\Kf{callK}{\sigma, D}$, the caller's frame $\sigma$ and the instructions
that move $\rax$ to where $D$ wants the result ($\mathit{Dest}$), then
continue as $D$:
\begin{align*}
  \halt &\approx \halt \\
  \Kf{callK}{\sigma, D} &\approx \key{retK}((\key{movq}\; \rax, d;\; \vec{\imath}), j, \vec{\beta}, \sigma, \Sigma)
    \Leftarrow \mathit{Dest}(D, d, \vec{\imath}, j, \vec{\beta}, \Sigma) \\[1ex]
  \mathit{Dest}(\Kf{seqK}{s, t, \vec{t}, \kappa}, s, \mathit{ins}(t), \mathit{end}(t), \pass{sel}(\vec{t}), \Sigma)
    &\Leftarrow \kappa \approx \Sigma \\
  \mathit{Dest}(\kappa, \rax, \epsilon, \key{retq}, \vec{\beta}, \Sigma) &\Leftarrow \kappa \approx \Sigma
\end{align*}
A CFG statement takes a step per operand, its instructions a few steps,
which must error exactly when the CFG does. So from each
synchronisation point, the start of a tail or a return ($s_1 \sim s_2$),
a script $\mathit{Fate}_i(s_1, s_2)$ says how the machines reach the
next one: both stop, the CFG steps while X86 waits (the index $i$ is the
length of the rest of the script), or both step:
\begin{align*}
  \exc{t}{\sigma}{\vec{t}}{\kappa} &\sim \xrun\;\mathit{ins}(t)\;\mathit{end}(t)\;\pass{sel}(\vec{t})\;\Sigma\;f\;\sigma\;\rho
    \Leftarrow \kappa \approx \Sigma \\
  \rti{v}{\sigma}{\kappa} &\sim \xrun\;\epsilon\;\key{retq}\;\vec{\beta}\;\Sigma\;f\;\sigma\;\rho
    \Leftarrow \kappa \approx \Sigma \land \rho(\rax) = v \\
  s_1 &\sim s_2 \Leftarrow \mathit{Fate}_i(s_1, s_2) \\[1ex]
  \mathit{Fate}_0(s_1, s_2) &\Leftarrow s_1 \mapsto \stp o \land s_2 \mapsto \stp o \\
  \mathit{Fate}_{i+1}(s_1, s_2) &\Leftarrow s_1 \lto_{\mathrm{CFG}} s_1' \land \mathit{Fate}_i(s_1', s_2) \\
  \mathit{Fate}_0(s_1, s_2) &\Leftarrow s_1 \lto_{\mathrm{CFG}} s_1' \land s_2 \lto^+_{\mathrm{X86}} s_2' \land s_1' \sim s_2'
\end{align*}
X86 runs an instruction at the CFG step that checks what it checks: an
operand as the CFG evaluates it, the destination slot when the CFG
assigns it. A call reads its operands in another order ($\rdi$ and
$\rsi$ before the code), so its moves are taken at the start of the
statement.

\begin{lemma}[Result to Its Destination]\label{lem:to-dest}
If $\mathit{Dest}(D, d, \vec{\imath}, j, \vec{\beta}, \Sigma)$ and $\rho(\rax) = v$, then
there is an $s_2$ with
\[
  \xrun\;(\key{movq}\; \rax, d;\; \vec{\imath})\;j\;\vec{\beta}\;\Sigma\;f\;\sigma\;\rho \lto^*_{\mathrm{X86}} s_2
  \quad\text{and}\quad \rti{v}{\sigma}{D} \sim s_2 .
\]
\end{lemma}

\begin{lemma}[Simple Expression Script]\label{lem:sel-simple-sim}
If $\mathit{Dest}(D, d, \vec{\imath}, j, \vec{\beta}, \Sigma)$ and
$s_1 \mapsto \mathit{res}(\mathit{first}(\sigma, r), \sigma, D)$, then
$\mathit{Fate}_i(s_1, \xrun\;(\pass{sel}(r, d); \vec{\imath})\;j\;\vec{\beta}\;\Sigma\;f\;\sigma\;\rho)$
for some $i$.
\end{lemma}

\begin{lemma}[Script]\label{lem:fate}
If $s_1 \sim s_2$, then $\mathit{Fate}_i(s_1, s_2)$ for some $i$.
\end{lemma}

\begin{lemma}[Select Simulation]\label{lem:select-fsim}
For every CFG program $P$,
\[
  \FSim(\Mach{\mathrm{CFG}}(P), \Mach{\mathrm{X86}}(\pass{select}(P)), \mathit{Fate}) .
\]
\end{lemma}

\begin{theorem}[Correctness of Select Instructions, Simulation]\label{thm:select-sim}
For every CFG program $P$,
\[
  \Equiv(\Mach{\mathrm{CFG}}(P), \Mach{\mathrm{X86}}(\pass{select}(P))) .
\]
\end{theorem}
There is no precondition.

\subsubsection{End to End}

\begin{theorem}[Compiler Correctness, Simulation]\label{thm:compiler-sim}
For every ISWIM term $M$,
\[
  \FV(M) = \emptyset \implies \Equiv(\Mach{\mathrm{ISWIM}}(M), \Mach{5}(\pass{compile}(M))) ,
\]
by the theorems for each pass, Lemma~\ref{lem:simplify-closed}, and Lemma~\ref{lem:equiv-trans}.
\end{theorem}

\begin{theorem}[Correctness of Compile-Mon, Simulation]\label{thm:compile-mon-sim}
For every ISWIM term $M$,
\[
  \FV(M) = \emptyset \implies \Equiv(\Mach{\mathrm{ISWIM}}(M), \Mach{5}(\pass{compile\text{-}mon}(M))) ,
\]
by Theorems~\ref{thm:compiler-sim} and~\ref{thm:rco-sim}, Lemma~\ref{lem:compile-scoped},
and Lemma~\ref{lem:equiv-trans}.
\end{theorem}

\begin{theorem}[Correctness of Compile-Imp, Simulation]\label{thm:compile-imp-sim}
For every ISWIM term $M$,
\[
  \FV(M) = \emptyset \implies \Equiv(\Mach{\mathrm{ISWIM}}(M), \Mach{\mathrm{Imp}}(\pass{compile\text{-}imp}(M))) ,
\]
by Theorems~\ref{thm:compile-mon-sim} and~\ref{thm:lift-sim} and Lemma~\ref{lem:equiv-trans}.
\end{theorem}

\begin{theorem}[Correctness of Compile-CFG, Simulation]\label{thm:compile-cfg-sim}
For every ISWIM term $M$,
\[
  \FV(M) = \emptyset \implies \Equiv(\Mach{\mathrm{ISWIM}}(M), \Mach{\mathrm{CFG}}(\pass{compile\text{-}cfg}(M))) ,
\]
by Theorems~\ref{thm:compile-imp-sim} and~\ref{thm:explicate-sim} and Lemma~\ref{lem:equiv-trans}.
\end{theorem}

\begin{theorem}[Correctness of Compile-X86, Simulation]\label{thm:compile-x86-sim}
For every ISWIM term $M$,
\[
  \FV(M) = \emptyset \implies \Equiv(\Mach{\mathrm{ISWIM}}(M), \Mach{\mathrm{X86}}(\pass{compile\text{-}x86}(M))) ,
\]
by Theorems~\ref{thm:compile-cfg-sim} and~\ref{thm:select-sim} and Lemma~\ref{lem:equiv-trans}.
\end{theorem}

%-------------------------------------------------------------------------------
\subsection{Backward Simulation for Delay}

The direct backward simulation for delay, which does not use determinism.
\begin{align*}
  \Div^\infty_L &\iff \exists g : \mathbb{N} \to S_L.\;
    g(0) = \start{L} \land \forall n.\; g(n) \lto_L g(n+1) \\[1ex]
  \BSim(L_1, L_2, \sim, \prec) &\iff
  \begin{array}[t]{l}
    \prec \text{ is well founded} \\
    \land\; \exists i.\; \start{L_1} \sim_i \start{L_2} \\
    \land\; \forall i\, s_1\, s_2\, o.\; s_1 \sim_i s_2 \land \Final_{L_2}(s_2, o) \implies
      \exists s_1'.\; s_1 \lto_{L_1}^* s_1' \land \Final_{L_1}(s_1', o) \\
    \land\; \forall i\, s_1\, s_2\, s_2'.\; s_1 \sim_i s_2 \land s_2 \lto_{L_2} s_2' \implies
      \exists i'\, s_1'.\; s_1' \sim_{i'} s_2' \\
    \qquad \land\; (s_1 \lto_{L_1}^+ s_1' \lor (s_1 \lto_{L_1}^* s_1' \land i' \prec i))
  \end{array}
\end{align*}

\begin{lemma}[Infinite Runs]\label{lem:div-inf}
If $\Div^\infty_L$, then $\Div_L$.
\end{lemma}

\begin{theorem}[Backward Simulation Reflects]\label{thm:bsim}
If $\BSim(L_1, L_2, \sim, \prec)$, then
\begin{align*}
  \Term_{L_2}(o) &\implies \Term_{L_1}(o) \\
  \Div^\infty_{L_2} &\implies \Div^\infty_{L_1} .
\end{align*}
\end{theorem}

The relations $\approx$, $\approxq$, and $\approxk$ are those of the forward
simulation for delay. The states are related with an index in $\mathbb{N}$,
the target steps left before the source moves.
\begin{align*}
  \ev{M}{\rho_1}{\kappa_1} &\sim_3 \ev{\pass{delay}(M)}{\rho_2}{\kappa_2} \\
  \rt{v}{\kappa_1} &\sim_9 \rt{w}{\kappa_2}
\end{align*}
when $\rho_1 \approxq \rho_2$, $\kappa_1 \approxk \kappa_2$, and $v \approx w$.
While the target builds a closure, with
$s_1 = \ev{(\clos{\bar{r}}{e\, y} N)}{\rho_1}{\kappa_1}$
and $\phi = \funV{e\, y.\, \pass{delay}(N)}$,
\begin{align*}
  s_1 &\sim_2 \ev{\fun{e\,y.\,\pass{delay}(N)}}{\rho_2}{\Kf{pairL}{\tuple{\bar{r}}, \rho_2, \kappa_2}} \\
  s_1 &\sim_1 \rt{\phi}{\Kf{pairL}{\tuple{\bar{r}}, \rho_2, \kappa_2}} \\
  s_1 &\sim_0 \ev{\tuple{\bar{r}}}{\rho_2}{\Kf{pairR}{\phi, \kappa_2}} .
\end{align*}
While the target calls a closure, with $L_M$ as in delay and
\[
  s_1 = \rt{v}{\Kf{appL}{M, \rho_1, \kappa_1}}, \qquad
  \rho_2' = \ext{\rho_2}{f}{w}, \qquad
  \kappa_F = \Kf{appF}{\SND f, \pass{delay}(M), \rho_2', \kappa_2},
\]
\begin{align*}
  s_1 &\sim_8 \ev{L_M}{\rho_2'}{\kappa_2} \\
  s_1 &\sim_7 \ev{(\FST f)}{\rho_2'}{\kappa_F} \\
  s_1 &\sim_6 \ev{f}{\rho_2'}{\Kf{fstK}{\kappa_F}} \\
  s_1 &\sim_5 \rt{w}{\Kf{fstK}{\kappa_F}}
\end{align*}
and, when moreover $v = \closV{e\, y.\, N}{v'}$ and $w = \pairV{\phi}{u}$ with $v' \approx u$,
\[
  \kappa_E = \Kf{appE}{\phi, \pass{delay}(M), \rho_2', \kappa_2},
\]
\begin{align*}
  s_1 &\sim_4 \rt{\phi}{\kappa_F} \\
  s_1 &\sim_3 \ev{(\SND f)}{\rho_2'}{\kappa_E} \\
  s_1 &\sim_2 \ev{f}{\rho_2'}{\Kf{sndK}{\kappa_E}} \\
  s_1 &\sim_1 \rt{w}{\Kf{sndK}{\kappa_E}} \\
  s_1 &\sim_0 \rt{u}{\kappa_E} .
\end{align*}

\begin{lemma}[Delay Backward Simulation]\label{lem:delay-bsim}
For every EnvClos term $M$,
\[
  \BSim(\Mach{3}(M), \Mach{4}(\pass{delay}(M)), \sim, <) .
\]
\end{lemma}

\begin{theorem}[Delay Reflects]\label{thm:delay-reflects}
For every EnvClos term $M$,
\begin{align*}
  \Term_{\Mach{4}(\pass{delay}(M))}(o) &\implies \Term_{\Mach{3}(M)}(o) \\
  \Div^\infty_{\Mach{4}(\pass{delay}(M))} &\implies \Div^\infty_{\Mach{3}(M)} .
\end{align*}
\end{theorem}

%===============================================================================
\section{Clocked Interpreter + Value Relations}\label{sec:u-interp}

\subsection{Results and Clocks}

\subsubsection{Results}

\begin{align*}
  r &::= \tmo \mid \error \mid \Ret v \\[1ex]
  \tmo \bind f &= \tmo \\
  \error \bind f &= \error \\
  \Ret v \bind f &= f(v) \\[1ex]
  \lift{v} &= \Ret v \\
  \lift{\bot} &= \error
\end{align*}

\subsubsection{Programs}

A clocked program gives a result for every clock.
\begin{align*}
  P &= (\obs_P,\; \run_P : \mathbb{N} \to \mathit{Result}) \\[1ex]
  \outc_P(\tmo) &= \bot \\
  \outc_P(\error) &= \error \\
  \outc_P(\Ret v) &= \obs_P(v) \\[1ex]
  \Term_P(o) &\iff \exists k.\; \outc_P(\run_P(k)) = o \\
  \Div_P &\iff \forall k.\; \run_P(k) = \tmo \\
  \Equiv(P_1, P_2) &\iff
    (\forall o.\; \Term_{P_1}(o) \iff \Term_{P_2}(o)) \land (\Div_{P_1} \iff \Div_{P_2})
\end{align*}
The observations $o$ are those of Section~\ref{sec:obs}, and $\Equiv$ has the
shape of the simulation one (Section~\ref{sec:sims}).

\subsubsection{Clock for Clock}

\begin{align*}
  \RelRes{R}(\tmo, \tmo), \quad \RelRes{R}(\error, \error) &\iff \mathsf{true} \\
  \RelRes{R}(\Ret v, \Ret w) &\iff v \mathrel{R} w \\
  \RelRes{R}(r_1, r_2) &\iff \mathsf{false} \quad \text{otherwise}
\end{align*}

\begin{lemma}[Related Binds]\label{lem:bind-rel}
If $\RelRes{R}(r_1, r_2)$ and $\RelRes{R'}(f(v), g(w))$ for every $v \mathrel{R} w$, then
\[
  \RelRes{R'}(r_1 \bind f,\; r_2 \bind g) .
\]
\end{lemma}

\begin{theorem}[Clock for Clock Gives Equivalence]\label{thm:clocked}
If $\obs_{P_1}(v) = \obs_{P_2}(w)$ for every $v \mathrel{R} w$ and
$\RelRes{R}(\run_{P_1}(k), \run_{P_2}(k))$ for every clock $k$, then
\[
  \Equiv(P_1, P_2) .
\]
\end{theorem}

\subsubsection{Clock Monotone}

A pass that removes calls is not clock for clock: its output needs no more
clock than its input, which may need more.
\begin{align*}
  \mathit{Le}_R(\tmo, r) &\iff \mathsf{true} \\
  \mathit{Le}_R(r_1, r_2) &\iff \RelRes{R}(r_1, r_2) \quad \text{for } r_1 \neq \tmo
\end{align*}

\begin{theorem}[Clock Monotone Gives Equivalence]\label{thm:clocked-le}
If $\obs_{P_1}(v) = \obs_{P_2}(w)$ for every $v \mathrel{R} w$, and for every clock $k$
\begin{align*}
  &\mathit{Le}_R(\run_{P_1}(k), \run_{P_2}(k)) , \\
  &\run_{P_2}(k) \neq \tmo \implies \forall k' \geq k.\; \run_{P_2}(k') = \run_{P_2}(k) , \\
  &\run_{P_2}(k) \neq \tmo \implies \exists k'.\; \run_{P_1}(k') \neq \tmo ,
\end{align*}
then
\[
  \Equiv(P_1, P_2) .
\]
\end{theorem}

\begin{lemma}[Clocked Equivalence Composes]\label{lem:equiv-trans-clk}
If $\Equiv(P_1, P_2)$ and $\Equiv(P_2, P_3)$, then
\[
  \Equiv(P_1, P_3) .
\]
\end{lemma}

%-------------------------------------------------------------------------------
\subsection{The Interpreters}

Each interpreter uses the values and environments of its language's machine
(Section~\ref{sec:machines}). The clock $k$ decreases only when a function is
called; every other construct evaluates its subterms with the same clock.
The do-notation is sugar for $\bind$, as in the Agda:
$\Do x \leftarrow e;\; e'$ means $e \bind \lambda x.\, e'$.
\begin{align*}
  \EV{\ell}{k} &: \mathit{Env} \to \mathit{Term} \to \mathit{Result}
    \qquad (\ell \in \{\mathrm{ISWIM}, 2, 3, 4\}) \\
  \EV{\vec{d}}{k} &: \mathit{Env} \to \mathit{Term} \to \mathit{Result}
    \qquad (\vec{d} = d_0, \ldots, d_{m-1}) \\[1ex]
  \Interp{\ell}(M) &= (\obs,\; \lambda k.\, \EV{\ell}{k}\,\emptyset\,M) \\
  \Interp{5}(\vec{d};\; M) &= (\obs,\; \lambda k.\, \EV{\vec{d}}{k}\,\emptyset\,M)
\end{align*}

\subsubsection{Clauses common to all interpreters}

\begin{align*}
  \EV{\ell}{k}\,\rho\,x &= \lift{\rho(x)} \\
  \EV{\ell}{k}\,\rho\,c &= \Ret c \\
  \EV{\ell}{k}\,\rho\,\tuple{M_1, \ldots, M_n} &=
    \Do \bar{v} \leftarrow \EVS{\ell}{k}\,\rho\,(M_1 \ldots M_n);\; \Ret \tupV{\bar{v}} \\
  \EV{\ell}{k}\,\rho\,(\prj{i} M) &= \EV{\ell}{k}\,\rho\,M \bind \key{get}^n_i \\
  \EV{\ell}{k}\,\rho\,(\inl M) &= \Do v \leftarrow \EV{\ell}{k}\,\rho\,M;\; \Ret (\inlV v) \\
  \EV{\ell}{k}\,\rho\,(\inr M) &= \Do v \leftarrow \EV{\ell}{k}\,\rho\,M;\; \Ret (\inrV v) \\
  \EV{\ell}{k}\,\rho\,(\Case{L}{x}{M}{y}{N}) &=
    \EV{\ell}{k}\,\rho\,L \bind \BR{\ell}{k}\,\rho\,(x.M)\,(y.N) \\
  \EV{\ell}{k}\,\rho\,\prim{p}{L}{M} &=
    \Do v \leftarrow \EV{\ell}{k}\,\rho\,L;\; w \leftarrow \EV{\ell}{k}\,\rho\,M;\; \key{prim}_p(v, w) \\[1ex]
  \EVS{\ell}{k}\,\rho\,\epsilon &= \Ret \epsilon \\
  \EVS{\ell}{k}\,\rho\,(M\,\bar{M}) &=
    \Do v \leftarrow \EV{\ell}{k}\,\rho\,M;\; \bar{v} \leftarrow \EVS{\ell}{k}\,\rho\,\bar{M};\; \Ret (v\,\bar{v}) \\[1ex]
  \key{get}^n_i(v) &=
    \begin{cases} \lift{\key{proj}^n_i(\bar{v})} & v = \tupV{\bar{v}} \\ \error & \text{otherwise} \end{cases} \\
  \key{prim}_p(v, w) &=
    \begin{cases} \Ret \res{\delta(p, n, m)} & v = n,\; w = m \\ \error & \text{otherwise} \end{cases} \\
  \BR{\ell}{k}\,\rho\,(x.M)\,(y.N)\,v &=
    \begin{cases}
      \EV{\ell}{k}\,\ext{\rho}{x}{u}\,M & v = \inlV u \\
      \EV{\ell}{k}\,\ext{\rho}{y}{u}\,N & v = \inrV u \\
      \error & \text{otherwise}
    \end{cases} \\[1ex]
  \CALL{\ell}{0}\,N\,\rho &= \tmo \\
  \CALL{\ell}{k+1}\,N\,\rho &= \EV{\ell}{k}\,\rho\,N
\end{align*}
Here $\key{proj}^n_i$ and $\res{-}$ are those of Section~\ref{sec:machine-shared}.

\subsubsection{ISWIM}

\begin{align*}
  \EV{\mathrm{ISWIM}}{k}\,\rho\,(\lam{x} N) &= \Ret \cloV{\lam{x} N}{\rho} \\
  \EV{\mathrm{ISWIM}}{k}\,\rho\,(L\;M) &= \EV{\mathrm{ISWIM}}{k}\,\rho\,L \bind \APP{\mathrm{ISWIM}}{k}\,\rho\,M \\[1ex]
  \APP{\mathrm{ISWIM}}{k}\,\rho\,M\,f &=
    \begin{cases}
      \Do w \leftarrow \EV{\mathrm{ISWIM}}{k}\,\rho\,M;\; \CALL{\mathrm{ISWIM}}{k}\,N\,\ext{\rho'}{x}{w}
        & f = \cloV{\lam{x} N}{\rho'} \\
      \error & \text{otherwise}
    \end{cases}
\end{align*}

\subsubsection{Clos}

\begin{align*}
  \EV{2}{k}\,\rho\,(\clos{z_1 \ldots z_n}{x_1 \ldots x_n\, y} N) &=
    \Do \bar{v} \leftarrow \EVS{2}{k}\,\rho\,(z_1 \ldots z_n);\; \Ret \closV{x_1 \ldots x_n\, y.\, N}{\bar{v}} \\
  \EV{2}{k}\,\rho\,(L\;M) &= \EV{2}{k}\,\rho\,L \bind \APP{2}{k}\,\rho\,M \\[1ex]
  \APP{2}{k}\,\rho\,M\,f &=
    \begin{cases}
      \Do w \leftarrow \EV{2}{k}\,\rho\,M;\; \CALL{2}{k}\,N\,\{\bar{x} {\mapsto} \bar{v}, y {\mapsto} w\}
        & f = \closV{\bar{x}\, y.\, N}{\bar{v}} \\
      \error & \text{otherwise}
    \end{cases}
\end{align*}

\subsubsection{EnvClos}

\begin{align*}
  \EV{3}{k}\,\rho\,(\clos{r_1 \ldots r_n}{e\, y} N) &=
    \Do \bar{v} \leftarrow \EVS{3}{k}\,\rho\,(r_1 \ldots r_n);\; \Ret \closV{e\, y.\, N}{\tupV{\bar{v}}} \\
  \EV{3}{k}\,\rho\,(L\;M) &= \EV{3}{k}\,\rho\,L \bind \APP{3}{k}\,\rho\,M \\[1ex]
  \APP{3}{k}\,\rho\,M\,f &=
    \begin{cases}
      \Do w \leftarrow \EV{3}{k}\,\rho\,M;\; \CALL{3}{k}\,N\,\{e {\mapsto} u, y {\mapsto} w\}
        & f = \closV{e\, y.\, N}{u} \\
      \error & \text{otherwise}
    \end{cases}
\end{align*}

\subsubsection{FlatClos}

As in the machine, the operator is checked after all three operands.
\begin{align*}
  \EV{4}{k}\,\rho\,\fun{e\,y.\,N} &= \Ret \funV{e\, y.\, N} \\
  \EV{4}{k}\,\rho\,L(M, N) &= \EV{4}{k}\,\rho\,L \bind \APP{4}{k}\,\rho\,M\,N \\
  \EV{4}{k}\,\rho\,\pair{M}{N} &=
    \Do v \leftarrow \EV{4}{k}\,\rho\,M;\; w \leftarrow \EV{4}{k}\,\rho\,N;\; \Ret \pairV{v}{w} \\
  \EV{4}{k}\,\rho\,(\FST M) &= \EV{4}{k}\,\rho\,M \bind \key{first} \\
  \EV{4}{k}\,\rho\,(\SND M) &= \EV{4}{k}\,\rho\,M \bind \key{second} \\
  \EV{4}{k}\,\rho\,(\Let x = M \In N) &=
    \Do v \leftarrow \EV{4}{k}\,\rho\,M;\; \EV{4}{k}\,\ext{\rho}{x}{v}\,N \\[1ex]
  \APP{4}{k}\,\rho\,M\,N\,f &=
    \Do u \leftarrow \EV{4}{k}\,\rho\,M;\; w \leftarrow \EV{4}{k}\,\rho\,N;\;
    \begin{cases}
      \CALL{4}{k}\,B\,\{e {\mapsto} u, y {\mapsto} w\} & f = \funV{e\, y.\, B} \\
      \error & \text{otherwise}
    \end{cases} \\
  \key{first}(v) &=
    \begin{cases} \Ret v_1 & v = \pairV{v_1}{v_2} \\ \error & \text{otherwise} \end{cases} \\
  \key{second}(v) &=
    \begin{cases} \Ret v_2 & v = \pairV{v_1}{v_2} \\ \error & \text{otherwise} \end{cases}
\end{align*}

\subsubsection{GlobalFun}

The interpreter of a program runs with its global definitions $\vec{d} = d_0, \ldots, d_{m-1}$.
Its clauses are those of FlatClos, with $\EV{\vec{d}}{k}$ in place of $\EV{4}{k}$, except
\begin{align*}
  \EV{\vec{d}}{k}\,\rho\,\funref{j} &= \Ret \codeV{j} \\[1ex]
  \APP{\vec{d}}{k}\,\rho\,M\,N\,f &=
    \Do u \leftarrow \EV{\vec{d}}{k}\,\rho\,M;\; w \leftarrow \EV{\vec{d}}{k}\,\rho\,N;\;
    \begin{cases}
      \CALL{\vec{d}}{k}\,j\,u\,w & f = \codeV{j} \\
      \error & \text{otherwise}
    \end{cases} \\
  \CALL{\vec{d}}{k}\,j\,u\,w &=
    \begin{cases}
      \error & j \geq m \\
      \tmo & j < m,\; k = 0 \\
      \EV{\vec{d}}{k'}\,\{e {\mapsto} u, y {\mapsto} w\}\,N & j < m,\; k = k' + 1,\; d_j = \key{fun}\; e\,y.\,N
    \end{cases}
\end{align*}

\subsubsection{Imp}

The interpreter of a program runs with its bodies $\vec{b} = b_0, \ldots, b_{m-1}$.
An expression runs in a frame $\sigma$, with $\sigma(a)$ and $\sigma[s := v]$ as in
the machine, and returns a value and the frame as it left it. The value
of a simple expression, given how to call an operator ($C$), its
injection test, and assignment are shared with CFG. As in GlobalFun, the
operator of a call is checked after the operands, and only calls tick.
\begin{align*}
  \EV{\vec{b}}{k} &: \mathit{Frame} \to \mathit{Exp} \to \mathit{Result}(\mathit{Val} \times \mathit{Frame}) \\
  \Interp{\mathrm{Imp}}(\vec{b};\; (n;\; e)) &= (\obs,\; \lambda k.\, \Do (v, \sigma) \leftarrow \EV{\vec{b}}{k}\,0^n\,e;\; \Ret v)
\end{align*}
\begin{align*}
  \key{simple}_C(\sigma, a) &= \lift{\sigma(a)} \\
  \key{simple}_C(\sigma, a_1(a_2, a_3)) &=
    \Do f \leftarrow \lift{\sigma(a_1)};\; u \leftarrow \lift{\sigma(a_2)};\; w \leftarrow \lift{\sigma(a_3)};\; C(f, u, w) \\
  \key{simple}_C(\sigma, \pair{a_1}{a_2}) &=
    \Do v \leftarrow \lift{\sigma(a_1)};\; w \leftarrow \lift{\sigma(a_2)};\; \Ret \pairV{v}{w} \\
  \key{simple}_C(\sigma, \FST a) &= \lift{\sigma(a)} \bind \key{first} \qquad
  \key{simple}_C(\sigma, \SND a) = \lift{\sigma(a)} \bind \key{second} \\
  \key{simple}_C(\sigma, \tuple{a_1, \ldots, a_l}) &=
    \Do v_1 \leftarrow \lift{\sigma(a_1)};\; \ldots;\; v_l \leftarrow \lift{\sigma(a_l)};\; \Ret \tupV{v_1, \ldots, v_l} \\
  \key{simple}_C(\sigma, \prj{i} a) &= \lift{\sigma(a)} \bind \key{get}^n_i \\
  \key{simple}_C(\sigma, \inl a) &= \Do v \leftarrow \lift{\sigma(a)};\; \Ret (\inlV v) \qquad
  \key{simple}_C(\sigma, \inr a) = \Do v \leftarrow \lift{\sigma(a)};\; \Ret (\inrV v) \\
  \key{simple}_C(\sigma, \prim{p}{a_1}{a_2}) &=
    \Do v \leftarrow \lift{\sigma(a_1)};\; w \leftarrow \lift{\sigma(a_2)};\; \key{prim}_p(v, w) \\[1ex]
  \key{inj?}(v) &= \begin{cases}
      \Ret \iota_1(u) & v = \inlV u \\
      \Ret \iota_2(u) & v = \inrV u \\
      \error & \text{otherwise}
    \end{cases} \qquad
  [G_1, G_2](\iota_i(u)) = G_i(u) \\
  \key{set}(\sigma, s, v, K) &= \lift{\sigma[s := v]} \bind K
\end{align*}
where $\key{simple}$ is the Agda's $\key{eval\text{-}simple}$.
\begin{align*}
  \EV{\vec{b}}{k}\,\sigma\,r &= \key{simple}_{\CALL{\vec{b}}{k}}(\sigma, r) \bind \key{here}_\sigma \\
  \EV{\vec{b}}{k}\,\sigma\,e &=
    \Do v \leftarrow \lift{\sigma(a)};\; \key{inj?}(v) \bind [\lambda u.\, \key{put}^k_{\vec{b}}\,\sigma\,s\,u\,e_1,\;
      \lambda u.\, \key{put}^k_{\vec{b}}\,\sigma\,s\,u\,e_2] \\
    &\qquad (e = \Casei{a}{s}{e_1}{e_2}) \\
  \EV{\vec{b}}{k}\,\sigma\,(\assign{s}{e_1}{e_2}) &=
    \Do (v, \sigma_1) \leftarrow \EV{\vec{b}}{k}\,\sigma\,e_1;\; \key{put}^k_{\vec{b}}\,\sigma_1\,s\,v\,e_2 \\[1ex]
  \key{here}_\sigma(v) &= \Ret (v, \sigma) \\
  \key{put}^k_{\vec{b}}\,\sigma\,s\,v\,e &= \key{set}(\sigma, s, v, \lambda \sigma'.\, \EV{\vec{b}}{k}\,\sigma'\,e) \\
  \CALL{\vec{b}}{k}(f, u, w) &=
    \begin{cases}
      \tmo & f = \codeV{j},\; j < m,\; k = 0 \\
      \Do (v, \sigma') \leftarrow \EV{\vec{b}}{k'}\,(u\, w\, 0^{n-2})\,e;\; \Ret v
        & f = \codeV{j},\; j < m,\; k = k' + 1, \\
        & \quad b_j = (n;\; e) \\
      \error & \text{otherwise}
    \end{cases}
\end{align*}

\subsubsection{CFG}

The interpreter of a program runs with its bodies $\vec{b} = b_0, \ldots, b_{m-1}$.
A tail runs in a frame with the blocks it may jump to, and returns a
value. The frames, $\key{simple}$, $\key{inj?}$, and $\key{set}$ are
those of Imp. Only calls tick; a jump does not, and runs its block with
the blocks older than it.
\begin{align*}
  \key{tail}^k_{\vec{b}} &: \mathit{Blocks} \to \mathit{Frame} \to \mathit{Tail} \to \mathit{Result}(\mathit{Val}) \\
  \Interp{\mathrm{CFG}}(\vec{b};\; (n;\; \vec{t};\; t)) &= (\obs,\; \lambda k.\, \key{tail}^k_{\vec{b}}\,\vec{t}\,0^n\,t)
\end{align*}
\begin{align*}
  \key{tail}^k_{\vec{b}}\,\vec{t}\,\sigma\,(\key{return}\; r) &= \key{simple}_{\CALL{\vec{b}}{k}}(\sigma, r) \\
  \key{tail}^k_{\vec{b}}\,\vec{t}\,\sigma\,(\assign{s}{r}{t}) &=
    \Do v \leftarrow \key{simple}_{\CALL{\vec{b}}{k}}(\sigma, r);\;
    \key{set}(\sigma, s, v, \lambda \sigma'.\, \key{tail}^k_{\vec{b}}\,\vec{t}\,\sigma'\,t) \\
  \key{tail}^k_{\vec{b}}\,\vec{t}\,\sigma\,(\key{goto}\; l) &= \key{jump}^k_{\vec{b}}\,\vec{t}\,\sigma\,l \\
  \key{tail}^k_{\vec{b}}\,\vec{t}\,\sigma\,(\key{case}\;a\;\key{of}_{s}\;\key{goto}\;l_1 \mid \key{goto}\;l_2) &=
    \Do v \leftarrow \lift{\sigma(a)};\; \key{inj?}(v) \bind [G_{l_1}, G_{l_2}] \\
  &\qquad\text{where } G_l(u) = \key{set}(\sigma, s, u, \lambda \sigma'.\, \key{jump}^k_{\vec{b}}\,\vec{t}\,\sigma'\,l) \\[1ex]
  \key{jump}^k_{\vec{b}}\,(t_0, \ldots, t_{j-1})\,\sigma\,l &=
    \begin{cases}
      \key{tail}^k_{\vec{b}}\,(t_0, \ldots, t_{l-1})\,\sigma\,t_l & l < j \\
      \error & \text{otherwise}
    \end{cases} \\
  \CALL{\vec{b}}{k}(f, u, w) &=
    \begin{cases}
      \tmo & f = \codeV{j},\; j < m,\; k = 0 \\
      \key{tail}^{k'}_{\vec{b}}\,\vec{t}\,(u\, w\, 0^{n-2})\,t
        & f = \codeV{j},\; j < m,\; k = k' + 1,\; b_j = (n;\; \vec{t};\; t) \\
      \error & \text{otherwise}
    \end{cases}
\end{align*}

\subsubsection{X86}

The interpreter of a program runs with its bodies $\vec{b} = b_0, \ldots, b_{m-1}$.
The state is a frame $\sigma$ and registers $\rho$: $\rax$, $\rdi$, $\rsi$
and the flag each hold a value or are unset ($\bot$); reading an unset
one errors. The frames and $\key{set}$ are those
of Imp, and $\key{prim}^?$, $\key{cmp}$, $\key{isinl}$, $\key{build}$, and
$\key{field}$ those of the X86 machine. A block runs its instructions,
each passing the state to the rest, then its jump. Only $\key{callq}$
ticks; it reads $\rdi$ and $\rsi$ and runs the start of the callee in a
fresh frame with them, $\rax$ unset and no flag, and the caller continues in its own frame with the
result in $\rax$ and $\rdi$, $\rsi$ and the flag unset.
\begin{align*}
  \rho &= (\rax{:}\; v^?,\; \rdi{:}\; v^?,\; \rsi{:}\; v^?,\; \key{flag}{:}\; b^?) \\
  \key{run}^k_{\vec{b}} &: \mathit{Blocks} \to \mathit{Frame} \times \mathit{Regs} \to \mathit{Instr}^* \to \mathit{Jump} \to \mathit{Result}(\mathit{Val}) \\
  \Interp{\mathrm{X86}}(\vec{b};\; (n;\; \vec{\beta};\; \vec{\imath};\, j)) &=
    (\obs,\; \lambda k.\, \key{run}^k_{\vec{b}}\,\vec{\beta}\,(0^n, (0, \bot, \bot, \bot))\,\vec{\imath}\,j)
\end{align*}
With $\sigma$ and $\rho$ those of the state $(\sigma, \rho)$:
\begin{align*}
  \mathit{arg}(s) &= \lift{\sigma(s)}, \quad
  \mathit{arg}(\%g) = \lift{\rho.\%g}, \\
  \mathit{arg}(c) &= \Ret c, \quad
  \mathit{arg}(\funref{k}) = \Ret \codeV{k} \\
  \key{write}(s, v, K) &= \key{set}(\sigma, s, v, \lambda \sigma'.\, K(\sigma', \rho)) \\
  \key{write}(\%g, v, K) &= K(\sigma, \rho[\%g := v]) \\[1ex]
  \key{step}(\key{movq}\; o, d, K) &= \Do v \leftarrow \mathit{arg}(o);\; \key{write}(d, v, K) \\
  \key{step}(\key{subq}\; o, d, K) &= \Do w \leftarrow \mathit{arg}(o);\; v \leftarrow \mathit{arg}(d);\;
    u \leftarrow \lift{\key{prim}^?_{\key{sub}}(v, w)};\; \key{write}(d, u, K) \\
  \key{step}(\key{cmpq}\; o_2, o_1, K) &= \Do v \leftarrow \mathit{arg}(o_1);\; w \leftarrow \mathit{arg}(o_2);\;
    b \leftarrow \lift{\key{cmp}(v, w)};\; K(\sigma, \rho[\key{flag} := b]) \\
  \key{step}(\key{sete}\; d, K) &= \Do b \leftarrow \lift{\rho.\key{flag}};\; \key{write}(d, \res{b}, K) \\
  \key{step}(\key{testinl}\; o, K) &= \Do v \leftarrow \mathit{arg}(o);\; b \leftarrow \lift{\key{isinl}(v)};\;
    K(\sigma, \rho[\key{flag} := b]) \\
  \key{step}(\key{alloc}\; \mathit{sh}\; o_1, \ldots, o_p, d, K) &=
    \Do v_1 \leftarrow \mathit{arg}(o_1);\; \ldots;\; v_p \leftarrow \mathit{arg}(o_p);\; \\
    &\qquad v \leftarrow \lift{\key{build}_{\mathit{sh}}(v_1, \ldots, v_p)};\; \key{write}(d, v, K) \\
  \key{step}(\key{load}\; \mathit{fl}\; o, d, K) &= \Do v \leftarrow \mathit{arg}(o);\;
    w \leftarrow \lift{\key{field}_{\mathit{fl}}(v)};\; \key{write}(d, w, K)
\end{align*}
Writing $\key{run}^k$, $\key{jump}^k$ and $\CALL{}{k}$ for those of $\vec{b}$:
\begin{align*}
  \key{run}^k\,\vec{\beta}\,(\sigma, \rho)\,(\key{callq}\; o;\; \vec{\imath})\,j &=
    \Do f \leftarrow \mathit{arg}(o);\; e \leftarrow \lift{\rho.\rdi};\; w \leftarrow \lift{\rho.\rsi};\;
    v \leftarrow \CALL{}{k}(f, e, w); \\
    &\qquad \key{run}^k\,\vec{\beta}\,(\sigma, (v, \bot, \bot, \bot))\,\vec{\imath}\,j \\
  \key{run}^k\,\vec{\beta}\,(\sigma, \rho)\,(i;\; \vec{\imath})\,j &=
    \key{step}(i, \lambda (\sigma', \rho').\, \key{run}^k\,\vec{\beta}\,(\sigma', \rho')\,\vec{\imath}\,j)
    \quad (i \neq \key{callq}\; o) \\
  \key{run}^k\,\vec{\beta}\,(\sigma, \rho)\,\epsilon\,\key{retq} &= \lift{\rho.\rax} \\
  \key{run}^k\,\vec{\beta}\,(\sigma, \rho)\,\epsilon\,(\key{jmp}\; l) &= \key{jump}^k\,\vec{\beta}\,(\sigma, \rho)\,l \\
  \key{run}^k\,\vec{\beta}\,(\sigma, \rho)\,\epsilon\,(\key{je}\; l_1, l_2) &=
    \Do b \leftarrow \lift{\rho.\key{flag}};\; \key{jump}^k\,\vec{\beta}\,(\sigma, \rho)\,(\text{$l_1$ if $b$, else $l_2$}) \\[1ex]
  \key{jump}^k\,(\beta_0, \ldots, \beta_{q-1})\,(\sigma, \rho)\,l &=
    \begin{cases}
      \key{run}^k\,(\beta_0, \ldots, \beta_{l-1})\,(\sigma, \rho)\,\vec{\imath}\,j & l < q,\; \beta_l = \vec{\imath};\, j \\
      \error & \text{otherwise}
    \end{cases} \\
  \CALL{}{k}(f, e, w) &=
    \begin{cases}
      \tmo & f = \codeV{j},\; j < m,\; k = 0 \\
      \key{run}^{k'}\,\vec{\beta}\,(0^{\max(2,n)}, (\bot, e, w, \bot))\,\vec{\imath}\,j'
        & f = \codeV{j},\; j < m,\; k = k' + 1, \\
        & b_j = (n;\; \vec{\beta};\; \vec{\imath};\, j') \\
      \error & \text{otherwise}
    \end{cases}
\end{align*}

%-------------------------------------------------------------------------------
\subsection{Relations}

Each pass relates source and target values by the relation $\approx$ of its
forward simulation (Section~\ref{sec:sim-relations}), and environments by the
relation used there for the states $\ev{M}{\rho_1}{\kappa_1} \sim \ev{M'}{\rho_2}{\kappa_2}$.
By Lemma~\ref{lem:obs}, if $v \approx w$, then $\obs(v) = \obs(w)$.

\subsubsection{Simplify}

The relations $\approx$ and $\models_E$ are those of
Section~\ref{sec:sim-simplify}. Simplify removes calls ($\beta$ of a small
value, inlining), so the proof is clock monotone (Theorem~\ref{thm:clocked-le}), from
the monotonicity of the ISWIM interpreter in the clock and two lemmas.

\begin{lemma}[Simplify, Forward at the Same Clock]
If $\rho \models_E \rho'$, then for every clock $k$,
\[
  \mathit{Le}_{\approx}(\EV{\mathrm{ISWIM}}{k}\,\rho\,M,\; \EV{\mathrm{ISWIM}}{k}\,\rho'\,\pass{simp}_E(M)) .
\]
\end{lemma}

\begin{lemma}[Simplify, Backward]
If $\rho \models_E \rho'$ and $r' = \EV{\mathrm{ISWIM}}{k}\,\rho'\,\pass{simp}_E(M) \neq \tmo$, then there are
$r$ and $k_0$ with
\[
  \RelRes{\approx}(r, r') \land \forall k' \geq k_0.\; \EV{\mathrm{ISWIM}}{k'}\,\rho\,M = r .
\]
\end{lemma}

\begin{theorem}[Correctness of Simplify, Interpreter]\label{thm:simplify-clk}
For every ISWIM term $M$,
\[
  \FV(M) = \emptyset \implies
  \Equiv(\Interp{\mathrm{ISWIM}}(M), \Interp{\mathrm{ISWIM}}(\pass{simplify}(M))) .
\]
\end{theorem}

\subsubsection{CC: Enclose}

The relations $\approx$ and $\approx_X$ are those of Section~\ref{sec:sim-enclose}.

\begin{lemma}[Enclose, Clock for Clock]\label{lem:enclose-clk}
If $\rho_1 \approx_{\FV(M)} \rho_2$, then for every clock $k$,
\[
  \RelRes{\approx}(\EV{\mathrm{ISWIM}}{k}\,\rho_1\,M,\; \EV{2}{k}\,\rho_2\,\pass{enclose}(M)) .
\]
\end{lemma}

\begin{theorem}[Correctness of Enclose, Interpreter]\label{thm:enclose-clk}
For every ISWIM term $M$,
\[
  \FV(M) = \emptyset \implies
  \Equiv(\Interp{\mathrm{ISWIM}}(M), \Interp{2}(\pass{enclose}(M))) .
\]
\end{theorem}

\subsubsection{CC: Optimize}

The relations $\approx$, $\approx_X$, $\mathit{Scoped}$, $\mathit{WS}$, and
$\mathit{Built}^j_N$ are those of Section~\ref{sec:sim-optimize}.

\begin{lemma}[Optimize, Clock for Clock]\label{lem:optimize-clk}
If $\mathit{Scoped}(M)$ and $\rho_1 \approx_{\FV(M)} \rho_2$, then for every clock $k$,
\[
  \RelRes{\approx}(\EV{2}{k}\,\rho_1\,M,\; \EV{2}{k}\,\rho_2\,\pass{optimize}(M)) .
\]
\end{lemma}

\begin{lemma}[Free Variables of a Closure]\label{lem:prune-clk}
If $\mathit{Scoped}(N)$ and $\rho_1 \approx_{\{z_1, \ldots, z_n\}} \rho_2$,
then for every clock $k$,
\[
  \RelRes{\mathit{Built}^{n+1}_N}(\EVS{2}{k}\,\rho_1\,(z_1 \ldots z_n),\;
    \EVS{2}{k}\,\rho_2\,\bar{z}'_{>0})
\]
for the closure $\clos{z_1 \ldots z_n}{x_1 \ldots x_n\, y} N$.
\end{lemma}

\begin{theorem}[Correctness of Optimize, Interpreter]\label{thm:optimize-clk}
For every Clos term $M$,
\[
  \mathit{WS}(M) \implies \Equiv(\Interp{2}(M), \Interp{2}(\pass{optimize}(M))) .
\]
\end{theorem}

\subsubsection{CC: Concretize}

The relations $\approx$ and $\approx_\varphi$ are those of Section~\ref{sec:sim-concretize}.

\begin{lemma}[Concretize, Clock for Clock]\label{lem:concretize-clk}
If $\rho_2 \approx_\varphi \rho_3$, then for every clock $k$,
\[
  \RelRes{\approx}(\EV{2}{k}\,\rho_2\,M,\; \EV{3}{k}\,\rho_3\,\pass{concretize}_\varphi(M)) .
\]
\end{lemma}

\begin{theorem}[Correctness of Concretize, Interpreter]\label{thm:concretize-clk}
For every Clos term $M$,
\[
  \Equiv(\Interp{2}(M), \Interp{3}(\pass{concretize}_{\mathit{id}}(M))) .
\]
\end{theorem}

\subsubsection{CC: Delay}

The relations $\approx$ and $\approxq$ are those of Section~\ref{sec:sim-delay}.

\begin{lemma}[Delay, Clock for Clock]\label{lem:delay-clk}
If $\rho_1 \approxq \rho_2$, then for every clock $k$,
\[
  \RelRes{\approx}(\EV{3}{k}\,\rho_1\,M,\; \EV{4}{k}\,\rho_2\,\pass{delay}(M)) .
\]
\end{lemma}

\begin{theorem}[Correctness of Delay, Interpreter]\label{thm:delay-clk}
For every EnvClos term $M$,
\[
  \Equiv(\Interp{3}(M), \Interp{4}(\pass{delay}(M))) .
\]
\end{theorem}

\subsubsection{CC: Globalize}

The relations $\leadsto$ and $\approx$, for the table $D$ that globalize makes
for the whole program, are those of Section~\ref{sec:sim-globalize}.

\begin{lemma}[Globalize, Clock for Clock]\label{lem:globalize-clk}
If $M \leadsto M'$ and $\rho_1 \approx \rho_2$, then for every clock $k$,
\[
  \RelRes{\approx}(\EV{4}{k}\,\rho_1\,M,\; \EV{D}{k}\,\rho_2\,M') .
\]
\end{lemma}

\begin{theorem}[Correctness of Globalize, Interpreter]\label{thm:globalize-clk}
For every FlatClos term $M$,
\[
  \Equiv(\Interp{4}(M), \Interp{5}(\pass{globalize}(M))) .
\]
\end{theorem}

\subsubsection{Remove Complex Operands}

Values are related by equality, as in Section~\ref{sec:sim-rco}, and
both programs run in the same environment: the fresh variables of rco
are not in $\rho$. The pass adds only lets, which do not tick, and atoms,
which call nothing.

\begin{lemma}[Let Chain, Clock for Clock]\label{lem:let-chain-clk}
Let $(\vec{b};\; a_1, \ldots, a_n) = \pass{atomize}(M_1, \ldots, M_n)$, and
let every operand $M_i$ that is a variable be in $\mathrm{dom}(\rho)$. If,
for every $\rho'$ that extends $\rho$ with the variables of $\vec{b}$,
$\EVS{\vec{d}}{k}\,\rho'\,(a_1 \ldots a_n) \bind K = \EV{\vec{d}}{k}\,\rho'\,X$, then
\[
  \EVS{\vec{d}}{k}\,\rho\,(M_1 \ldots M_n) \bind K = \EV{\vec{d}}{k}\,\rho\,\pass{lets}(\vec{b}, X) .
\]
\end{lemma}

\begin{lemma}[RCO, Clock for Clock]\label{lem:rco-clk}
If $\FV(d_j) = \emptyset$ for every $j$ and $\FV(M) \subseteq \mathrm{dom}(\rho)$,
then for every clock $k$,
\[
  \EV{\vec{d}}{k}\,\rho\,M = \EV{\pass{rco}(\vec{d})}{k}\,\rho\,\pass{rco}(M) .
\]
\end{lemma}

\begin{theorem}[Correctness of RCO, Interpreter]\label{thm:rco-clk}
For every GlobalFun program $P$,
\[
  \mathit{Scoped}_5(P) \implies \Equiv(\Interp{5}(P), \Interp{5}(\pass{rco}(P))) .
\]
\end{theorem}

\subsubsection{Lift Locals}

Values are related by equality, as in Section~\ref{sec:sim-lift}, and frames
to environments by $\approx^n_\varphi$. A result of the target, a value and
the frame it left, is related to a value of the source by
\begin{align*}
  \mathit{Out}^n_{\rho,\varphi}(v, (v', \sigma')) &\iff v = v' \land \rho \approx^n_\varphi \sigma' .
\end{align*}
The pass maps each call to one call, so the interpreters tick in lock step.

\begin{lemma}[Lift, Clock for Clock]\label{lem:lift-clk}
Let $M$ be a monadic GlobalFun term and $\vec{d}$ monadic GlobalFun definitions.
If $|\rho| + \mathit{height}(M) \leq n$ and $\rho \approx^n_\varphi \sigma$, then for every clock $k$,
\[
  \RelRes{\mathit{Out}^n_{\rho,\varphi}}(\EV{\vec{d}}{k}\,\rho\,M,\;
    \EV{\pass{lift}(\vec{d})}{k}\,\sigma\,\pass{lift}_{n,\varphi,|\rho|}(M)) .
\]
\end{lemma}

\begin{theorem}[Correctness of Lift Locals, Interpreter]\label{thm:lift-clk}
For every monadic GlobalFun program $P$,
\[
  \Equiv(\Interp{5}(P), \Interp{\mathrm{Imp}}(\pass{lift}(P))) .
\]
\end{theorem}

\subsubsection{Explicate Control}

The pass keeps every call, so the interpreters tick in lock step and
return equal results; $\codeV{j}$ names body $b_j$ in the source and
$\pass{explicate}(b_j)$ in the target. As in
Section~\ref{sec:den-explicate}, the pass threads the list of blocks, and
\begin{align*}
  \vec{t} \sqsubseteq \vec{t}' &\iff \vec{t} \text{ is a prefix of } \vec{t}' .
\end{align*}
Fix the bodies $\vec{b}$ and a clock $k$ at which calls agree:
$\CALL{\vec{b}}{k}(f, u, w) = \CALL{\pass{explicate}(\vec{b})}{k}(f, u, w)$
for all $f$, $u$, and $w$. Write $\key{tail}^k$ and $\key{jump}^k$ for
those of $\pass{explicate}(\vec{b})$, and $\EV{}{k}$ for $\EV{\vec{b}}{k}$.

\begin{lemma}[Jump to a Block]\label{lem:jump-at}
If $\vec{t}, t \sqsubseteq \vec{t}'$, then
$\key{jump}^k\,\vec{t}'\,\sigma\,|\vec{t}| = \key{tail}^k\,\vec{t}\,\sigma\,t$.
\end{lemma}

A tail made with the blocks $\vec{t}$ means $K$ when it returns
$K(\sigma)$ with any newer blocks, and so does the output of the pass:
\begin{align*}
  \mathit{Means}_{\vec{t}}(t, K) &\iff \forall \vec{t}' \sqsupseteq \vec{t}.\; \forall \sigma.\;
    \key{tail}^k\,\vec{t}'\,\sigma\,t = K(\sigma) \\
  \mathit{Ok}_{\vec{t}}((t;\, \vec{t}'), K) &\iff \vec{t} \sqsubseteq \vec{t}' \land \mathit{Means}_{\vec{t}'}(t, K)
\end{align*}

\begin{lemma}[Assignment, Clock for Clock]\label{lem:assign-ok}
If $\mathit{Means}_{\vec{t}}(t', K)$, then
\[
  \mathit{Ok}_{\vec{t}}(\pass{assign}(e, s, t', \vec{t}),\;
    \lambda \sigma.\, \Do (v, \sigma_1) \leftarrow \EV{}{k}\,\sigma\,e;\; \key{set}(\sigma_1, s, v, K)) .
\]
\end{lemma}

\begin{lemma}[Tail, Clock for Clock]\label{lem:tail-ok}
For every Imp expression $e$ and blocks $\vec{t}$,
\[
  \mathit{Ok}_{\vec{t}}(\pass{tail}(e, \vec{t}),\;
    \lambda \sigma.\, \Do (v, \sigma') \leftarrow \EV{}{k}\,\sigma\,e;\; \Ret v) .
\]
\end{lemma}
Calls agree at every clock, by induction on the clock and
Lemma~\ref{lem:tail-ok}.

\begin{theorem}[Correctness of Explicate Control, Interpreter]\label{thm:explicate-clk}
For every Imp program $P$,
\[
  \Equiv(\Interp{\mathrm{Imp}}(P), \Interp{\mathrm{CFG}}(\pass{explicate}(P))) .
\]
\end{theorem}
There is no precondition.

\subsubsection{Select Instructions}

The pass keeps every call, so the interpreters tick in lock step and
return equal results. The X86 frame is the CFG frame, and the
registers are arbitrary. Fix the bodies $\vec{b}$ and a clock $k$ at
which calls agree:
$\CALL{\vec{b}}{k}(f, u, w) = \CALL{\pass{select}(\vec{b})}{k}(f, u, w)$
for all $f$, $u$, and $w$. Write $\key{tail}^k$ and
$\key{jump}^k$ for those of $\vec{b}$, and $\key{run}^k_{\mathrm{X}}$ and
$\key{jump}^k_{\mathrm{X}}$ for those of $\pass{select}(\vec{b})$, and
$\pass{sel}(t) = \mathit{ins}(t);\, \mathit{end}(t)$ as in
Section~\ref{sec:sim-select}. The
instructions of $\pass{sel}(r, d)$ nest around the rest of the block;
they compute the value of $r$ and write it to $d$, with any registers
after. A call reads its operator after its operands, which matters only
when reading it errors.

\begin{lemma}[Simple Expression, Clock for Clock]\label{lem:simple-ok}
If $K(v) = \key{write}((\sigma, \rho'), d, v, \lambda s.\, \key{run}^k_{\mathrm{X}}\,\vec{\beta}\,s\,\vec{\imath}\,j)$
for every $v$ and $\rho'$, then
\[
  \key{simple}_{\CALL{\vec{b}}{k}}(\sigma, r) \bind K =
  \key{run}^k_{\mathrm{X}}\,\vec{\beta}\,(\sigma, \rho)\,(\pass{sel}(r, d); \vec{\imath})\,j .
\]
\end{lemma}

\begin{lemma}[Tail and Jump, Clock for Clock]\label{lem:sel-tail-ok}
For every tail $t$, blocks $\vec{t}$, label $l$, frame $\sigma$, and registers $\rho$,
\begin{align*}
  \key{tail}^k\,\vec{t}\,\sigma\,t &= \key{run}^k_{\mathrm{X}}\,\pass{sel}(\vec{t})\,(\sigma, \rho)\,\mathit{ins}(t)\,\mathit{end}(t) , \\
  \key{jump}^k\,\vec{t}\,\sigma\,l &= \key{jump}^k_{\mathrm{X}}\,\pass{sel}(\vec{t})\,(\sigma, \rho)\,l .
\end{align*}
\end{lemma}
Calls agree at every clock, by induction on the clock and
Lemma~\ref{lem:sel-tail-ok}: the prologue moves $\rdi$ and $\rsi$ to
slots $0$ and $1$, giving the frame $u\, w\, 0^{n-2}$ of CFG.

\begin{theorem}[Correctness of Select Instructions, Interpreter]\label{thm:select-clk}
For every CFG program $P$,
\[
  \Equiv(\Interp{\mathrm{CFG}}(P), \Interp{\mathrm{X86}}(\pass{select}(P))) ,
\]
with equal results for every clock.
\end{theorem}
There is no precondition.

\subsubsection{End to End}

\begin{theorem}[Compiler Correctness, Interpreter]\label{thm:compiler-clk}
For every ISWIM term $M$,
\[
  \FV(M) = \emptyset \implies \Equiv(\Interp{\mathrm{ISWIM}}(M), \Interp{5}(\pass{compile}(M))) ,
\]
by the theorems for each pass, Lemma~\ref{lem:simplify-closed}, and Lemma~\ref{lem:equiv-trans-clk}.
\end{theorem}

\begin{theorem}[Correctness of Compile-Mon, Interpreter]\label{thm:compile-mon-clk}
For every ISWIM term $M$,
\[
  \FV(M) = \emptyset \implies \Equiv(\Interp{\mathrm{ISWIM}}(M), \Interp{5}(\pass{compile\text{-}mon}(M))) ,
\]
by Theorems~\ref{thm:compiler-clk} and~\ref{thm:rco-clk}, Lemma~\ref{lem:compile-scoped},
and Lemma~\ref{lem:equiv-trans-clk}.
\end{theorem}

\begin{theorem}[Correctness of Compile-Imp, Interpreter]\label{thm:compile-imp-clk}
For every ISWIM term $M$,
\[
  \FV(M) = \emptyset \implies \Equiv(\Interp{\mathrm{ISWIM}}(M), \Interp{\mathrm{Imp}}(\pass{compile\text{-}imp}(M))) ,
\]
by Theorems~\ref{thm:compile-mon-clk} and~\ref{thm:lift-clk} and Lemma~\ref{lem:equiv-trans-clk}.
\end{theorem}

\begin{theorem}[Correctness of Compile-CFG, Interpreter]\label{thm:compile-cfg-clk}
For every ISWIM term $M$,
\[
  \FV(M) = \emptyset \implies \Equiv(\Interp{\mathrm{ISWIM}}(M), \Interp{\mathrm{CFG}}(\pass{compile\text{-}cfg}(M))) ,
\]
by Theorems~\ref{thm:compile-imp-clk} and~\ref{thm:explicate-clk} and Lemma~\ref{lem:equiv-trans-clk}.
\end{theorem}

\begin{theorem}[Correctness of Compile-X86, Interpreter]\label{thm:compile-x86-clk}
For every ISWIM term $M$,
\[
  \FV(M) = \emptyset \implies \Equiv(\Interp{\mathrm{ISWIM}}(M), \Interp{\mathrm{X86}}(\pass{compile\text{-}x86}(M))) ,
\]
by Theorems~\ref{thm:compile-cfg-clk} and~\ref{thm:select-clk} and Lemma~\ref{lem:equiv-trans-clk}.
\end{theorem}

%===============================================================================
\section{Reduction Semantics + Simulation}\label{sec:u-red}

\subsection{Reduction Semantics}\label{sec:red}

The languages with binders (ISWIM, Clos--GlobalFun) get a reduction semantics
by substitution: a state is a term, and a step reduces a redex in an
evaluation context, substituting values for bound variables. A variable
is never looked up: one that is reached is unbound, and the term is stuck.
Imp, CFG, and X86 have no binders, so they keep their machines
(Section~\ref{sec:machines}), whose proofs of explicate control and select
instructions are shared with the CEK machine + simulation development.

\subsubsection{Reductions as machines}

A reduction says of a term whether it is a value, is stuck, or reduces.
\begin{align*}
  \mathit{Red}(T) &::= \key{val} \mid \error \mid \key{to}\; M \qquad (M \in T) \\
  \mathit{toNext}(o, \key{val}) &= \stp o \qquad
  \mathit{toNext}(o, \error) = \stp \error \qquad
  \mathit{toNext}(o, \key{to}\; M) = \nxt M \\
  \mathcal{R}^{\ell}(M) &= (\mathit{Term}^{\ell},\; N \mapsto \mathit{toNext}(\obs(N), \mathit{red}^{\ell}(N)),\; M)
\end{align*}
so termination, divergence, and $\Equiv$ are those of Section~\ref{sec:sims}.
On paper we give $\mathit{red}^{\ell}$ with evaluation contexts $E$: $M = E[r]$
for a unique context and redex or value $r$, and
\begin{align*}
  \mathit{red}(V) &= \key{val} \\
  \mathit{red}(E[r]) &= \key{to}\; E[r'] \quad \text{if } r \to r' \\
  \mathit{red}(E[r]) &= \error \quad \text{otherwise (a variable, or a redex of the wrong shape)}
\end{align*}
The Agda writes $\mathit{red}$ as one function with congruence rules:
$\key{step\text{-}in}(f, \mathit{red}(X), k)$ reduces inside the context $f$ if
the subterm $X$ reduces, is stuck if $X$ is, and continues with $k$ if $X$ is
a value. This needs no decomposition into a context and a redex, no plugging,
and no unique-decomposition lemma, and is deterministic and total by
construction, as the machines are; a run of a subterm lifts through a
context by one generic lemma (frames, below).
Values are observed as in Section~\ref{sec:obs}: $\obs(\lam{x} N) = \key{fun}$,
and so on.

\subsubsection{ISWIM}

\[
\begin{array}{rcl}
  V, W &::=& \lam{x} N \mid c \mid \tuple{V_1, \ldots, V_n} \mid \inl V \mid \inr V \\
  E &::=& \square \mid E\;M \mid (\lam{x} N)\;E \mid \tuple{V_1, \ldots, V_{i-1}, E, M_{i+1}, \ldots, M_n}
    \mid \prj{i} E \mid \inl E \mid \inr E \\
  &\mid& \Case{E}{x}{M}{y}{N} \mid \prim{p}{E}{M} \mid \prim{p}{V}{E}
\end{array}
\]
\[
\begin{array}{@{}r@{\;\to\;}l@{\quad}l@{}}
  (\lam{x} N)\;V & N[x := V] \\
  \prj{i} \tuple{V_1, \ldots, V_m} & V_i & m = n,\; \prj{i} \text{ of arity } n \\
  \Case{\inl V}{x}{M}{y}{N} & M[x := V] \\
  \Case{\inr V}{x}{M}{y}{N} & N[y := V] \\
  \prim{p}{n}{m} & \res{\delta(p, n, m)}
\end{array}
\]
The context $(\lam{x} N)\;E$ checks that the operator is a function before
the argument is evaluated, as the ISWIM machine does: $V\;M$ with $V$ not a
function is stuck.

\subsubsection{Clos and EnvClos}

A closure is a value once its free variables are values.
\[
\begin{array}{rcl}
  V &::=& \clos{V_1\ldots V_n}{x_1\ldots x_n\, y} N \mid c \mid \tuple{\bar{V}} \mid \inl V \mid \inr V \\
  E &::=& \clos{V_1 \ldots V_{i-1}\, E\, M_{i+1} \ldots M_n}{x_1\ldots x_n\, y} N \mid E\;M \mid V_\lambda\;E \mid \cdots
\end{array}
\]
where $V_\lambda$ is a closure value, and in Clos
\[
  (\clos{V_1\ldots V_n}{x_1\ldots x_n\, y} N)\;W \to N[x_1 := V_1, \ldots, x_n := V_n, y := W]
\]
and in EnvClos, with the tuple of the free variables for $e$,
\[
  (\clos{V_1\ldots V_n}{e\, y} N)\;W \to N[e := \tuple{V_1, \ldots, V_n}, y := W] .
\]
The other contexts and redexes are those of ISWIM.

\subsubsection{FlatClos and GlobalFun}

The operator of an application is checked to be code only at the call,
after its three operands.
\[
\begin{array}{rcl}
  V &::=& \fun{e\,y.\,N} \mid \pair{V}{W} \mid c \mid \tuple{\bar{V}} \mid \inl V \mid \inr V \\
  E &::=& E(M, N) \mid V(E, N) \mid V(W, E) \mid \pair{E}{N} \mid \pair{V}{E} \mid \FST E \mid \SND E
    \mid \Let x = E \In N \mid \cdots
\end{array}
\]
\[
\begin{array}{@{}r@{\;\to\;}l@{}}
  (\fun{e\,y.\,N})(U, W) & N[e := U, y := W] \\
  \FST \pair{V}{W} & V \qquad \SND \pair{V}{W} \to W \\
  \Let x = V \In N & N[x := V]
\end{array}
\]
In GlobalFun, $\funref{k}$ is a value, and a program $(d_0, \ldots, d_{m-1};\; M)$
reduces $M$ with the definitions:
$\funref{k}(U, W) \to N[e := U, y := W]$ when $k < m$ and $d_k = \key{fun}\; e\,y.\,N$.

%-------------------------------------------------------------------------------
\subsection{Simulations Between Reductions}\label{sec:red-sims}

The forward simulations of Section~\ref{sec:sims} answer a source that stops
by a target that has stopped. Here the target may first take administrative
steps (an extra projection, a let, the evaluation of an atom), so a stop
is answered by a run to a final state:
\[
  \mathit{Answers}_i(\stp o, s_2) \iff \exists t.\; s_2 \lto_{L_2}^* t \land \Final_{L_2}(t, o) .
\]
Theorem~\ref{thm:equivalent} still holds, by determinism.

Between two reductions, a relation $\approx$ on terms gives a simulation when
every source term's reduction is answered by every related target term:
\begin{align*}
  \mathit{Ans}(M, \key{val}, M') &\iff \exists V'.\; M' \lto^* V' \land V' \text{ is a value} \land M \approx V' \\
  \mathit{Ans}(M, \error, M') &\iff \exists W'.\; M' \lto^* W' \land \mathit{red}(W') = \error \\
  \mathit{Ans}(M, \key{to}\; M_1, M') &\iff \exists M_1'.\; M' \lto^+ M_1' \land M_1 \approx M_1'
\end{align*}

\begin{lemma}[Congruence]\label{lem:ans-in}
If $g$ is a frame of the target ($\mathit{red}(g(X)) = \key{step\text{-}in}(g, \mathit{red}(X), \kappa(X))$),
$Y \approx Y' \implies f(Y) \approx g(Y')$, $\mathit{Ans}(X, r, X')$, and
$\mathit{Ans}(f(X), k, g(V'))$ for every value $V'$ with $X \approx V'$, then
$\mathit{Ans}(f(X), \key{step\text{-}in}(f, r, k), g(X'))$.
\end{lemma}

\begin{theorem}[Reduction Simulation Gives Equivalence]\label{thm:red-equivalent}
If $M_0 \approx M_0'$, $\mathit{Ans}(M, \mathit{red}^1(M), M')$ whenever $M \approx M'$, and
$\obs(V) = \obs(V')$ whenever $V \approx V'$ are values, then
$\Equiv(\mathcal{R}^1(M_0), \mathcal{R}^2(M_0'))$.
\end{theorem}

\begin{lemma}[Closure Under Substitution]\label{lem:red-subst}
For each relation $\approx$ below that is closed under substitution, if
$M \approx M'$ and $\sigma(x) \approx \sigma'(x)$ for every $x$, then
$M[\sigma] \approx M'[\sigma']$. It is proved once for renamings, then for
every substitution, from one proof for any relation between substitutions
that is preserved under binders.
\end{lemma}

%-------------------------------------------------------------------------------
\subsection{Relations}\label{sec:red-relations}

Each relation is a congruence: $x \approx x$, $c \approx c$, and each
construct is related to the same construct with related subterms; we show
only the other clauses. Where a pass compiles code that no substitution
enters (the body of a Clos--FlatClos closure, a global definition), the
relation on that code is the compilation function, and applying related
closures substitutes related values into it, which a fundamental lemma
relates.

\subsubsection{Simplify}

The target, simplified, takes fewer steps than the source, so the simulation
goes from the target to the source: each target step is answered by one or
more source steps, and $\Equiv$ is symmetric. Simplify reduces under binders,
so the source catches up by steps $\to$ on open terms, in which a variable
stands for a value (in a closed term they are steps of $\mathit{red}$). Writing
the source on the left, besides the congruence,
\begin{align*}
  M \approx M' &\iff M \to M_1 \land M_1 \approx M' \\
  (\lam{x} N)\;W \approx (\lam{x} N')\;V' &\iff W \text{ is a value} \land W \approx V' \land N[x := W] \approx N'[x := V'] ,
\end{align*}
the second for a value that simplify kept bound and known, whose bodies are
related only after $\beta$. The relation is closed under substitution by values.
\begin{lemma}[Fundamental Lemma, Simplify]
If $\sigma$ and $\sigma'$ map variables to values, and for every $x$ with
$E(x) = (T, K)$, $\sigma(x) \approx T[\sigma']$ and, if $K \neq \bot$,
$\sigma(x) \approx K[\sigma']$ for a known value $K$, and
$\sigma(x) \to^* (\lam{y} N)[\sigma_f]$ with $\sigma_f$, $E_f$ satisfying the same
for a known function $K = \langle \lam{y} N, E_f \rangle$, then
$M[\sigma] \approx \pass{simp}_E(M)[\sigma']$, and $\mathcal{K}_E(M)$ is true of
$M[\sigma]$ in the same way.
\end{lemma}
A known value is reached: if $M \approx V'$ for a value $V'$, then $M \to^* W$
for a value $W \approx V'$.
\begin{theorem}[Correctness of Simplify, Reduction]\label{thm:simplify-red}
For every ISWIM term $M$,
\[
  \FV(M) = \emptyset \implies
  \Equiv(\mathcal{R}^{\mathrm{ISWIM}}(M), \mathcal{R}^{\mathrm{ISWIM}}(\pass{simplify}(M))) .
\]
\end{theorem}

\subsubsection{CC: Enclose}

The bodies of ISWIM functions are substituted into, so a function is related
to the closure of the compiled body, through a substitution $\sigma$ of the
body's free variables by terms related to what the closure captures:
\[
  \lam{y} N[\sigma] \approx \clos{F_1\ldots F_n}{x_1\ldots x_n\, y} \pass{enclose}(N)
  \iff \forall i.\; x_i \in \FV(N) \implies \sigma(x_i) \approx F_i ,
\]
where each $F_i$ is a variable or a value.

\begin{lemma}[Fundamental Lemma, Enclose]
If $\sigma(x) \approx \sigma'(x)$ for every $x \in \FV(M)$, and each
$\sigma'(x)$ is a variable or a value, then $M[\sigma] \approx \pass{enclose}(M)[\sigma']$.
\end{lemma}

The target is closed, so a closure that is reduced has captured values,
and is a value as the ISWIM function is.

\begin{theorem}[Correctness of Enclose, Reduction]\label{thm:enclose-red}
For every ISWIM term $M$,
\[
  \FV(M) = \emptyset \implies
  \Equiv(\mathcal{R}^{\mathrm{ISWIM}}(M), \mathcal{R}^{2}(\pass{enclose}(M))) .
\]
\end{theorem}

\subsubsection{CC: Optimize}

A closure is related to its optimization when the code is the compiled code
and the free variables are related as $\pass{prune}$ makes them: the ones the
code uses are kept and related, the others may be dropped if they are
variables or values ($\mathit{Pr}$). In a closed program a dropped free
variable that is reduced has become a value, so the two reductions run in
lock step.
\begin{lemma}[Fundamental Lemma, Optimize]
If $\sigma(x) \approx \sigma'(r(x))$ for every $x \in \FV(M)$, and each
$\sigma(x)$ is a variable or a value, then $M[\sigma] \approx \pass{optimize}_r(M)[\sigma']$.
\end{lemma}

\begin{theorem}[Correctness of Optimize, Reduction]\label{thm:optimize-red}
For every Clos term $M$,
\[
  \FV(M) = \emptyset \implies \Equiv(\mathcal{R}^{2}(M), \mathcal{R}^{2}(\pass{optimize}(M))) .
\]
\end{theorem}

\subsubsection{CC: Concretize}

Applying a EnvClos closure substitutes the tuple of its free variables, so where
the source has a value, the target has a projection from that tuple, which
takes a step:
\[
  V \approx \prj{i} \tuple{V_1', \ldots, V_n'} \iff V \approx V_i' \land V_1', \ldots, V_n' \text{ are values} ,
\]
and closures are related by their free variables and by their code, the
target's being $\pass{concretize}$ of the source's.
\begin{lemma}[Fundamental Lemma, Concretize]
If $\sigma(x) \approx r_\varphi(x)[\sigma']$ for every $x \in \FV(M)$, then
$M[\sigma] \approx \pass{concretize}_\varphi(M)[\sigma']$, where $r_\varphi(x)$
is the term of $\varphi(x)$ (a variable or a projection).
\end{lemma}

\begin{theorem}[Correctness of Concretize, Reduction]\label{thm:concretize-red}
For every Clos term $M$,
\[
  \FV(M) = \emptyset \implies \Equiv(\mathcal{R}^{2}(M), \mathcal{R}^{3}(\pass{concretize}_{\mathit{id}}(M))) .
\]
\end{theorem}
The precondition is new: a variable of a closure's code beyond its
parameters becomes, in the two languages, two different variables, and
substitution moves both differently; the machines look both up in the
closure's environment, where both fail.

\subsubsection{CC: Delay}

\begin{align*}
  \clos{\bar{F}}{e\, y} N \approx \pair{\fun{e\,y.\,N'}}{\tuple{\bar{F}'}} &\iff N \approx N' \land \bar{F} \approx \bar{F}' \\
  L\;M \approx \Let f = L' \In (\FST f)(\SND f, M') &\iff L \approx L' \land M \approx M' \\
  (\clos{\bar{V}}{e\, y} N)\;M \approx \fun{e\,y.\,N'}(\tuple{\bar{V}'}, M') &\iff
    N \approx N' \land \bar{V} \approx \bar{V}' \land M \approx M'
\end{align*}
The last clause relates the source, which waits for the argument, to the
target after its let, $\FST$, and $\SND$. Code is related by the relation
itself, so $M \approx \pass{delay}(M)$ for every $M$, and $\approx$ is closed
under substitution.
\begin{theorem}[Correctness of Delay, Reduction]\label{thm:delay-red}
For every EnvClos term $M$,
\[
  \Equiv(\mathcal{R}^{3}(M), \mathcal{R}^{4}(\pass{delay}(M))) .
\]
\end{theorem}

\subsubsection{CC: Globalize}

For the table $D$ of the whole program,
\[
  \fun{e\,y.\,N} \approx_D \funref{k} \iff d_k = \key{fun}\; e\,y.\,N' \text{ in } D \land N \approx_D N' ,
\]
and a term is related to its globalization when the table it makes grows
into $D$, as in the other proofs.
\begin{theorem}[Correctness of Globalize, Reduction]\label{thm:globalize-red}
For every FlatClos term $M$,
\[
  \Equiv(\mathcal{R}^{4}(M), \mathcal{R}^{5}(\pass{globalize}(M))) .
\]
\end{theorem}

\subsubsection{Remove Complex Operands}

The values of GlobalFun contain no code, so source and target compute the same
values, and the relation is not closed under substitution but built from the
code under a substitution $\sigma$ of values. A complex operand $M$ in the
context $h$ of its construct is related to its let in the target:
\[
  h[M] \approx \Let x = M' \In B \iff M \approx M' \land \forall V.\; h[V] \approx B[x := V] ,
\]
for every value $V$ (the same on both sides, so this is not a logical
relation). Values are related to themselves, and so are the constructs whose
operands are values, $\key{case}$ and $\key{let}$ with the code of their
branches and body compiled under $\sigma$.
\begin{lemma}[Fundamental Lemma, RCO]
If $\sigma$ maps the variables of $M$ to values, then $M[\sigma] \approx \pass{rco}(M)[\sigma]$.
\end{lemma}
\begin{theorem}[Correctness of RCO, Reduction]\label{thm:rco-red}
For every GlobalFun program $P$,
\[
  \mathit{Scoped}_5(P) \implies \Equiv(\mathcal{R}^{5}(P), \mathcal{R}^{5}(\pass{rco}(P))) .
\]
\end{theorem}

\subsubsection{Lift Locals}

The source substitutes values for the variables; the target keeps them in its
frame. A source term $C[M[\sigma]]$, for code $M$ of the program in an
evaluation context $C$ of lets, is related to the Imp state that runs
$\pass{lift}(M)$ in a frame holding the values of $\sigma$, with a continuation
that does what $C$ does (assigning the let's slot and running its body,
and returning from calls), and a value term to the machine value with the
same data.
\begin{theorem}[Correctness of Lift Locals, Reduction]\label{thm:lift-red}
For every monadic GlobalFun program $P$,
\[
  \Equiv(\mathcal{R}^{5}(P), \Mach{\mathrm{Imp}}(\pass{lift}(P))) .
\]
\end{theorem}

\subsubsection{Explicate Control and Select Instructions}

Theorems~\ref{thm:explicate-sim} and~\ref{thm:select-sim}, between the
machines of Imp, CFG, and X86.

\subsubsection{End to End}

\begin{theorem}[Compiler Correctness, Reduction]\label{thm:compiler-red}
For every ISWIM term $M$, if $\FV(M) = \emptyset$, then
\begin{align*}
  &\Equiv(\mathcal{R}^{\mathrm{ISWIM}}(M), \mathcal{R}^{5}(\pass{compile}(M))) \\
  &\Equiv(\mathcal{R}^{\mathrm{ISWIM}}(M), \mathcal{R}^{5}(\pass{compile\text{-}mon}(M))) \\
  &\Equiv(\mathcal{R}^{\mathrm{ISWIM}}(M), \Mach{\mathrm{Imp}}(\pass{compile\text{-}imp}(M))) \\
  &\Equiv(\mathcal{R}^{\mathrm{ISWIM}}(M), \Mach{\mathrm{CFG}}(\pass{compile\text{-}cfg}(M))) \\
  &\Equiv(\mathcal{R}^{\mathrm{ISWIM}}(M), \Mach{\mathrm{X86}}(\pass{compile\text{-}x86}(M))) ,
\end{align*}
by the theorems for each pass, Lemmas~\ref{lem:simplify-closed} and~\ref{lem:compile-scoped}, and
Lemma~\ref{lem:equiv-trans}.
\end{theorem}

%===============================================================================
\section{The Graph Model with Errors}\label{sec:graph-errors}

In the graph model, a program that stops with an error and one that
diverges both denote $\emptyset$. The graph model with errors tells them
apart: a trapped error denotes a set that holds the error value $\omega$
(the value that $\rho_0$ already gives to a variable not in scope, and
which is consistent only with itself), and divergence still denotes
$\emptyset$. It is defined for every language of the compiler
(\texttt{agda/Semantics/GraphErr/}); its operators are the graph model's
(Section~\ref{sec:den-ops}), under an error monad on sets of values
(\texttt{agda/Framework/GraphModel/ErrorMonad.agda}). Each semantics errs
exactly where its machine does (Section~\ref{sec:machines}).

\subsection{Operators}

A set $D$ is a computation: it errs when $\omega \in D$, its values are
$\lfloor D \rfloor = \{ d \in D \mid d \neq \omega \}$, and it diverges
when it has neither. The bind runs $F$ on the values of $D$:
\begin{align*}
  \mathit{val}(D) &\iff \lfloor D \rfloor \NE \\
  D \bind F &= \{ \omega \mid \omega \in D \} \cup \{ w \mid \mathit{val}(D),\; w \in F(\lfloor D \rfloor) \} \\
  \key{checked}_k(X, A) &= \{ \omega \mid \exists d \in X.\; \key{kind}(d) \neq k \} \cup A
\end{align*}
and $(D_1, \ldots, D_n) \bind K$ binds $n$ operands, left to right. So an
operand that diverges ($\emptyset$) before one that errs gives no error,
and $F$ never sees $\omega$. $\key{kind}(d)$ is the function ($\nu$,
$V \mapsto w$), the pair, the injection, the tuple of $n$ components or
the base type $B$ of a value; $\key{checked}_k$ adds the error of a value
of another kind. Each operator is the graph model's operator (its
\emph{kernel}) under binds of its operands, in the machines' order:
\begin{align*}
  D \app^{\omega} E &= D \bind \lambda F.\, \key{checked}_{\mathit{fun}}(F,\; E \bind \lambda X.\, F \app X) \\
  \Lambda^{\omega}(F) &= \{ e \in \Lambda(F) \mid e = V \mapsto w \Rightarrow \omega \notin V \} \\
  \mathcal{T}^{\omega}(D_1,\ldots,D_n) &= (D_1, \ldots, D_n) \bind \mathcal{T} \\
  \key{proj}^{\omega}_i(D) &= D \bind \lambda X.\, \key{checked}_{\mathit{tup}_n}(X, \key{proj}_i(X)) \qquad
  \key{inl}^{\omega}(D) = D \bind \key{inl} \\
  \mathcal{C}^{\omega}(D, F, G) &= D \bind \lambda X.\, \key{checked}_{\mathit{inj}}(X,\; \mathcal{C}(X^{\checkmark}, F, G)) \\
  \key{prim}^{\omega}_p(D, E) &= D \bind \lambda X.\, E \bind \lambda Y.\,
    \key{checked}_{\mathit{int}}(X, \key{checked}_{\mathit{int}}(Y, \key{prim}_p(X, Y))) \\
  \key{pair}^{\omega}(D, E) &= (D, E) \bind \key{pair} \qquad
  \key{car}^{\omega}(D) = D \bind \lambda X.\, \key{checked}_{\mathit{pair}}(X, \key{car}(X)) \\
  \key{call}^{\omega}(L, M, N) &= (L, M, N) \bind \lambda F\, E\, X.\, \key{checked}_{\mathit{fun}}(F,\; (F \app E) \app X)
\end{align*}
where $X^{\checkmark}$ keeps the injections of $X$ whose payload is not
$\omega$; $\key{inr}^{\omega}$ and $\key{cdr}^{\omega}$ likewise.
Application checks its operator before it binds its argument, as the
ISWIM, Clos and EnvClos machines do: applying a non-function is an error
even when the argument diverges. $\key{call}^{\omega}$ is the call of
FlatClos and the later languages: code, environment and argument are bound
left to right, and only then is the code checked (the late check of
their machines). The binders $\Lambda^{\omega}$ and $\mathcal{C}^{\omega}$
bind no $\omega$; operands never are $\omega$ (the bind sees to that),
but without the restriction $\lam{x}\lam{y} x$ would have the entry
$[\omega] \mapsto \nu$ in ISWIM, where its closure in Clos has
$[\omega] \mapsto \omega$. For the global definitions and the
first-order languages there are also
\begin{align*}
  \Lambda\Lambda^{\omega}(F) &= \{ e \in \Lambda\Lambda(F) \mid \text{no argument of } e \text{ holds } \omega \} \qquad
  \key{errcode} = \Lambda\Lambda^{\omega}(\lambda X\, Y.\, \{\omega\}) \\
  \key{strict}^{\omega}(D, E) &= D \bind \lambda \_.\, E \qquad
  \key{payload}^{\omega}(D) = D \bind \lambda X.\, \key{checked}_{\mathit{inj}}(X, \key{payload}(X^{\checkmark})) \\
  \key{isinl}^{\omega}(D) &= D \bind \lambda X.\, \key{checked}_{\mathit{inj}}(X, \key{isinl}(X))
\end{align*}
$\key{errcode}$ is the meaning of a reference to a definition that is
not there: code whose call errs, as the machines err at the call.
$\key{strict}^{\omega}(D, E)$ is a write of $D$ to a register or the flag
followed by $E$. The monad laws, and the monotonicity, consistency and
continuity of the bind and the check, are proved once; each operator's
are derived from its kernel's
(\texttt{agda/proof/Framework/GraphModel/ErrorMonad.agda}).

\subsection{Semantics}

$\SEM{-}^{\mathrm{ISWIM},\omega}$ and $\SEM{-}^{2,\omega}$ are the
clauses of the graph model's semantics of ISWIM and Clos with each
operator replaced by its error-propagating counterpart:
\begin{align*}
  \SEM{x}^{\ell,\omega}\rho &= \rho(x) \qquad
  \SEM{c}^{\ell,\omega}\rho = \mathcal{B}(c) \\
  \SEM{\lam{x} N}^{\mathrm{ISWIM},\omega}\rho &= \Lambda^{\omega}(\lambda X.\, \SEM{N}^{\mathrm{ISWIM},\omega}\ext{\rho}{x}{X}) \\
  \SEM{L\;M}^{\ell,\omega}\rho &= \SEM{L}^{\ell,\omega}\rho \app^{\omega} \SEM{M}^{\ell,\omega}\rho \\
  \SEM{\Let x = M \In N}^{\ell,\omega}\rho &= \Lambda^{\omega}(\lambda X.\, \SEM{N}^{\ell,\omega}\ext{\rho}{x}{X}) \app^{\omega} \SEM{M}^{\ell,\omega}\rho \\
  \SEM{\clos{z_1\ldots z_n}{x_1\ldots x_n\, y} N}^{2,\omega}\rho &=
    (\rho(z_1), \ldots, \rho(z_n)) \bind \lambda X_1 \ldots X_n.\,
       \Lambda^{\omega}(\lambda Y.\, \SEM{N}^{2,\omega}\rho_0(x_1{:=}X_1,\ldots,x_n{:=}X_n, y{:=}Y))
\end{align*}
and tuples, projections, injections, case and primitives likewise with
$\mathcal{T}^{\omega}$, $\key{proj}^{\omega}_i$, $\key{inl}^{\omega}$,
$\key{inr}^{\omega}$, $\mathcal{C}^{\omega}$ and $\key{prim}^{\omega}_p$.
For EnvClos and FlatClos,
\begin{align*}
  \SEM{\clos{r_1\ldots r_n}{e\, y} N}^{3,\omega}\rho &=
    \Lambda\Lambda^{\omega}(\lambda X\, Y.\, \SEM{N}^{3,\omega}\rho_0(e{:=}X, y{:=}Y))
    \app^{\omega} \mathcal{T}^{\omega}(\SEM{r_1}^{3,\omega}\rho, \ldots, \SEM{r_n}^{3,\omega}\rho) \\
  \SEM{\fun{e\,y.\,N}}^{4,\omega}\rho &= \Lambda\Lambda^{\omega}(\lambda X\, Y.\, \SEM{N}^{4,\omega}\rho_0(e{:=}X, y{:=}Y)) \\
  \SEM{L(M, N)}^{4,\omega}\rho &= \key{call}^{\omega}(\SEM{L}^{4,\omega}\rho, \SEM{M}^{4,\omega}\rho, \SEM{N}^{4,\omega}\rho) \\
  \SEM{\pair{M}{N}}^{4,\omega}\rho &= \key{pair}^{\omega}(\SEM{M}^{4,\omega}\rho, \SEM{N}^{4,\omega}\rho) \qquad
  \SEM{\FST M}^{4,\omega}\rho = \key{car}^{\omega}(\SEM{M}^{4,\omega}\rho)
\end{align*}
and $\key{snd}$ by $\key{cdr}^{\omega}$; EnvClos's application is
$\app^{\omega}$, with the early check. All of these semantics are
monotone and consistent, and all but Clos's (whose continuity no proof
needs) continuous.

\paragraph{GlobalFun.} As FlatClos, with the global definitions of the graph
model's GlobalFun, except that a
reference to a definition that is not there is $\key{errcode}$, not
$\emptyset$:
\begin{align*}
  \SEM{\funref{k}}^{5,\omega}_T\rho &= T(k) \qquad
  \key{code}^{\omega}_T(\key{fun}\; e\,y.\,N) = \Lambda\Lambda^{\omega}(\lambda X\, Y.\, \SEM{N}^{5,\omega}_T\rho_0(e{:=}X, y{:=}Y)) \\
  \key{table}^{\omega}(d_0,\ldots,d_{m-1})(k) &=
    \begin{cases}
      \key{code}^{\omega}_{\key{table}^{\omega}(d_0,\ldots,d_{k-1})}(d_k) & k < m \\
      \key{errcode} & \text{otherwise}
    \end{cases}
\end{align*}

\paragraph{Imp and CFG.} The atoms and simple expressions are those of
the graph model with the
error-propagating operators ($\mathcal{S}_\sigma(a_1(a_2, a_3)) =
\key{call}^{\omega}(\mathcal{A}_\sigma(a_1), \mathcal{A}_\sigma(a_2),
\mathcal{A}_\sigma(a_3))$), and the table gives $\key{errcode}$ for a
missing definition. Reading a slot beyond the frame gives $\{\omega\}$,
as before. Assigning one errs once the value assigned is computed:
\[
  \mathit{checked}^{\omega}_n(s, D, E) =
    \begin{cases} E & s < n \\ D \bind \lambda \_.\, \{ \omega \} & \text{otherwise.} \end{cases}
\]
A body runs in a fresh frame of $\max(2, n)$ slots when called (and of
$n$ for main), as in the machines. Writing $\SEM{e}\sigma\,\kappa$ for
Imp's continuation-passing semantics and $\SEM{t}\sigma$ for CFG's:
\begin{align*}
  \SEM{\Casei{a}{s}{e_1}{e_2}}\sigma\,\kappa &= \mathit{checked}^{\omega}_n(s,
    \key{payload}^{\omega}(\mathcal{A}_\sigma(a)),\;
    \mathcal{C}^{\omega}(\mathcal{A}_\sigma(a),\; \lambda X.\, \SEM{e_1}\ext{\sigma}{s}{X}\,\kappa,\;
      \lambda X.\, \SEM{e_2}\ext{\sigma}{s}{X}\,\kappa)) \\
  \SEM{\assign{s}{e_1}{e_2}}\sigma\,\kappa &= \SEM{e_1}\sigma\,(\lambda D\, \sigma'.\,
    \mathit{checked}^{\omega}_n(s, D, \Lambda^{\omega}(\lambda X.\, \SEM{e_2}\ext{\sigma'}{s}{X}\,\kappa) \app^{\omega} D)) \\
  \SEM{\assign{s}{r}{t}}\sigma &= \mathit{checked}^{\omega}_n(s, \mathcal{S}_\sigma(r),
    \Lambda^{\omega}(\lambda X.\, \SEM{t}\ext{\sigma}{s}{X}) \app^{\omega} \mathcal{S}_\sigma(r)) \\
  \SEM{\key{case}\;a\;\key{of}_{s}\;\key{goto}\;l_1 \mid \key{goto}\;l_2}\sigma &=
    \mathit{checked}^{\omega}_n(s, \key{payload}^{\omega}(\mathcal{A}_\sigma(a)),\; \mathcal{C}^{\omega}(\mathcal{A}_\sigma(a),\;
      \lambda X.\, B(l_1)(\ext{\sigma}{s}{X}),\; \lambda X.\, B(l_2)(\ext{\sigma}{s}{X})))
\end{align*}
and a jump to a block that is not there is $\{\omega\}$. A case checks
the tag first and then writes the slot, so a slot beyond the frame errs
once the payload is computed.

\paragraph{X86.} As in the graph model, with these differences. An unset
register and the absent flag hold $\{\omega\}$, so that reading them
errs: $\rdi$ and $\rsi$ in main, $\rax$ at the start of a definition,
$\rdi$, $\rsi$ and the flag after a call. A write to a slot is
$\mathit{checked}^{\omega}$ and $\Lambda^{\omega}$, as in CFG; a write to a
register or the flag is $\key{strict}^{\omega}$:
\begin{align*}
  \key{put}_{\%g}(D, \kappa)(\sigma, \rho, \varphi) &= \key{strict}^{\omega}(D,\; \kappa(\sigma, \ext{\rho}{\%g}{D}, \varphi)) \qquad
  \key{flag}(F, \kappa)(\sigma, \rho) = \key{strict}^{\omega}(F,\; \kappa(\sigma, \rho, F)) \\
  \SEM{\key{callq}\; o}\kappa\,\sigma\,\rho\,\varphi &= \key{strict}^{\omega}(R,\; \kappa(\sigma, \key{called}(R), \{\omega\}))
    \quad\text{where } R = \key{call}^{\omega}(\mathcal{O}(o), \rho(\rdi), \rho(\rsi)) \\
  \key{field}_{\key{payload}} &= \key{payload}^{\omega} \qquad \text{the flag of } \key{testinl} = \key{isinl}^{\omega}
\end{align*}
with $\key{called}(R)(\%g) = \{\omega\}$ and $\key{args}(X, Y)(\rax) =
\{\omega\}$; $\key{alloc}$ of the wrong number of fields is $\{\omega\}$;
$\key{cmpq}$ is $\key{prim}^{\omega}_{\key{eq}}$; $\key{je}$ is
$\mathcal{C}^{\omega}$ on the flag; a missing block is $\{\omega\}$; and
main runs with the flag $\{\omega\}$.

\subsection{Correctness of the Passes}

\begin{theorem}[Each Pass, with Errors]
For every closed ISWIM term $M$ (and, for optimize and rco, every
well-scoped program, which the compiler produces),
\begin{align*}
  \SEM{M}^{\mathrm{ISWIM},\omega}\rho_0 &\simeq \SEM{\pass{simplify}(M)}^{\mathrm{ISWIM},\omega}\rho_0
  \qquad
  \SEM{M}^{\mathrm{ISWIM},\omega}\rho_0 \simeq \SEM{\pass{enclose}(M)}^{2,\omega}\rho_0 \\
  \SEM{M}^{2,\omega}\rho_0 &\simeq \SEM{\pass{optimize}(M)}^{2,\omega}\rho_0 \qquad
  \SEM{M}^{2,\omega}\rho_0 \simeq \SEM{\pass{concretize}(M)}^{3,\omega}\rho_0 \\
  \SEM{M}^{4,\omega}\rho_0 &\simeq \SEM{\pass{globalize}(M)}^{5,\omega} \qquad
  \SEM{\pass{rco}(P)}^{5,\omega} \simeq \SEM{P}^{5,\omega} \qquad
  \SEM{\pass{lift}(P)}^{\mathrm{ImpK},\omega} \simeq \SEM{P}^{5,\omega} \\
  \SEM{\pass{explicate}(P)}^{\mathrm{CFG},\omega} &\simeq \SEM{P}^{\mathrm{ImpK},\omega} \qquad
  \SEM{\pass{select}(P)}^{\mathrm{X86},\omega} \simeq \SEM{P}^{\mathrm{CFG},\omega}
\end{align*}
(lift locals for a program in monadic normal form). For every EnvClos
term $M$, writing $D = \SEM{M}^{3,\omega}\rho_0$ and
$D' = \SEM{\pass{delay}(M)}^{4,\omega}\rho_0$,
\[
  \omega \in D \iff \omega \in D', \qquad
  c \in D \iff c \in D', \qquad
  \mathit{val}(D) \iff \mathit{val}(D') .
\]
\end{theorem}
The proofs are those of the graph model with environments that hold
values and no $\omega$, and with the monad's laws for the errors: an
operand that is a set of values binds as the left unit, a let is the bind
of its body, binds associate, and a case and a check commute with a bind. Simplify duplicates and drops only values, and a
variable stands for a value, so the theorem is for closed terms (for an
open one, $\Let y = x \In 0$ becomes $0$, though $x$ errs in $\rho_0$).
A Clos closure is an error when one of its free variables is, so
enclose needs the variables a closure captures to have values and no
errors, and optimize, which drops the free variables a closure does not
use, needs its input well scoped. Rco binds every operand by a
$\key{let}$, so it needs that the variables of a well-scoped program
hold values. Delay changes the representation of closures, so, as in the
graph model, it is proved in two directions, with every target error
traced to a source error. Select needs no precondition: both CFG and
X86 run a called body in a frame of $\max(2, n)$ slots.

\begin{theorem}[The Compiler, with Errors]
For every closed ISWIM term $M$, writing $D = \SEM{M}^{\mathrm{ISWIM},\omega}\rho_0$
and $D'$ for the meaning of $\pass{compile}(M)$ (GlobalFun),
$\pass{compile\text{-}mon}(M)$, $\pass{compile\text{-}imp}(M)$,
$\pass{compile\text{-}cfg}(M)$ or $\pass{compile\text{-}x86}(M)$,
\[
  \omega \in D \iff \omega \in D', \qquad
  c \in D \iff c \in D', \qquad
  \mathit{val}(D) \iff \mathit{val}(D') .
\]
\end{theorem}

\subsection{Adequacy}

\begin{theorem}[Adequacy, with Errors]
For every closed ISWIM term $M$, writing $D$ for
$\SEM{M}^{\mathrm{ISWIM},\omega}\rho_0$ and $L$ for $\Mach{\mathrm{ISWIM}}(M)$
(or $D$ for $\SEM{\pass{compile}(M)}^{5,\omega}$ and $L$ for
$\Mach{\mathrm{GlobalFun}}(\pass{compile}(M))$, or $D$ for
$\SEM{\pass{compile\text{-}x86}(M)}^{\mathrm{X86},\omega}$ and $L$ for
$\Mach{\mathrm{X86}}(\pass{compile\text{-}x86}(M))$),
\begin{align*}
  \omega \in D &\iff \Term_L(\error) \\
  c \in D &\iff \Term_L(c) \\
  \mathit{val}(D) &\iff \exists o \neq \error.\; \Term_L(o) \\
  D = \emptyset &\iff \Div_L .
\end{align*}
\end{theorem}
Soundness (a run shrinks the denotation of the machine's state) and
adequacy (a logical relation between values of the domain and of the
machine, which relates $\omega$ to a run that stops with an error) are as
for the graph model.

\begin{theorem}[End to End, with Errors]
For every closed ISWIM term $M$, with $L = \Mach{\mathrm{ISWIM}}(M)$ and
$L' = \Mach{\mathrm{GlobalFun}}(\pass{compile}(M))$ or
$L' = \Mach{\mathrm{X86}}(\pass{compile\text{-}x86}(M))$,
\[
  \Term_L(\error) \iff \Term_{L'}(\error), \quad
  \Term_L(c) \iff \Term_{L'}(c), \quad
  (\exists o \neq \error.\; \Term_L(o)) \iff (\exists o \neq \error.\; \Term_{L'}(o)), \quad
  \Div_L \iff \Div_{L'} .
\]
\end{theorem}
This is $\Equiv$ for the observations of errors and constants, and for
divergence; an answer that is not a constant (a function, a tuple, an
injection) is observed only as an answer, since delay and the later
passes change the representation of closures.

\end{document}
